\documentclass[11pt,letterpaper]{article}
\usepackage[T1]{fontenc}
\usepackage[utf8]{inputenc}
\usepackage{lmodern}
\usepackage[margin=1in,headheight=15pt]{geometry}
\usepackage{amsmath,amssymb,amsthm,mathtools,mathrsfs}
\usepackage{microtype,booktabs,array,longtable,enumitem}
\usepackage[dvipsnames]{xcolor}
\usepackage{fancyhdr,aliascnt}
\usepackage[colorlinks=true,linkcolor=MidnightBlue,citecolor=MidnightBlue,urlcolor=MidnightBlue]{hyperref}
\usepackage[nameinlink,noabbrev]{cleveref}
\hypersetup{%
 pdftitle={Surcomplex Dynamics: Linearization Thresholds, Periods, and Normal Forms},
 pdfauthor={Merged research report, Surreal project},
 pdfsubject={Hahn-supported dynamics: linearization, resonance, periods, normal forms}}
\pagestyle{fancy}\fancyhf{}
\fancyhead[L]{\small Surcomplex dynamics: thresholds, periods, normal forms}
\fancyhead[R]{\small\thepage}
\renewcommand{\headrulewidth}{0.3pt}
\setlength{\parskip}{0.25em}
\setlength{\parindent}{1.3em}
\setlength{\emergencystretch}{2em}
\setlist{nosep,leftmargin=2em}
\allowdisplaybreaks[2]

\newtheorem{theorem}{Theorem}[section]
\newaliascnt{lemma}{theorem}
\newtheorem{lemma}[lemma]{Lemma}
\aliascntresetthe{lemma}
\newaliascnt{proposition}{theorem}
\newtheorem{proposition}[proposition]{Proposition}
\aliascntresetthe{proposition}
\newaliascnt{corollary}{theorem}
\newtheorem{corollary}[corollary]{Corollary}
\aliascntresetthe{corollary}
\theoremstyle{definition}
\newaliascnt{definition}{theorem}
\newtheorem{definition}[definition]{Definition}
\aliascntresetthe{definition}
\newaliascnt{example}{theorem}
\newtheorem{example}[example]{Example}
\aliascntresetthe{example}
\newaliascnt{question}{theorem}
\newtheorem{question}[question]{Question}
\aliascntresetthe{question}
\theoremstyle{remark}
\newaliascnt{remark}{theorem}
\newtheorem{remark}[remark]{Remark}
\aliascntresetthe{remark}
\newaliascnt{warning}{theorem}
\newtheorem{warning}[warning]{Warning}
\aliascntresetthe{warning}
\newaliascnt{convention}{theorem}
\newtheorem{convention}[convention]{Convention}
\aliascntresetthe{convention}

\crefname{lemma}{lemma}{lemmas}
\Crefname{lemma}{Lemma}{Lemmas}
\crefname{proposition}{proposition}{propositions}
\Crefname{proposition}{Proposition}{Propositions}
\crefname{corollary}{corollary}{corollaries}
\Crefname{corollary}{Corollary}{Corollaries}
\crefname{definition}{definition}{definitions}
\Crefname{definition}{Definition}{Definitions}
\crefname{example}{example}{examples}
\Crefname{example}{Example}{Examples}
\crefname{question}{question}{questions}
\Crefname{question}{Question}{Questions}
\crefname{remark}{remark}{remarks}
\Crefname{remark}{Remark}{Remarks}
\crefname{warning}{warning}{warnings}
\Crefname{warning}{Warning}{Warnings}
\crefname{convention}{convention}{conventions}
\Crefname{convention}{Convention}{Conventions}

\newcommand{\C}{\mathbb C}
\newcommand{\R}{\mathbb R}
\newcommand{\Q}{\mathbb Q}
\newcommand{\Z}{\mathbb Z}
\newcommand{\N}{\mathbb N}
\newcommand{\No}{\mathbf{No}}
\newcommand{\SC}{\mathbf{SC}}
\newcommand{\HH}{\mathscr H}
\newcommand{\OO}{\mathcal O}
\newcommand{\MM}{\mathcal M}
\newcommand{\GG}{\mathscr G}
\newcommand{\AAA}{\mathscr A}
\newcommand{\BB}{\mathscr B}
\newcommand{\EE}{\mathscr E}
\newcommand{\PP}{\mathscr P}
\newcommand{\LL}{\mathscr L}
\newcommand{\m}{\mathfrak m}
\newcommand{\id}{\operatorname{id}}
\newcommand{\supp}{\operatorname{supp}}
\newcommand{\st}{\operatorname{st}}
\newcommand{\Cent}{\operatorname{Cent}}
\newcommand{\Fix}{\operatorname{Fix}}
\newcommand{\ord}{\operatorname{ord}}
\newcommand{\Res}{\operatorname{Res}}
\newcommand{\res}{\operatorname{res}}
\newcommand{\Log}{\operatorname{Log}}
\newcommand{\ExpH}{\operatorname{Exp}_{\!H}}
\newcommand{\rad}{\operatorname{rad}}
\newcommand{\radiso}{\operatorname{rad}_{\mathrm{iso}}}
\newcommand{\word}{\mathcal W}
\newcommand{\Li}{\mathcal L}
\newcommand{\NN}{\mathcal N}
\newcommand{\ad}{\operatorname{ad}}
\newcommand{\Disc}{\operatorname{Disc}}
\newcommand{\eps}{\varepsilon}
\newcommand{\dd}{\,\mathrm d}
\newcommand{\sdz}{\mathsf{SD}_{0}}
\newcommand{\sddr}{\mathsf{SD}_{\mathrm{dr}}}
\newcommand{\sdex}{\mathsf{SD}_{\mathrm{ex}}}
\newcommand{\sdfr}{\mathsf{SD}_{\mathrm{fr}}}
\newcommand{\sdno}{\mathsf{SD}_{\varnothing}}
\newcommand{\resq}{\mathsf{RES}_q}
\newcommand{\mres}{\mathsf{MR}}
\newcommand{\nres}{\mathsf{NR}}
\newcommand{\fres}{\mathsf{FR}}
\newcommand{\pares}{\mathsf{PAR}}
\newcommand{\qres}{\mathsf{QR}}
\newcommand{\sdss}{\mathsf{SD}_{\mathrm{ss}}}
\newcommand{\TT}{\mathbb T}
\newcommand{\strip}{\OO(\TT^{d}_{\rho})}
\newcommand{\Cups}{\mathcal C_{\Upsilon}}
\newcommand{\Gups}{\mathcal G_{\Upsilon}}
\newcommand{\Tups}{\mathsf T_{\Upsilon}}
\newcommand{\qlat}{\mathsf N_{\Upsilon}}
\newcommand{\kups}{\kappa_{\Upsilon}}
\newcommand{\fla}{\mathfrak a}
\newcommand{\av}{\operatorname{av}}
\newcommand{\Qz}{\mathsf Q_{0}}
\newcommand{\diver}{\operatorname{div}}
\newcommand{\rank}{\operatorname{rank}}
\newcolumntype{P}[1]{>{\raggedright\arraybackslash}p{#1}}
\numberwithin{equation}{section}

\begin{document}
\hypersetup{pageanchor=false}
\begin{titlepage}
\centering
\vspace*{0.35in}
{\large\scshape A merged research report of the Surreal project\par}
\vspace{0.40in}
{\Huge\bfseries Surcomplex Dynamics\par}
\vspace{0.20in}
{\LARGE\bfseries Linearization Thresholds, Periods,\par}
\vspace{0.08in}
{\LARGE\bfseries and Normal Forms\par}
\vspace{0.35in}
{\large 21 September 2026\par}
\vspace{0.40in}
\begin{minipage}{0.92\textwidth}
\begin{abstract}
Seven independently written manuscripts developed local and global dynamics
over complex Hahn fields $K_\Gamma=\C((t^\Gamma))$ realized inside
$\No[i]$, with $\Gamma$ a set-sized ordered abelian group of arbitrary rank.
This report is their \emph{union}. Every result any one of them reached
survives; results all of them reached are printed once; results proved by
genuinely different routes keep each route, marked; and every honest
limitation is preserved.

Six of the seven are germ and fixed-point dynamics: multipliers, return
iterates, periods, Hamiltonians at the origin of a phase halo. The seventh
is the \emph{quasi-periodic} half of the same subject --- a cohomological
equation on a real torus whose frequency vector has surreal components,
with divisors indexed by an integer lattice --- and it is not a duplicate
of any of the six: none of them contains a resonance lattice, a
cohomological equation, a Fourier divisor or a conjugacy to a constant
vector field.

Five inequivalent conditions travelled under the single word
\emph{resonance} across the seven sources, and six inequivalent conditions
travelled under the words \emph{small divisor}; worse, two families of
sources used the word \emph{radius} with opposite monotonicity.
\Cref{dyn:sec:notions} therefore opens the report with a notion table that
gives each variant its own name and symbol, and every threshold statement
below names which variant its hypothesis uses and in which coefficient
category it holds. A threshold proved under one notion is never quoted
under another.

The mathematical content: a sharp five-row coefficient-category
classification of universal linearization for an exact nonresonant diagonal
unitary multiplier, keyed to one exponential divisor rate, with necessity
witnessed at a single infinitesimal weight; an exact radius-collapse law
$\rad(h_{k,1})=R\,e^{-(k+1)\sigma}$ in the \emph{drifted}-multiplier
category, which settles the common-domain germ question there and leaves it
open for a fixed multiplier; an exact maximal centered valuation ball read
off a \emph{finite return iterate}, with its shell cycles, its
cancellation phase diagram and its common-domain lifting; an algebraicity
dichotomy and an exact coefficientwise monodromy exponent for parabolic
Euler maps, whose infinitesimal part no finite branched cover removes; a
complete period classification, explicit moduli and beyond-all-orders
invisible moduli for global near-identity maps; a gauge-normalized
symplectic Hamiltonian normal form, its full Poisson centralizer, and
support-certified domains reaching infinite phase magnitudes; and, on the
quasi-periodic side, the theorem that a frequency vector with $d$ surreal
components has \textbf{at most $d$} arithmetically visible divisor
valuations, identified by a canonical descending flag of integer resonance
lattices, together with the exact arithmetic criterion for universal
solvability of the cohomological equation on one fixed analytic strip, an
arbitrary-order analytic lifting obstruction, a unique infinitesimal
conjugacy to a constant vector field proved by marked support words rather
than by sequential valuation convergence, and its canonical invariant
density.
\end{abstract}
\end{minipage}
\vfill
\begin{minipage}{0.92\textwidth}\small
\textbf{Status.} This is a merged report assembled from seven AI-assisted,
unrefereed research drafts. The full report is not proof-assistant verified.
The repository's Lean development checks the finite-word support lemma and
the input-relative ancestry bounds; its coverage ledger is
\texttt{docs/FORMALIZATION.md}. Priority is not certified, and no named published
conjecture is claimed solved. The seven sources' own verification programs
supply finite exact symbolic checks only; \cref{dyn:sec:verification}
records exactly what each one does and does not establish, and
\cref{dyn:sec:nonclaims} carries every limitation stated by every source.
\end{minipage}
\end{titlepage}
\hypersetup{pageanchor=true}
\pagenumbering{roman}
{\small\setlength{\parskip}{0pt}\tableofcontents}
\clearpage
\pagenumbering{arabic}

\section{What this report is, and how to read it}
\label{dyn:sec:read}

\subsection{Seven manuscripts, one body of machinery}
The seven source manuscripts were all
written independently on 21 September 2026 against a pinned repository
snapshot: sources 01--06 against commit
\texttt{39f2be66\allowbreak 67ade51b\allowbreak ca2b45da\allowbreak
a47e289d\allowbreak 69c09764} and source 07 against commit
\texttt{aa846271\allowbreak b4dcae2c\allowbreak 055b2161\allowbreak
26a87210\allowbreak 292ec19b}. They share a single body
of machinery --- the set-sized Hahn workspace inside $\No[i]$, strong
summability by well-ordered support and local finiteness, the
positive-support lemma of Neumann in finite-word form, halo evaluation by
coefficientwise double Taylor sums, Taylor substitution and its
near-identity inverse --- and each of the seven states that machinery from
scratch. The positive-support lemma is stated in all seven. Between $45$
and $55$ per cent of the combined text of sources 01--06 is the same
theory, and source 07 restates the same foundation again.

Sources 01--06 are germ and fixed-point dynamics and were merged first;
source 07 is the quasi-periodic half of the subject and was integrated
afterwards, as \cref{dyn:sec:flag,dyn:sec:qplinear,dyn:sec:qpnonlinear}.
It is genuinely additive rather than a seventh treatment of the same
material: grepping sources 01--06 for a resonance lattice, a cohomological
equation, a Fourier divisor or a mean-zero conjugacy to a constant vector
field returns nothing, and conversely source 07 contains no multiplier, no
return iterate and no period. What the two halves do share is this
machinery, which is therefore stated once for both, and the
coefficient-category architecture of \cref{dyn:rem:oneinvariant}, which is
stated once and instantiated twice
(\cref{dyn:rem:qponeinvariant}).

They also disagree, in the specific and dangerous way that independent
runs at one \emph{foundation} disagree: not by contradicting each other's
theorems, but by attaching the same words to different hypotheses. That is
what \cref{dyn:sec:notions} exists to prevent.

\begin{convention}[The merge rule]\label{dyn:conv:merge}
A merge is a union, not a selection.
\begin{enumerate}[label=(\roman*)]
\item A statement proved by every source is printed once.
\item A statement proved by genuinely different routes keeps every route,
explicitly marked, because the routes carry different resources and
different by-products.
\item Anything one source reached that the others did not survives in full.
\item Every honest limitation survives. \Cref{dyn:sec:nonclaims} is the
consolidated record; dropping one of its items would silently upgrade the
report, and that is the worst failure available in an assembly of this
kind.
\item Nothing is strengthened beyond what a source proves. Where a source
records a hypothesis as load-bearing, the hypothesis travels with the
statement everywhere it is quoted.
\end{enumerate}
\end{convention}

\subsection{Discipline inherited from companion reports}
Three constructions here are inherited rather than restated, and are named
by the title and repository path of the companion report that owns them.
\Cref{dyn:sec:companions} gives the full inheritance ledger. In brief:
\begin{itemize}
\item the no-topological-convergence discipline --- every set of surcomplex
numbers is closed and discrete in the fine order topology, so
$\sum_{n}t^{n}=(1-t)^{-1}$ is a Hahn identity about a summable family and
never a limit of partial sums --- is the foundation fixed by
\emph{Surcomplex analysis: holomorphic functions on $K=\No[i]$}
(\texttt{docs/surcomplex/analysis/}). All seven sources restate it; we cite
it and state it once, in \cref{dyn:def:summable}.
\item the coefficient-ring taxonomy is the four-ring table of
\emph{Infinitesimal Analytic Geometry over the Surcomplex Numbers}
(\texttt{docs/surcomplex/analytic-geometry/}), whose names we adopt
verbatim. We add exactly two rows to it, and they are new \emph{rows},
not new rings: see \cref{dyn:subsec:rings}.
\item the warning that a Hahn constant for a coordinate derivative need not
be a constant for the scalar Berarducci--Mantova derivation is the
cross-category warning of \emph{Differential Algebra and Differential
Equations over the Surreal and Surcomplex Numbers}
(\texttt{docs/surcomplex/differential-equations/}). All seven sources
derive it independently --- source 07 in the form that its torus
derivatives are external derivations annihilating $K_\Gamma$ --- and we
adopt that report's wording once, in \cref{dyn:warn:derivative}.
\end{itemize}

\subsection{What is offered as new, and what is not}
Not new, and credited as classical throughout: the positive-support lemma
(Neumann, via Higman's word well-quasi-order) \cite{Neumann,Higman};
algebraic closure of a maximally complete Hahn field \cite{Poonen}; formal
logarithms of tangent-to-identity maps, formal flow embedding and
semisimple--unipotent splitting \cite{Ribon,GV}; the exponential
correspondence for strongly summable derivations of generalized-series
algebras \cite{BKKPS}; the ordinary homological equation and the
Brjuno--Yoccoz theory of small divisors
\cite{CarlettiMarmi,Yoccoz,BuffCheritat,Marmi}; rank-one ultrametric
linearization and its disk estimates \cite{Lindahl,LindahlPadic}; the
classical iterative residue and modified-equation coefficients
\cite{Hairer,AB}; and classical formal Birkhoff normalization with its
convergence--divergence theory \cite{PerezMarco,Krikorian}.

Offered for scrutiny as candidate contributions, in the exact form proved
below and with no priority claim: the coefficient-category classification
\cref{dyn:thm:main} with its first-weight necessity mechanism; the exact
drifted radius-collapse law \cref{dyn:thm:collapse} and the trichotomy
\cref{dyn:thm:trichotomy} it yields; the exact finite-return threshold
\cref{dyn:thm:exact} with \cref{dyn:thm:shell,dyn:thm:phase,dyn:thm:coherent};
the algebraicity dichotomy \cref{dyn:thm:dichotomy} with the monodromy
formula \cref{dyn:thm:monodromy} and its torsion obstruction; the period
classification \cref{dyn:thm:classification} with
\cref{dyn:thm:moduli,dyn:thm:flat,dyn:thm:slowtime}; the Hamiltonian
package \cref{dyn:thm:normalform,dyn:thm:universal,dyn:thm:pcentralizer,dyn:thm:rescale};
and, from source 07, the finite resonance flag with its single common
reciprocal support \cref{dyn:thm:flag}, the exact fixed-strip solvability
criterion \cref{dyn:thm:qplinear}, the arbitrary-order analytic lifting
obstruction \cref{dyn:thm:qpjet}, and the torus normal-form package
\cref{dyn:thm:qpnormalform,dyn:thm:qpnecessity,dyn:prop:qpboundary,dyn:thm:qpdensity}.
Each source's own bounded novelty assessment is reproduced in
\cref{dyn:sec:nonclaims}. Source 07 adds to the classical list the ordinary
Fourier division mechanism with its subexponential criterion, the Liouville
construction method, and positive-support Hahn inversion
\cite{Bounemoura,Kamienski,Neumann,Poonen}, none of which it claims.

\section{The notion table: five resonances, six small divisors, two radii}
\label{dyn:sec:notions}

This section is the load-bearing apparatus of the merge. Read it before
anything else. Every threshold statement in this report carries tags from
the tables below, and \cref{dyn:subsec:theoremtags} lists, theorem by
theorem, which tags it actually uses.

\subsection{Why unifying the notation without this table would invert results}
The seven sources use the word \emph{resonance} for five inequivalent
conditions and the words \emph{small divisor} for six, and two groups of
sources use the word \emph{radius} with opposite monotonicity. A merge that
harmonized the symbols without separating the notions would produce
statements that read correctly and are false: it would apply an exact
maximal-ball formula to a multiplier its source explicitly excludes,
attach an arithmetic threshold to a theorem proved with no arithmetic
hypothesis at all, and reverse the direction of every domain inequality in
half the report.

\subsection{Four notions of resonance}
\label{dyn:subsec:resonance}

\begin{longtable}{@{}P{0.13\textwidth}P{0.40\textwidth}P{0.40\textwidth}@{}}
\toprule
Tag & Condition & Where it is the hypothesis, and what it is not\\
\midrule
\endhead
$\resq$\newline \emph{residue-resonant of order $q$}
&
$\lambda\in K_\Gamma^\times$ with $v(\lambda)=0$, $\lambda$ is
\textbf{never} a root of unity, and its residue
$\alpha=\bar\lambda\in\C^\times$ is a root of unity of exact order
$q\ge1$. Then $c=\lambda^{q}-1\ne0$ and $\delta=v(c)>0$.
&
This is the hypothesis of the exact finite-return threshold
\cref{dyn:thm:exact}, the shell theorem \cref{dyn:thm:shell}, the phase
diagram \cref{dyn:thm:phase} and the common-domain lifting
\cref{dyn:thm:coherent}. It is \textbf{not} $\lambda$ itself being a root
of unity --- that case is excluded --- and it is \textbf{not} a divisor
size condition of any kind.\\[4pt]
$\mres$ / $\nres$\newline \emph{multiplier-resonant}
&
For $\Lambda=\operatorname{diag}(\lambda_1,\dots,\lambda_d)$ with
$|\lambda_j|=1$: $\lambda^{\alpha}=\lambda_j$ for some $|\alpha|\ge2$ and
some $j$. $\nres$ is its negation, nonresonance. In one variable
$\nres$ reads $\lambda^{n}\ne\lambda$ for $n\ge2$.
&
This is the hypothesis of \cref{dyn:thm:main,dyn:thm:lifting,dyn:thm:trichotomy}
and of everything built on the homological operator $\Li$ of
\cref{dyn:subsec:homological}. The multiplier is \emph{exactly} $\Lambda$
in \cref{dyn:thm:main} and \emph{drifted} in \cref{dyn:thm:collapse}; see
\cref{dyn:warn:drift}.\\[4pt]
$\fres$\newline \emph{frequency-resonant}
&
For a frequency vector $\omega\in\C^{d}$ of ordinary complex constants:
$\omega\cdot k=0$ for some $0\ne k\in\Z^{d}$; equivalently
$\omega\cdot(\alpha-\beta)=0$ for a monomial $q^{\alpha}p^{\beta}$. Its
negation is the hypothesis of the Hamiltonian package.
&
$\fres$ is the exact Hamiltonian analogue of $\mres$, and its negation is
what makes the kernel of $L=\{\cdot,H_0\}$ exactly the diagonal monomials
$\alpha=\beta$. It is a condition on \emph{ordinary complex} numbers,
never on Hahn or surreal data.\\[4pt]
$\pares$\newline \emph{parabolic}
&
The leading multiplier is exactly $1$, perturbed infinitesimally:
$F=z+h$ with $v(h)>0$, or $G_\eps=w+\eps V(w)$ with $\eps=t^{\delta}$,
$\delta>0$.
&
This is the hypothesis of the whole of
\cref{dyn:sec:parabolic-monodromy,dyn:sec:periods}. \textbf{No divisor
hypothesis of any kind is imposed anywhere under $\pares$.} Source 04
says plainly that the word in its own title is not a small-divisor
condition, and source 05 says it has no resonance and no small-divisor
condition at all because its maps are unipotent. A reader who reads
$\pares$ as a small-divisor hypothesis will misattribute every theorem in
those two sections.\\[4pt]
$\qres$\newline \emph{lattice-resonant, for a Hahn frequency}
&
For a frequency vector $\Upsilon\in F_\Gamma^{d}=\R((t^{\Gamma}))^{d}$ whose
components are \textbf{Hahn (hence surreal) scalars}: the exact resonance
lattice $\qlat=\{k\in\Z^{d}:k\cdot\Upsilon=0\}$ is nonzero. Its negation
$\lnot\qres$ is $\qlat=\{0\}$.
&
The flag \cref{dyn:thm:flag} and the linear theorem
\cref{dyn:thm:qplinear} \emph{permit} $\qres$; the torus normal form
\cref{dyn:thm:qpnormalform} and all of \cref{dyn:sec:qpnonlinear} require
$\lnot\qres$. This is \textbf{exact vanishing of a Hahn scalar}, never
smallness of anything. It is \textbf{not} $\fres$: that one is a condition
on ordinary complex $\omega\in\C^{d}$ with divisors graded by Taylor
degree, whereas $\qres$ concerns surreal components with divisors
$k\cdot\Upsilon$ graded by the Fourier index $|k|$ and of \emph{nonzero}
valuation (\cref{dyn:warn:twofrequencies}).\\
\bottomrule
\end{longtable}

\begin{warning}[The four germ notions do not intersect in the way the word
suggests]
\label{dyn:warn:resonance}
$\resq$ requires $\lambda$ \emph{not} to be a root of unity while its
residue is one; $\mres$ is about exact equality of ordinary unit
multipliers; $\fres$ is a linear condition on an integer vector; $\pares$
is about the multiplier being $1$ up to an infinitesimal. In particular:
\begin{itemize}
\item the exact maximal-ball formula of \cref{dyn:thm:exact} needs
$\resq$ and must never be applied under $\pares$ or under $\mres$;
\item the resonant projection and the retained resonant Hamiltonian
algebra of \cref{dyn:sec:hamiltonian} need $\fres$;
\item the parabolic results of \cref{dyn:sec:parabolic-monodromy,dyn:sec:periods}
need none of the three others, and acquire no arithmetic hypothesis by
being printed next to them;
\item the finite-order splitting of \cref{dyn:thm:finite-order} needs
$\Lambda^{q}=I$ --- the multiplier itself of finite order --- which is
\emph{neither} $\resq$ (where $\lambda$ is not a root of unity) nor mere
$\mres$ (where resonances exist but need not have finite order).
\end{itemize}
The fifth notion $\qres$ intersects none of these four: it is a condition
on an integer lattice attached to a vector of \emph{Hahn} scalars, it has
no multiplier and no fixed point, and no theorem stated under $\resq$,
$\mres$, $\fres$ or $\pares$ is quoted under it or vice versa.
\end{warning}

\subsection{Six notions of small divisor}
\label{dyn:subsec:smalldivisor}

The first five concern \emph{ordinary complex} divisors. Every nonzero ordinary
complex number, however tiny, has Hahn valuation zero
(\cref{dyn:rem:invisible}), so none of the five is visible to the Hahn
valuation, and a proof that keeps only the valuation cannot distinguish
them. The sixth, $\sdss$, is different in one respect and identical in
another: its divisors $k\cdot\Upsilon$ are Hahn scalars of \emph{nonzero}
valuation, but the arithmetic condition it imposes is again a condition on
ordinary real numbers only --- the leading real forms --- so that here too
the valuation cannot see the arithmetic
(\cref{dyn:warn:qptwosizes}).

\begin{longtable}{@{}P{0.12\textwidth}P{0.41\textwidth}P{0.40\textwidth}@{}}
\toprule
Tag & Condition & Scope\\
\midrule
\endhead
$\sdz$\newline \emph{support only}
&
No arithmetic condition whatsoever. Control is by well-ordered positive
support and local finiteness alone.
&
Sources 01, and the polynomial and formal rows of \cref{dyn:thm:main},
and the polynomial-layer case (a) of \cref{dyn:thm:normalform}. Source 01
insists on the point: the Schr\"oder denominators $\nu-\nu^{n}$ have
nonzero ordinary constant terms $\alpha-\alpha^{n}$, which may be
arbitrarily small in the ordinary modulus and still all have Hahn
valuation zero, so no Diophantine, Brjuno or Siegel condition on surreal
data appears in \cref{dyn:sec:return} at all.\\[4pt]
$\sddr$\newline \emph{drifted exponential rate}
&
$\displaystyle
\sigma(\lambda)=\limsup_{n\to\infty}
\frac{\log^{+}|\lambda^{n}-\lambda|^{-1}}{n}\in[0,\infty]$,
\textbf{together with multiplier drift} $F'(0)=\lambda+u$,
$v(u)>0$.
&
Source 02. The drift is not decoration: it is what creates the higher
reciprocal jets $c_{n,k}$, $k\ge1$, whose poles of order $k+1$ drive the
collapse law. \Cref{dyn:warn:drift} records that removing drift removes
the mechanism.\\[4pt]
$\sdex$\newline \emph{exact-multiplier exponential rate}
&
$\displaystyle
\tau(\Lambda)=\limsup_{m\to\infty}\frac1m\log M_{m}(\Lambda)$,
$M_{m}=\max\bigl(1,\max_{|\alpha|=m,j}|\lambda^{\alpha}-\lambda_j|^{-1}\bigr)$,
with the multiplier \textbf{exactly} $\Lambda$: every
$f_\gamma$ has $f_\gamma(0)=0$ and $\mathrm df_\gamma(0)=0$.
&
Source 03. No infinitesimal linear part is permitted in
\cref{dyn:thm:main}; infinitesimal linear terms \emph{are} permitted in
the identity-reducing flow theorem \cref{dyn:thm:exp-log}.\\[4pt]
$\sdfr$\newline \emph{frequency exponential rate}
&
$\displaystyle \Theta(\omega)=\limsup_{D\to\infty}M_{D}(\omega)^{1/D}$,
$M_{D}=\max\bigl(1,\max_{|\alpha|+|\beta|=D,\ \alpha\ne\beta}
|\omega\cdot(\alpha-\beta)|^{-1}\bigr)$, so
$\Theta\in[1,\infty]$.
&
Source 06, where this invariant is called $\Lambda(\omega)$; we write
$\Theta(\omega)$ to free the symbol $\Lambda$ for the diagonal multiplier
matrix of sources 02 and 03. Note the \textbf{different normalization}:
$\sigma$ and $\tau$ are logarithmic rates in $[0,\infty]$ while $\Theta$
is their exponential. $\Theta(\omega)<\infty$ is the exact analogue of
$\tau<\infty$; $\Theta(\omega)=1$ would be the analogue of $\tau=0$ and
is \textbf{nowhere required}, because source 06's entire coefficients
already have the common domain $\C^{2d}$ and so lose nothing to a finite
radius factor.\\[4pt]
$\sdno$\newline \emph{none}
&
No divisor condition, and no divisor family to impose one on.
&
Sources 04 and 05. Under $\pares$ the multiplier is $1+\mu\eps$ or
$1$, so $\lambda^{n}-\lambda$ is automatically nonzero for every
$n\ge2$, formal uniqueness is free, and the entire difficulty is
well-ordered support.\\[4pt]
$\sdss$\newline \emph{stratified subexponential}
&
For each flag stratum $\Sigma_j$ of \cref{dyn:thm:flag} and every real
$\eps>0$ there is $C_{j,\eps}>0$ with
$|\fla_j(k)|^{-1}\le C_{j,\eps}e^{\eps|k|}$ for $k\in\Sigma_j$, where
$\fla_j(k)=k\cdot\varpi_{\gamma_j}$ is the leading \emph{real} form of the
divisor $k\cdot\Upsilon$ on that stratum; equivalently
$\limsup_{k\in\Sigma_j}\log^{+}(1/|\fla_j(k)|)/|k|=0$ for every $j$.
&
Source 07, where it is called $(\mathrm{SS})$. It is the vanishing-rate
condition of \cref{dyn:rem:oneinvariant} on a \emph{fourth} divisor family
(\cref{dyn:rem:qponeinvariant}), imposed stratum by stratum on the
finitely many leading real forms. It therefore corresponds to $\tau=0$ and
$\sigma=0$, \textbf{not} to $\tau<\infty$; and it says nothing whatever
about the surreal size $v(k\cdot\Upsilon)=\gamma_j$ of the divisor itself
(\cref{dyn:warn:qptwosizes}).\\
\bottomrule
\end{longtable}

\begin{remark}[One invariant, three names, one category conclusion]
\label{dyn:rem:oneinvariant}
Fix a family $(d_i)_{i\in I}$ of nonzero ordinary complex divisors graded
by a degree $\deg i\in\N$, and set
\[
 \kappa=\limsup_{m\to\infty}\frac1m
 \log\max\Bigl(1,\max_{\deg i=m}|d_i|^{-1}\Bigr).
\]
Then $\sigma$, $\tau$ and $\log\Theta$ are this one quantity evaluated on
three different divisor families: $\lambda^{n}-\lambda$ graded by $n$;
$\lambda^{\alpha}-\lambda_j$ graded by $|\alpha|$; and
$\omega\cdot(\alpha-\beta)$ graded by $|\alpha|+|\beta|$. The
\emph{category conclusion} is the same in all three instances, and is the
single most quotable fact in this report:
\begin{center}
\begin{tabular}{@{}ll@{}}
\toprule
Coefficient category & Universal criterion\\
\midrule
polynomial coefficients & nonresonance only ($\sdz$)\\
formal coefficients & nonresonance only ($\sdz$)\\
radius-free convergent germs & $\kappa<\infty$\\
entire coefficients & $\kappa<\infty$\\
one fixed unchanged polydisk & $\kappa=0$\\
\bottomrule
\end{tabular}
\end{center}
The three sources reach this table on three divisor families, in three
notations, with three different proofs of necessity. What they do
\emph{not} share is the sixth category, common-domain germs; see
\cref{dyn:sec:drift} and \cref{dyn:q:commongerm}.

A \emph{fourth} divisor family, $k\cdot\Upsilon$ for a Hahn frequency
vector graded by the Fourier index $|k|$, is added by source 07, and the
first row of the table above is instantiated on it as
\cref{dyn:thm:qplinear}: ``one fixed unchanged polydisk'' becomes ``one
fixed unchanged strip'' and ``$\kappa=0$'' becomes $(\sdss)$, the same
vanishing rate imposed on each of the finitely many strata. The table is
stated here once and instantiated exactly twice; the two instantiations,
and the rows that are \textbf{not} instantiated on the quasi-periodic
side, are set out in \cref{dyn:rem:qponeinvariant}. Reading an
uninstantiated row across would assert a theorem no source proves.
\end{remark}

\begin{remark}[Arithmetic invisible to the Hahn valuation]
\label{dyn:rem:invisible}
For every nonzero ordinary divisor, $v(\lambda^{\alpha}-\lambda_j)=0$ and
$v(\omega\cdot(\alpha-\beta))=0$. The distinction among $\sdz$, $\sddr$,
$\sdex$ and $\sdfr$ therefore comes entirely from the \emph{ordinary
complex Taylor growth inside each coefficient function}, not from the
valuation. This is exactly why formal and polynomial coefficients always
permit normalized linearization under nonresonance while the analytic
categories have sharp arithmetic thresholds, and it is why the negative
results below never prohibit a conjugacy on an infinitesimal monad: the
formal-coefficient conjugacy exists and can be evaluated there
(\cref{dyn:prop:monad}).
\end{remark}

\subsection{Two notions of radius, running in opposite directions}
\label{dyn:subsec:radius}

\begin{longtable}{@{}P{0.16\textwidth}P{0.37\textwidth}P{0.40\textwidth}@{}}
\toprule
Object & Definition & Monotonicity and sources\\
\midrule
\endhead
valuation ball $B_{r}$, shell $S_{r}$
&
$B_{r}=\{z\in K_\Gamma:v(z)>r\}=t^{r}\m$ and
$S_{r}=\{z:v(z)=r\}$, for $r\in\Gamma$.
&
\textbf{Increasing $r$ makes the ball smaller.} Sources 01
($B_{r}$, $r=\max(\delta-v(b_n))/(n-1)$) and 04
($D=\{v(z)>\delta/q\}$, so larger $\delta/q$ is a smaller disk). $S_{r}$
is called a bounding shell and is \textbf{not} the topological boundary
(valuation balls are clopen) and \textbf{not} a circle for the
ordered-surreal modulus.\\[4pt]
valuation polydisk $D_{\rho}$
&
$D_{\rho}=(t^{\rho}\OO_\Gamma)^{2d}=\{z:v(z_j)\ge\rho\ \forall j\}$.
&
\textbf{Increasing $\rho$ makes the polydisk smaller}; $\rho<0$ gives
domains containing phase points of \emph{infinite} magnitude. Source 06.
Same direction as $B_r$.\\[4pt]
ordinary disk / polydisk $\mathbb D_{R}$, $\mathbb D_{R}^{d}$, $D$
&
$\mathbb D_R^{d}=\{z\in\C^{d}:|z_j|<R\}$; $D$ an ordinary connected open
subset of $\C$.
&
\textbf{Increasing $R$ makes the domain larger.} Sources 02, 03, 05, and
the coefficient domains of 04 and 06. This is a domain for the ordinary
\emph{coefficient functions}, never a surcomplex ball.\\[4pt]
Taylor radius $\rad$, $\radiso$
&
$\rad(g)^{-1}=\limsup_n|g_n|^{1/n}$;
$\radiso(g)^{-1}=\limsup_m\bigl(\max_{|\alpha|=m,j}|g_{j,\alpha}|\bigr)^{1/m}$.
&
\textbf{Larger is better.} Sources 02 ($\rad$) and 03 (written $\rho$
there; we write $\radiso$, freeing $\rho$ for source 06's rescaling
exponent).\\
\bottomrule
\end{longtable}

\begin{warning}[Do not identify a valuation ball with an ordinary disk]
\label{dyn:warn:ball}
$B_{r}$ and $\mathbb D_{R}$ are objects of different kinds and their
inequalities run in opposite directions. Source 01 states the point
sharply: all points of the normalized monad have ordinary residue zero,
whereas the ordinary disk $D$ of \cref{dyn:thm:coherent} contains many
nonzero ordinary points; $D$ is a domain of holomorphy for coefficient
\emph{functions} and not a claim of a larger centered surcomplex
linearization ball. Likewise a Lindahl-type ultrametric disk radius in a
fixed rank-one workspace and source 01's maximal centered valuation ball
at arbitrary rank are different categories; see
\cref{dyn:subsec:berkovich}.
\end{warning}

\begin{warning}[Support orientation]\label{dyn:warn:orientation}
Throughout this report supports are well ordered in the \emph{increasing}
order of $\Gamma$ and $v=\min\supp$, so that $t^{\gamma}=\omega^{-\gamma}$
converts increasing Hahn exponents to decreasing surreal monomials. A
threshold rewritten using the Conway growth exponent $\alpha=-\gamma$
requires negating its exponent and reversing the corresponding exponent
inequality. Ordinary coefficient estimates are unchanged. Silently
transferring a valuation threshold without this conversion inverts it.
\end{warning}

\subsection{Tag ledger: which notion each threshold statement uses}
\label{dyn:subsec:theoremtags}

\begin{longtable}{@{}P{0.24\textwidth}P{0.10\textwidth}P{0.11\textwidth}P{0.22\textwidth}P{0.20\textwidth}@{}}
\toprule
Statement & Resonance & Divisor & Coefficient category & Domain notion\\
\midrule
\endhead
\cref{dyn:thm:exact} exact finite-return ball
& $\resq$ & $\sdz$ & formal, $\C[[X]]((t^\Gamma))$ & $B_{r}$ (valuation)\\
\cref{dyn:thm:shell} shell cycles
& $\resq$ & $\sdz$ & ordinary polynomial $P\in\C[X^{q}]$ & $S_{r}$ (valuation shell)\\
\cref{dyn:thm:phase} phase diagram
& $\resq$ & $\sdz$ & formal & $B_{r}$\\
\cref{dyn:thm:coherent} common-domain lifting
& $\resq$ & $\sdz$ & fixed-polydisk $\OO(D)((t^\Gamma))$ & $\mathbb D$ (ordinary)\\
\cref{dyn:prop:noncoherent} no common domain
& $\nres$ & $\sdz$ hypothesis, Brjuno used externally & fails radius-free germ & ordinary\\
\cref{dyn:thm:stability} return perturbation certificate
& $\resq$ & $\sdz$ & return coefficients only & $B_{r}$\\
\cref{dyn:thm:main} five-row classification
& $\nres$, exact $\Lambda$ & $\sdex$ & all five rows & $\mathbb D_R^{d}$, germ, entire, poly, formal\\
\cref{dyn:lem:radius-loss} single-step loss
& $\nres$ & $\sdex$ & isotropic Taylor radius & $\radiso$\\
\cref{dyn:lem:sharp-divisor} sharpness
& $\nres$ & $\sdex$ & $\OO(\mathbb D_R^d)$ and entire & $\radiso$\\
\cref{dyn:thm:radius-depth} word-depth radius bound
& $\nres$ & $\sdex$ & germ / fixed polydisk & $\radiso$, \emph{lower bound only}\\
\cref{dyn:thm:lifting} support-controlled lifting
& $\nres$ & $\sddr$ allowed (drift permitted) & four rows & $\mathbb D_R$\\
\cref{dyn:thm:collapse} exact radius collapse
& $\nres$ & $\sddr$, \textbf{drift required} & germ yes, common domain no & $\rad$\\
\cref{dyn:thm:trichotomy} trichotomy
& $\nres$ & $\sddr$ & common domain vs germ & $\mathbb D_R$\\
\cref{dyn:thm:multicollapse} several-variable collapse
& $\nres$ & $\sddr$, multiplicative detuning & polydisk germ & $\radiso$\\
\cref{dyn:thm:dichotomy} algebraicity dichotomy
& $\pares$ & $\sdno$ & $\C(w)((t^\Gamma))$ base & ordinary universal cover\\
\cref{dyn:thm:monodromy} finite-cover obstruction
& $\pares$ & $\sdno$ & rational-coefficient & ordinary loops\\
\cref{dyn:thm:sharp} sharp summability disk
& $\pares$ & $\sdno$ & formal Taylor family & $B_{\delta/q}$ (valuation)\\
\cref{dyn:thm:classification} period classification
& $\pares$ & $\sdno$ & $\OO(D)((t^\Gamma))$, no shrinking & ordinary $D$\\
\cref{dyn:thm:moduli} moduli $\m_K^{r}$
& $\pares$ & $\sdno$ & $\OO(D_r)((t^\Gamma))$ & ordinary punctured plane\\
\cref{dyn:thm:flat} invisible moduli
& $\pares$ & $\sdno$ & $\OO(D_r)((t^\Gamma))$ & ordinary; needs $J_\delta\ne0$\\
\cref{dyn:thm:homological} sharp homological criterion
& $\lnot\fres$ & $\sdfr$ & entire / germ & $\C^{2d}$ / germ radius\\
\cref{dyn:thm:normalform}(a) polynomial layers
& $\lnot\fres$ & $\sdz$ & $\PP_{2d}$ polynomial & finite halo $\OO_\Gamma^{2d}$\\
\cref{dyn:thm:normalform}(b) entire layers
& $\lnot\fres$ & $\sdfr$ & $\EE_{2d}=\OO(\C^{2d})$ & finite halo\\
\cref{dyn:thm:universal} universal criterion
& $\lnot\fres$ & $\sdfr$ & entire; germ at first order & finite halo\\
\cref{dyn:thm:pcentralizer} Poisson centralizer
& $\lnot\fres$ & $\sdfr$ (through 06(b)) & entire-coefficient & finite halo\\
\cref{dyn:thm:rescale} certified domains
& $\lnot\fres$ & $\sdz$ (polynomial input) & polynomial & $D_{\rho}$ (valuation)\\
\cref{dyn:thm:finite-order} finite-order splitting
& $\Lambda^{q}=I$ & $\sdz$ & all five rows & $\mathbb D_R^{d}$\\
\cref{dyn:thm:exp-log} exp/log correspondence
& $\pares$ & $\sdno$ & $\OO(U)$, germ, entire, poly, formal & arbitrary open $U$\\
\cref{dyn:thm:centralizer} centralizer
& $\pares$ & $\sdno$ & $\OO(U)$, $U\subset\C$ connected & ordinary $U$\\
\cref{dyn:thm:fixed-ideal} fixed-point ideals
& $\pares$ & $\sdno$ & any of the five & finite halo\\
\cref{dyn:thm:flag} finite resonance flag
& $\qres$ allowed & none used & \emph{no coefficient functions}: a statement
about $F_\Gamma^{d}$ and $\Z^{d}$ & none\\
\cref{dyn:prop:qpjets} inverse coefficients
& $\qres$ allowed & none used & scalar, in $K_\Gamma$ & none\\
\cref{dyn:prop:qpvalbound} arithmetic-free bound
& $\qres$ allowed & \textbf{none}: no divisor estimate at all
& fixed strip $\HH_\Gamma(\strip)$ & $\TT^{d}_\rho$ (ordinary)\\
\cref{dyn:thm:qplinear} fixed-strip solvability
& $\qres$ allowed & $\sdss$, \textbf{iff} & fixed strip
& $\TT^{d}_\rho$ (ordinary)\\
\cref{dyn:cor:qpsharp} loss $\kups$ attained
& $\qres$ allowed & $\sdss$ & fixed strip & $\TT^{d}_\rho$\\
\cref{dyn:cor:qptail} higher tails
& $\lnot\qres$ for the last clause & $\sdss$ preserved & fixed strip
& $\TT^{d}_\rho$\\
\cref{dyn:thm:qpjet} arbitrary-order lifting obstruction
& $\lnot\qres$ & $\sdss$ \emph{fails}; Liouville & entire input, fixed-strip
target & $\TT^{2}_\rho$, all $\rho$\\
\cref{dyn:thm:qpnormalform} torus normal form
& $\lnot\qres$ & $\sdss$; also $v(\mathcal V)>\kups$ & fixed strip
& $\TT^{d}_\rho$\\
\cref{dyn:lem:qpcoordinverse} coordinate inverses
& --- & none used & fixed strip & $\TT^{d}_\rho$\\
\cref{dyn:thm:qpnecessity} necessity beyond any scale
& $\lnot\qres$ & $\sdss$ \emph{fails} & fixed strip & $\TT^{d}_\rho$\\
\cref{dyn:prop:qpboundary} failure at $v(\mathcal V)=\kups$
& $\lnot\qres$ & $\sdss$ \emph{holds} & fixed strip & $\TT^{d}_\rho$\\
\cref{dyn:thm:qpdensity} invariant density
& $\lnot\qres$ & $\sdss$ (through \cref{dyn:thm:qpnormalform})
& fixed strip & $\TT^{d}_\rho$\\
\bottomrule
\end{longtable}

\noindent Three entries of this ledger deserve to be read twice.

\begin{warning}[The drift hypothesis travels with the radius law]
\label{dyn:warn:drift}
\Cref{dyn:thm:collapse} and its consequences
\cref{dyn:cor:finitefailure,dyn:thm:trichotomy,dyn:thm:multicollapse}
require $\sddr$: the multiplier drifts, $F'(0)=\lambda+u$ with
$v(u)>0$, and two rationally independent positive exponents are used
($\Gamma_{*}=\Q+\sqrt2\,\Q$, $u=t$, $v=t^{\sqrt2}$). Source 02 records
that eliminating the drift removes the direct reciprocal-jet mechanism and
that it does not settle the fixed-multiplier problem for
$0<\sigma<\infty$. The exact law $\rad(h_{k,1})=R\,e^{-(k+1)\sigma}$ must
therefore never be quoted without the drift hypothesis, and in particular
must not be quoted as an answer in the exact-multiplier category $\sdex$
of \cref{dyn:thm:main}. See \cref{dyn:sec:drift} for the resolution of
this apparent tension with \cref{dyn:q:commongerm}.
\end{warning}

\begin{warning}[$\resq$ and equal characteristic zero]
\label{dyn:warn:charzero}
The divisor cancellation behind \cref{dyn:thm:exact} leaves a factor
$n-1$ in the denominator (\cref{dyn:lem:displacement}). In equal
characteristic zero every such nonzero integer is a valuation unit; in
mixed characteristic it is not. Source 01 states outright that the exact
finite-return formula and the shell-periodicity conclusion must not be
transported to $\C_p$, and cites $p$-adic quadratic-like examples whose
linearization bounding sphere carries no periodic point at all
\cite[Cor.~C]{LindahlPadic}. That contrast identifies a hypothesis; it is
not a contradiction.
\end{warning}

\begin{warning}[$\pares$ acquires no arithmetic by proximity]
\label{dyn:warn:parabolic}
\Cref{dyn:sec:parabolic-monodromy,dyn:sec:periods} sit between two
sections full of divisor thresholds. They have none. Source 04's title
word ``resonance'' means only that the leading multiplier is $1$ perturbed
by an infinitesimal; source 05 has no eigenvalue ratios at all. Their
analogue of a linearization domain is a fixed common ordinary domain with
no shrinking whatever, and their analogue of a normal form is a rational
time form. A merge of either with a framework carrying a resonance or
small-divisor hypothesis would identify two different frameworks.
\end{warning}

\begin{warning}[Two frequency vectors, two divisor families, two $\kappa$'s]
\label{dyn:warn:twofrequencies}
This report now contains two frequency vectors and they are objects of
different kinds. Conflating them would attach a Hahn-scale statement to an
ordinary-complex hypothesis, or the reverse.
\begin{itemize}
\item $\omega\in\C^{d}$ of \cref{dyn:sec:hamiltonian} is a vector of
\textbf{ordinary complex constants}. Its divisors are
$\omega\cdot(\alpha-\beta)$ for Taylor multi-indices, graded by the
\emph{Taylor degree} $|\alpha|+|\beta|$, and they all have Hahn valuation
zero (\cref{dyn:rem:invisible}). Its resonance notion is $\fres$ and its
divisor notion $\sdfr$.
\item $\Upsilon\in F_\Gamma^{d}=\R((t^{\Gamma}))^{d}$ of
\cref{dyn:sec:flag,dyn:sec:qplinear,dyn:sec:qpnonlinear} is a
vector of \textbf{Hahn (hence surreal) scalars}. Its divisors are
$k\cdot\Upsilon$ for $k\in\Z^{d}$, graded by the \emph{Fourier index}
$|k|=|k_1|+\cdots+|k_d|$, and they have \emph{nonzero} Hahn valuations ---
that is the entire point of \cref{dyn:thm:flag}. Its resonance notion is
$\qres$ and its divisor notion $\sdss$
(\cref{dyn:subsec:resonance,dyn:subsec:smalldivisor}).
\end{itemize}
Further, $\kappa$ in \cref{dyn:rem:oneinvariant} is a generic
\emph{logarithmic divisor rate} in $[0,\infty]$, whereas $\kups$ in
\cref{dyn:sec:flag,dyn:sec:qplinear,dyn:sec:qpnonlinear} is an
\emph{element of $\Gamma$}: the valuation of the last visible flag level. They are never compared. The symbol $\Omega$ continues to
denote an infinite element of $\Gamma$
(\cref{dyn:ex:ranktwo,dyn:ex:rank2}) and $\Omega_F=\dd z/W_F$ a time form
(\cref{dyn:sec:periods}); the frequency vector is $\Upsilon$ and never
$\Omega$. The coordinate change $\Phi_h=\id+h$ of
\cref{dyn:thm:qpnormalform} is a change of torus coordinates, \emph{not}
the flow map $\Phi_X$ or $\Phi_c$ of \cref{dyn:thm:exp-log,dyn:cor:flow}.
\end{warning}

\begin{warning}[The visible levels are not a spectrum]
\label{dyn:warn:qpspectrum}
The finitely many visible divisor valuations $\gamma_1<\cdots<\gamma_r$ of
\cref{dyn:thm:flag} read like a spectrum of scales and are not one in any
sense used in this collection: not an eigenvalue list, not a multiset of
purely infinite phases, and not the prime spectrum of a ring. They are
values of the valuation $v$ on a family of Hahn scalars. The word
\emph{spectrum} is not used for them anywhere below.
\end{warning}

\begin{warning}[Arithmetic size and valuation size are incomparable]
\label{dyn:warn:qptwosizes}
In \eqref{dyn:eq:qpleading} the quantity $|\fla_j(k)|$ is an \emph{ordinary
real absolute value} while $v(k\cdot\Upsilon)=\gamma_j$ lies in $\Gamma$.
Neither bounds the other, and no inequality in one kind of size is used as
an estimate for the other anywhere below without the explicit
leading-coefficient formula \eqref{dyn:eq:qpleading}. A divisor can be
infinitesimal (large $\gamma_j$) with completely harmless arithmetic, and a
divisor of scale one ($\gamma_j=0$) can have disastrous arithmetic. This is
the torus form of \cref{dyn:rem:invisible}: the arithmetic lives in the
ordinary complex coefficient, not in the valuation --- with the one
difference that here the \emph{divisors themselves} do have nonzero
valuations, whereas the divisors of \cref{dyn:rem:invisible} all have
valuation zero.
\end{warning}

\section{Workspaces, supports, coefficient categories, and evaluation}
\label{dyn:sec:framework}

All seven sources set up this material independently. It is stated once
here.

\subsection{The set-sized workspace inside \texorpdfstring{$\No[i]$}{No[i]}}
Let $\Gamma$ be a nonzero ordered abelian group which is a \emph{set};
divisibility is assumed only where it is used and is flagged there. A Hahn
series over an ordinary commutative $\C$-algebra $B$ is a formal sum
$A=\sum_{\gamma\in S}t^{\gamma}A_\gamma$ with $A_\gamma\in B$ and
$S\subset\Gamma$ well ordered in the increasing order
(\cref{dyn:warn:orientation}), with the standard convolution product. For
$A\ne0$ put $v(A)=\min\supp A$ and $v(0)=+\infty$; for a vector or matrix
take the minimum over entries. Write
\[
 \HH_\Gamma(B)=B((t^{\Gamma})),\qquad
 K_\Gamma=\HH_\Gamma(\C)=\C((t^{\Gamma})),\qquad
 \OO_\Gamma=\{v\ge0\},\qquad \m_\Gamma=\{v>0\},
\]
and append a subscript $+$ or $>0$ for strictly positive support. Positive
support is a condition on \emph{exponents}, never on the signs of the
ordinary complex coefficients. The residue field of $K_\Gamma$ is
canonically $\C$, and a bar denotes reduction modulo $\m_\Gamma$; in
particular every nonzero ordinary complex number, including every nonzero
integer, has valuation zero. When $\Gamma$ is divisible, $K_\Gamma$ is
algebraically closed \cite[Cor.~4]{Poonen}, and algebraic closure is used
only in \cref{dyn:lem:root-count,dyn:thm:workspace}; the Hamiltonian
package of \cref{dyn:sec:hamiltonian} needs neither closure nor
divisibility.

For a set-sized ordered subgroup $\Gamma\subset\No$ the Conway normal form
and the substitution
\begin{equation}\label{dyn:eq:embedding}
 t^{\gamma}\longmapsto\omega^{-\gamma}
\end{equation}
embed $K_\Gamma$ as a subfield of $\SC=\No[i]$, using Gonshor's normal-form
arithmetic \cite{Gonshor}; the minus sign is what converts increasing Hahn
supports to decreasing surreal monomial supports. Every \emph{set} of
surcomplex constants lies in one such workspace: collect the exponents of
the real and imaginary Conway normal forms of those constants, change sign,
and take the additive (if needed $\Q$-linear) span. The proper-class field
$\No[i]$ is never formed as a single object; every conclusion about it is
read pointwise through a workspace enlarged to contain the exponents of the
point and of the input, and all constructions commute with such
enlargements (\cref{dyn:prop:workspace,dyn:thm:workspace}). No sum in this
report is indexed by a proper class.

\begin{definition}[Strong summability]\label{dyn:def:summable}
A \emph{set}-indexed family $(A_i)_{i\in I}$ of Hahn series is
\emph{strongly summable} if $\bigcup_i\supp A_i$ is well ordered and each
exponent lies in the support of only finitely many members; its sum is then
formed coefficientwise by finite sums. A map is \emph{strongly linear} if
it is linear and preserves such families and their sums.
\end{definition}

There is no limiting operation in \cref{dyn:def:summable}. The identity
$\sum_{j\ge0}c^{j}=(1-c)^{-1}$ for $c\in\m_\Gamma$ is a statement about a
summable family and multiplication, not about convergence of partial sums
in the fine order topology of $\No$; this is the discipline established by
the companion report \emph{Surcomplex analysis: holomorphic functions on
$K=\No[i]$} (\texttt{docs/surcomplex/analysis/}), which we inherit rather
than re-derive. No Banach contraction, Cauchy limit, compactness principle
or surcomplex path occurs in any proof below. Ordinary Taylor series of
coefficient functions, by contrast, use ordinary complex convergence, and
the two summations are never interchanged without a support argument.
Cancellation may not rescue an inadmissible family: infinitely many nonzero
terms at one exponent already violates \cref{dyn:def:summable} even if some
other summation method would assign a value to their complex coefficients.

\begin{warning}[The derivative in this report; the cross-category pair]
\label{dyn:warn:derivative}
The symbols $\partial_j$, $\partial_z$, $\partial_w$ differentiate the
\emph{ordinary} coordinate and act coefficientwise; every $t^{\gamma}$ is a
constant for them. They are not a derivation of the surreal coefficient
field, and in particular not the Berarducci--Mantova derivation
\cite{BM}. No compatibility with a global composition operator on $\No$ is
asserted. We adopt here, once, the wording of the cross-category warning of
the companion report \emph{Differential Algebra and Differential Equations
over the Surreal and Surcomplex Numbers}
(\texttt{docs/surcomplex/differential-equations/}): a Hahn constant for a
coordinate derivative need not be a constant for the scalar
Berarducci--Mantova derivation, so $\eps=\omega^{-1}$ is permitted as a
Hahn coefficient of a coordinate problem and prohibited as an infinitesimal
phase residue of the scalar derivation. All seven sources derive this
distinction independently; it is not restated again below.
\end{warning}

\subsection{The coefficient categories, in the collection's ring vocabulary}
\label{dyn:subsec:rings}

The companion report \emph{Infinitesimal Analytic Geometry over the
Surcomplex Numbers}
(\texttt{docs/\allowbreak surcomplex/\allowbreak analytic-\allowbreak
geometry/}) fixes a
four-ring table and a ring discipline: every numbered statement names the
ring over which it is asserted, and no theorem proved of one ring is
transferred to another. We adopt its names:
\begin{equation}\label{dyn:eq:rings}
 \OO(D)((t^{\Gamma}))\longrightarrow
 \mathcal R_n\subsetneq\mathcal A_n\subsetneq\mathcal F_n ,
\end{equation}
where
\begin{itemize}
\item $\OO(D)((t^{\Gamma}))$ is the \emph{fixed-polydisk} ring: one
specified, unchanged ordinary domain for every coefficient;
\item $\mathcal R_n=\varinjlim_{U\ni0}\HH_\Gamma(\OO(U))$ is the
\emph{common-domain germ} ring: one common domain, then allowed to shrink;
\item $\mathcal A_n=\C\{z_1,\dots,z_n\}((t^{\Gamma}))$ is the
\emph{radius-free} ring: individually convergent germs, with no common
radius required;
\item $\mathcal F_n=\C[[z_1,\dots,z_n]]((t^{\Gamma}))$ is the
\emph{formal-coefficient} ring.
\end{itemize}
We add exactly two rows to that table, and they are new \emph{rows}, not
new rings: the \emph{entire-coefficient} row $\OO(\C^{n})((t^{\Gamma}))$
(source 06's $\EE_n$, also the third row of \cref{dyn:thm:main}) and the
\emph{polynomial-coefficient} row $\C[z_1,\dots,z_n]((t^{\Gamma}))$
(source 06's $\PP_n$, also the fourth row of \cref{dyn:thm:main}). Both sit
between the fixed-polydisk ring and $\mathcal F_n$ in the obvious way, and
neither carries a uniform bound: in the entire row \emph{no} uniform norm
bound is imposed as $\gamma$ varies, and in the polynomial row \emph{no}
uniform degree bound.

We do not reprint the witness that separates the rings; it is the element
$F(z)=\sum_{r\ge1}t^{r\eta}/(1-rz)$ of the analytic-geometry report, whose
coefficient at $r\eta$ has a pole at $1/r$. Sources 02, 03 and 04 each
re-erect that separation from scratch, with $\eta=1$ in source 02
($\sum_{m\ge1}t^{m}/(1-mz)$) and $\eta=\delta$ in source 04
($\sum_{n\ge1}t^{n\delta}/(1-nw)$). Those are the same series.

\begin{convention}[Ring discipline, inherited]\label{dyn:conv:rings}
Every numbered statement below names its coefficient category in the
vocabulary of \eqref{dyn:eq:rings} and the two added rows, and the tag
ledger \cref{dyn:subsec:theoremtags} records it. No statement proved of one
category is transferred to another. In particular the fixed-polydisk row
of \cref{dyn:thm:main} concerns preservation of a \emph{specified
unchanged} polydisk, the radius-free row concerns separate germs, and
neither row, nor their combination, classifies $\mathcal R_n$ in the
intermediate regime; that is \cref{dyn:q:commongerm}.
\end{convention}

\subsection{The support certificate: finite words, not sequential limits}
\label{dyn:sec:support}

For well ordered $S\subset\Gamma_{>0}$ let $S^{*}=M(S)$ be its additive
monoid, including the empty sum $0$, and for $\gamma\in\Gamma$ let
\[
 \word_S(\gamma)=\{(s_1,\dots,s_k):k\ge0,\ s_i\in S,\
 s_1+\cdots+s_k=\gamma\}
\]
be its set of \emph{ordered} words of total weight $\gamma$, permutations
counted separately.

\begin{lemma}[Positive-support lemma of Neumann, in finite-word form]
\label{dyn:lem:neumann}
Let $S\subset\Gamma_{>0}$ be well ordered. Then $S^{*}$ is well ordered,
and $\word_S(\gamma)$ is finite for every $\gamma\in\Gamma$; in particular
\begin{equation}\label{dyn:eq:depth}
 \ell_S(\gamma)=\max\{k:(s_1,\dots,s_k)\in\word_S(\gamma)\}
\end{equation}
is finite for $\gamma\in S^{*}$, with $\ell_S(0)=0$. Moreover, if
$A\subset\Gamma$ is well ordered then $A+S^{*}$ is well ordered, and for
fixed $\gamma$ there are only finitely many pairs consisting of $a\in A$
and a word in $S$ with $a+s_1+\cdots+s_k=\gamma$.
\end{lemma}
\begin{proof}
Order finite words over $S$ by subsequence embedding with coordinatewise
domination. Higman's lemma says this is a well-quasi-order: every infinite
sequence of words contains indices $i<j$ with the earlier word embedded in
the later \cite{Higman}; the special case needed is reproved from scratch
in \cref{dyn:app:higman}. Because all letters are positive, an embedding
weakly increases the sum of the letters, and increases it strictly unless
the two words are identical: an unused positive letter contributes
positively, and a strictly increased matched letter increases the sum.
An infinite strictly decreasing sequence in $S^{*}$, after choosing one
word per element, would contradict the weak increase; infinitely many
distinct words of one fixed weight would contradict the strict increase.
This proves the first two assertions.

For the last, assume $A\ne\emptyset$, subtract $a_0=\min A$, and apply the
first part to the well-ordered positive set
$T=S\cup\bigl((A-a_0)\setminus\{0\}\bigr)$ at weight $\gamma-a_0$: a pair
in the assertion is encoded by a word with one distinguished initial letter
from $A-a_0$, or none if $a=a_0$. Well-ordering follows from
$A+S^{*}\subset a_0+T^{*}$.
\end{proof}

This is classical Hahn--Neumann mathematics \cite{Neumann,Poonen} and is not
claimed as a contribution. All seven sources state it; sources 02 and 05
each reproduce a minimal-bad-sequence proof of the underlying word order,
which is the same argument as \cref{dyn:app:higman} and is therefore not
printed again. Source 07 uses the \emph{full} finiteness assertion, across
all word lengths, together with a \emph{labelled} refinement of it
(\cref{dyn:subsec:qpmarked}); both are already contained in
\cref{dyn:lem:neumann}.

For $\gamma\in S^{*}$ let $D_S(\gamma)$ be the finite set of letters
occurring in the words of $\word_S(\gamma)$.

For a well-ordered input support $B\subset\Gamma$, retain the starting
exponent as well as the positive word. Define
\begin{align*}
 \word_{B,S}(\gamma)
 &=\{(b,w):b\in B,\ w\in\word_S(\gamma-b)\},\\
 \ell_{B,S}(\gamma)
 &=\max\bigl(\{0\}\cup\{|w|:(b,w)\in\word_{B,S}(\gamma)\}\bigr).
\end{align*}
Here $|w|$ is the number of letters of $w$. Let $E_{B,S}(\gamma)$ be the
set of starting exponents $b$ in these pairs, and $D_{B,S}(\gamma)$ the set
of positive letters occurring in their words. All three sets are finite by
\cref{dyn:lem:neumann}; the depth is zero if there is no pair. When
$B=\{0\}$ and $\gamma\in S^{*}$, this depth is $\ell_S(\gamma)$.

\begin{corollary}[Finite ancestry relative to the input support]
\label{dyn:cor:ancestry}
Let $S\subset\Gamma_{>0}$ be well ordered, and suppose a $\C$-linear
operator $T$ on Hahn coefficient families is given coefficientwise by
\[
 (TA)_\gamma
 =\sum_{b\in\supp A}\;
   \sum_{\substack{w\in\word_S(\gamma-b)\\ |w|\ge1}}
       T_{b,w}(A_b).
\]
Each $T_{b,w}$ is a fixed linear coefficient operator, expressible by
finitely many operations in the chosen coefficient category. For any Hahn
input $A$ with $B=\supp A$ and any ordinary scalars $c_k$, the family
$(c_kT^kA)_{k\ge0}$ is strongly summable and its sum has support in
$B+S^{*}$. At an output exponent $\gamma$, contributions use at most
$\ell_{B,S}(\gamma)$ applications of $T$, only input coefficients $A_b$
with $b\in E_{B,S}(\gamma)$, and only positive support letters in
$D_{B,S}(\gamma)$. Only finitely many coefficient operations contribute.

These conclusions hold for finite-dimensional vector and matrix
coefficient families and for operators with Hahn coefficients whenever
the displayed positive-word hypothesis holds. If $(A_i)_{i\in I}$ is
strongly summable, the family $(c_kT^kA_i)_{(i,k)\in I\times\N}$ is
strongly summable too, using $B=\bigcup_i\supp A_i$. Formal power-series
identities in one operator may therefore be evaluated coefficientwise;
this applies to the geometric, logarithmic and exponential identities used
below.
\end{corollary}
\begin{proof}
The displayed definition of $T$ is meaningful: at a fixed output exponent,
\cref{dyn:lem:neumann} gives only finitely many input-exponent/word pairs.
Its support is contained in $B+(S^{*}\setminus\{0\})$, a subset of the
well-ordered set $B+S^{*}$, so iteration is legitimate.

Expand a contribution to $(T^kA)_\gamma$ by its successive nonempty
words $w_1,\ldots,w_k$, starting at $b\in B$. Their concatenation $w$
satisfies $b+\sum w=\gamma$, hence $(b,w)\in\word_{B,S}(\gamma)$.
Conversely, once $(b,w)$ and its decomposition into $k$ consecutive
nonempty blocks are specified, all intermediate exponents are determined
by the partial sums. A word of length $n\ge1$ has only $2^{n-1}$ such
block decompositions across all possible $k$, with $k\le n$; the empty
word occurs only for $k=0$. Since $\word_{B,S}(\gamma)$ is finite and
every block uses finitely many coefficient operations, only finitely many
expanded contributions occur at $\gamma$, across all iterates. This also
proves the stated depth and input-data bounds. Their supports all lie in
$B+S^{*}$, proving both requirements of strong summability.

For a strongly summable input family, its union of supports $B$ is well
ordered. There are finitely many relevant pairs $(b,w)$, and for each such
$b$ only finitely many $i$ have $(A_i)_b\ne0$. The same argument thus
proves the joint assertion and permits commuting $T$ with the input sum.
Finite-dimensional vector and matrix operations introduce only finitely
many component choices and do not change the argument.

Finally, fix an input and an output exponent. All contributing powers of
$T$ have bounded degree; in a product or a composition of scalar formal
power series, each such degree has only finitely many decompositions.
Consequently the formal identity reduces to a finite coefficient
identity. Compositions are understood in their usual formal domain (in
particular the inner series has zero constant term), as for
$\log(1+T)$ and $\exp(T)-I$.
\end{proof}

\noindent\textbf{Why the input support cannot be omitted.}
Take $\Gamma=\Z$, $S=\{1\}$ and $T(A)=tA$. For every integer $n\ge0$,
\[
 T^n(t^{-n})=1,\qquad
 \ell_S(0)=0,\qquad \ell_{\{-n\},S}(0)=n.
\]
Thus no bound depending only on $S$ and the output exponent can control
all Hahn inputs, and input exponents need not be letters of $S$.
The earlier version of \cref{dyn:cor:ancestry} used $\ell_S(\gamma)$ and
$D_S(\gamma)$ for arbitrary input; the relative sets above correct that
statement. The sharper unshifted bound remains valid in
\cref{dyn:prop:linearization}: its initial input $\Li^{-1}f$ is already
supported in $S$, so its starting exponent itself supplies a positive
letter. The original source manuscripts remain available in repository history.

\begin{warning}[Why a valuation-contraction argument would be incomplete]
\label{dyn:warn:contraction}
Writing $v(T^{k}A)\ge v(A)+k\min S$ does \emph{not} imply that these lower
bounds become cofinal in $\Gamma$: for example $k<\Omega$ for every integer
$k$. The finite-word certificate of
\cref{dyn:lem:neumann,dyn:cor:ancestry} establishes summability even when
that sequence of bounds never passes a specified higher rank, and every
infinite operator series in this report rests on the stronger certificate.
Likewise ``only finitely many terms below each cutoff'' is \emph{not} the
criterion: in $\Gamma=\Z^{2}$ lexicographic with $\eta=(0,1)$ and
$\theta=(1,0)$ one has $n\eta<\theta$ for every integer $n$, yet
$\sum_{n\ge0}t^{n\eta}$ is a valid Hahn series with
$(1-t^{\eta})\sum_{n\ge0}t^{n\eta}=1$, while the partial sums do not
converge in the valuation topology. Sources 03, 05 and 06 each record this
independently.
\end{warning}

\begin{example}[Infinitely many earlier weights, one relevant word]
\label{dyn:ex:ranktwo}
Let $\Gamma=\Q+\Q\Omega$ with $\Omega$ greater than every rational, and
$S=\{1\}\cup\{2-1/n:n\ge2\}\cup\{\Omega\}$. This $S$ is well ordered and
lies in no finitely generated additive grid, and infinitely many of its
members are below $\Omega$. Nevertheless
$\word_S(\Omega)=\{(\Omega)\}$, $\ell_S(\Omega)=1$ and
$\ell_S(\Omega+2)=3$, the words at $\Omega+2$ being the three permutations
of $(\Omega,1,1)$: any finite sum of rational letters is rational, and any
sum involving two $\Omega$'s exceeds $\Omega+2$. An algorithm computing a
coefficient at $\Omega$ needs the input at $\Omega$, not infinitely many
rational predecessors. \Cref{dyn:ex:Omega} computes with this $S$.
\end{example}

\subsection{A uniform-monoid support algebra}
\label{dyn:subsec:BM}

The finite-return arguments of \cref{dyn:sec:return} need a slightly
stronger object than \cref{dyn:cor:ancestry}: one monoid uniform across all
Taylor degrees.

\begin{definition}\label{dyn:def:BM}
For $\mathsf M=\langle S\rangle$ with $S\subset\Gamma_{>0}$ well ordered,
put
\[
 \BB_{\mathsf M}=\Bigl\{\textstyle\sum_{n\ge0}h_nX^{n}\in K_\Gamma[[X]]:
 \supp(h_n)\subseteq\mathsf M\ \text{for \emph{every}}\ n\Bigr\},
\]
a subring of $\mathcal F_1=\C[[X]]((t^{\Gamma}))$ after regrouping as
$\sum_{\gamma\in\mathsf M}H_\gamma(X)t^{\gamma}$ with
$H_\gamma\in\C[[X]]$. Coefficients need not converge as ordinary complex
power series.
\end{definition}

The uniformity of the single monoid across all Taylor degrees is
load-bearing: an arbitrary member of $\OO_\Gamma[[X]]$ need not lie in any
such algebra.

\begin{lemma}[Evaluation, inversion, isometry]\label{dyn:lem:evaluation}
Every $H\in\BB_{\mathsf M}$ is strongly Hahn-evaluable at every
$u\in\m_\Gamma$. If $H=X+O(X^{2})$ then its formal inverse also lies in
$\BB_{\mathsf M}$, and the evaluated maps are mutually inverse isometric
bijections of $\m_\Gamma$, isometry meaning $v(H(u)-H(w))=v(u-w)$. Formal
composition identities between zero-fixing series with integral unit linear
coefficients and a common positive-support certificate remain valid after
evaluation.
\end{lemma}
\begin{proof}
For $u\ne0$ with $U=\supp u\subset\Gamma_{>0}$, the supports of $h_nu^{n}$
lie in $\langle S\cup U\rangle$, and every monomial of $u^{n}$ uses at
least $n$ generators from $U$, so \cref{dyn:lem:neumann} leaves finitely
many $n$ at a fixed exponent: both clauses of \cref{dyn:def:summable} hold.
The coefficients of the formal inverse of $X+\sum_{n\ge2}h_nX^{n}$ are
integer polynomials in finitely many $h_n$, hence supported in $\mathsf M$,
and the substitutions in the inverse identities are summable by the same
count. For $u\ne w$ in $\m_\Gamma$,
\[
 \frac{H(u)-H(w)}{u-w}=1+\sum_{n\ge2}h_n\sum_{j=0}^{n-1}u^{j}w^{n-1-j},
\]
whose nonconstant summands all have positive valuation and whose total has
residue one and valuation zero; this is the isometry. For compositions,
expand both levels into products: at a fixed exponent the number of
positive generators is finite, and at a fixed formal degree the number of
zero-exponent contributions is finite because the substituted series has
zero constant term. A unit linear coefficient and its inverse are adjoined
by separating their nonzero complex constant terms and expanding the
inverse in its positive part.
\end{proof}

\begin{example}[Coefficient-valuation bounds are insufficient at rank $>1$]
\label{dyn:ex:bad-support}
Let $\Gamma=\Q+\Q\eps_0$ with $0<\eps_0<1/n$ for every positive integer
$n$. All coefficients of $A(X)=\sum_{n\ge2}t^{1/n}X^{n}$ have positive
valuation. Yet at $u=t^{\eps_0}$ the proposed Taylor sum has exponents
$1/n+n\eps_0$, which \emph{strictly decrease} as $n$ grows: it is not a
Hahn series. Hence an assertion that ``all normalized coefficients are
integral'' would leave a real gap in an arbitrary-rank linearization proof.
The set $\{1/n:n\ge2\}$ is not well ordered, so this example does not
satisfy the uniform-support hypothesis of \cref{dyn:lem:evaluation}; it is
the reason that hypothesis is stated. This is source 01's observation, and
no other source records it.
\end{example}

\subsection{Halos, monads, and evaluation as a ring homomorphism}
\label{dyn:subsec:halo}

For an ordinary open $U\subset\C^{d}$ put
\[
 U^{\#}_\Gamma=\{a+\xi:a\in U,\ \xi\in\m_\Gamma^{d}\},
\]
the \emph{finite halo}; the representation is unique and $a=\st(a+\xi)$ is
the standard part. For $A=\sum_\gamma t^{\gamma}A_\gamma$ with all
$A_\gamma$ holomorphic on $U$, define
\begin{equation}\label{dyn:eq:halo-evaluation}
 A^{\#}(a+\xi)=\sum_{\gamma,\alpha}t^{\gamma}
 \frac{\partial^{\alpha}A_\gamma(a)}{\alpha!}\,\xi^{\alpha}.
\end{equation}

\begin{proposition}[Faithful analytic realization]\label{dyn:prop:eval}
The family in \eqref{dyn:eq:halo-evaluation} is strongly summable, and
\eqref{dyn:eq:halo-evaluation} defines a \emph{faithful} $K_\Gamma$-algebra
homomorphism from $\HH_\Gamma(\OO(U))$ to functions on $U^{\#}_\Gamma$. It
respects coordinate derivatives, substitution by positive near-identity
Hahn maps, and infinitesimal Taylor expansions, with the expected
coefficients $(\partial^{\alpha}A)^{\#}(z)/\alpha!$. For germ or formal
coefficients use the Taylor coefficients at $0$ and restrict the domain to
$\m_\Gamma^{d}$; polynomial and entire coefficients evaluate on every
finite halo, but not automatically at infinite surcomplex points.
\end{proposition}
\begin{proof}
Let $T=\supp A$ and let $S$ be the union of the positive supports of the
components of $\xi$, a finite union of well-ordered positive sets. Every
monomial contribution has exponent in $T+S^{*}$, and at a fixed exponent
\cref{dyn:lem:neumann} permits only finitely many choices of $\gamma$ and
of a word from $S$; the word length bounds $|\alpha|$ and a finite word has
finitely many coordinate assignments. Hence the expanded family is
strongly summable. Ordinary Taylor identities for products and derivatives
then apply, every regrouping at a Hahn exponent being finite; adding the
positive support of a near-identity substitution to $S$ gives the
composition assertion from the ordinary Taylor chain rule. For
faithfulness, let $\gamma_0=\min\supp A$ for $A\ne0$ and pick an ordinary
$a\in U$ with $A_{\gamma_0}(a)\ne0$: the coefficient of $t^{\gamma_0}$ in
$A^{\#}(a)$ is then nonzero.
\end{proof}

At a fixed ordinary Taylor center the numbers $n!$, or extremely large
inverse small divisors, remain \emph{ordinary complex constants} of
valuation zero. Strong summability imposes no real growth condition on
them. This single fact explains both the strength of the theorems below and
the necessity of stating their domains exactly.

\begin{lemma}[Faithfulness of infinitesimal evaluation]\label{dyn:lem:faithful}
If $0\ne A\in\mathcal F_1$ and $\Gamma\ne\{0\}$ is divisible, then
$A(u)\ne0$ for some $u\in\m_\Gamma$; evaluation on the whole monad is
injective on the formal-coefficient ring.
\end{lemma}
\begin{proof}
Let $\gamma_0=\min\supp A$ and let $n$ be the least degree of the nonzero
formal series $A_{\gamma_0}$. If another Hahn exponent occurs, let
$\Delta>0$ be the gap from $\gamma_0$ to the next one and choose $s>0$ with
$ns<\Delta$; divisibility supplies such an $s$ when $n>0$, and any positive
$s$ works when $n=0$. If no other exponent occurs, take any $s>0$. At
$u=t^{s}$ the exponent $\gamma_0+ns$ appears with nonzero coefficient,
while higher degrees at $\gamma_0$ and all other exponents contribute only
larger exponents. Evaluation is legitimate by \cref{dyn:lem:evaluation}
after factoring out $t^{\gamma_0}$.
\end{proof}

\Cref{dyn:lem:faithful} blocks the evasion of representing a divergent
conjugacy by a \emph{different} common-domain family agreeing on all
infinitesimals; it is used in \cref{dyn:prop:noncoherent}.

\section{Positive substitution, logarithms, flows, and centralizers}
\label{dyn:sec:substitution}

\subsection{The substitution group}
For $U\subset\C^{d}$ open set
\[
 \GG_+(U)=\{F(z)=z+u(z):u\in\HH_\Gamma(\OO(U))_+^{d}\}.
\]
No condition at $z=0$ is imposed, and $0$ need not lie in $U$. The pullback
of $F=\id+u$ on $A\in\HH_\Gamma(\OO(U))$ is
\begin{equation}\label{dyn:eq:pullback}
 C_FA=A\circ F=\sum_{\alpha\in\N^{d}}\frac1{\alpha!}(\partial^{\alpha}A)u^{\alpha},
\end{equation}
every coefficient of which is a finite sum of products and derivatives of
functions holomorphic on $U$, hence holomorphic on $U$ regardless of its
ordinary norm; the support lies in $\supp A+S^{*}$ for $S$ the union of the
supports of the components of $u$.

\begin{proposition}[Substitution group]\label{dyn:prop:subgroup}
\eqref{dyn:eq:pullback} defines an associative group law on $\GG_+(U)$.
Every inverse again lies in $\GG_+(U)$, on the same $U$, with support in
$S^{*}\setminus\{0\}$. Evaluation gives mutually inverse bijections of
$U^{\#}_\Gamma$ preserving standard parts and commutes with composition.
The same statements hold in the common-domain germ, radius-free germ,
entire, polynomial and formal categories, with the evaluation domains of
\cref{dyn:subsec:halo}. The order convention is
\begin{equation}\label{dyn:eq:pullback-order}
 C_FC_G=C_{G\circ F}:
\end{equation}
pullback reverses composition.
\end{proposition}
\begin{proof}
Strong summability is \cref{dyn:lem:neumann,dyn:cor:ancestry}, applied
with the support of each Hahn input retained. At a fixed Hahn exponent
only finitely many input coefficients, derivatives and products contribute.
Keep those input coefficients separately, introduce a formal parameter for
each relevant positive support letter, and discard parameter monomials
beyond the corresponding relative word depth. Taylor composition in that
finite nilpotent parameter algebra is associative and multiplicative, and
each asserted coefficient identity is
one of those identities. For the inverse put $\Delta=C_F-I$: every term of
$\Delta$ adds a nonempty word from $S$, so $C_F^{-1}=\sum_{k\ge0}(-\Delta)^{k}$
is defined coefficientwise and is a two-sided multiplicative inverse. Let
$G=C_F^{-1}z$, identity plus positive support with support in $S^{*}$. On
coordinate polynomials multiplicativity gives $C_F^{-1}=C_G$; at each Hahn
weight both operators are finite-order differential operators on an
ordinary input, so agreement on polynomial jets extends to all holomorphic
inputs, and strong linearity extends it to Hahn inputs. For formal
coefficients the finite-jet argument is coefficientwise in $z$.
Substituting $z=a+\xi$ and adding $\supp\xi$ to the control set preserves
every identity, and $\st F(a+\xi)=a$ keeps the evaluated map in
$U^{\#}_\Gamma$. Iterates of a single map are unaffected by
\eqref{dyn:eq:pullback-order}, and commuting maps have commuting pullbacks;
the convention is retained explicitly to avoid a sign or order ambiguity in
logarithms.
\end{proof}

For maps $F=\Lambda z+f$ on $\mathbb D_R^{d}$ with $\Lambda$ diagonal
unitary and $f$ positive, use instead
\begin{equation}\label{dyn:eq:rotation-substitution}
 A\circ F=\sum_{\alpha\in\N^{d}}\frac1{\alpha!}(\partial^{\alpha}A)(\Lambda z)f^{\alpha},
\end{equation}
legitimate because $\Lambda\mathbb D_R^{d}=\mathbb D_R^{d}$. For formal
coefficients or germs, $F(0)=0$ is imposed whenever composition with the
nonidentity reduction is used. \Cref{dyn:lem:inverse} records the
near-identity inverse in the form source 02 uses.

\begin{lemma}[Near-identity inverses]\label{dyn:lem:inverse}
Let $H=z+U$ be normalized with coefficients in any one of
$\OO(\mathbb D_R)$, $\C\{z\}$, $\C[[z]]$, $\OO(\C)$, $\C[z]$, and positive
Hahn support $B$. There is a unique normalized inverse $J=z+V$ with
coefficient support in $B^{*}\setminus\{0\}$ and coefficients in the same
algebra, with $H\circ J=J\circ H=\id$.
\end{lemma}
\begin{proof}
In $H(J)=J+U(J)$ the coefficient $v_\gamma$ occurs with coefficient one,
and every other contribution at exponent $\gamma$ uses only $v_\beta$ with
$\beta<\gamma$; set $v_\gamma$ to minus that finite sum and recurse on
$B^{*}$. Derivatives, products and finite sums preserve each algebra and
the order-$z^{2}$ normalization. The identical triangular argument gives a
left inverse, and associativity, which is coefficientwise a finite Taylor
identity, identifies the two.
\end{proof}

\subsection{Logarithms are actual coefficientwise vector fields}
\label{dyn:sec:flows}

A \emph{positive Hahn vector field} on $U$ is
$X=a_1\partial_1+\cdots+a_d\partial_d$ with
$a\in\HH_\Gamma(\OO(U))_+^{d}$, and
\begin{equation}\label{dyn:eq:explog}
 \exp X=\sum_{k\ge0}\frac{X^{k}}{k!},\qquad
 \log C_F=\sum_{k\ge1}\frac{(-1)^{k+1}}{k}(C_F-I)^{k}
\end{equation}
are operators on the Hahn algebra, not numerical logarithms of values of
$F$.

\begin{theorem}[Same-domain exponential--logarithm correspondence]
\label{dyn:thm:exp-log}
For every open $U\subset\C^{d}$ the assignments
$F\mapsto X_F=\log C_F$ and $X\mapsto\Phi_X=(\exp X)z$ are inverse
bijections between $\GG_+(U)$ and positive Hahn vector fields on $U$. All
coefficient functions stay holomorphic on the \emph{original} $U$: no
shrinking, and no uniform bound on the coefficient family. If $S$ contains
the positive support of the input then the output support lies in
$S^{*}\setminus\{0\}$, and each output coefficient has a finite word
certificate. For $F\ne\id$,
\begin{equation}\label{dyn:eq:valXF}
 v(X_F)=v(F-\id),
\end{equation}
and the leading coefficient of $X_Fz$ equals that of $F-\id$. The inverse
of $F$ is $(\exp(-X_F))z$. The same assertions hold with common-domain
germ, radius-free germ, entire, polynomial or formal coefficients.
\end{theorem}
\begin{proof}
\emph{Summability.} Put $\Delta=C_F-I$ and $S=\supp(F-\id)$. By
\eqref{dyn:eq:pullback}, $\Delta$ adds nonempty words from $S$, so both
series in \eqref{dyn:eq:explog} are strongly summable by
\cref{dyn:cor:ancestry} for each fixed Hahn input $A$, with depth
$\ell_{\supp A,S}(\gamma)$ at output $\gamma$. Every output coefficient is
a finite differential expression with coefficients holomorphic on $U$.
For the coordinate input $z$, the input support is $\{0\}$, so the
unshifted depth $\ell_S(\gamma)$ applies for $\gamma\in S^{*}$.

\emph{The derivation identity.} Introduce an ordinary indeterminate $s$ and
put $C_F^{s}=\sum_{k\ge0}\binom{s}{k}\Delta^{k}$. For fixed inputs and a
fixed Hahn exponent this is a \emph{polynomial} in $s$, not a power series.
At every nonnegative integer $s$ it is the usual integer power of the
multiplicative operator $C_F$, so the coefficientwise polynomial identity
$C_F^{s}(AB)=C_F^{s}(A)C_F^{s}(B)$ holds identically in $s$.
Differentiating at $s=0$ and using
$\frac{\mathrm d}{\mathrm ds}\binom{s}{k}\big|_{s=0}=(-1)^{k+1}/k$ for
$k\ge1$ gives $X_F(AB)=X_F(A)B+AX_F(B)$.

\emph{Why the derivation is a vector field.} At each Hahn weight $\gamma$
the restriction of $X_F$ to ordinary coefficient functions is a
finite-order differential operator $D_\gamma$, and the preceding identity
makes it a derivation. Put $a_{j,\gamma}=D_\gamma z_j$; by Leibniz,
$D_\gamma$ agrees with $\sum_ja_{j,\gamma}\partial_j$ on every coordinate
polynomial, and at an ordinary point a finite-order differential operator
is determined by a finite Taylor jet, so the two agree on every holomorphic
function. For germs or formal series use finite jets at the origin; for
polynomials it is immediate. This finite-order step is essential: no
classification of arbitrary abstract algebraic derivations of a function
ring is invoked.

\emph{Conversely.} If $X=a\cdot\partial$ is positive, each application adds
a positive letter, so $\exp X$ is strongly summable, and repeated Leibniz
$X^{k}(AB)=\sum_j\binom kjX^{j}(A)X^{k-j}(B)$ makes it multiplicative with
inverse $\exp(-X)$. Put $\Phi_X=(\exp X)z=\id+\text{positive}$. On
coordinate polynomials multiplicativity gives $\exp X=C_{\Phi_X}$, and the
finite-jet argument extends the equality to the whole coefficient algebra
and then to Hahn inputs.

\emph{Inverse identities and valuation.} $\exp(\log(I+T))=I+T$ and
$\log(\exp T)=T$ hold in a nilpotent operator algebra, and at a specified
Hahn coefficient only finitely many words occur, so they hold
coefficientwise; this is the bijection. Finally $X_Fz=(F-z)+$ terms with at
least two support letters, whose valuation is at least $2v(F-\id)$, which
gives \eqref{dyn:eq:valXF} and the leading-coefficient assertion.
Derivatives and finite products preserve every stated category.
\end{proof}

\begin{remark}[Three sources, one route, three by-products]
\label{dyn:rem:threelogs}
Sources 03, 04 and 05 each prove \cref{dyn:thm:exp-log}, and all three use
the same essential device: the binomial operator family
$\sum_k\binom sk\Delta^{k}$, polynomiality in $s$ at each fixed Hahn
exponent, and a finite-jet identification of the resulting derivation with
a coordinate vector field. Being one route, it is printed once, in source
03's several-variable form on an arbitrary open $U\subset\C^{d}$. The
routes differ only in by-products, and each by-product survives separately
below: source 05 adds that $v(W_F)=v(h)$ \emph{with equal leading
coefficient} and that inverses are $\exp(-W_F\partial_z)z$ (both included
in \cref{dyn:thm:exp-log}); source 04 adds the exact fixed-point identity
\cref{dyn:lem:fixedderivative}, which is what determines the complete
monodromy exponent of \cref{dyn:sec:parabolic-monodromy}; and source 05
adds the Julia equation \cref{dyn:cor:julia} and the intertwining criterion
\cref{dyn:prop:intertwine}, which are what make
\cref{dyn:sec:periods} possible. Source 05 also notes explicitly that the
general algebraic mechanism for exponentiating strongly summable
derivations of generalized-series algebras is established work
\cite{BKKPS} and that the existence of a formal logarithm is not claimed
new.
\end{remark}

\begin{lemma}[Derivative of the logarithmic generator at an ordinary fixed
point]\label{dyn:lem:fixedderivative}
Let $F=w+\Delta$ with $\Delta\in\HH_\Gamma(\OO(U))_{>0}$ in one variable,
$X_F=A(w)\partial_w$. If an ordinary $a\in U$ satisfies $F(a)=a$, then
\begin{equation}\label{dyn:eq:fixedderivative}
 A(a)=0,\qquad A'(a)=\Log F'(a),
\end{equation}
where $\Log(1+u)=\sum_{n\ge1}(-1)^{n+1}u^{n}/n$ for $v(u)>0$.
\end{lemma}
\begin{proof}
Evaluation at $a$ annihilates $N\Phi=(C_F-I)\Phi$, hence annihilates
$A=Dw$. With $\lambda_a=F'(a)$, the chain rule for Taylor substitution
gives $(N\Phi)'(a)=(\lambda_a-1)\Phi'(a)$, so
$(N^{r}w)'(a)=(\lambda_a-1)^{r}$; apply the logarithm series, which is a
strong sum because $v(\lambda_a-1)>0$ (the identity is immediate if
$\lambda_a=1$).
\end{proof}

This \emph{exact} identity, and not a truncated modified differential
equation, is what fixes the complete monodromy exponent in
\cref{dyn:sec:parabolic-monodromy}.

\begin{corollary}[Julia equation]\label{dyn:cor:julia}
For every $F\in\GG_+(U)$ in one variable, with $W_F=X_Fz$,
$W_F(F(z))=F'(z)W_F(z)$.
\end{corollary}
\begin{proof}
$C_F$ commutes with its logarithm; apply
$C_F(W_F\partial_z)=(W_F\partial_z)C_F$ to $z$.
\end{proof}

\begin{proposition}[Conjugacy and vector fields]\label{dyn:prop:intertwine}
For $F_1,F_2,H\in\GG_+(U)$ in one variable,
$H\circ F_1=F_2\circ H$ if and only if $W_{F_2}(H)=H'W_{F_1}$.
\end{proposition}
\begin{proof}
The conjugacy relation is $C_{F_1}C_H=C_HC_{F_2}$; conjugating the summable
logarithm gives $(W_{F_1}\partial_z)C_H=C_H(W_{F_2}\partial_z)$, and
evaluation on $z$ gives the displayed identity. Conversely that identity
and the chain rule give the operator intertwining on the whole algebra, and
exponentiation recovers the conjugacy.
\end{proof}

\subsection{Times, flows, and unique roots}

\begin{corollary}[Exact flow law and admissible Hahn times]\label{dyn:cor:flow}
Let $X=a\cdot\partial$ be positive. For every ordinary $s\in\C$,
$\Phi_s=(\exp(sX))z\in\GG_+(U)$ and
\begin{equation}\label{dyn:eq:flow-law}
 \Phi_0=\id,\qquad\Phi_{s+r}=\Phi_s\circ\Phi_r,\qquad
 \frac{\mathrm d}{\mathrm ds}\Phi_s=a\circ\Phi_s,
\end{equation}
each Hahn coefficient of $\Phi_s$ being a \emph{polynomial} in $s$. For
$c\in K_\Gamma$ with $c=0$ or
\begin{equation}\label{dyn:eq:admissible-time}
 v(c)+v(a)>0,
\end{equation}
the map $\Phi_c=(\exp(cX))z$ is also defined and positive, and the additive
flow law holds for such times.
\end{corollary}
\begin{proof}
All ordinary scalars have valuation zero, each coefficient involves
finitely many powers of $s$, and commuting exponentials of multiples of
$X$ satisfy the addition law. Coefficientwise differentiation and
commuting $X$ with its exponential give the last identity. For a Hahn time,
\eqref{dyn:eq:admissible-time} says exactly that $cX$ has positive support,
since $c$ is a scalar independent of $z$; apply \cref{dyn:thm:exp-log} to
$cX$, and note that the positive supports of two admissible multiples have
well-ordered union.
\end{proof}

Condition \eqref{dyn:eq:admissible-time} characterizes when this scalar
multiple of a nonzero vector field is positive; it is \emph{not} asserted
to be a maximal possible lifetime in a larger class of maps (a constant
vector field, for instance, has translations beyond the positive subgroup).
Likewise \eqref{dyn:eq:flow-law} is coefficientwise analytic time
dependence, not a continuous path in the fine surreal topology.

\begin{corollary}[Unique fractional iteration; torsion-freeness]
\label{dyn:cor:roots}
For every $F\in\GG_+(U)$ and integer $N\ge1$ there is a unique
$G\in\GG_+(U)$ with $G^{\circ N}=F$, namely $G=(\exp(X_F/N))z$. The group
$\GG_+(U)$ is uniquely divisible and torsion-free.
\end{corollary}
\begin{proof}
Existence is \cref{dyn:cor:flow}. If $G^{\circ N}=F$ then $C_G^{N}=C_F$,
and taking logarithms gives $NX_G=X_F$, whence uniqueness. If
$F^{\circ N}=\id$ then $NX_F=0$, and characteristic zero gives $F=\id$.
\end{proof}

\subsection{Fixed-point ideals, and two routes to iteration invariance}
\label{dyn:sec:fixed}

For a coefficient map $F$ let $I_{\Fix(F)}=(F_1-z_1,\dots,F_d-z_d)$ be its
fixed-point ideal in the chosen Hahn coefficient algebra. No Noetherianity
and no Nullstellensatz is needed to compare such ideals, and none is
invoked: ideals are compared through an explicit invertible matrix.

\begin{theorem}[Fixed-point ideal of a positive flow]\label{dyn:thm:fixed-ideal}
Let $X=a\cdot\partial$ be a positive Hahn vector field on $U$ and let
$c\ne0$ be an admissible Hahn time. Then
\begin{equation}\label{dyn:eq:fixed-ideal}
 I_{\Fix(\Phi_c)}=(a_1,\dots,a_d)\qquad\text{in }\HH_\Gamma(\OO(U)),
\end{equation}
and likewise in the other coefficient categories and after localization. In
particular all nonzero integer iterates of $F\in\GG_+(U)$ have identical
fixed-point ideals; on the evaluated halo there are no nonfixed points of
finite period; and whenever a fixed-point quotient or local algebra has
finite length, that length is unchanged by nonzero integer iteration.
\end{theorem}
\begin{proof}
Set $E=cX=b\cdot\partial$ with $b=ca$ positive, so that
$\Phi_c-z=b+\sum_{k\ge2}E^{k}z/k!$. Since
$E^{k}z_j=\sum_\ell b_\ell\partial_\ell(E^{k-1}z_j)$, define
\begin{equation}\label{dyn:eq:fixedmatrix}
 B_{j\ell}=\delta_{j\ell}+\sum_{k\ge2}\frac{\partial_\ell(E^{k-1}z_j)}{k!},
\end{equation}
strongly summable by \cref{dyn:cor:ancestry}, of the form
$I+\text{positive}$, hence invertible by its matrix geometric series, and
\begin{equation}\label{dyn:eq:fixedfactor}
 \Phi_c-z=Bb.
\end{equation}
The two coordinate ideals are therefore equal, and since the scalar $c$ is
a unit of $K_\Gamma$, the ideal generated by $b$ is that generated by $a$.
On evaluation at any halo point $B$ remains invertible by
\cref{dyn:prop:subgroup}, so a point is fixed by $\Phi_c$ exactly when all
components of $a$ vanish there; apply this with $c=1$ and any nonzero
integer $c=n$. Identical ideals give identical quotient algebras.
\end{proof}

This theorem concerns maps whose ordinary reduction is the identity. It
makes no claim that a general surcomplex map lacks periodic points; in
particular rotations of finite order lie outside $\GG_+(U)$ unless they are
the identity, and \cref{dyn:thm:shell} exhibits genuine period-$q$ points
for a map that is \emph{not} in $\GG_+(U)$.

\begin{remark}[Two routes]\label{dyn:rem:tworoutes-fixed}
Source 03 proves iteration invariance by the explicit matrix
\eqref{dyn:eq:fixedmatrix}--\eqref{dyn:eq:fixedfactor}, obtaining ideal
equality --- scheme structure and finite length --- rather than inferring
it from equality of zero sets. Source 05 obtains the same invariance as a
corollary of unique $n$-th roots and the centralizer
(\cref{dyn:cor:roots,dyn:thm:centralizer}), which gives less (no ideal
statement) by a shorter argument but also gives torsion-freeness of the
whole group. Both are kept: the ideal statement is not available from the
root argument, and the group statement is not available from the matrix.
\end{remark}

\begin{theorem}[The full positive centralizer in one variable]
\label{dyn:thm:centralizer}
Let $U\subset\C$ be connected, let $F\in\GG_+(U)$ be nonidentity, and write
$X_F=a(z)\partial_z$ with $\alpha=v(a)>0$. Then
\begin{equation}\label{dyn:eq:centralizer}
 \Cent_{\GG_+(U)}(F)=\bigl\{(\exp(ca\partial_z))z:
 c\in K_\Gamma,\ c=0\ \text{or}\ v(c)+\alpha>0\bigr\},
\end{equation}
and the parametrization is an isomorphism from that additive group of Hahn
times onto the centralizer. Zeros of $a$ need \emph{not} be removed from
$U$. No nowhere-zero hypothesis on the leading coefficient of $a$ is
required.
\end{theorem}
\begin{proof}
Let $G$ be positive with $X_G=b\partial_z$. Commuting maps have commuting
pullbacks, hence commuting logarithms, all double regroupings being valid
over the union of the two positive supports; conversely commuting
logarithms have commuting exponentials. Thus $FG=GF$ iff
$[a\partial_z,b\partial_z]=0$ iff $ab'-ba'=0$. Embed the coefficient
algebra into $\MM(U)((t^{\Gamma}))$ with $\MM(U)$ the field of ordinary
meromorphic functions on the connected $U$; the nonzero series $a$ is
invertible there (factor its leading term and use a positive geometric
series), so $(b/a)'=(ab'-ba')/a^{2}=0$. Every Hahn coefficient of $b/a$ is
a meromorphic function with zero derivative, hence an ordinary complex
constant by connectedness, so $b/a=c\in K_\Gamma$. Conversely $b=ca$
commutes with $a$ for every scalar $c$, and has positive support exactly
when $c=0$ or $v(c)+\alpha>0$; \cref{dyn:thm:exp-log} gives
\eqref{dyn:eq:centralizer}. Equal maps give $(c_1-c_2)a=0$ and $a\ne0$, so
$c_1=c_2$; \cref{dyn:cor:flow} gives the group assertion. Meromorphic
division was used only to identify the ratio: the resulting coefficient
$ca$ is holomorphic on the original $U$, including at the zeros of $a$.
\end{proof}

Sources 03 and 05 state this theorem with the same hypotheses, the same
conclusion --- source 05 writes $v(c)+\delta>0$ for $\delta=v(W_F)$, which
is \eqref{dyn:eq:centralizer} verbatim --- the same meromorphic-ratio proof
and the same remark that zeros of $a$ need not be removed. It is therefore
printed once, per \cref{dyn:conv:merge}(i). Two consequences belong to
source 05 and survive separately: if the leading coefficient of $W_F$ is
nonzero at $z_{*}$, no nonidentity centralizer element fixes $z_{*}$, so
the basepoint normalization in \cref{dyn:thm:classification} removes
exactly the centralizer freedom, and the unnormalized conjugacies between
two conjugate maps form a torsor under either centralizer. The admissible
times $c$ may have negative valuation provided $cW_F$ stays positive; the
positivity threshold is part of the conclusion, not an unrestricted
large-time claim.

\begin{remark}[Limits of the centralizer statement]
\label{dyn:rem:centralizer-limits}
Connectedness is necessary: otherwise one obtains independent scalar series
on different ordinary components. In dimension $d>1$ the correct general
statement is only that commuting positive maps have commuting positive
logarithmic vector fields, and those need \emph{not} be scalar multiples of
one another; source 03 states that the one-variable meromorphic-ratio
argument cannot classify higher-dimensional Lie centralizers and that
additional first integrals or commuting geometric structures can occur, and
source 05 states that in several variables there is no single reciprocal
time one-form encoding a vector field, so the centralizer proof and the
finite-period normal form must not be transported. The Hamiltonian
centralizer \cref{dyn:thm:pcentralizer} is a genuinely different
several-variable statement, proved by a different argument in a different
algebra.
\end{remark}

\begin{proposition}[Workspace invariance]\label{dyn:prop:workspace}
Let $\Gamma\hookrightarrow\Gamma'$ be an ordered group inclusion. The
substitution inverse, the positive logarithm and exponential, and the
normalized nonresonant linearizer of \cref{dyn:thm:main} are unchanged when
their inputs are regarded in the larger workspace; their support
certificates remain valid; a normalized positive solution constructed in
the larger workspace acquires no exponents outside $\Gamma$; and evaluation
identities remain valid after adjoining the supports of new halo points.
\end{proposition}
\begin{proof}
The support words use only input exponents and finite addition, and the
ordinary coefficient operations are independent of $\Gamma$, so the
explicit formulas compute the same coefficients in $\Gamma$ and $\Gamma'$,
with support already inside $\Gamma$. Uniqueness of substitution inverses,
of the exponential--logarithm correspondence and of normalized
linearization identifies these with the corresponding positive solutions in
the larger workspace. Evaluation follows by including the new positive
point supports in \cref{dyn:lem:neumann}.
\end{proof}

\begin{warning}[One recorded exception: the centralizer can grow]
\label{dyn:warn:workspace}
\Cref{dyn:prop:workspace} does \emph{not} say that the centralizer
\eqref{dyn:eq:centralizer} is unchanged by scalar extension: a larger
workspace can supply new admissible times $c$. Source 03 records this
exception explicitly. What is invariant is the particular normalized
construction from the original data.
\end{warning}

\section{The homological operator and the radius of one inversion}
\label{dyn:sec:divisors}
\label{dyn:subsec:homological}

Fix $\Lambda=\operatorname{diag}(\lambda_1,\dots,\lambda_d)$ with
$|\lambda_j|=1$ and assume $\nres$: $\lambda^{\alpha}\ne\lambda_j$ for
$|\alpha|\ge2$ and all $j$. Let $B_2^{d}$ denote coefficient vectors whose
Taylor series have no terms of total degree $0$ or $1$, and put
\begin{equation}\label{dyn:eq:homological}
 \Li h=h(\Lambda z)-\Lambda h(z),\qquad
 (\Li h)_{j,\alpha}=(\lambda^{\alpha}-\lambda_j)h_{j,\alpha},
\end{equation}
so that under $\nres$ the inverse on formal coefficients is
\begin{equation}\label{dyn:eq:Linverse}
 (\Li^{-1}g)_j(z)=\sum_{|\alpha|\ge2}
 \frac{g_{j,\alpha}}{\lambda^{\alpha}-\lambda_j}z^{\alpha}.
\end{equation}
It preserves total polynomial degree and acts coefficientwise on Hahn
families without changing Hahn exponents. Recall
$\radiso(g)^{-1}=\limsup_m(\max_{|\alpha|=m,j}|g_{j,\alpha}|)^{1/m}$, the
radius of absolute convergence on centered equal-radius polydisks (the
count $\binom{m+d-1}{d-1}$ of multi-indices of degree $m$ is polynomial in
$m$ and does not affect the root test).

\begin{lemma}[Single-step analytic loss]\label{dyn:lem:radius-loss}
If $\tau=\tau(\Lambda)<\infty$ then
\begin{equation}\label{dyn:eq:loss}
 \radiso(\Li^{-1}g)\ge e^{-\tau}\radiso(g).
\end{equation}
Hence $\Li^{-1}$ preserves $\OO(\mathbb D_R^{d})_2^{d}$ when $\tau=0$,
preserves radius-free convergent germs and entire coefficients when
$\tau<\infty$, and always preserves polynomial and formal coefficients
under $\nres$ alone.
\end{lemma}
\begin{proof}
By \eqref{dyn:eq:Linverse} the maximal degree-$m$ coefficient of
$\Li^{-1}g$ is at most $M_m$ times that of $g$; take $m$-th roots and
limsups. For $\radiso(g)=\infty$ the root growth of $g$ is zero while
$\limsup M_m^{1/m}=e^{\tau}<\infty$, so the product again has root growth
zero. Polynomial support in the Taylor degree is not changed at all, and
for formal coefficients \eqref{dyn:eq:Linverse} divides by nonzero
constants.
\end{proof}

\begin{lemma}[Sharpness and analytic failure]\label{dyn:lem:sharp-divisor}
If $\tau>0$ and $0<R<\infty$, there is
$\varphi\in\OO(\mathbb D_R^{d})_2^{d}$ for which $\Li^{-1}\varphi$ is not
holomorphic on $\mathbb D_R^{d}$; it can be chosen with Taylor radius
exactly $Re^{-\tau}$, read as zero when $\tau=\infty$. If $\tau=\infty$
there is even an \emph{entire} $\varphi\in\OO(\C^{d})_2^{d}$ for which
$\Li^{-1}\varphi$ has zero radius.
\end{lemma}
\begin{proof}
Choose increasing degrees $m_k$ with $m_k^{-1}\log M_{m_k}\to\tau$ and for
each $k$ a pair $(\alpha_k,j_k)$, $|\alpha_k|=m_k$, realizing the maximal
inverse divisor; discard finitely many degrees so that $M_{m_k}>1$. Then
\begin{equation}\label{dyn:eq:finite-radius-witness}
 \varphi(z)=\sum_{k\ge1}R^{-m_k}z^{\alpha_k}e_{j_k}
\end{equation}
is holomorphic on $\mathbb D_R^{d}$, and its inverse has degree-$m_k$
coefficient of modulus $R^{-m_k}M_{m_k}$ and no other nonzero coefficient,
so its radius is exactly $Re^{-\tau}$. For the stronger assertion choose
$m_k$ with $M_{m_k}\ge\exp(k^{2}m_k)$ and put
\begin{equation}\label{dyn:eq:entire-witness}
 \varphi(z)=\sum_{k\ge1}\exp(-km_k)z^{\alpha_k}e_{j_k},
\end{equation}
which converges normally on every closed polydisk once $R\exp(-k)\to0$ and
is therefore entire, while its inverse has a degree-$m_k$ coefficient of
modulus at least $\exp((k^{2}-k)m_k)$ and Taylor radius zero.
\end{proof}

These witnesses use only ordinary analytic coefficient functions; no
infinitesimal support manipulation is hidden in them. Sources 02, 03 and 06
each construct such a witness; in one variable source 02's forcing is the
explicit rational function
\begin{equation}\label{dyn:eq:PR}
 P_R(z)=\frac{z^{2}}{1-z/R}=\sum_{n\ge2}R^{2-n}z^{n},
\end{equation}
for which $\rad(\Li^{-1}P_R)=Re^{-\sigma}$ exactly --- all its Taylor
coefficients being $R^{2-n}$, so no choice of a sparse subsequence is
needed --- and source 02's entire witness sets
$q_{n_j}=|d_{n_j}|^{1/2}$ along a sequence with
$|d_{n_j}|^{-1/n_j}\to\infty$. Source 06's version,
\cref{dyn:thm:homological}, is on frequency divisors and is stated as a
sharp \emph{iff}; source 06's explicit super-Liouville instance is
\cref{dyn:prop:explicitbad}.

\begin{proposition}[A compact-disk estimate, offered only as a starting point]
\label{dyn:prop:norm}
Suppose $d=1$ and $\sigma=0$, and let $0<r<s<R$. Then
\begin{equation}\label{dyn:eq:compactestimate}
 \|\Li^{-1}q\|_{r}\le\|q\|_{s}\sum_{n\ge2}\frac{(r/s)^{n}}{|\lambda^{n}-\lambda|},
 \qquad \|q\|_{s}=\max_{|z|\le s}|q(z)|,
\end{equation}
and the numerical series converges.
\end{proposition}
\begin{proof}
Cauchy's estimate gives $|q_n|\le\|q\|_ss^{-n}$; sum the absolutely
convergent Taylor majorant on $|z|\le r$, convergence following from
$\limsup|\lambda^{n}-\lambda|^{-1/n}=e^{\sigma}=1$.
\end{proof}

Source 02 offers \eqref{dyn:eq:compactestimate} explicitly as a
\emph{starting point} for normed questions and not as part of any theorem:
it shows the inverse is continuous for the compact-open Fr\'echet topology,
at the cost of a slightly larger compact disk, while the domain of
holomorphy remains the full open $\mathbb D_R$. None of the theorems below
imposes or produces such an estimate.

\section{The coefficient-category classification}
\label{dyn:sec:categories}

\subsection{The sharp five-row table, for an exact multiplier}
Consider
\begin{equation}\label{dyn:eq:Fintro}
 F(z)=\Lambda z+f(z),\qquad f\in\HH_\Gamma(B)_+^{d},\qquad
 f_\gamma(0)=0,\quad \mathrm df_\gamma(0)=0 ,
\end{equation}
the last condition being the exactness of the multiplier: no infinitesimal
linear part is permitted at any Hahn exponent
(tag $\sdex$, \cref{dyn:subsec:smalldivisor}).

\begin{theorem}[Universal coefficient-category classification]
\label{dyn:thm:main}
For every nonzero $\Gamma$ and every $\Lambda$ satisfying $\nres$, the
following table gives necessary and sufficient conditions for
\emph{every} map \eqref{dyn:eq:Fintro} in the indicated category to admit a
normalized conjugacy
\begin{equation}\label{dyn:eq:Schroeder}
 H(z)=z+h(z),\quad h\in\HH_\Gamma(B)_+^{d},\quad
 h_\gamma(0)=0,\ \mathrm dh_\gamma(0)=0,\qquad H\circ F=\Lambda H .
\end{equation}
Whenever it exists, that conjugacy is unique among normalized
positive-support conjugacies and its inverse lies in the same category.
{\normalfont
\begin{center}
\begin{tabular}{@{}P{0.62\textwidth}P{0.25\textwidth}@{}}
\toprule
Coefficient category $B$ & Exact condition\\
\midrule
$\OO(\mathbb D_R^{d})$, on the \emph{same fixed} polydisk
 (fixed-polydisk row) & $\tau(\Lambda)=0$\\
$\C\{z_1,\dots,z_d\}$, no common radius required (radius-free row)
 & $\tau(\Lambda)<\infty$\\
$\OO(\C^{d})$, entire coefficients (new row) & $\tau(\Lambda)<\infty$\\
$\C[z_1,\dots,z_d]$, polynomial coefficients (new row) & $\nres$ only\\
$\C[[z_1,\dots,z_d]]$, formal coefficients & $\nres$ only\\
\bottomrule
\end{tabular}
\end{center}}
If $S=\supp f\subset\Gamma_{>0}$ then $\supp h\subseteq S^{*}\setminus\{0\}$.
Each $h_\gamma$ depends on only finitely many coefficient functions of $f$
--- those at letters in $D_S(\gamma)$ --- and on at most $\ell_S(\gamma)$
ordinary homological inversions. The result is compatible with enlargement
of the Hahn workspace (\cref{dyn:prop:workspace}).
\end{theorem}

Necessity is already visible in an input with support $\{\eta\}$ for any
single $\eta>0$. Sufficiency is a genuinely nonlinear arbitrary-support
statement: no countability, finite rank, or finitely generated grid
assumption is made on $S$.

\subsection{Sufficiency: a support-controlled inverse of the nonlinear operator}
Write $F=\Lambda z+f$, seek $H=z+h$, and set
\begin{equation}\label{dyn:eq:Noperator}
 \NN_fh=h(\Lambda z+f(z))-h(\Lambda z)
 =\sum_{|\beta|\ge1}\frac{(\partial^{\beta}h)(\Lambda z)}{\beta!}f^{\beta},
\end{equation}
so that Schr\"oder's equation becomes
\begin{equation}\label{dyn:eq:linear-Schroeder}
 (\Li+\NN_f)h=-f .
\end{equation}
The map $F$ is nonlinear, but with $F$ fixed this equation is \emph{linear
in $h$}, and in the Hahn setting that observation is decisive: $\Li$ does
not alter Hahn weights while $\NN_f$ adds positive support words, and both
preserve the absence of coordinate terms of degrees zero and one. Whenever
$\Li^{-1}$ acts on the category, put $T=\Li^{-1}\NN_f$ and
\begin{equation}\label{dyn:eq:linearizer-series}
 h=-\sum_{k\ge0}(-T)^{k}\Li^{-1}f .
\end{equation}

\begin{proposition}[Support-controlled inverse; route A]
\label{dyn:prop:linearization}
Let $B_2^{d}$ be one of the five categories of \cref{dyn:thm:main} and
suppose $\Li^{-1}$ preserves it. Then \eqref{dyn:eq:linearizer-series}
defines the unique normalized positive conjugacy; its support lies in
$S^{*}\setminus\{0\}$; at weight $\gamma$ only terms with
\begin{equation}\label{dyn:eq:word-term-bound}
 k+1\le\ell_S(\gamma)
\end{equation}
contribute, and only input functions $f_\sigma$ with $\sigma\in D_S(\gamma)$
are involved.
\end{proposition}
\begin{proof}
Each term of $\NN_f$ contains at least one coefficient of $f$, hence adds a
nonempty word in $S$, and the initial input $\Li^{-1}f$ already contains one
letter; a term with $k$ copies of $T$ therefore contains at least $k+1$
letters. More explicitly, prefix the initial support exponent to the
concatenated operator words: the resulting word belongs to
$\word_S(\gamma)$ and has length at least $k+1$. This gives
\eqref{dyn:eq:word-term-bound} and puts every participating coefficient of
$f$ at a letter of $D_S(\gamma)$. \Cref{dyn:cor:ancestry}, with input
support $\supp(\Li^{-1}f)\subseteq S$, gives strong summability and the
support inclusion. This use of the unshifted depth relies on the initial
positive letter, not on an input-independent operator bound. Every
contributing expression is assembled from finitely many derivatives,
products, compositions with the fixed diagonal rotation, and applications
of $\Li^{-1}$, all of which preserve the assumed category, so $h_\gamma\in
B_2^{d}$; the operator geometric identity gives $(I+T)h=-\Li^{-1}f$, which
is \eqref{dyn:eq:linear-Schroeder}, hence $H\circ F=\Lambda H$. For
uniqueness let $h,\widetilde h$ be normalized positive solutions, not
necessarily with the constructed support, and let $\gamma_0$ be the least
weight of $q=h-\widetilde h\ne0$; since $\NN_f$ adds positive weights, the
coefficient of $\NN_fq$ at $\gamma_0$ vanishes, so $\Li q_{\gamma_0}=0$ and
$\nres$ gives $q_{\gamma_0}=0$, a contradiction. Finally $H$ has a positive
inverse in the same category by \cref{dyn:prop:subgroup}, normalized by
differentiating the inverse identity.
\end{proof}

Applying \cref{dyn:lem:radius-loss} gives sufficiency in all five rows.
Note that the polynomial row requires no bound on the inverse divisors at
all: a polynomial uses only finitely many of them at any one Hahn
coefficient, and the infinitely many nonlinear corrections are separated
into different positive Hahn weights.

\begin{theorem}[Support-controlled lifting; route B, drift permitted]
\label{dyn:thm:lifting}
Let $\lambda$ be an irrational unit multiplier ($d=1$), let
$A\subset\Gamma_{>0}$ be well ordered, and write
\begin{equation}\label{dyn:eq:data}
 F(z)=\lambda z+P(z),\quad P=\sum_{a\in A}t^{a}f_a,\quad
 \delta_a=f_a'(0),\quad \delta=\sum_at^{a}\delta_a,\quad
 \Lambda_{\mathrm{dr}}=\lambda+\delta,\quad N=P-\delta z,
\end{equation}
with $f_a\in z\C[[z]]$, so that every coefficient of $N$ vanishes to order
at least two and the multiplier is permitted to \emph{drift}. There is a
unique normalized positive-support formal solution $H=z+U$ of
$H(F(z))=\Lambda_{\mathrm{dr}}H(z)$, with
\begin{equation}\label{dyn:eq:supportcertificate}
 \supp(H-z)\subseteq A^{*}\setminus\{0\},
\end{equation}
and a normalized compositional inverse with the same bound. Moreover:
(i) if $\sigma=0$ and all $f_a$ are holomorphic on $\mathbb D_R$, all
coefficients of $H$ and $H^{-1}$ are holomorphic on that same
$\mathbb D_R$; (ii) if $\sigma<\infty$ and all $f_a$ are convergent germs,
the coefficients are convergent germs; (iii) if $\sigma<\infty$ and all
$f_a$ are entire, the coefficients are entire; (iv) if all $f_a$ are
polynomials, all coefficients are polynomials, with no arithmetic
assumption beyond $\nres$. No uniform coefficient bound, no lower
Archimedean bound on the positive exponents, and no convergence in a
complex parameter is assumed.
\end{theorem}
\begin{proof}
Substituting $H=z+U$ and expanding at $\lambda z$ gives, after cancelling
$\Lambda_{\mathrm{dr}}z$,
\begin{equation}\label{dyn:eq:masterrecursion}
 \Li U=-N+\delta U-\sum_{\beta>0}\sum_{k\ge1}t^{\beta}
 \frac{h_\beta^{(k)}(\lambda z)}{k!}P(z)^{k},
\end{equation}
whose coefficient at $\gamma\in A^{*}\setminus\{0\}$ is
\begin{equation}\label{dyn:eq:recursion}
 \Li h_\gamma=-N_\gamma
 +\sum_{\substack{a\in A,\ \beta\in A^{*}_{>0}\\ a+\beta=\gamma}}\delta_ah_\beta
 -\sum_{\substack{k\ge1,\ \beta\in A^{*}_{>0}\\ a_1,\dots,a_k\in A\\
 \beta+a_1+\cdots+a_k=\gamma}}
 \frac{h_\beta^{(k)}(\lambda z)}{k!}f_{a_1}(z)\cdots f_{a_k}(z).
\end{equation}
Every sum is finite even for a non-Archimedean order: represent $\beta$ by
a word in $A$ and append the $a_i$, then apply \cref{dyn:lem:neumann}; and
every occurring $\beta$ is strictly less than $\gamma$. The right side
vanishes to order at least two --- $N_\gamma$ by definition, $\delta_ah_\beta$
by induction, and in the last line $h_\beta^{(k)}$ vanishes to order at
least $\max(2-k,0)$ while the product of $k$ functions $f_{a_i}$ vanishes
to order at least $k$. Since $A^{*}$ is well ordered and $\Li$ is
invertible on $z^{2}\C[[z]]$, \eqref{dyn:eq:recursion} defines $h_\gamma$
by well-founded recursion, with all coefficients outside $A^{*}$ set to
zero. For (i)--(iv), induct on $\gamma$: differentiation, rotation,
multiplication and finite addition preserve each stated algebra, and
\cref{dyn:lem:radius-loss} does the same for the homological inverse ---
and in case (i) the induction retains \emph{exactly} $\mathbb D_R$ at every
exponent rather than replacing it by a sequence of smaller disks. For
uniqueness compare two normalized positive-support solutions at the least
exponent where they differ, as in \cref{dyn:prop:linearization}.
\Cref{dyn:lem:inverse} constructs the inverse in the same algebra, and
$H\circ F\circ H^{-1}(z)=\Lambda_{\mathrm{dr}}z$ is then a Hahn identity.
\end{proof}

\begin{remark}[Two genuinely different routes to a support-controlled linearizer]
\label{dyn:rem:tworoutes-lin}
\Cref{dyn:prop:linearization,dyn:thm:lifting} prove overlapping statements
by different means, and both are kept because each carries something the
other does not. Route A (source 03) inverts the nonlinear Schr\"oder
operator as an \emph{operator geometric series}
\eqref{dyn:eq:linearizer-series}, which is what yields the finite
word-depth bound \eqref{dyn:eq:word-term-bound} and hence the radius and
degree certificates \cref{dyn:thm:radius-depth,dyn:thm:degrees}; it
requires the multiplier to be exactly $\Lambda$ and works in several
variables. Route B (source 02) is an exponent-by-exponent well-founded
recursion \eqref{dyn:eq:recursion}, which is what accommodates the
\emph{drift} term $\delta_ah_\beta$; source 02 remarks that omitting that
term gives the wrong recurrence and misses the obstruction of
\cref{dyn:sec:drift} entirely. A third route, source 01's Taylor-degree
recurrence with a pre-cancelled infinitesimal divisor, appears as
\cref{dyn:lem:nonresonant,dyn:lem:displacement} and is what makes the
\emph{finite-return} threshold \cref{dyn:thm:exact} exact. Source 02 also
records the route-A statement in category-free form:
\cref{dyn:prop:reusable}.
\end{remark}

\begin{proposition}[Reusable abstract form]\label{dyn:prop:reusable}
Let $B$ be any ordinary coefficient algebra closed under differentiation,
finite products and rotation by $\lambda$, and such that the nonresonant
homological inverse preserves $z^{2}B$. Then every positive Hahn
perturbation $F=\lambda z+\sum_at^{a}f_a$ with $f_a\in zB$ admits a unique
normalized linearizer with coefficients in $B$ and support in the additive
monoid generated by the input support, and its inverse has the same
properties. No norm on $B$ is required.
\end{proposition}
\begin{proof}
This is the proof of \cref{dyn:thm:lifting} with the coefficient algebra
left abstract.
\end{proof}

The common-disk, germ, entire and polynomial rows are then obtained by
checking one closure property. \Cref{dyn:thm:collapse} shows that such a
closure check is sufficient but that passing to an \emph{inductive limit}
of domains cannot simply be interchanged with taking a full Hahn family:
for $0<\sigma<\infty$ every coefficient belongs to the germ inductive
limit while the entire Hahn family does not belong to the common-domain
inductive limit. That noncommutation is the algebraic form of the exact
radius-collapse phenomenon.

\subsection{Necessity: the first-weight obstruction}

\begin{proposition}[The first-weight obstruction]\label{dyn:prop:first-weight}
Fix $\eta>0$ and an ordinary coefficient vector $\varphi$ vanishing to
order at least two, and let $F=\Lambda z+t^{\eta}\varphi(z)$. Any
normalized positive conjugacy in any of the stated analytic categories has
no nonzero coefficients of weight below $\eta$, and necessarily
\begin{equation}\label{dyn:eq:first-weight}
 h_\eta=-\Li^{-1}\varphi .
\end{equation}
Later Hahn weights, additional scales, and enlargement of $\Gamma$ cannot
repair this.
\end{proposition}
\begin{proof}
If a coefficient below $\eta$ were nonzero, take the least nonzero weight
of $h$: there the right side of \eqref{dyn:eq:linear-Schroeder} vanishes
and $\NN_fh$ also vanishes because it adds at least $\eta$, leaving
$\Li h_\gamma=0$, contradicting $\nres$. At weight $\eta$ the same
reasoning eliminates $\NN_fh$, whose input would need weight at most zero,
and the equation is exactly \eqref{dyn:eq:first-weight}.
\end{proof}

\begin{proof}[Completion of the proof of \cref{dyn:thm:main}]
For the fixed-polydisk row with $\tau>0$, take
\eqref{dyn:eq:finite-radius-witness} and $f=t^{\eta}\varphi$:
\cref{dyn:prop:first-weight} forces a coefficient not holomorphic on the
specified $\mathbb D_R^{d}$, so universal preservation of that polydisk
fails. For the radius-free and entire rows with $\tau=\infty$, take the
entire witness \eqref{dyn:eq:entire-witness}: the forced $h_\eta$ has zero
complex radius, so it is not even a convergent germ and certainly not
entire. Higher Hahn coefficients, extra scales and enlargement of $\Gamma$
cannot change the equation at this first weight. Polynomial and formal
sufficiency was proved without arithmetic restrictions beyond $\nres$, and
the support and uniqueness claims are those of
\cref{dyn:prop:linearization}; workspace compatibility is
\cref{dyn:prop:workspace}.
\end{proof}

\begin{remark}[Nonresonance is not cosmetic]\label{dyn:rem:nres}
If $\lambda^{\alpha}=\lambda_j$ for some $|\alpha|\ge2$, take
$\varphi=z^{\alpha}e_j$: \eqref{dyn:eq:first-weight} would require $0=-1$
in that Taylor coefficient. A universal normalized linearization statement
therefore fails under $\mres$ even for a polynomial with one infinitesimal
weight. What is needed is a different target equation, not division by a
zero divisor; \cref{dyn:thm:finite-order} supplies one for multipliers of
finite order, and \cref{dyn:q:resonant} records that the general
infinite-order resonant case is open. The several-variable version of the
same remark is source 02's: a first-order forcing $tz^{\alpha}e_j$ lies
outside the range of $\Li$ at that monomial and cannot be removed by a
normalized near-identity change of variable.
\end{remark}

\begin{corollary}[The entire-coefficient threshold; polynomial universality]
\label{dyn:cor:entirethreshold}
Universal linearization of entire-coefficient perturbations by
entire-coefficient normalized coordinate changes holds exactly when the
relevant divisor rate is finite. Universal polynomial-coefficient
linearization holds for \emph{every} nonresonant multiplier.
\end{corollary}
\begin{proof}
Positive assertions: \cref{dyn:thm:main} and \cref{dyn:thm:lifting}(iii),
(iv). For the negative direction with an infinite rate, take the entire
forcing $q$ of \cref{dyn:lem:sharp-divisor} and $F=\lambda z+tq(z)$: the
coefficient of $t$ in the normalized linearizer is $-\Li^{-1}q$, which is
not even a convergent germ.
\end{proof}

\subsection{Quantitative certificates at a specified exponent}

\begin{theorem}[Radius at a specified Hahn exponent]\label{dyn:thm:radius-depth}
Suppose $\tau<\infty$ and $f=\sum_{\sigma\in S}t^{\sigma}f_\sigma$ has
convergent-germ coefficients. For $\gamma\in S^{*}\setminus\{0\}$ choose
positive $r_\sigma\le\radiso(f_\sigma)$ for the finite set
$\sigma\in D_S(\gamma)$ and put
$r_{\mathrm{in}}(\gamma)=\min_{\sigma\in D_S(\gamma)}r_\sigma$. Then
\begin{equation}\label{dyn:eq:radius-depth}
 \radiso(h_\gamma)\ge r_{\mathrm{in}}(\gamma)
 \exp\bigl(-\ell_S(\gamma)\,\tau\bigr),
\end{equation}
which for inputs on one $\mathbb D_R^{d}$ reads
$\radiso(h_\gamma)\ge R\exp(-\ell_S(\gamma)\tau)$; entire input gives
entire $h_\gamma$.
\end{theorem}
\begin{proof}
Consider one term of \eqref{dyn:eq:linearizer-series} contributing at
$\gamma$: its input letters lie in $D_S(\gamma)$; derivatives, products and
composition with a diagonal unitary map do not reduce the minimum ordinary
Taylor radius; each application of $\Li^{-1}$ costs at most $e^{-\tau}$ by
\cref{dyn:lem:radius-loss}; and a term with $k$ copies of $T$ uses $k+1$
inversions and has at least $k+1$ letters, hence at most $\ell_S(\gamma)$
inversions. Finitely many terms occur at this weight, so the minimum is
preserved under their sum.
\end{proof}

\begin{warning}[\eqref{dyn:eq:radius-depth} is a lower bound only]
\label{dyn:warn:lowerbound}
The single-inversion factor is sharp by \cref{dyn:lem:sharp-divisor}, but
it is \emph{not} asserted that the exponent $\ell_S(\gamma)\tau$ is ever
attained by a nonlinear example; cancellations can improve the radius
substantially. In particular \eqref{dyn:eq:radius-depth} alone does
\textbf{not} decide whether one common positive radius survives when
$0<\tau<\infty$ and word depths are unbounded, and source 03 states
explicitly that its coefficientwise lower bound is not a proof that radius
collapse occurs. That question is \cref{dyn:q:commongerm}; the source that
answers it does so in a different category, and \cref{dyn:sec:drift} keeps
the two apart.
\end{warning}

\begin{theorem}[Polynomial coefficient degrees]\label{dyn:thm:degrees}
Suppose every $f_\sigma$ is polynomial and put
$d_\sigma=\max_j\deg(f_{\sigma,j})\ge2$. Then
\begin{equation}\label{dyn:eq:degree-bound}
 \deg h_\gamma\le 1+\max_{(\sigma_1,\dots,\sigma_k)\in\word_S(\gamma)}
 \sum_{i=1}^{k}(d_{\sigma_i}-1),
\end{equation}
and if all $d_\sigma\le D$ then
$\deg h_\gamma\le 1+\ell_S(\gamma)(D-1)$ uniformly. No small-divisor bound
is assumed. For polynomial and formal coefficients the theorem and the
linearization construction also hold for any diagonal invertible
nonresonant $\Lambda$, \emph{dropping} $|\lambda_j|=1$.
\end{theorem}
\begin{proof}
$\Li^{-1}$ preserves total degree; the initial term at one letter $\sigma$
has degree at most $d_\sigma=1+(d_\sigma-1)$. If a term for the word $w$
has degree at most $1+\sum_{\sigma\in w}(d_\sigma-1)$, a term of $\NN_f$
differentiates it $|\beta|$ times and multiplies by $|\beta|$ further
coefficient polynomials, giving degree at most
$1+\sum_{\sigma\in w}(d_\sigma-1)+\sum_{i\le|\beta|}(d_{\sigma_i}-1)$; a
vanishing derivative contributes no term. Induct through
\eqref{dyn:eq:linearizer-series} and take the finite maximum at $\gamma$.
For polynomial and formal coefficients, composition with any invertible
diagonal matrix preserves the category and the same coefficient equations
apply.
\end{proof}

\begin{warning}[What ``finite dependence'' does and does not mean]
\label{dyn:warn:finitedep}
\Cref{dyn:prop:linearization,dyn:thm:radius-depth,dyn:thm:degrees} bound
the number of input \emph{coefficient functions} and analytic operators
needed at a specified Hahn weight. This is not an unconditional finite-time
algorithm for arbitrary analytic inputs: one application of $\Li^{-1}$ to
an arbitrary analytic function can involve infinitely many of its ordinary
Taylor coefficients. Only with polynomial input and an effective support
certificate does each requested Hahn coefficient become a finite exact
algebraic calculation; \cref{dyn:app:recipe} states exactly what must then
be certified. Source 02 adds the practical corollary that a domain
certificate cannot be a single global Boolean flag: for the $u^{k}v$
coefficient of \cref{dyn:thm:collapse} the appropriate certificate is the
exact radius $Re^{-(k+1)\sigma}$, and reporting merely ``each coefficient
converges'' would erase the decisive failure.
\end{warning}

\subsection{Actual conjugacies on halos and on the monad}

\begin{proposition}[Actual conjugacy on a finite halo]\label{dyn:prop:halo}
Under \cref{dyn:thm:main} or \cref{dyn:thm:lifting}(i), $H^{\#}$ is a
bijection of $(\mathbb D_R^{d})^{\#}_\Gamma$ with inverse
$(H^{-1})^{\#}$; both $F^{\#}$ and multiplication by the multiplier are
bijections of that halo; and
$H^{\#}(F^{\#}(w))=\Lambda H^{\#}(w)$ there. With entire or polynomial
coefficients the same holds on $\C^{\#}=\C^{d}+\m_\Gamma^{d}$.
\end{proposition}
\begin{proof}
Hahn--Taylor evaluation respects composition whenever the reductions map
the base domains into one another, by expanding both sides and regrouping
at a fixed exponent (\cref{dyn:lem:neumann} reduces this to finite Taylor
identities). The reductions of $H,H^{-1}$ are the identity and preserve
$\mathbb D_R^{d}$; the reduction of $F$ and of the linear normal form is
$z\mapsto\Lambda z$, which preserves $\mathbb D_R^{d}$ because
$|\lambda_j|=1$. Positive perturbations of the invertible $\Lambda$ have a
Hahn inverse, so the conjugate maps are bijections of the halo.
\end{proof}

\begin{proposition}[The monad needs no arithmetic]\label{dyn:prop:monad}
For every irrational unit multiplier, the formal-coefficient solution of
\cref{dyn:thm:main}/\cref{dyn:thm:lifting} evaluates on $\m_\Gamma^{d}$ and
yields a \emph{bijective} conjugacy there, with \textbf{no} arithmetic
hypothesis, even when some ordinary coefficient series have radius zero.
\end{proposition}
\begin{proof}
At $\xi\in\m_\Gamma^{d}$ use the formal Taylor coefficients at $0$ in
\eqref{dyn:eq:halo-evaluation}: a coefficient may be an arbitrary element
of $\C[[z]]$, since its scalar coefficients all have valuation zero while
powers of $\xi$ carry positive-support factors, and
\cref{dyn:lem:neumann} makes the family summable. The condition $F(0)=0$
keeps $F(\xi)\in\m_\Gamma^{d}$, and the normalized inverse and conjugacy
identities survive evaluation.
\end{proof}

\Cref{dyn:prop:monad} is the exact scope limit of every negative result in
this report: they concern analyticity of the ordinary coefficient
\emph{functions on neighborhoods}, not the absence of every form of
surcomplex linearization at infinitesimal arguments. There is no assertion
that every algebraic conjugacy of evaluated functions must arise from a
coefficient family, and failure in one analytic category prohibits no
set-theoretic conjugacy.

\subsection{Several variables}
\label{dyn:subsec:several}

For $d\ge1$, $L=\operatorname{diag}(\lambda_1,\dots,\lambda_d)$ with
$|\lambda_j|=1$ and $\nres$, put
\begin{equation}\label{dyn:eq:multisigma}
 \Delta_n=\min_{|\alpha|=n,\ 1\le j\le d}|\lambda^{\alpha}-\lambda_j|>0,
 \qquad
 \sigma_d(L)=\limsup_{n\to\infty}\frac{\log^{+}\Delta_n^{-1}}{n}.
\end{equation}
All several-variable assertions below are about the specified
\emph{equal-radius} polydisk geometry $\mathbb D_R^{d}$; no claim about an
optimal general Reinhardt domain is made or needed.

\begin{lemma}[Multivariable homological inverse]\label{dyn:lem:multihomological}
$\Li_Lh=h(Lz)-Lh(z)$ is invertible on vector formal series of order at
least two, with $(\Li_L^{-1}q)_{\alpha,j}=q_{\alpha,j}/(\lambda^{\alpha}-\lambda_j)$.
If $\sigma_d=0$ it preserves holomorphy on every $\mathbb D_R^{d}$; if
$\sigma_d<\infty$ it preserves radius-free germs and entire functions, a
finite nonzero isotropic radius losing at most the factor $e^{-\sigma_d}$;
and it preserves polynomial vectors under $\nres$ alone.
\end{lemma}
\begin{proof}
$z^{\alpha}e_j$ is an eigenvector with eigenvalue
$\lambda^{\alpha}-\lambda_j$. Taking maxima at total degree $n$ gives
$\max_{|\alpha|=n,j}|(\Li_L^{-1}q)_{\alpha,j}|\le\Delta_n^{-1}
\max_{|\alpha|=n,j}|q_{\alpha,j}|$; apply the root test, the count of
degree-$n$ monomials having only polynomial growth.
\end{proof}

\begin{theorem}[Multivariable support-controlled lifting]\label{dyn:thm:multilifting}
Let $F(z)=Lz+\sum_{a\in A}t^{a}f_a(z)$ with $f_a(0)=0$, $A\subset\Gamma_{>0}$
well ordered, and $M=\mathrm DF(0)=L+\sum_at^{a}\mathrm Df_a(0)$. Under
$\nres$ there is a unique normalized formal coordinate change $H=z+U$ with
$\mathrm DU(0)=0$ coefficientwise and
\begin{equation}\label{dyn:eq:multischroeder}
 H(F(z))=MH(z),\qquad \supp(H-z)\subseteq A^{*}\setminus\{0\}.
\end{equation}
If $\sigma_d=0$ and all $f_a$ are holomorphic on $\mathbb D_R^{d}$, then
$H$ and its inverse have coefficients holomorphic on that \emph{same}
polydisk; for finite $\sigma_d$, radius-free germs and entire coefficients
are preserved; polynomial coefficients are preserved for every nonresonant
$L$. The infinitesimal matrix $M-L$ need \emph{not} be diagonal.
\end{theorem}
\begin{proof}
With $P=F-Lz$, $E=M-L$, $N=P-Ez$, the vector form of
\eqref{dyn:eq:masterrecursion} is
\begin{equation}\label{dyn:eq:multimaster}
 \Li_LU=-N+EU-\sum_{\beta>0}\sum_{k\ge1}\frac{t^{\beta}}{k!}
 \mathrm D^{k}h_\beta(Lz)[P(z),\dots,P(z)] ,
\end{equation}
with $\mathrm D^{k}h$ the symmetric $k$-linear derivative. At a fixed Hahn
exponent only finitely many terms occur by \cref{dyn:lem:neumann}, every
$h_\beta$ occurring has strictly smaller exponent, and the right side has
total Taylor order at least two by the order count of
\cref{dyn:thm:lifting}. Apply \cref{dyn:lem:multihomological} by
well-founded recursion; the operations preserve the common polydisk because
$L$ does. The least-exponent argument proves uniqueness, and the
inverse-coordinate recursion is vectorial with coefficient one on the new
unknown, so \cref{dyn:lem:inverse} applies with multilinear derivatives.
\end{proof}

As in \cref{dyn:prop:halo}, common-polydisk solutions evaluate to actual
conjugacies on $(\mathbb D_R^{d})^{\#}_\Gamma$.

\begin{theorem}[A common coordinate for the centralizer]
\label{dyn:thm:commoncoord}
Suppose $d=1$, $\sigma(\lambda)=0$, and $F$ satisfies
\cref{dyn:thm:lifting}(i) on $\mathbb D_R$ with normalized linearizer $H$.
Let $G(z)=\mu z+\sum_{b\in B}t^{b}g_b(z)$ with $|\mu|=1$,
$g_b\in z\OO(\mathbb D_R)$, $B\subset\Gamma_{>0}$ well ordered, and suppose
$G\circ F=F\circ G$. Then, with $M=G'(0)$,
\begin{equation}\label{dyn:eq:simultaneous}
 H\circ G\circ H^{-1}(z)=Mz,
\end{equation}
with \textbf{no} nonresonance or small-divisor assumption on $\mu$.
Consequently every such $G$ commuting with $F$ is linearized by the same
$H$, any two such maps commute with each other, and $G'(0)$ determines $G$
uniquely --- in particular $G'(0)=1$ forces $G=\id$.
\end{theorem}
\begin{proof}
All compositions are legitimate on $\mathbb D_R$ by
\cref{dyn:lem:neumann} and $|\lambda|=|\mu|=1$. Put
$K=H\circ G\circ H^{-1}$: it commutes with $z\mapsto\Lambda_{\mathrm{dr}}z$,
fixes zero and has derivative $M$ there. Expanding
$K(z)=\sum_{n\ge1}k_nz^{n}$ and comparing coefficients in
$K(\Lambda_{\mathrm{dr}}z)=\Lambda_{\mathrm{dr}}K(z)$ gives
$(\Lambda_{\mathrm{dr}}^{n}-\Lambda_{\mathrm{dr}})k_n=0$ for $n\ge2$; the
factor is a nonzero Hahn-field element because its exponent-zero
coefficient is $\lambda^{n}-\lambda\ne0$. Hence $k_n=0$ for $n\ge2$ and
$k_1=M$; at each Hahn exponent the coefficient function has the Taylor
series of a linear function, so the identity theorem on $\mathbb D_R$
yields $K(z)=Mz$ as a common-domain section. Conjugating back and using
that two scalar multiplications commute gives the remaining assertions.
\end{proof}

The same proof works in the radius-free germ category when
$\sigma<\infty$, and in the polynomial setting whenever the constructed
linearizers have polynomial coefficients; in those statements equality is
interpreted in the indicated category rather than on a disk that has not
been shown to exist. A set-indexed commuting family is treated one member
at a time with the same $H$; no unsupported sum over the family is taken.

\section{The drifted category: exact radius collapse and a trichotomy}
\label{dyn:sec:drift}

This section is the one place in the report where the common-domain germ
category is decided, and it is decided \emph{only} in the drifted category
$\sddr$. \Cref{dyn:warn:drift} applies to everything in it.

\subsection{Two independent monomials in one Hahn field}
Work in the divisible ordered group
\[
 \Gamma_{*}=\Q+\sqrt2\,\Q\subset\R,\qquad u=t,\qquad v=t^{\sqrt2}.
\]
The map $(k,\ell)\mapsto k+\ell\sqrt2$ is injective on $\N^{2}$ and only
finitely many such pairs lie below any fixed real bound, so the image of
$\N^{2}$ is a well-ordered additive monoid and a double formal series
$\sum_{k,\ell\ge0}b_{k,\ell}u^{k}v^{\ell}$ is an unambiguous Hahn series
with that support. This is a use of two independent \emph{exponents}, not
of two ordinary complex parameter values. For a fixed finite $R>0$ put
\begin{equation}\label{dyn:eq:familyintro}
 F_{u,v}(z)=(\lambda+u)z+vP_R(z),\qquad P_R \text{ as in }\eqref{dyn:eq:PR},
\end{equation}
so that both perturbation coefficients are holomorphic on $\mathbb D_R$ and
the multiplier is $\Lambda_{\mathrm{dr}}=\lambda+u$.
\Cref{dyn:thm:lifting} gives a unique formal linearizer with support in
$\N+\sqrt2\,\N$. At $v=0$ the map is already linear and its normalized
linearizer is the identity --- either by setting $v=0$ in
\eqref{dyn:eq:recursion} or by ordinary formal uniqueness for the
multiplier $\lambda+u$, every $(\lambda+u)^{n}-(\lambda+u)$ having nonzero
constant term in $u$. Hence its first nontrivial $v$-block is
\begin{equation}\label{dyn:eq:Qblock}
 H(z)=z+vQ(u,z)+O(v^{2}),\qquad Q(u,z)=\sum_{k\ge0}u^{k}h_{k,1}(z),
\end{equation}
where $O(v^{2})$ means only the blocks of $v$-degree at least two and is
\emph{not} an asymptotic assertion in an ordinary parameter. Extracting the
coefficient of $v$ in $H(F)=\Lambda_{\mathrm{dr}}H$ gives
$Q(u,(\lambda+u)z)-(\lambda+u)Q(u,z)=-P_R(z)$ and hence
\begin{equation}\label{dyn:eq:blockformula}
 h_{k,1}(z)=-\sum_{n\ge2}R^{2-n}c_{n,k}z^{n},\qquad
 c_{n,k}=[u^{k}]\frac1{(\lambda+u)^{n}-(\lambda+u)} .
\end{equation}
The radius problem is thereby reduced to an exact estimate of the
\emph{reciprocal jets} $c_{n,k}$ for fixed $k$ and increasing $n$.

\subsection{Reciprocal jets and dominance of the highest pole}
With $a_{n,1}=n\lambda^{n-1}-1$ and $a_{n,j}=\binom nj\lambda^{n-j}$ for
$j\ge2$ (and $\binom nj=0$ for $j>n$), one has
$(\lambda+u)^{n}-(\lambda+u)=d_n+\sum_{j\ge1}a_{n,j}u^{j}$ with
$d_n=\lambda^{n}-\lambda$, so that for $k\ge1$
\begin{equation}\label{dyn:eq:reciprocaljet}
 c_{n,k}=\sum_{r=1}^{k}(-1)^{r}d_n^{-r-1}
 \sum_{\substack{j_1+\cdots+j_r=k\\ j_i\ge1}}a_{n,j_1}\cdots a_{n,j_r},
 \qquad c_{n,0}=d_n^{-1},
\end{equation}
and the term with $r=k$ is unique and equals
\begin{equation}\label{dyn:eq:highestpole}
 (-1)^{k}a_{n,1}^{k}d_n^{-k-1}.
\end{equation}
It is this \emph{highest} pole in the small divisor, and not merely its
first power, that drives the successive radius losses. These are exactly the
objects source 02's program checks.

\begin{lemma}[Uniform reciprocal-jet estimates]\label{dyn:lem:jets}
For each fixed $k\ge1$ there is $C_k>0$ independent of $n$ with
\begin{align}
 |c_{n,k}|&\le C_kn^{k}\max\{1,|d_n|^{-k-1}\},\label{dyn:eq:jetupper}\\
 \bigl|c_{n,k}-(-1)^{k}a_{n,1}^{k}d_n^{-k-1}\bigr|
 &\le C_kn^{k}|d_n|^{-k}\qquad\text{when }|d_n|\le1.\label{dyn:eq:jetremainder}
\end{align}
Also $|a_{n,1}|\ge n-1$. Consequently, along any sequence with
$n\to\infty$ and $d_n\to0$,
\begin{equation}\label{dyn:eq:jetasymptotic}
 c_{n,k}=(-1)^{k}a_{n,1}^{k}d_n^{-k-1}\bigl(1+o(1)\bigr),
\end{equation}
which in particular \emph{excludes} cancellation of the highest-pole
contribution along the subsequences that control exponential growth.
\end{lemma}
\begin{proof}
The reverse triangle inequality gives $|n\lambda^{n-1}-1|\ge n-1$ and
$|a_{n,1}|\le n+1$; for $j\ge2$, $|a_{n,j}|\le n^{j}/j!$. Any product
indexed by a composition $j_1+\cdots+j_r=k$ is therefore bounded by a
constant depending only on $k$ times $n^{k}$, and there are finitely many
compositions of $k$; summing \eqref{dyn:eq:reciprocaljet} gives
\eqref{dyn:eq:jetupper}, since $|d_n|\le2$ and every pole order is at most
$k+1$. Removing the $r=k$ term leaves pole orders at most $k$, giving
\eqref{dyn:eq:jetremainder} for $|d_n|\le1$. The removed term has modulus
at least $(n-1)^{k}|d_n|^{-k-1}$, and the ratio of the remainder bound to
that lower bound is at most $C_k(n/(n-1))^{k}|d_n|\to0$ when $d_n\to0$.
\end{proof}

\subsection{The exact collapse law and the trichotomy}

\begin{theorem}[Exact coefficient-radius collapse; requires drift]
\label{dyn:thm:collapse}
For the family \eqref{dyn:eq:familyintro} with $\sigma=\sigma(\lambda)>0$
($\sigma=\infty$ allowed), for every $k\ge0$,
\begin{equation}\label{dyn:eq:jetroot}
 \limsup_{n\to\infty}|c_{n,k}|^{1/n}=e^{(k+1)\sigma},
\end{equation}
and the coefficient of $u^{k}v$ in the normalized formal linearizer has
\emph{exact} radius
\begin{equation}\label{dyn:eq:collapse}
 \boxed{\;\rad(h_{k,1})=R\,e^{-(k+1)\sigma}.\;}
\end{equation}
If $0<\sigma<\infty$ the full normalized linearizer lies in the radius-free
ring $\mathcal A_1$ over $\Gamma_{*}$ but \textbf{not} in the
common-domain germ ring $\mathcal R_1$; if $\sigma=\infty$ it does not even
lie in $\mathcal A_1$.
\end{theorem}
\begin{proof}
For $k=0$, \eqref{dyn:eq:jetroot} is
$\limsup|d_n|^{-1/n}=e^{\sigma}$, which follows from the definition of
$\sigma$ because $|d_n|\le2$, so replacing $\log|d_n|^{-1}$ by its positive
part does not change the limsup after division by $n$. For fixed $k\ge1$
and finite positive $\sigma$, \eqref{dyn:eq:jetupper} gives the upper
bound after $n$-th roots since $(C_kn^{k})^{1/n}\to1$. Choose a
subsequence $n_j$ with $n_j^{-1}\log|d_{n_j}|^{-1}\to\sigma$; since
$\sigma>0$, $d_{n_j}\to0$, so \eqref{dyn:eq:jetasymptotic} applies and
\[
 |c_{n_j,k}|^{1/n_j}=|a_{n_j,1}|^{k/n_j}|d_{n_j}|^{-(k+1)/n_j}
 |1+o(1)|^{1/n_j}\longrightarrow e^{(k+1)\sigma},
\]
using $n_j-1\le|a_{n_j,1}|\le n_j+1$. If $\sigma=\infty$, choose $n_j$ with
$|d_{n_j}|^{-1/n_j}\to\infty$: again $d_{n_j}\to0$ and the same asymptotic
gives $|c_{n_j,k}|^{1/n_j}\to\infty$. The Cauchy--Hadamard formula in
\eqref{dyn:eq:blockformula} now gives \eqref{dyn:eq:collapse}. For
$0<\sigma<\infty$, \cref{dyn:thm:lifting}(ii) says every coefficient of the
full nonlinear linearizer is a convergent germ, while the radii
\eqref{dyn:eq:collapse} have infimum zero, so its coefficients cannot all
be holomorphic on any $\mathbb D_r$ with $r>0$. For $\sigma=\infty$ already
$h_{0,1}$ has radius zero. Formal uniqueness in \cref{dyn:thm:lifting}
shows that no alternative normalized solution, even with additional
positive Hahn exponents, can avoid these coefficients.
\end{proof}

\begin{remark}[Why two exponents, and why the first one does not suffice]
\label{dyn:rem:twoexponents}
The first coefficient $h_{0,1}$ alone loses the radius only by
$e^{-\sigma}$, a loss compatible with the existence of \emph{some} smaller
common disk; it is the higher multiplier jets $h_{k,1}$ that force losses
by every power $e^{-(k+1)\sigma}$ and thereby exclude such a disk. The
argument is therefore strictly stronger than a first-order homological
counterexample. The rational independence of $1$ and $\sqrt2$ prevents
different pairs $(k,\ell)$ from merging into the same Hahn coefficient and
lets the $v$-linear block be isolated without contamination by higher
$v$-powers. Source 02 records explicitly that replacing $v$ by an integer
power of $u$ would merge blocks and could introduce cancellation, so that
the proof would no longer establish the claimed radius for a particular
Hahn coefficient, and that it does not silently make that replacement. It
also records that the same necessity statement is \emph{not} asserted for a
fixed cyclic value group such as $\Z$, where different parameter blocks can
collide at one Hahn exponent; resolving the effect of those collisions is a
separate problem.
\end{remark}

\begin{corollary}[A finite-support instability]\label{dyn:cor:finitefailure}
For every irrational multiplier with $0<\sigma<\infty$, \emph{two} nonzero
positive Hahn coefficients already suffice for a holomorphic perturbation
of a rotation to have a radius-free germ linearizer and no common-domain
linearizer; the input coefficients can be taken to be rational functions of
$z$, namely $z$ at exponent $1$ and $P_R(z)$ at exponent $\sqrt2$.
\end{corollary}

\begin{theorem}[The common-domain / radius-free trichotomy]
\label{dyn:thm:trichotomy}
Let $\mathsf{CD}(\lambda)$ assert that for every set-sized ordered abelian
$\Gamma$, every $R>0$ and every normalized input \eqref{dyn:eq:data} with
coefficients on $\mathbb D_R$ there is a normalized linearizer in
$\mathcal R_1$; let $\mathsf{CG}(\lambda)$ assert that every such input
with radius-free germ coefficients has a normalized linearizer in
$\mathcal A_1$. Then for an irrational unit multiplier
\begin{equation}\label{dyn:eq:equiv}
 \mathsf{CD}(\lambda)\iff\sigma(\lambda)=0,
 \qquad
 \mathsf{CG}(\lambda)\iff\sigma(\lambda)<\infty,
\end{equation}
and both equivalences survive replacing the universal quantifier over
$\Gamma$ by the single group $\Gamma_{*}$. In the positive common-domain
case the linearizer and its inverse preserve the \emph{original} disk, not
merely an unspecified smaller one.
\end{theorem}
\begin{proof}
Positive implications: \cref{dyn:thm:lifting}(i), (ii). For
$0<\sigma<\infty$, \cref{dyn:thm:collapse} disproves $\mathsf{CD}$ within
$\Gamma_{*}$ while the radius-free implication still holds; for
$\sigma=\infty$ the same example has a nonconvergent coefficient and
disproves both.
\end{proof}

\begin{table}[ht]
\centering\small
\begin{tabular}{@{}P{0.17\textwidth}P{0.25\textwidth}P{0.24\textwidth}P{0.21\textwidth}@{}}
\toprule
Arithmetic ($\sddr$) & Common ordinary disk ($\mathcal R_1$)
& Radius-free germ / entire & Polynomial coefficients\\
\midrule
$\sigma=0$ & Always, with the \emph{same} $\mathbb D_R$
& Always / always & Always\\[3pt]
$0<\sigma<\infty$ & Not universally; explicit radii collapse to zero
& Always / always & Always\\[3pt]
$\sigma=\infty$ & Not universally & Neither universally & Always\\
\bottomrule
\end{tabular}
\caption{Universal normalized linearization in the drifted category.
Negative common-domain assertions use $\Gamma_{*}=\Q+\sqrt2\,\Q$; positive
assertions hold for every set-sized ordered abelian group. The polynomial
column means that each Hahn coefficient is a polynomial, not that the whole
coordinate change is a finite polynomial over the Hahn field.}
\label{dyn:tab:trichotomy}
\end{table}

The distinction between the middle two columns of
\cref{dyn:tab:trichotomy} is not a contradiction: for entire input
coefficients a finite exponential loss of Taylor coefficient growth leaves
every coefficient entire, whereas for a finite input radius repeated losses
can exhaust every positive common radius. The fixed coefficient algebra,
and not the mere existence of Hahn sums, determines the answer.

\subsection{The several-variable collapse}

\begin{theorem}[Exact multivariable radius collapse]\label{dyn:thm:multicollapse}
Put $P_R^{(d)}(z)=\prod_{j=1}^{d}(1-z_j/R)^{-1}-1-\sum_jz_j/R
=\sum_{|\alpha|\ge2}R^{-|\alpha|}z^{\alpha}$, let
$\mathbf1=(1,\dots,1)^{T}$, and in the same two-monomial workspace consider
\begin{equation}\label{dyn:eq:multiF}
 F_{u,v}(z)=(1+u)Lz+vP_R^{(d)}(z)\mathbf1 ,
\end{equation}
the \emph{multiplicative} detuning $(1+u)L$ being what makes the highest
reciprocal pole uniform over every component and every multi-index. Assume
$\nres$ and $\sigma_d(L)>0$. Then for every $k\ge0$
\begin{equation}\label{dyn:eq:multicollapse}
 \radiso(h_{k,1})=R\,e^{-(k+1)\sigma_d(L)} .
\end{equation}
Consequently the universal common-polydisk-germ property holds exactly when
$\sigma_d=0$, and the universal radius-free germ property exactly when
$\sigma_d<\infty$, when the workspaces include $\Gamma_{*}$.
\end{theorem}
\begin{proof}
At $v=0$ the map is linear with identity normalized coordinate change, so
the coefficient of $v$ satisfies the homological equation for the diagonal
multiplier $(1+u)L$: in component $j$ at multi-index $\alpha$ with
$n=|\alpha|$ it is $-R^{-n}/\bigl((1+u)^{n}\lambda^{\alpha}-(1+u)\lambda_j\bigr)$.
Write $D_{\alpha,j}=\lambda^{\alpha}-\lambda_j$,
$A_{\alpha,j,1}=n\lambda^{\alpha}-\lambda_j$ and
$A_{\alpha,j,r}=\binom nr\lambda^{\alpha}$ for $r\ge2$;
\eqref{dyn:eq:reciprocaljet} remains valid with these replacements, and
since all phases have modulus one, $|A_{\alpha,j,1}|\ge n-1$ and
$|A_{\alpha,j,r}|\le n^{r}/r!$, so the estimates of \cref{dyn:lem:jets}
hold with constants independent not only of $n$ but also of $\alpha$ and
$j$. The upper estimate gives
$\max_{|\alpha|=n,j}|C_{\alpha,j,k}|\le C_kn^{k}\max\{1,\Delta_n^{-k-1}\}$;
for the lower bound choose $n$ along a subsequence realizing the limsup in
\eqref{dyn:eq:multisigma} and at each such degree a pair $(\alpha,j)$
attaining the finite minimum $\Delta_n$, positivity of $\sigma_d$ ensuring
$\Delta_n\to0$, so that the highest-pole term dominates its remainder
uniformly in the chosen multi-indices. Multiplication by $R^{-n}$ and the
definition of $\radiso$ give \eqref{dyn:eq:multicollapse}. For finite
positive $\sigma_d$, \cref{dyn:thm:multilifting} makes every coefficient a
convergent germ while the displayed isotropic radii have infimum zero, and
every neighborhood of $0$ in $\C^{d}$ contains an equal-radius polydisk;
for infinite $\sigma_d$ the first $v$ coefficient already fails to be a
germ.
\end{proof}

\subsection{Continued fractions: every regime occurs, and none is Brjuno}
\label{dyn:subsec:arithmetic}

Let $\theta$ be irrational with continued-fraction convergents $p_j/q_j$,
and $\|x\|_{\Z}$ the distance to the nearest integer. The classical
estimates
\begin{equation}\label{dyn:eq:cfbounds}
 \frac1{q_j+q_{j+1}}<\|q_j\theta\|_{\Z}<\frac1{q_{j+1}},
 \qquad
 \|m\theta\|_{\Z}\ge\|q_j\theta\|_{\Z}\quad(q_j\le m<q_{j+1})
\end{equation}
are standard \cite{Marmi,CarlettiMarmi}.

\begin{proposition}[Continued-fraction formula]\label{dyn:prop:cf}
For $\lambda=e^{2\pi i\theta}$,
\begin{equation}\label{dyn:eq:tau-cf}
 \sigma(\lambda)=\tau(\lambda)
 =\limsup_{j\to\infty}\frac{\log q_{j+1}}{q_j}.
\end{equation}
\end{proposition}
\begin{proof}
For real $x$, $4\|x\|_{\Z}\le|1-e^{2\pi ix}|\le2\pi\|x\|_{\Z}$ --- the
upper bound from $|\sin y|\le|y|$, the lower from concavity of $\sin$ on
$[0,\pi/2]$ --- and constants do not affect a logarithmic limsup divided by
$m$. Hence both invariants equal
$\limsup_m m^{-1}\log\|m\theta\|_{\Z}^{-1}$, with $m=n-1$ and
$m/(m+1)\to1$. At $m=q_j$, \eqref{dyn:eq:cfbounds} gives the lower bound
$\log q_{j+1}/q_j$; for $q_j\le m<q_{j+1}$ the other half gives
$m^{-1}\log\|m\theta\|_{\Z}^{-1}\le\log(q_j+q_{j+1})/q_j
\le\log q_{j+1}/q_j+(\log2)/q_j$, and $(\log2)/q_j\to0$.
\end{proof}

\Cref{dyn:eq:tau-cf} is the point at which \emph{one} continued-fraction
quantity governs both of the one-variable divisor rates $\sddr$ and
$\sdex$; sources 02 and 03 prove it independently with the same estimates,
so it is printed once. It does \emph{not} identify the two notions: the
hypotheses that accompany them, drift versus exactness, remain different
(\cref{dyn:warn:drift}).

The classical Brjuno condition is
\begin{equation}\label{dyn:eq:brjuno}
 B(\theta)=\sum_{j\ge0}\frac{\log q_{j+1}}{q_j}<\infty ,
\end{equation}
strictly stronger than the mere convergence of those summands to zero,
which is exactly the vanishing of the rate.

\begin{theorem}[Explicit arithmetic separations]\label{dyn:thm:nonbrjuno}
The following continued fractions are explicitly specified, not
numerically approximated.
\begin{enumerate}[label=(\alph*)]
\item \emph{Rate zero, not Brjuno (source 03's form).} With $a_1=16$ and
$a_{k+1}=\lceil\exp(q_k/k)\rceil$, $q_{k+1}=a_{k+1}q_k+q_{k-1}$,
$q_0=1$, $q_1=16$, one has rate $0$ and $B(\theta)=\infty$.
\item \emph{Rate zero, not Brjuno (source 02's form).} With $a_1=2$ and
$a_{j+1}=\lceil\exp(q_j/(j+1))\rceil$, the same conclusion holds.
\item \emph{Any prescribed finite rate $s>0$.}
$a_{j+1}=\lceil e^{sq_j}\rceil$ gives rate exactly $s$.
\item \emph{Infinite rate.} $a_{j+1}=\lceil e^{q_j^{2}}\rceil$, or
$a_{k+1}=\lceil\exp(kq_k)\rceil$, gives rate $\infty$.
\item \emph{Source 01's non-Brjuno witness.} With $q_{-1}=0$, $q_0=1$ and
$a_{n+1}=2^{q_n^{2}}$, one has $(\log q_{n+1})/q_n\ge q_n\log2$, so
$B(\theta)=\infty$ (this one is used in \cref{dyn:prop:noncoherent}, where
only divergence of the Brjuno sum is needed).
\end{enumerate}
Consequently every row of \cref{dyn:thm:main} and every regime of
\cref{dyn:thm:trichotomy} is a nonempty arithmetic possibility, and the
rate-zero rows apply to rotations that are not Brjuno.
\end{theorem}
\begin{proof}
In (a), every partial quotient is a finite positive integer, so $\theta$ is
a well-specified irrational and $q_k\to\infty$; from the recursion
$q_k\exp(q_k/k)\le q_{k+1}\le q_k(\exp(q_k/k)+2)$, whence
\[
 \frac1k+\frac{\log q_k}{q_k}\le\frac{\log q_{k+1}}{q_k}
 \le\frac1k+\frac{\log q_k}{q_k}+\frac{\log(1+2e^{-q_k/k})}{q_k}.
\]
Both outer expressions tend to zero, so the rate is zero by
\cref{dyn:prop:cf}, while the middle expression is at least $1/k$, so
\eqref{dyn:eq:brjuno} diverges. Case (b) is identical with $j+1$ in place
of $k$ and $\log3$ in place of $\log(1+2e^{-q_k/k})$. In (c) and (d) the
same estimates give the quotient tending to $s$, respectively to
$q_j+o(1)$ or $k+o(1)$. In (e) the recursion gives
$q_{n+1}\ge a_{n+1}q_n=2^{q_n^{2}}q_n$, so
$\log q_{n+1}\ge q_n^{2}\log2$.
\end{proof}

\begin{remark}[No contradiction with Brjuno--Yoccoz]\label{dyn:rem:nocontra}
By the Brjuno--Yoccoz theorem the ordinary quadratic
$Q_\lambda(w)=\lambda w+w^{2}$ is locally analytically linearizable exactly
for Brjuno rotations \cite{Yoccoz,BuffCheritat}. For the multipliers of
\cref{dyn:thm:nonbrjuno}(a), (b) it is therefore \emph{not} analytically
linearizable, while every positive Hahn perturbation with coefficients on a
fixed disk \emph{is} linearizable on that same disk. This weakens nothing:
the ordinary nonlinear part in $Q_\lambda$ is of size one, whereas the
nonlinear part in \cref{dyn:thm:main} has strictly positive Hahn support
and zero ordinary reduction, and the resulting conjugacy is a Hahn family,
not a convergent ordinary family under a numerical specialization of $t$.
Sources 01, 02 and 03 all invoke Brjuno--Yoccoz for exactly this
separation, as an external classical theorem which none of them reproves.
\end{remark}

\subsection{Resolution: which category answers which question}
\label{dyn:subsec:resolution}

\begin{question}[Common-domain germs for a finite positive rate; \emph{open}
in the exact-multiplier category]\label{dyn:q:commongerm}
Assume $0<\tau(\Lambda)<\infty$ and an \emph{exact} multiplier
(\eqref{dyn:eq:Fintro}, tag $\sdex$). Characterize the positive
perturbations whose coefficients are holomorphic on one common ordinary
polydisk and whose unique radius-free linearizer has coefficients
holomorphic on \emph{some} common, possibly smaller, polydisk.
\end{question}

\begin{warning}[The two sources answer different questions; do not average
them]\label{dyn:warn:resolution}
Source 03 poses \cref{dyn:q:commongerm} as open and proves only the
word-depth lower bound \eqref{dyn:eq:radius-depth}, stating explicitly that
this does not prove radius collapse (\cref{dyn:warn:lowerbound}). Source 02
proves outright \emph{failure}: $\mathsf{CD}(\lambda)$ holds iff
$\sigma(\lambda)=0$, with the exact law \eqref{dyn:eq:collapse}, so for
$0<\sigma<\infty$ the linearizer is a radius-free coefficient germ with no
common domain whatsoever.

This is not a contradiction, and each source says so. Source 02's negative
direction \textbf{requires} multiplier drift $F'(0)=\lambda+u$ and two
rationally independent positive exponents, both of which source 03
excludes: source 03's multiplier is \emph{exactly} $\Lambda$, with no
infinitesimal linear part at any Hahn exponent. The resolution adopted
here, and the only one consistent with both proofs:
\begin{itemize}
\item \Cref{dyn:thm:collapse,dyn:thm:trichotomy} are the answer
\textbf{in the drifted category} $\sddr$. There
\cref{dyn:q:commongerm} is settled, negatively, and the law
\eqref{dyn:eq:collapse} is exact.
\item \Cref{dyn:q:commongerm} remains \textbf{open} in the
fixed-multiplier category $\sdex$. Source 02 itself records that
eliminating drift removes the direct reciprocal-jet mechanism and that a
different counterexample, or a positive cancellation theorem, would be
needed.
\item The two must not be averaged into a single ``intermediate case''
claim, and \eqref{dyn:eq:collapse} must never be quoted without its drift
hypothesis.
\end{itemize}
For the weaker requirement of retaining the \emph{original} disk, one
positive parameter is already enough: \cref{dyn:lem:sharp-divisor} and the
first-order equation give failure whenever the rate is positive. The
unresolved issue concerns the existence of \emph{some} smaller common
neighborhood for every input.
\end{warning}

\section{The exact finite-return threshold and its shell cycles}
\label{dyn:sec:return}

Everything in this section is under $\resq$ and $\sdz$: $\lambda$ is
\emph{never} a root of unity, its residue $\alpha=\bar\lambda$ has exact
finite order $q$, and \textbf{no} arithmetic condition of any kind is
imposed. \Cref{dyn:warn:resonance,dyn:warn:charzero,dyn:warn:ball} apply
throughout. The ambient $\Gamma$ is divisible here, so that $K_\Gamma$ is
algebraically closed and rational division of exponents is available.

\subsection{Two support-controlled linearization lemmas; route C}
For $T(X)=\nu X+\sum_{j=2}^{D}d_jX^{j}$ and
$H(X)=X+\sum_{n\ge2}h_nX^{n}$ with $\nu\ne0$ not a root of unity,
comparison of degree $n$ in $H(T(X))=\nu H(X)$ gives the triangular
recurrence
\begin{equation}\label{dyn:eq:recurrence}
 (\nu-\nu^{n})h_n=\sum_{j=1}^{n-1}h_j\,[X^{n}]T(X)^{j},\qquad h_1=1,
\end{equation}
whose denominators are nonzero, so \eqref{dyn:eq:recurrence} defines a
unique normalized formal conjugacy over any field of characteristic zero.
What remains is to control all of its coefficients \emph{simultaneously}.
This is route C of \cref{dyn:rem:tworoutes-lin}.

\begin{lemma}[Nonresonant residue]\label{dyn:lem:nonresonant}
Suppose $T\in\OO_\Gamma[X]$, $v(\nu)=0$, and $\alpha=\bar\nu$ is \emph{not}
a root of unity in $\C$. Then the normalized formal conjugacy and its
inverse belong to $\BB_{\mathsf M}$ for a positive-support monoid
$\mathsf M$ depending only on the finitely many coefficients of $T$.
\end{lemma}
\begin{proof}
Write $\nu=\alpha+e$ with $e\in\m_\Gamma$ and let $S$ be the finite union
of the positive supports of $e$ and of the nonlinear coefficients $d_j$.
For $n\ge2$ the element $\nu-\nu^{n}$ has nonzero constant term
$\alpha-\alpha^{n}$, so its inverse is a formal power series in $e$ with
ordinary complex coefficients and support in $\langle\supp e\rangle$. Then
\eqref{dyn:eq:recurrence} proves inductively that every $h_n$ is integral
with support in $\mathsf M=\langle S\rangle$, and the formal inverse has
the same property by \cref{dyn:lem:evaluation}.
\end{proof}

\begin{remark}[No Diophantine bound, and why that is a real statement]
\label{dyn:rem:nodiophantine}
There is \emph{no} lower bound on the ordinary complex numbers
$\alpha-\alpha^{n}$ anywhere in \cref{dyn:lem:nonresonant}. Those constants
may be arbitrarily small in the ordinary modulus and still all have Hahn
valuation zero (\cref{dyn:rem:invisible}), which is exactly what makes this
a $\sdz$ statement. It is not the assertion that ``all small divisors are
units, so all analytic problems disappear'': \cref{dyn:prop:noncoherent}
shows what is lost, and \cref{dyn:ex:bad-support} shows that
coefficient-valuation bounds alone are insufficient at rank $>1$.
\end{remark}

\begin{lemma}[Divisible displacement: cancelling a common infinitesimal
divisor]\label{dyn:lem:displacement}
Let $0\ne c\in\m_\Gamma$, $A\in X^{2}\OO_\Gamma[X]$ and
\begin{equation}\label{dyn:eq:T-displacement}
 T(X)=(1+c)X+cA(X)=X+cV(X),\qquad V=X+A .
\end{equation}
Then $1+c$ is not a root of unity, and the normalized formal conjugacy of
$T$ to $(1+c)X$, together with its inverse, has all coefficients supported
in one positive-support monoid generated by $\supp c$ and the positive
parts of the coefficients of $A$.
\end{lemma}
\begin{proof}
For a positive integer $n$ the first nonzero term of $(1+c)^{n}-1$ is
$nc$, all remaining terms having strictly larger valuation, so the
multiplier is not a root of unity. For $n\ge2$,
\begin{equation}\label{dyn:eq:denomcancel}
 (1+c)-(1+c)^{n}=-c(n-1)u_n(c),\qquad u_n(0)=1,\quad u_n\in\Q[c].
\end{equation}
For $j<n$ the coefficient $[X^{n}]T(X)^{j}$ contains at least one nonlinear
term, hence is divisible by $c$ as a polynomial in $c$ and the coefficients
of $A$. Substituting this and \eqref{dyn:eq:denomcancel} into
\eqref{dyn:eq:recurrence} cancels one factor $c$ \emph{before any
valuation estimate is made}. The remaining division is by the nonzero
ordinary integer $n-1$ --- a valuation unit in equal characteristic zero,
which is where \cref{dyn:warn:charzero} enters --- and by $u_n(c)$, whose
inverse is supported in $\langle\supp c\rangle$. Induction on $n$ gives one
common monoid, a finite union of well-ordered positive supports; inversion
of $H$ uses only polynomials in its coefficients.
\end{proof}

This cancellation is the whole reason the exact threshold is read from a
\emph{return} map: multiplying all the small divisors in
\eqref{dyn:eq:recurrence} together would lose the common factor present in
every nonlinear numerator. Source 01 states this explicitly, and no other
source performs this cancellation.

\begin{lemma}[Roots on the unit shell]\label{dyn:lem:root-count}
Let $Q\in\OO_\Gamma[X]$ with $Q(0)=1$ and suppose $P=\bar Q$ has degree
$M>0$. Then $Q$ has no root in $\m_\Gamma$ and exactly $M$ roots of
valuation zero, counted with algebraic multiplicity; for each nonzero root
$a\in\C$ of $P$, the roots of $Q$ with residue $a$ have total multiplicity
equal to the multiplicity of $a$ in $P$.
\end{lemma}
\begin{proof}
A root of positive valuation is impossible because the constant term would
strictly dominate. Factor in the algebraically closed $K_\Gamma$ as
$Q(X)=\prod_j(1-X/\beta_j)$ with roots repeated by multiplicity; every
$v(\beta_j)\le0$ so all factors are integral, a factor with
$v(\beta_j)<0$ reduces to $1$, and one with $v(\beta_j)=0$ reduces to
$1-X/\bar\beta_j$. Reducing this finite factorization proves all
assertions; a simple root of $P$ corresponds to exactly one simple root of
$Q$ in that residue class.
\end{proof}

\subsection{The exact maximal ball}
Let
\begin{equation}\label{dyn:eq:f-main}
 f(z)=\lambda z+\sum_{n=2}^{d}a_nz^{n}\in K_\Gamma[z],\qquad d\ge2,\ a_d\ne0,
 \ v(\lambda)=0,
\end{equation}
with $\lambda$ not a root of unity, and put
$\eta(f)=\max_{a_n\ne0}\bigl(-v(a_n)/(n-1)\bigr)$, $\alpha=\bar\lambda$.
All maxima are finite maxima over nonzero polynomial coefficients.

\begin{theorem}[Exact finite-return linearization]\label{dyn:thm:exact}
Define $r=r(f)$ as follows. If $\alpha$ is \emph{not} a root of unity, set
$r=\eta(f)$. If $\alpha$ has exact order $q\ge1$ (tag $\resq$), write
\begin{equation}\label{dyn:eq:g-main}
 g(z)=f^{\circ q}(z)=\lambda^{q}z+\sum_{n=2}^{d^{q}}b_nz^{n},
 \qquad c=\lambda^{q}-1,\qquad \delta=v(c)>0,
\end{equation}
and set
\begin{equation}\label{dyn:eq:exact-r}
 \boxed{\;r=\max_{b_n\ne0}\frac{\delta-v(b_n)}{n-1}\;}
\end{equation}
--- read off the \emph{actual coefficients of the finite $q$-th return
iterate}, including all cancellations --- in which case $r>\eta(f)$.

Let $H(z)=z+\sum_{n\ge2}h_nz^{n}$ be the unique normalized formal solution
of $H\circ f=\lambda H$. Then $H$ and its formal inverse are strongly
Hahn-evaluable on $B_r$ and their evaluated maps are mutually inverse
isometries of $B_r$; in particular $f:B_r\to B_r$ is conjugate to
multiplication by $\lambda$. More strongly, with $\rho=t^{r}$,
\begin{equation}\label{dyn:eq:H-scaled}
 H_r(X)=\rho^{-1}H(\rho X)=X+\sum_{n\ge2}h_n\rho^{n-1}X^{n}
\end{equation}
and its inverse have a uniform positive-support certificate, so in
particular $v(h_n)+(n-1)r\ge0$ for $n\ge2$ --- though the certificate is
strictly stronger than that inequality.

No strictly larger \emph{centered valuation ball} admits an injective
conjugacy of $f$ to multiplication by $\lambda$. In the nonresonant-residue
case the obstruction is a nonzero $p\in S_r$ with $f(p)=0$; in the
$\resq$ case it is a nonzero $p\in S_r$ of exact period $q$.
\end{theorem}
\begin{proof}
\emph{Nonresonant residue.} Take $\rho=t^{\eta}$ and
$F(X)=\rho^{-1}f(\rho X)$; its nonlinear coefficients have valuations
$v(a_n)+(n-1)\eta\ge0$ with at least one equal to zero, so
$F\in\OO_\Gamma[X]$ with nonlinear reduction. By
\cref{dyn:lem:nonresonant,dyn:lem:evaluation} its normalized conjugacy and
inverse are strongly summable mutually inverse isometries of $\m_\Gamma$,
and scaling back gives every existence and support assertion. For
maximality apply \cref{dyn:lem:root-count} to
$Q(X)=F(X)/(\lambda X)=1+\sum_{n\ge2}(a_n\rho^{n-1}/\lambda)X^{n-1}$,
whose reduction is nonconstant: $Q$ has a unit root $u$, and
$p=\rho u\in S_\eta$ satisfies $f(p)=0=f(0)$. Any centered ball strictly
containing $B_\eta$ contains that whole shell, and an injective conjugacy
there would give $\lambda H(p)=H(f(p))=H(f(0))=\lambda H(0)$.

\emph{The resonant threshold lies inside the injectivity scale.} With
$\alpha$ of order $q$, the multiplier $\lambda^{q}$ has residue one and is
not one because $\lambda$ is not a root of unity, so $c\ne0$ and
$\delta>0$. Scaling by $\sigma_0=t^{\eta}$, the polynomial
$F_\eta(X)=\sigma_0^{-1}f(\sigma_0X)$ is integral with reduction of degree
at least two; reduction commutes with polynomial iteration, so
$\overline{\sigma_0^{-1}g(\sigma_0X)}=\bar F_\eta^{\circ q}(X)$ also has
degree at least two, and some $n\ge2$ satisfies $v(b_n)+(n-1)\eta=0$. For
that index $(\delta-v(b_n))/(n-1)=\eta+\delta/(n-1)>\eta$, so
$r>\eta$. At the final scale $\rho=t^{r}$ the polynomial
$F(X)=\rho^{-1}f(\rho X)$ has reduction exactly $\alpha X$: every nonlinear
coefficient acquires positive valuation. It maps $\m_\Gamma$ to itself and
preserves the residue direction of a unit up to multiplication by $\alpha$.

\emph{Linearizing the return and recovering $f$.} At $\rho=t^{r}$ write
\begin{equation}\label{dyn:eq:G}
 G(X)=\rho^{-1}g(\rho X)=(1+c)X+cA(X),\qquad
 A(X)=\sum_{n\ge2}\frac{b_n\rho^{n-1}}{c}X^{n}\in X^{2}\OO_\Gamma[X],
\end{equation}
integrality being exactly \eqref{dyn:eq:exact-r}.
\Cref{dyn:lem:displacement} constructs the normalized conjugacy $J$ of $G$
to $(1+c)X$ with a uniform support certificate. It remains to see that $J$
linearizes $F$ and not merely its $q$-th iterate: in the formal
power-series group $T=J\circ F\circ J^{-1}$ commutes with multiplication by
$\mu=1+c$ because $F$ commutes with $G=F^{\circ q}$, so writing
$T(X)=\sum_{n\ge1}\tau_nX^{n}$ the commutation says
$\tau_n(\mu^{n}-\mu)=0$; since $\mu$ is not a root of unity all $\tau_n$
with $n\ge2$ vanish, and the linear coefficient is $\lambda$, so
$T(X)=\lambda X$ and $J$ is the unique normalized conjugacy of $F$.
Adjoining the positive supports of the finitely many coefficients of $F$ to
the certificate, all compositions are legitimate by
\cref{dyn:lem:evaluation}, and scaling by $\rho$ gives the existence,
isometry and coefficient assertions.

\emph{The exact-period obstruction.} Define the integral polynomial
\begin{equation}\label{dyn:eq:Q-shell}
 Q(X)=\frac{G(X)-X}{cX}=1+\sum_{n\ge2}\frac{b_n\rho^{n-1}}{c}X^{n-1}.
\end{equation}
At least one nonlinear coefficient of $G$ attains equality in
\eqref{dyn:eq:exact-r}, so $P=\bar Q$ is nonconstant, and
\cref{dyn:lem:root-count} gives a unit root $u$; then $p=\rho u\in S_r$ and
$g(p)=p$. The least period $k$ of $p$ divides $q$; every point of its
finite orbit lies on $S_r$ and reduction of the normalized orbit multiplies
by $\alpha$ at each step, so if $a=\bar u\ne0$ then $\alpha^{k}a=a$ and
$q\mid k$, whence $k=q$. On any larger centered ball an injective conjugacy
would give $H(p)=H(f^{\circ q}(p))=\lambda^{q}H(p)$, and $\lambda^{q}\ne1$
forces $H(p)=0=H(0)$.
\end{proof}

\begin{remark}[Exactly what is exact]\label{dyn:rem:exact-scope}
\Cref{dyn:thm:exact} determines the maximal \emph{centered valuation ball}
only: not the full pointwise summability locus of the formal series
(exceptional evaluations on the bounding shell are allowed), not a maximal
non-ball invariant domain, and not every analytic continuation of the
conjugacy. The formal-coefficient family lands in $\mathcal F_1$;
\cref{dyn:thm:coherent} is what reaches the fixed-polydisk row, and only
under $\resq$.
\end{remark}

\subsection{The shell polynomial and the closest cycles}
Retain $r,\rho,c,Q$ and put
\begin{equation}\label{dyn:eq:P}
 P(X)=\bar Q(X),\qquad
 M=\deg P=\max\{n-1:v(b_n)+(n-1)r=\delta\},
\end{equation}
a finite ordinary complex polynomial with $P(0)=1$: the initial polynomial
on the first face of the return equation seen from the origin.

\begin{theorem}[Resonant shell cycles]\label{dyn:thm:shell}
\begin{enumerate}[label=(\roman*)]
\item $f^{\circ q}(z)=z$ has exactly $M$ nonzero solutions on $S_r$,
counted with algebraic multiplicity; every such solution has exact period
$q$; and there are no nonzero periodic points in $B_r$.
\item $P(\alpha X)=P(X)$, i.e.\ $P\in\C[X^{q}]$; in particular $q\mid M$.
\item If $P$ is squarefree there is one simple point
$p_a=\rho(a+u_a)$, $u_a\in\m_\Gamma$, for every root $a$ of $P$; these form
exactly $M/q$ distinct cycles, and $f(p_a)=p_{\alpha a}$.
\item For each such simple point,
\begin{equation}\label{dyn:eq:multiplier}
 \res\left(\frac{(f^{\circ q})'(p_a)-1}{c}\right)=aP'(a),
 \qquad v\bigl((f^{\circ q})'(p_a)-1\bigr)=\delta .
\end{equation}
\end{enumerate}
\end{theorem}
\begin{proof}
Counting and residue classes follow from \cref{dyn:lem:root-count} applied
to $Q$; exact period $q$ was proved in \cref{dyn:thm:exact}; a nonzero
periodic point in $B_r$ would be sent by the injective conjugacy to a
nonzero periodic point of multiplication by $\lambda$, which has none.
For the symmetry write $F=\rho^{-1}f(\rho X)$ and
$G=F^{\circ q}=X+cW(X)$ with $W=XQ$. The identity $F\circ G=G\circ F$ gives
$\bigl(F(X+cW(X))-F(X)\bigr)/c=W(F(X))$, an identity of integral
polynomials whose reduction leaves only the first Taylor term on the left
since $c\in\m_\Gamma$; as $\bar F=\alpha X$ this is
$\alpha\bar W(X)=\bar W(\alpha X)$, and substituting $\bar W=XP$ gives
$P(\alpha X)=P(X)$, so only degrees divisible by $q$ occur. When $P$ is
squarefree, uniqueness in each residue class gives $f(p_a)=p_{\alpha a}$,
and the action of $\alpha$ on $\C^{\times}$ has orbits of length $q$.
Finally $G'(X)=1+cW'(X)$, and at $p_a=\rho u$ with $\bar u=a$ the
derivative of the unscaled iterate equals $G'(u)$, whose reduction is
$\overline{W'(u)}=P(a)+aP'(a)=aP'(a)\ne0$.
\end{proof}

\begin{remark}[Multiplicity is not the number of distinct points]
\label{dyn:rem:multiplicity}
If $P$ has a multiple root, \cref{dyn:thm:shell} gives a multiplicity count
only. It does \emph{not} say that several actual periodic points have
collided: higher Hahn coefficients can split one repeated initial root into
distinct points with the same residue direction, and in that situation
$v((f^{\circ q})'(p)-1)=\delta$ need not hold, because the displayed
residue may vanish. \Cref{dyn:ex:kappa} exhibits the degeneracy
concretely.
\end{remark}

\begin{corollary}[A finite obstruction certificate]\label{dyn:cor:finite-certificate}
In the $\resq$ case the data $q$, $c$, $r$, $P$ certify exact maximality
without computing a single coefficient of the infinite conjugacy. At most
$d^{q}-1$ shell points occur with multiplicity, and their count is a
positive multiple of $q$.
\end{corollary}
\begin{proof}
The return polynomial has degree $d^{q}$ and $M=\deg P\le d^{q}-1$;
exactness and positivity are \cref{dyn:thm:exact,dyn:thm:shell}.
\end{proof}

\Cref{dyn:app:certificate} states the accompanying four-step finite
algebraic procedure and \emph{exactly} what it is relative to.

\subsection{A finite-return perturbation certificate}

\begin{theorem}[Stability above the first return face]\label{dyn:thm:stability}
Let $f$ satisfy the $\resq$ hypotheses of \cref{dyn:thm:exact} with data
$q,c,\delta,r,P$, and let $\widetilde f$ be another nonlinear polynomial
with the \emph{same} multiplier $\lambda$ and return coefficients
$\widetilde b_n$ (pad either list with zeros). Suppose
\begin{equation}\label{dyn:eq:stable-assumption}
 v(\widetilde b_n-b_n)+(n-1)r>\delta
 \qquad\text{for every nonzero difference.}
\end{equation}
Then $r(\widetilde f)=r(f)$ and the shell polynomial of $\widetilde f$ is
also $P$. If $P$ is squarefree and at least one coefficient changes, put
\begin{equation}\label{dyn:eq:epsilon}
 \epsilon=\min_{\widetilde b_n\ne b_n}
 \bigl(v(\widetilde b_n-b_n)+(n-1)r-\delta\bigr)>0 ;
\end{equation}
the closest cycles have a canonical matching by residue direction, matched
points satisfy
\begin{equation}\label{dyn:eq:stable-distance}
 v(\widetilde p_a-p_a)\ge r+\epsilon ,
\end{equation}
the matching intertwines $f$ and $\widetilde f$ on these finite cycle sets,
and their leading multiplier ratios in \eqref{dyn:eq:multiplier} agree.
\end{theorem}
\begin{proof}
At $\rho=t^{r}$ the two normalized displacement polynomials are $Q$ as in
\eqref{dyn:eq:Q-shell} and $\widetilde Q$ with $\widetilde b_n$ in place of
$b_n$. The hypothesis says every coefficient of $\widetilde Q-Q$ has
positive valuation, so $\widetilde Q$ is integral with the same nonconstant
reduction $P$; the inequalities defining the maximum hold at $r$ for all
$\widetilde b_n$ and at least one remains an equality, so
$r(\widetilde f)=r$. For a simple root $a$ of $P$ let $u=p_a/\rho$ and
$\widetilde u=\widetilde p_a/\rho$, both of residue $a$. When
$u\ne\widetilde u$ the divided difference
$(Q(u)-Q(\widetilde u))/(u-\widetilde u)$ is integral with nonzero residue
$P'(a)$; moreover $Q(u)=0$ and
$Q(\widetilde u)=(Q-\widetilde Q)(\widetilde u)$, of valuation at least
$\epsilon$, whence $v(u-\widetilde u)\ge\epsilon$, which is
\eqref{dyn:eq:stable-distance} (it is automatic if the two roots coincide).
The direction rule $a\mapsto\alpha a$ and \cref{dyn:thm:shell} give the
rest.
\end{proof}

\begin{warning}[The certificate is sufficient, not necessary; and it is
about the return]\label{dyn:warn:stability}
A perturbation may change the first return face and nevertheless preserve
its maximal slope, so \eqref{dyn:eq:stable-assumption} is sufficient only.
It is also a condition on the \emph{return} polynomial: it must not be
silently replaced by the same inequalities on the original coefficients of
$f$. Where coefficients are available only to a specified precision,
\cref{dyn:thm:stability} says which strict inequalities certify that the
computed first face is already correct; at a tie or a possible cancellation
more precision may be necessary, and treating a truncated zero as an exact
zero would invalidate the maximum.
\end{warning}

\section{Common-domain lifting: one leading coefficient, two instantiations}
\label{dyn:sec:coherent}

The support arguments of \cref{dyn:sec:return} establish more than formal
existence but less than ordinary analyticity of every coefficient function.
Under $\resq$ a further argument supplies exactly that regularity, and it
lands in the fixed-polydisk row of \cref{dyn:subsec:rings}.

\begin{theorem}[Coherent realization of resonant linearization]
\label{dyn:thm:coherent}
Under the $\resq$ hypotheses of \cref{dyn:thm:exact}, the scaled conjugacy
$H_r$ of \eqref{dyn:eq:H-scaled} has an expansion
\begin{equation}\label{dyn:eq:coherent-expansion}
 H_r(X)=\sum_{\gamma\in\mathsf M}H_\gamma(X)t^{\gamma}
 \quad\text{in }\OO(D)((t^{\Gamma}))
\end{equation}
for a \emph{single} ordinary complex disk $D$ centered at zero, with
$\mathsf M$ a positive-support monoid, $H_0(0)=0$, $H_0'(0)=1$, and
$H_\gamma\in X^{2}\OO(D)$ for $\gamma>0$. One may take any sufficiently
small disk on which the shell polynomial $P$ has no zero, and the leading
coefficient is
\begin{equation}\label{dyn:eq:H0}
 \boxed{\;H_0(X)=X\exp\left(\int_0^{X}
 \Bigl(\frac1{sP(s)}-\frac1s\Bigr)\dd s\right),\;}
\end{equation}
whose integrand is holomorphic at zero. After an ordinary shrinking, the
formal inverse of $H_r$ also has all Hahn coefficient functions holomorphic
on one common ordinary neighborhood.
\end{theorem}
\begin{proof}
Use $G=X+cV$ from \eqref{dyn:eq:G}, so $V=X+A=XQ$ and
$V_0:=\bar V=XP(X)$; all coefficients of the normalized conjugacy lie in a
common monoid $\mathsf M$ by \cref{dyn:lem:displacement}. Taylor expansion
of the formal identity $H_r(X+cV)=(1+c)H_r(X)$, followed by cancellation of
$c$, gives
\begin{equation}\label{dyn:eq:differential-schroeder}
 VH_r'+\sum_{k\ge2}\frac{c^{k-1}}{k!}V^{k}H_r^{(k)}=H_r .
\end{equation}
This is first read coefficientwise in $X$, where only finitely many terms
contribute at each degree because $V$ has order one; it may also be
regrouped at a fixed Hahn exponent, because the term of index $k$ uses at
least $k-1$ positive generators from $c$, so \cref{dyn:lem:neumann} makes
that regrouping finite. At exponent zero,
\eqref{dyn:eq:differential-schroeder} becomes
\begin{equation}\label{dyn:eq:ode0}
 V_0H_0'=H_0 .
\end{equation}
Choose an ordinary disk $D$ centered at zero on which $P$ has no zero;
since $P(0)=1$ the function $1/(XP(X))-1/X$ has a removable singularity at
zero, and \eqref{dyn:eq:H0} defines the unique solution of
\eqref{dyn:eq:ode0} with $H_0'(0)=1$, with $H_0/X$ nonvanishing on $D$.

Now let $\gamma>0$ in the well-ordered monoid $\mathsf M$ and assume
$H_\beta$ constructed on this same disk for every $\beta<\gamma$.
Isolating the terms containing $H_\gamma$ in
\eqref{dyn:eq:differential-schroeder} gives
\begin{equation}\label{dyn:eq:ode-gamma}
 LH_\gamma:=V_0H_\gamma'-H_\gamma=R_\gamma,\qquad
 R_\gamma=-[t^{\gamma}]\Bigl((V-V_0)H_r'
 +\sum_{k\ge2}\frac{c^{k-1}}{k!}V^{k}H_r^{(k)}\Bigr).
\end{equation}
Every appearance of a coefficient of $H_r$ on the right is multiplied by a
positive exponent, so only $H_\beta$ with $\beta<\gamma$ occurs, and
\cref{dyn:lem:neumann} makes $R_\gamma$ a \emph{finite} expression in
earlier coefficient functions and their derivatives --- hence holomorphic
on the same $D$, with no new shrinking. Moreover $V-V_0$ vanishes to order
at least two, since the coefficient of $X$ in $V$ is exactly one, and the
terms with $k\ge2$ likewise, so $R_\gamma\in X^{2}\OO(D)$. Variation of the
constant in \eqref{dyn:eq:ode0} gives the unique normalized solution
\begin{equation}\label{dyn:eq:gamma-integral}
 H_\gamma(X)=H_0(X)\int_0^{X}\frac{R_\gamma(s)}{V_0(s)H_0(s)}\dd s ,
\end{equation}
whose apparent singularity at zero is removable (numerator of order at
least two, denominator of order exactly two) and whose denominator is
nonzero elsewhere on $D$; thus $H_\gamma\in X^{2}\OO(D)$. Transfinite
induction on $\mathsf M$ constructs all coefficients on one fixed disk, and
they have the same formal Taylor expansions as the unique solution of
\eqref{dyn:eq:recurrence}, so this analytic family is the regrouping of the
$H_r$ already constructed, not a different conjugacy. For the inverse,
first shrink to a neighborhood on which $H_0$ is biholomorphic; the
exponent-zero coefficient of the inverse is $H_0^{-1}$, and at any positive
exponent the coefficient equation for composition is linear in the unknown
with coefficient $H_0'(H_0^{-1}(Y))$, which never vanishes on a
sufficiently small fixed neighborhood, all other terms depending on earlier
exponents through finitely many ordinary derivatives and products.
\end{proof}

\begin{remark}[Why the differential recursion is genuinely triangular]
\label{dyn:rem:triangular}
It would be \emph{incorrect} to justify \eqref{dyn:eq:ode-gamma} by saying
that $(k-1)v(c)\to+\infty$: that can fail to be cofinal in an
arbitrary-rank value group (\cref{dyn:warn:contraction}). The correct
statement is that at a \emph{specified} exponent $\gamma$ infinitely many
decompositions into generators from one well-ordered positive support are
impossible, which simultaneously controls the Taylor index, the Hahn
products and the number of earlier coefficients entering the right side.
The construction is transfinite, but each individual coefficient equation
has finite dependence.
\end{remark}

\begin{warning}[No uniform norm bound on the coefficient family]
\label{dyn:warn:nouniform}
\Cref{dyn:thm:coherent} asserts \emph{no} uniform ordinary norm bound on
$(H_\gamma)$: the coefficient functions may be arbitrarily large with
arbitrarily large derivatives. What is uniform is their ordinary domain of
holomorphy and the Hahn support certificate. The same disclaimer is made
independently by sources 02, 03, 05 and 06 for their own coefficient
families.
\end{warning}

\subsection{The same formula, twice, and why both instantiations survive}
\label{dyn:subsec:sameformula}

\begin{proposition}[The leading-coefficient exponential integral]
\label{dyn:prop:leading}
Let $U$ be a simply connected ordinary domain containing $0$, and let
$V_0\in\OO(U)$ satisfy $V_0(0)=0$, $V_0'(0)=\mu\ne0$ and $V_0(w)\ne0$ for
$w\in U\setminus\{0\}$. The unique solution of $V_0h_0'=\mu h_0$ on $U$
with $h_0(0)=0$, $h_0'(0)=1$ is
\begin{equation}\label{dyn:eq:leadingH}
 h_0(w)=w\exp\left(\int_0^{w}
 \Bigl(\frac{\mu}{V_0(\zeta)}-\frac1\zeta\Bigr)\dd\zeta\right),
\end{equation}
whose integrand is holomorphic at $0$ and whose quotient $h_0/w$ is
nonvanishing on $U$.
\end{proposition}
\begin{proof}
$\mu/V_0-1/w$ has a removable singularity at $0$ because
$(V_0/w)(0)=\mu$; on a simply connected $U$ its primitive vanishing at $0$
exists and is unique, and exponentiating gives a solution with the stated
normalization, unique because the logarithmic derivative determines $h_0$
up to a constant fixed at $0$.
\end{proof}

Formula \eqref{dyn:eq:leadingH} is \emph{one} formula. It is source 01's
\eqref{dyn:eq:H0} with $V_0=XP(X)$ and $\mu=1$, and it is source 04's
leading Koenigs coordinate with $V_0=V$ the Euler polynomial and
$\mu=V'(0)$, and both sources prove it by the same transfinite
variation-of-constants recursion --- source 01's
\eqref{dyn:eq:gamma-integral}, source 04's
\eqref{dyn:eq:Hconstruction}. It is therefore printed once, per
\cref{dyn:conv:merge}(i). What does \emph{not} merge is the pair of
instantiations, and both survive:
\begin{itemize}
\item under $\resq$, $V_0=XP(X)$ where $P$ is the \emph{shell polynomial}
of \eqref{dyn:eq:P}, an ordinary polynomial determined by the first face of
the return equation. Its nonzero roots are the closest periodic directions
in the valued field, and its \emph{absence} of zeros on a small ordinary
disk is what permits \eqref{dyn:eq:ode0} to be solved there. One finite
polynomial thus plays two different roles, and \eqref{dyn:eq:H0} is the
precise bridge. This identification of $P$ is unavailable under $\pares$.
\item under $\pares$, $V_0=V$ is the perturbing field itself and $\mu$ is
its derivative at $0$, which is what makes the \emph{global} factorization
\cref{dyn:cor:factorization} and the monodromy exponents
\cref{dyn:lem:residue} available. That factorization is unavailable under
$\resq$, where $V_0$ is not the input.
\end{itemize}
Concretely, source 01's three regimes give
\begin{equation}\label{dyn:eq:phase-H0}
 H_0(X)=\begin{cases}
 X(1+X^{2})^{-1/2}, & \sigma_0<\delta/2,\\[2pt]
 X(1-X^{2})^{-1/6}(1+2X^{2})^{-1/3}, & \sigma_0=\delta/2,\\[2pt]
 X(1-2X^{4})^{-1/4}, & \sigma_0>\delta/2,
 \end{cases}
\end{equation}
with branches equal to one at zero (see \cref{dyn:thm:phase});
differentiating their logarithms verifies $XP(X)H_0'(X)=H_0(X)$ directly.
Source 04's instantiation gives $h_0=w/(1+w)$ for $V=w+w^{2}$ and
$h_0=w(1+w^{q})^{-1/q}$ for $V=w(1+w^{q})$; see
\cref{dyn:sec:parabolic-monodromy}.

\section{Nonresonant Hahn linearization need not be coherent}
\label{dyn:sec:noncoherent}

The extra conclusion of \cref{dyn:thm:coherent} cannot be imported into
\cref{dyn:lem:nonresonant}. The obstruction is ordinary small-divisor
divergence in the exponent-zero coefficient, \emph{not} a failure of Hahn
summability.

\begin{proposition}[A divergent ordinary linearizer survives on the monad]
\label{dyn:prop:noncoherent}
There is an ordinary complex polynomial $f(X)=\alpha X+X^{2}$ with $\alpha$
not a root of unity whose normalized formal conjugacy is strongly
Hahn-evaluable and isometric on $\m_\Gamma$ in \emph{every} nontrivial
workspace $K_\Gamma$, yet does not lie in the radius-free ring
$\mathcal A_1=\C\{X\}((t^{\Gamma}))$ and so has no common-domain
representative as a normalized formal conjugacy.
\end{proposition}
\begin{proof}
Choose $\alpha=e^{2\pi i\theta}$ with $\theta$ irrational and non-Brjuno
--- for instance \cref{dyn:thm:nonbrjuno}(e). By Yoccoz's quadratic theorem
the germ $\alpha X+X^{2}$ is analytically linearizable exactly when
$\theta$ is Brjuno \cite{Yoccoz}, so its unique normalized formal
linearizer $H\in\C[[X]]$ has zero ordinary radius of convergence. All its
coefficients lie in the constant field, so its Hahn support is $\{0\}$ and
\cref{dyn:lem:evaluation} makes $H$ and its formal inverse strongly
evaluable isometries of $\m_\Gamma$, the exact maximal ball being $B_0$ by
\cref{dyn:thm:exact}. But the coefficient at Hahn exponent zero is the
divergent ordinary series $H$ itself, so the family is not in
$\mathcal A_1$, let alone in the fixed-polydisk or common-domain germ
rings.
\end{proof}

\Cref{dyn:lem:faithful} rules out an evasion by change of representation: a
different common-domain family cannot represent the divergent conjugacy as
the same function on all infinitesimals. Likewise a normalized
common-domain conjugacy satisfying the same functional identity would
satisfy it formally, and uniqueness in \eqref{dyn:eq:recurrence} would
force the same divergent exponent-zero coefficient.

\begin{warning}[What \cref{dyn:prop:noncoherent} does and does not say]
\label{dyn:warn:noncoherent}
There is no conflict with Yoccoz's theorem: an ordinary divergent series can
be Hahn-summable at an infinitesimal argument because each successive power
introduces a new positive exponent. That is a change of analytic category,
not a proof of classical analytic linearizability. Moreover the coherent
conclusion of \cref{dyn:thm:coherent} is \textbf{resonant-only}, and
\cref{dyn:prop:noncoherent} establishes failure of a \emph{universal}
common-domain assertion in the nonresonant case --- it is explicitly
\textbf{not} an if-and-only-if classification, and not every nonresonant
case is noncoherent.
\end{warning}

\section{Workspace invariance and the full surcomplex class}
\label{dyn:sec:workspace}

\begin{theorem}[Workspace invariance for the finite-return package]
\label{dyn:thm:workspace}
Let $\Gamma\subseteq\Gamma'$ be divisible ordered abelian groups and view
$K_\Gamma\subset K_{\Gamma'}$. For $f\in K_\Gamma[z]$ satisfying
\cref{dyn:thm:exact}, the threshold $r$, the normalized formal conjugacy
and the shell polynomial are unchanged by the extension. The conjugacy and
its inverse strongly evaluate at every point of $B_r(K_{\Gamma'})$,
including points absent from $K_\Gamma$. The finite shell root set,
algebraic multiplicities, exact periods and multiplier formulas gain no new
points and do not change, and the common-domain coefficient functions of
\cref{dyn:thm:coherent} are also unchanged.
\end{theorem}
\begin{proof}
The embedding preserves valuation, residue and all finite coefficient
arithmetic, so the return polynomial and the finite maximum
\eqref{dyn:eq:exact-r} are identical, and \eqref{dyn:eq:recurrence} gives
the same formal coefficients in both fields. For a new point
$z\in B_r(K_{\Gamma'})$, the normalized input $z/t^{r}$ has positive
well-ordered support in $\Gamma'$, whose union with the constructed
certificate is still well ordered, so \cref{dyn:lem:evaluation} applies
unchanged. Every relevant root is a root of a polynomial over the
algebraically closed $K_\Gamma$, which already factors completely into
linear factors in $K_\Gamma[z]$, and the same finite factorization remains
valid in $K_{\Gamma'}[z]$: no additional roots, no change of multiplicity.
The remaining assertions concern the same formal coefficients, the same
ordinary residue polynomial and the same uniquely normalized differential
recursion.
\end{proof}

\begin{corollary}[Transfer to the surcomplex numbers]
\label{dyn:cor:surcomplex}
Let $f\in\SC[z]=\No[i][z]$ be a nonlinear polynomial fixing zero, with unit
multiplier for the natural Hahn valuation and that multiplier not a root of
unity. Then the formulas and conclusions of
\cref{dyn:thm:exact,dyn:thm:shell,dyn:thm:coherent,dyn:thm:stability} hold
on the corresponding class-sized ball in $\SC$, interpreted pointwise by
set-sized Hahn sums, and the threshold and finite shell cycles are
independent of the workspace used to calculate them.
\end{corollary}
\begin{proof}
Place the finitely many coefficients of $f$ in one divisible workspace and
construct its sequence of formal conjugacy coefficients there. For a
particular argument, enlarge the workspace to include it and apply
\cref{dyn:thm:workspace}; two choices agree in their common larger
workspace and so define one class function pointwise. The finite
obstructing root set already lies in the initial algebraically closed
workspace. No sum in this construction is indexed by a proper class.
\end{proof}

\Cref{dyn:cor:surcomplex} may be read in a class theory such as NBG or with
definable classes in ZFC; it requires no proper-class polynomial ring as a
set and does not replace Hahn summation by fine-topological convergence. The
analogous statements for the Hamiltonian package are
\cref{dyn:cor:hamsurcomplex}, and for the several-variable linearizer
\cref{dyn:prop:workspace}.

\section{Cancellation: a tuned cubic and a two-parameter phase diagram}
\label{dyn:sec:phase}

The exact threshold \eqref{dyn:eq:exact-r} is read off the \emph{actual}
return coefficients $b_n$, including the cancellations that occur when
$f^{\circ q}$ is composed. This subsection exhibits a family in which that
distinction is strictly consequential: any bound depending only on
$a=\max_{n\ge2}|a_n|^{1/(n-1)}$, $q$ and $|1-\lambda^{q}|$ --- the shape of
the general lower radius of the rank-one antecedent
\cite[Thm.~1.1]{Lindahl} --- throws the improvement away.

\subsection{A tangent binomial, for calibration}
For $f(z)=(1+c)z+az^{m}$ with $0\ne c\in\m_\Gamma$, $a\ne0$, $m\ge2$, one
has $q=1$ and $r=(v(c)-v(a))/(m-1)$. The nonzero fixed points are exactly
$p^{m-1}=-c/a$, all $m-1$ of them simple and on $S_r$, and at each the
multiplier is known \emph{exactly}, not merely to leading order:
\begin{equation}\label{dyn:eq:binomial-mult}
 f'(p)=1+c+map^{m-1}=1-(m-1)c .
\end{equation}
With $A=\res(a\rho^{m-1}/c)$ the shell polynomial is $P(X)=1+AX^{m-1}$ and
$H_0(X)=X(1+AX^{m-1})^{-1/(m-1)}$ on the ordinary branch normalized to one
at zero. This already shows why a degree-independent small-divisor loss can
be wasteful.

\subsection{A cubic with two suppressed return terms}
\label{dyn:subsec:tuned}
Let $\delta>0$, $e=t^{\delta}$, and
\begin{equation}\label{dyn:eq:tuned-cubic}
 f_e(z)=(-1+e)z+z^{2}-z^{3} .
\end{equation}
Its residue multiplier has order $q=2$, $\eta(f_e)=0$, and
$c=\lambda^{2}-1=e(e-2)$ has valuation $\delta$. The nonlinearity of the
second iterate is not guessed from the original quadratic term; its exact
coefficients are
\begin{center}
\begin{tabular}{c l c l}
\toprule
$n$ & $b_n$ in $f_e^{\circ2}$ & $n$ & $b_n$ in $f_e^{\circ2}$\\
\midrule
$2$ & $e(e-1)$ & $6$ & $6(e-1)$\\
$3$ & $-e(e-2)(e-1)$ & $7$ & $-3(e-2)$\\
$4$ & $-e(3e-4)$ & $8$ & $-3$\\
$5$ & $3e^{2}-9e+4$ & $9$ & $1$\\
\bottomrule
\end{tabular}
\end{center}
At $e=0$ the return displacement factors as
\begin{equation}\label{dyn:eq:tuned-factor}
 f_0^{\circ2}(z)-z=z^{5}(z^{2}-2z+2)(z^{2}-z+2),
\end{equation}
so that the cubic \emph{and} quartic terms have both disappeared; their
actual valuations, not the separate sizes of the terms that cancelled, are
what enter \eqref{dyn:eq:exact-r}. The degree-five coefficient is a unit
and the exact maximum is
\begin{equation}\label{dyn:eq:tuned-r}
 \boxed{\;r(f_e)=\delta/4 .\;}
\end{equation}
For $\rho=t^{\delta/4}$ one gets $P(X)=1-2X^{4}$: four simple closest
period-two points in two cycles, with residue directions satisfying
$a^{4}=1/2$; since $aP'(a)=-4$,
\begin{equation}\label{dyn:eq:tuned-mult}
 (f_e^{\circ2})'(p_a)=1-4c+o_v(c)=1+8e+o_v(e),
\end{equation}
where $R=o_v(c)$ means $v(R)>v(c)$, and
$H_0(X)=X(1-2X^{4})^{-1/4}$.

For comparison, the untuned quadratic $(-1+e)z+z^{2}$ has threshold
$\delta/2$ and shell polynomial $1+X^{2}$. The cubic
\eqref{dyn:eq:tuned-cubic} therefore admits an \emph{exactly larger}
linearization ball with the same multiplier and the same basic coefficient
scale $\eta=0$; the difference is produced entirely by cancellation in a
finite return.

\subsection{The two-parameter transition}
The cancellation can be broken on a second valuation scale. Let
\begin{equation}\label{dyn:eq:phase-family}
 f_{e,s}(z)=(-1+e)z+z^{2}+(-1+s)z^{3},
 \qquad e=t^{\delta},\ s=t^{\sigma_0},\ \delta,\sigma_0>0,
\end{equation}
with no Archimedean relationship assumed between the two positive
exponents.

\begin{theorem}[Cancellation phase diagram]\label{dyn:thm:phase}
For \eqref{dyn:eq:phase-family},
\begin{equation}\label{dyn:eq:phase-r}
 \boxed{\;r(f_{e,s})=\max\Bigl\{\frac{\delta-\sigma_0}{2},\frac{\delta}{4}\Bigr\},\;}
\end{equation}
and the shell polynomial, cycle count and leading return-multiplier ratios
are
\begin{center}
\small
\begin{tabular}{@{}llll@{}}
\toprule
Regime & $P(X)$ & Two-cycles & $aP'(a)$\\
\midrule
$\sigma_0<\delta/2$ & $1+X^{2}$ & one & $-2$\\
$\sigma_0=\delta/2$ & $1+X^{2}-2X^{4}=(1-X^{2})(1+2X^{2})$ & two & $-6$, $-3$\\
$\sigma_0>\delta/2$ & $1-2X^{4}$ & two & $-4$\\
\bottomrule
\end{tabular}
\end{center}
Every listed initial root is simple. In the middle row the ratio $-6$
occurs for $a^{2}=1$ and $-3$ for $a^{2}=-1/2$. The three leading coherent
coordinates are \eqref{dyn:eq:phase-H0}.
\end{theorem}
\begin{proof}
For the general cubic $f(z)=(-1+e)z+z^{2}+bz^{3}$, direct composition gives
\begin{align}
 b_2&=e(e-1), & b_3&=(e-1)(be^{2}-2be+2b+2),\nonumber\\
 b_4&=3be^{2}-4be+b+1, & b_5&=b(3be^{2}-6be+3b+3e-1),\label{dyn:eq:general-cubic}\\
 b_6&=b(6be-5b+1), & b_7&=3b^{2}(be-b+1),\nonumber\\
 b_8&=3b^{3}, & b_9&=b^{4}.\nonumber
\end{align}
Substituting $b=-1+s$, the degree-three coefficient has leading terms
$-2e-2s$ and the degree-four coefficient has leading terms $4e+s$; because
$e$ and $s$ have leading coefficient one, no cancellation occurs between
those pairs even when $\delta=\sigma_0$. Hence $v(b_2)=\delta$,
$v(b_3)=v(b_4)=\min\{\delta,\sigma_0\}$ and $v(b_n)=0$ for
$5\le n\le9$, so the maximum in \eqref{dyn:eq:exact-r} is the larger of
$(\delta-\min\{\delta,\sigma_0\})/2$ and $\delta/4$, which is
\eqref{dyn:eq:phase-r}. If $\sigma_0<\delta/2$, take
$r=(\delta-\sigma_0)/2$: only the normalized degree-three return
coefficient has nonzero residue in $Q-1$, and that residue is one, so
$P=1+X^{2}$. If $\sigma_0>\delta/2$, take $r=\delta/4$: only degree five
survives, with residue $4/(-2)=-2$, so $P=1-2X^{4}$. At equality both terms
survive and $P=1+X^{2}-2X^{4}$. All three are squarefree, so the cycle
counts follow from \cref{dyn:thm:shell}, and reduction of $XP'(X)$ at the
roots gives $-2$, $-4$ and, in the middle case,
$2a^{2}-8a^{4}=-6$ for $a^{2}=1$ and $-3$ for $a^{2}=-1/2$.
\end{proof}

\begin{example}[A repeated initial root at the transition]\label{dyn:ex:kappa}
Replace $s$ in the equality regime by $\kappa t^{\delta/2}$ with
$\kappa\in\C^{\times}$. The same calculation gives
$P_\kappa(X)=1+\kappa X^{2}-2X^{4}$, and as a polynomial in $Y=X^{2}$,
\[
 \Disc_Y(1+\kappa Y-2Y^{2})=\kappa^{2}+8 .
\]
For $\kappa^{2}+8\ne0$ the four shell points are simple in distinct residue
directions. For $\kappa^{2}+8=0$ the initial polynomial has repeated roots
--- an \emph{initial-direction degeneracy}, and not by itself an assertion
of an actual multiple periodic point (\cref{dyn:rem:multiplicity});
determining the finer splitting there requires higher Hahn terms of the
return equation. The exact threshold remains $\delta/4$, since the
degree-four term of $P_\kappa$ cannot vanish.
\end{example}

\begin{example}[A genuinely higher-rank example]\label{dyn:ex:rank2}
Let $\Gamma=\Q+\Q\Omega$ with $\Omega>n$ for every ordinary positive
integer $n$, and
\begin{equation}\label{dyn:eq:rank2}
 f(z)=(1+t^{\Omega})z+t^{3}z^{2}+z^{4}.
\end{equation}
Here $q=1$, $\eta=0$ and $r=\max\{\Omega-3,\Omega/3\}=\Omega-3$. At
$\rho=t^{\Omega-3}$ the shell polynomial is $P(X)=1+X$, and there is one
closest nonzero fixed point $p=\rho(-1+u)$, $u\in\m_\Gamma$, with
$f'(p)=1-t^{\Omega}+o_v(t^{\Omega})$.

The other two nonzero fixed points are much farther from zero: all three
solve $z^{3}+t^{3}z+t^{\Omega}=0$, and factoring off $p$ gives
$z^{3}+t^{3}z+t^{\Omega}=(z-p)(z^{2}+pz+p^{2}+t^{3})$; scaling the
quadratic factor by $z=t^{3/2}Y$ and dividing by $t^{3}$ gives reduction
$Y^{2}+1$, because $v(p)=\Omega-3$ dominates every ordinary rational.
Those two roots therefore have valuation $3/2$ and normalized residue
directions $\pm i$.

This example \emph{cannot} be faithfully described by a real-valued
valuation preserving the order of all of $\Gamma$: sending $\Omega$ to a
finite real number would fail to preserve $\Omega>n$ for all $n$. The exact
calculation uses only ordered-group comparisons and
\cref{dyn:lem:neumann}. Finally, the same threshold persists if
$t^{\Omega}$ in the linear term is replaced by
$c=t^{\Omega}\bigl(1+\sum_{j\ge2}t^{\,2-1/j}\bigr)$, a legitimate Hahn
coefficient whose positive support is well ordered although infinitely many
exponents lie below $2$ before the shift by $\Omega$: the leading valuation
remains $\Omega$ and the normalized shell polynomial remains $1+X$. The
theorem is therefore \emph{not} restricted to coefficients with discrete or
finitely generated exponent support. Both halves of this example are
specific to source 01.
\end{example}

\section{Parabolic dynamics I: algebraicity and hidden monodromy}
\label{dyn:sec:parabolic-monodromy}

Everything in this section is under $\pares$ and $\sdno$
(\cref{dyn:warn:parabolic}): the leading multiplier is exactly $1$
perturbed by an infinitesimal, and no divisor condition is imposed
anywhere. Here $\Gamma$ is divisible and set-sized --- divisibility is
load-bearing, because it is what keeps the resonant rescaling
$s=t^{\delta/q}$ inside the chosen workspace --- and $\eps=t^{\delta}$ with
$\delta>0$ is a single Hahn monomial, \textbf{not} a small nonzero complex
number. Ordinary powers of $\eps$ form the Hahn subring $\C[[\eps]]$, and
every $O(\eps^{k})$ below means \emph{divisibility} by $\eps^{k}$ in a
formal series ring, never an analytic estimate uniform in a numerical
parameter.

\subsection{The common-domain coefficient ring and the algebraicity base}
Fix an ordinary connected open $U\subset\C$ and work in
$\AAA(U)=\OO(U)((t^{\Gamma}))$, the fixed-polydisk row of
\cref{dyn:subsec:rings} in one variable: one well-ordered exponent support,
every coefficient holomorphic on that \emph{same} $U$, derivatives and
primitives coefficientwise with no estimate uniform in the exponent. A
series with nowhere-zero leading coefficient on $U$ is invertible in
$\AAA(U)$ (\cref{dyn:lem:unit}). Write $\AAA^{>0}(U)$ for the
positive-valuation part. The algebraicity base is
\begin{equation}\label{dyn:eq:base}
 \BB=\C(w)((t^{\Gamma})),
\end{equation}
Hahn series with \emph{arbitrary rational-function} coefficients.

\begin{warning}[The base field is deliberately large]\label{dyn:warn:base}
$\BB$ is \textbf{not} $K_\Gamma(w)$. Its elements may have coefficient
functions whose extra poles vary with the Hahn exponent, so there is no
common pole-free subdomain for all of $\BB$, and proving transcendence over
$\BB$ is strictly stronger than proving it over $K_\Gamma(w)$. A merge must
not weaken \eqref{dyn:eq:base} to $K_\Gamma(w)$. The mechanism that makes
the larger base available is that an ordinary deck transformation fixes
\emph{every rational coefficient function separately}, so no degree bound,
no finite support and no common bound on pole sets is needed.
\end{warning}

The ambient field for algebraicity is
\begin{equation}\label{dyn:eq:ambientfield}
 \LL=\MM(\widetilde X)((t^{\Gamma})),\qquad
 \BB\hookrightarrow\LL,
\end{equation}
where $X=\C\setminus Z(V)$, $p:\widetilde X\to X$ is its \emph{ordinary}
universal cover and rational functions are pulled back along $p$. A deck
transformation acts on each coefficient by ordinary analytic continuation
and leaves Hahn exponents fixed, hence is a strong field automorphism of
$\LL$ fixing $\BB$ pointwise. No covering-space topology on $\No[i]$ is
introduced, and no topological convergence in the surcomplex plane is used.

\begin{lemma}[A sufficient unit criterion]\label{dyn:lem:unit}
If $f=t^{\lambda_0}(a+u)$ with $a\in\OO(U)$ nowhere zero and
$u\in\AAA^{>0}(U)$, then $f$ is invertible in $\AAA(U)$ with
$f^{-1}=t^{-\lambda_0}a^{-1}\sum_{n\ge0}(-u/a)^{n}$, and every coefficient
of the inverse is holomorphic on the \emph{original} $U$.
\end{lemma}
\begin{proof}
The family is summable by \cref{dyn:lem:neumann}, multiplication by
$1+u/a$ telescopes coefficientwise to $1$, and powers of $a^{-1}$ are
holomorphic on $U$, so no new singularity appears.
\end{proof}

No converse characterization of the units of $\AAA(U)$ is claimed; only the
stated criterion is used. Every nonzero series is invertible in
$\MM(U)((t^{\Gamma}))$.

\subsection{Common-domain Hahn linearization}

\begin{theorem}[Common-domain Hahn linearization]\label{dyn:thm:common}
Let $U\subset\C$ be simply connected with $0\in U$, let $V\in\OO(U)$
satisfy $V(0)=0$, $V'(0)=\mu\ne0$ and $V(w)\ne0$ on $U\setminus\{0\}$, and
let $E\in\AAA^{>0}(U)$ satisfy $E(0)=0$. Put
\begin{equation}\label{dyn:eq:generalF}
 F(w)=w+t^{\delta}\bigl(V(w)+E(w)\bigr),\quad \lambda=F'(0),\quad
 \mathsf M=\bigl\langle\{\delta\}\cup\supp E\bigr\rangle .
\end{equation}
There is a unique $H\in\AAA(U)$ with
\begin{equation}\label{dyn:eq:generalH}
 H\circ F=\lambda H,\qquad H(0)=0,\qquad H'(0)=1,
\end{equation}
its support is contained in $\mathsf M$, every coefficient is holomorphic
on $U$, and its coefficient at exponent zero is $h_0$ of
\eqref{dyn:eq:leadingH}. All constructions commute with enlargement of the
set-sized value group.
\end{theorem}
\begin{proof}
\Cref{dyn:thm:exp-log} gives $\log C_F=A\partial_w$ with
$A\in\AAA^{>0}(U)$. Every word in $\Delta=t^{\delta}(V+E)$ has an exponent
which, after removing one initial $\delta$, lies in $\mathsf M$, so
$A=t^{\delta}W$ with $W=V+W_+$, $\supp W\subseteq\mathsf M$, $W_+$ of
positive valuation and all coefficients holomorphic on $U$. Put
$c=\Log\lambda$ and $\alpha_{\mathrm c}=t^{-\delta}c$;
\cref{dyn:lem:fixedderivative} gives $W(0)=0$ and
$W'(0)=\alpha_{\mathrm c}$, whose leading coefficient is $\mu$.
Coefficientwise division by $w$ is possible since $W(0)=0$, and the leading
coefficient of $W/w$ is $V/w$, nowhere zero on $U$ including at $0$, so
$W/w$ is invertible in $\AAA(U)$ by \cref{dyn:lem:unit}. Hence
\begin{equation}\label{dyn:eq:regularR}
 R(w):=\frac{\alpha_{\mathrm c}}{W(w)}-\frac1w
 =\frac{\alpha_{\mathrm c}/(W(w)/w)-1}{w}
\end{equation}
belongs to $\AAA(U)$, its numerator vanishing coefficientwise at $0$
because $(W/w)(0)=\alpha_{\mathrm c}$, with support in $\mathsf M$. Every
coefficient of $R$ has a unique primitive vanishing at $0$ since $U$ is
simply connected; write $B(w)=\int_0^{w}R$, $B=B_0+B_+$ with $B_0$ the
coefficient at exponent zero, and set
\begin{equation}\label{dyn:eq:Hconstruction}
 H(w)=we^{B_0(w)}\sum_{n\ge0}\frac{B_+(w)^{n}}{n!},
\end{equation}
whose exponential is certified by \cref{dyn:lem:neumann} and preserves the
support bound $\mathsf M$, using the same $U$ at every stage. Then
$H(0)=0$, $H'(0)=1$ and $H'/H=\alpha_{\mathrm c}/W$ away from $0$, with the
cross-multiplied identity everywhere, so $DH=cH$; as $D$ fixes Hahn
constants, exponentiation gives $C_FH=\ExpH(c)H=\lambda H$, and the
exponent-zero part of \eqref{dyn:eq:Hconstruction} is
\eqref{dyn:eq:leadingH}. For uniqueness, if $\widetilde H$ satisfies
\eqref{dyn:eq:generalH} then $N^{r}\widetilde H=(\lambda-1)^{r}\widetilde H$
so $D\widetilde H=c\widetilde H$; the quotient
$\widetilde H/H=(\widetilde H/w)/(H/w)$ lies in $\AAA(U)$ because the
leading coefficient of $H/w$ is the nowhere-zero $e^{B_0}$, the
differential equations force its derivative to vanish ($\AAA(U)$ being an
integral domain for connected $U$), every coefficient is therefore
constant, and its value at $0$ is $1$. Finally all outputs are formed from
input supports by finite words, coefficientwise differentiation and
uniquely normalized integration, none of which depends on the ambient group
beyond containing the indicated supports.
\end{proof}

\begin{warning}[What kind of embedding was proved]\label{dyn:warn:embedding}
\Cref{dyn:thm:common} constructs a strong logarithm and solves a
coefficientwise ordinary differential equation. It does \textbf{not} assert
convergence of the series at a nonzero ordinary complex value of
$t^{\delta}$, nor analytic embedding of every numerical near-identity map
in an autonomous flow; those are different questions, and the
Gelfreich--Vieiro reference is cited precisely to mark that line
\cite{GV}. Nor does it assert that the global evaluated map $H^{\#}$ is
injective on all of $U^{\#}_\Gamma$; injectivity is proved only on the
infinitesimal disk of the resonant model (\cref{dyn:thm:sharp}). There is
also no ordinary invariant-domain hypothesis: $C_F$ denotes the Taylor
operation \eqref{dyn:eq:pullback}, not a numerical self-map of $U$, and at
a point $a+\eta\in U^{\#}_\Gamma$ the positive displacement preserves the
standard part.
\end{warning}

\begin{corollary}[Universal leading profile at a resonant scale]
\label{dyn:cor:universal}
Let $q\ge1$, $a\in\C^{*}$, and let
$\phi(z)=az^{q+1}+\sum_{n>q+1}a_nz^{n}$ be an ordinary convergent germ.
Choose $c_0\in\C^{*}$ with $ac_0^{q}=1$, set $s=c_0t^{\delta/q}$, and let
$f(z)=(1+\eps)z+\phi^{\#}(z)$ at infinitesimal $z$. After $z=sw$, the
normalized Hahn coordinate on any simply connected ordinary $U$ containing
$0$ but no root of $1+w^{q}$ has leading coefficient
\begin{equation}\label{dyn:eq:universalprofile}
 w(1+w^{q})^{-1/q},
\end{equation}
with the branch normalized at $0$.
\end{corollary}
\begin{proof}
Taylor evaluation of the convergent germ at $sw$ is strongly meaningful for
every finite $w$, and direct rescaling gives
\[
 s^{-1}f(sw)=w+\eps\Bigl(w+w^{q+1}
 +\sum_{n>q+1}a_nc_0^{\,n-1}t^{\delta(n-q-1)/q}w^{n}\Bigr),
\]
whose final sum has positive Hahn support and polynomial coefficients on
such a $U$. Apply \cref{dyn:thm:common} with $V=w(1+w^{q})$; integration of
$1/V-1/w$ gives $-q^{-1}\log(1+w^{q})$.
\end{proof}

\Cref{dyn:cor:universal} asserts a leading profile and a common-domain
construction. It explicitly does \textbf{not} extend the exact monodromy
formula or the sharp disk theorem below to arbitrary higher-order
perturbations.

\subsection{Polynomial Euler maps and their exact logarithmic residues}
Specialize to
\begin{equation}\label{dyn:eq:Euler}
 G_\eps(w)=w+\eps V(w),\qquad V\in\C[w],\quad V(0)=0,\quad V'(0)=\mu\ne0,
\end{equation}
so \cref{dyn:thm:common} with $E=0$ constructs
$H_\eps\in\OO(U)[[\eps]]$. The polynomial hypothesis supplies structure the
general theorem does not have.

\begin{proposition}[The generator gains no new pole]\label{dyn:prop:factorV}
The logarithmic generator of \eqref{dyn:eq:Euler} has the form
\begin{equation}\label{dyn:eq:Wfactor}
 \log C_{G_\eps}=\eps W_\eps(w)\partial_w,\qquad
 W_\eps=VB_\eps,\qquad B_\eps\in1+\eps\C[w][[\eps]] .
\end{equation}
Setting
$\alpha_\eps=\Log(1+\mu\eps)/\eps$ and
$\Omega_\eps=\alpha_\eps/W_\eps$, every coefficient of $\Omega_\eps$ is a
rational function of $w$ with poles \emph{only} at zeros of $V$, of order
at most the multiplicity in $V$.
\end{proposition}
\begin{proof}
For a polynomial $p$,
$(C_{G_\eps}-I)p=\sum_{k\ge1}\eps^{k}V^{k}p^{(k)}/k!$, every coefficient
of which is polynomial and divisible by $V$; the same holds of each
$N^{r}w$, and at a fixed power of $\eps$ only finitely many $r$ occur, so
the logarithm has polynomial coefficients all divisible by $V$, its first
term being $\eps V$. The inverse of $B_\eps$ has polynomial coefficients by
the formal geometric series, so $\Omega_\eps$ is a formal series of
polynomials divided by the single polynomial $V$.
\end{proof}

The divisor assertion is stronger than rationality of each coefficient:
higher terms of the modified vector field cannot introduce \emph{new}
spatial pole locations into the logarithmic derivative of the coordinate.
Finite expansion of the operator logarithm gives, for the record,
\begin{equation}\label{dyn:eq:modifiedvector}
 W_\eps=V-\frac\eps2VV'
 +\eps^{2}\Bigl(\frac1{12}V^{2}V''+\frac13V(V')^{2}\Bigr)+O(\eps^{3}),
\end{equation}
obtained from $Nw=\eps V$,
$N^{2}w=\eps^{2}VV'+\frac{\eps^{3}}2V^{2}V''+O(\eps^{4})$,
$N^{3}w=\eps^{3}(V(V')^{2}+V^{2}V'')+O(\eps^{4})$ substituted in
$N-N^{2}/2+N^{3}/3$. \textbf{No assertion about all coefficients is
inferred from this truncation}; the exact identity
\cref{dyn:lem:residue} is proved separately.

\begin{lemma}[Exact residue at a simple fixed point]\label{dyn:lem:residue}
At a simple zero $a$ of $V$, with $\nu_a=V'(a)$,
\begin{equation}\label{dyn:eq:exactresidue}
 \Res_{w=a}\bigl(\Omega_\eps(w)\dd w\bigr)=\rho_a(\eps)
 =\frac{\Log(1+\mu\eps)}{\Log(1+\nu_a\eps)},\qquad \rho_0=1 .
\end{equation}
\end{lemma}
\begin{proof}
\Cref{dyn:lem:fixedderivative} at $a$ gives
$\eps W_\eps'(a)=\Log G_\eps'(a)=\Log(1+\nu_a\eps)$. Since $a$ is a simple
zero of $V$, the factorization $W_\eps=VB_\eps$ makes it a simple zero over
$\C[[\eps]]$ as well, so the residue of $\alpha_\eps/W_\eps$ is
$\alpha_\eps/W_\eps'(a)$.
\end{proof}

These are \emph{coefficientwise ordinary} residues of a rational
differential. They are not local residues at displaced surcomplex poles;
the fixed points of this polynomial family are exactly the ordinary zeros
of $V$, so no displacement ambiguity arises.

\begin{proposition}[The first correction is an explicit logarithm]
\label{dyn:prop:firstcorrection}
On the common domain $U$,
\begin{equation}\label{dyn:eq:firstcorrection}
 H_\eps=h_0\left[1+\frac{\mu\eps}{2}
 \log\Bigl(\frac{V}{\mu h_0}\Bigr)+O(\eps^{2})\right],
\end{equation}
the quotient $V/(\mu h_0)$ being holomorphic, nowhere zero and equal to $1$
at $0$ after removing the common zero, so that its logarithm is uniquely
fixed by vanishing at $0$.
\end{proposition}
\begin{proof}
Since $\alpha_\eps=\mu-\mu^{2}\eps/2+O(\eps^{2})$,
\eqref{dyn:eq:modifiedvector} gives
$\Omega_\eps=\mu/V+\frac{\mu\eps}2(V'-\mu)/V+O(\eps^{2})$, whose
coefficient at $\eps$ is holomorphic at $0$ and equals
$\frac\mu2\bigl(\log(V/(\mu h_0))\bigr)'$ because $h_0'/h_0=\mu/V$.
Integrate with the normalization at $0$ and exponentiate.
\end{proof}

\begin{corollary}[Separation of the branching factors]
\label{dyn:cor:factorization}
Suppose $V$ has degree $d\ge2$ with distinct zeros $0,a_1,\dots,a_{d-1}$.
Partial fractions taken separately at each $\eps$ coefficient give
\begin{equation}\label{dyn:eq:partialfractions}
 \Omega_\eps(w)=Q_\eps(w)+\frac1w+\sum_{j=1}^{d-1}
 \frac{\rho_{a_j}(\eps)}{w-a_j},\qquad Q_\eps\in\eps^{2}\C[w][[\eps]],
\end{equation}
and with $P_\eps(w)=\int_0^{w}Q_\eps$, on branches normalized near $0$,
\begin{equation}\label{dyn:eq:fullfactorization}
 H_\eps(w)=w\ExpH(P_\eps(w))\prod_{j=1}^{d-1}(1-w/a_j)^{\rho_{a_j}(\eps)},
 \qquad P_\eps\in\eps^{2}\C[w][[\eps]],\ P_\eps(0)=0 .
\end{equation}
The factor $\ExpH(P_\eps)$ and its inverse have polynomial coefficients and
belong to $\BB$; in particular \emph{all} nontrivial coefficientwise
monodromy is carried by the displayed power factors.
\end{corollary}
\begin{proof}
The coefficients at $\eps^{0}$ and $\eps^{1}$ have numerator degree smaller
than $d$ by \eqref{dyn:eq:modifiedvector} and
\cref{dyn:prop:firstcorrection}, so $Q_\eps\in\eps^{2}\C[w][[\eps]]$. For
$\rho=\rho_0+\rho_+$ define the power on a chosen ordinary logarithm branch
by $(1-w/a)^{\rho}=\exp(\rho_0\log(1-w/a))\ExpH(\rho_+\log(1-w/a))$. The
logarithmic derivative of the right side of
\eqref{dyn:eq:fullfactorization} is \eqref{dyn:eq:partialfractions} and its
value and derivative at $0$ have the required normalization, so uniqueness
in \cref{dyn:thm:common} proves the identity; at a fixed power of $\eps$
the exponential of $P_\eps$ uses only finitely many products of
polynomials.
\end{proof}

\subsection{The torsion test and the finite-cover obstruction}

\begin{lemma}[An infinitesimal phase has no torsion]\label{dyn:lem:torsion}
For a finite Hahn constant $c=c_0+c_+$ with $c_0\in\C$ and $v(c_+)>0$,
\[
 \ExpH(c)=1\iff c_+=0\ \text{and}\ c_0\in2\pi i\Z .
\]
Consequently, for $\rho\in\C+\m_\Gamma$, the element $\ExpH(2\pi i\rho)$
has finite multiplicative order if and only if $\rho$ is an ordinary
\emph{rational} number.
\end{lemma}
\begin{proof}
The coefficient at exponent zero in $\ExpH(c)$ is $e^{c_0}$, so equality to
$1$ first requires $c_0\in2\pi i\Z$; if $c_+\ne0$ then
$\ExpH(c_+)-1=c_++c_+^{2}/2+\cdots$ has the same leading term as $c_+$, all
higher powers having strictly larger valuation, so it cannot vanish. The
last assertion follows from
$\ExpH(2\pi i\rho)^{m}=\ExpH(2\pi im\rho)$.
\end{proof}

This single lemma converts ``the exponent has a nonzero infinitesimal
part'' into ``no finite branched cover works''.

\begin{theorem}[Exact monodromy and obstruction to finite ramification]
\label{dyn:thm:monodromy}
At every simple zero $a$ of $V$, continuation of $H_\eps$ around a positive
simple loop about $a$ is multiplication by
\begin{equation}\label{dyn:eq:monodromy}
 M_a=\ExpH(2\pi i\rho_a),\qquad
 \rho_a=\frac{\Log(1+\mu\eps)}{\Log(1+V'(a)\eps)} .
\end{equation}
The multiplier has finite order exactly when $V'(a)=\mu$, in which case it
is $1$. When $V'(a)\ne\mu$, \emph{no} finite ordinary branched cover at $a$
makes the continued coordinate single-valued.
\end{theorem}
\begin{proof}
Near $a$ the rational-coefficient differential decomposes as
$\Omega_\eps\dd w=\rho_a\dd w/(w-a)+$ a holomorphic-coefficient
differential; a positive circuit changes a coefficientwise primitive by
$2\pi i\rho_a$, so exponentiation multiplies $H_\eps$ by $M_a$. Writing
$\nu=V'(a)$, formal division of the two logarithms gives
\begin{equation}\label{dyn:eq:rhofirst}
 \rho_a=\frac\mu\nu+\frac{\mu(\nu-\mu)}{2\nu}\eps+O(\eps^{2}),
\end{equation}
so if $\nu\ne\mu$ the coefficient of $\eps$ is nonzero and
\cref{dyn:lem:torsion} makes $M_a$ of infinite order, while if $\nu=\mu$
the ratio is exactly $1$. A finite branched cover has a positive integer
ramification index $m$ at a point above $a$, and in suitable local
coordinates a circuit upstairs projects to $m$ circuits downstairs, so its
monodromy multiplier is $M_a^{m}\ne1$ in the infinite-order case. This
applies in particular to the covers used to make an ordinary algebraic
function single-valued.
\end{proof}

The obstruction is already visible at order $\eps^{1}$ whenever
$\nu\ne\mu$: it is not necessary to inspect infinitely many unknown
correction terms to decide whether finite ramification will work.

\begin{lemma}[A nonlinear polynomial cannot have one slope at all its
zeros]\label{dyn:lem:slopes}
If a polynomial of degree $d\ge2$ has $d$ distinct zeros then
$\sum_{V(a)=0}1/V'(a)=0$. In particular, if $V'(0)=\mu\ne0$, not all slopes
at its zeros can equal $\mu$.
\end{lemma}
\begin{proof}
Partial fractions give $1/V(w)=\sum_a(V'(a)(w-a))^{-1}$; the coefficient of
$1/w$ at infinity on the left is zero because $d\ge2$, and on the right it
is the displayed sum. If every derivative were $\mu$ the sum would be
$d/\mu\ne0$.
\end{proof}

\subsection{The Hahn-algebraicity dichotomy}

\begin{lemma}[Reduction of an algebraic Hahn element]\label{dyn:lem:reduction}
If $Y\in\LL$ has valuation zero with leading coefficient $y_0$ and $Y$ is
algebraic over $\BB$, then $y_0$ is algebraic over $\C(w)$.
\end{lemma}
\begin{proof}
Take a nonzero relation $\sum_{j=0}^{m}b_jY^{j}=0$ with $b_j\in\BB$, let
$\gamma$ be the least valuation among the nonzero $b_j$, and multiply by
$t^{-\gamma}$: all coefficients then have nonnegative valuation and at
least one has nonzero coefficient at zero, and taking the coefficient at
zero yields a nonzero polynomial over $\C(w)$ annihilating $y_0$ ---
multiplication reduces correctly because all series in the normalized
relation have nonnegative valuation.
\end{proof}

\begin{lemma}[A higher-order logarithmic pole forces transcendence]
\label{dyn:lem:essential}
Let $r$ be a rational function with a pole of order at least two at $a$.
Any nonzero local solution of $y'/y=r$ is transcendental over $\C(w)$.
\end{lemma}
\begin{proof}
On a small punctured sector, integration gives
$y(w)=u(w)(w-a)^{\beta}\exp\bigl(\sum_{j=1}^{k-1}c_j(w-a)^{-j}\bigr)$ with
$c_{k-1}\ne0$, $u$ holomorphic and nonzero and $k\ge2$ the pole order.
Along a suitable ray ending at $a$ the real part of the leading expression
$c_{k-1}(w-a)^{-(k-1)}$ is positive and dominates the lower terms, so
$|y(w)|$ grows faster than every power of $|w-a|^{-1}$ along that ray. An
algebraic function cannot: after dividing a polynomial relation by its
leading coefficient, its coefficients are meromorphic and hence bounded by
a fixed power of $|w-a|^{-1}$ on a small punctured disk, and the elementary
bound $|y|\le1+\max_{j<m}|b_j|$ for a root of $y^{m}+\sum_{j<m}b_jy^{j}$
gives a polynomial growth bound for every root.
\end{proof}

\begin{theorem}[Polynomial Euler maps: the algebraicity dichotomy]
\label{dyn:thm:dichotomy}
Let $V\in\C[w]$ with $V(0)=0$, $V'(0)=\mu\ne0$. On every simply connected
ordinary $U$ containing $0$ and no other zero of $V$, equation
\eqref{dyn:eq:schroederintro} below has a unique solution
$H_\eps\in\OO(U)[[\eps]]$, whose coefficients continue to the universal
cover of $X=\C\setminus Z(V)$, and
\begin{equation}\label{dyn:eq:dichotomy}
 H_\eps\ \text{is algebraic over }\BB=\C(w)((t^{\Gamma}))
 \iff \deg V=1 ,
\end{equation}
where
\begin{equation}\label{dyn:eq:schroederintro}
 H_\eps\bigl(w+\eps V(w)\bigr)=(1+\mu\eps)H_\eps(w),
 \qquad H_\eps(0)=0,\quad H_\eps'(0)=1,
\end{equation}
composition being Taylor substitution in the positive Hahn displacement
$\eps V(w)$. At a simple zero $a$ of $V$, positive continuation once around
$a$ multiplies $H_\eps$ by $M_a$ of \eqref{dyn:eq:monodromy}. If $V$ is
nonlinear and all its zeros are simple, at least one zero has
infinite-order monodromy, which no finite branched cover removes.
\end{theorem}
\begin{proof}
Existence, uniqueness and the common-domain assertion are
\cref{dyn:thm:common} with $E=0$; continuation and the monodromy formula
are \cref{dyn:prop:factorV,dyn:lem:residue,dyn:thm:monodromy}. If $V$ is
linear it is $\mu w$ and $H_\eps=w$ solves the equation, an element of
$\BB$. Assume $\deg V\ge2$.

\emph{Route 1: a growth bound (needed for multiple zeros).} If $V$ has a
multiple zero then $\mu/V$ has a pole of order at least two there, and
$h_0'/h_0=\mu/V$, so \cref{dyn:lem:essential} makes $h_0$ transcendental
over $\C(w)$; since $H_\eps$ has valuation zero in $\LL$,
\cref{dyn:lem:reduction} excludes algebraicity over $\BB$.

\emph{Route 2: an infinite monodromy orbit (needed for simple zeros).} If
every zero is simple, \cref{dyn:lem:slopes} supplies a zero $a$ with
$V'(a)\ne\mu$. Let $\tau_a$ be the corresponding deck action on $\LL$; it
fixes $\BB$ and satisfies $\tau_a^{n}(H_\eps)=M_a^{n}H_\eps$ for $n\ge0$.
By \cref{dyn:thm:monodromy} these elements are all distinct, $H_\eps\ne0$
and $M_a$ of infinite order; if $H_\eps$ were algebraic over $\BB$ they
would all be roots of one nonzero polynomial over $\BB$, impossible in a
field.
\end{proof}

\begin{remark}[Two routes, neither substituting for the other]
\label{dyn:rem:tworoutes-alg}
Source 04 states explicitly that its two proof branches --- the elementary
growth bound for algebraic functions, and the infinite-monodromy-orbit
argument --- meet in the resonant model but do not logically substitute for
one another. Both are therefore kept in \cref{dyn:thm:dichotomy}: route 1
is what covers multiple zeros, where there is no simple-pole residue to
work with, and route 2 is what covers simple zeros, where $\mu/V$ has only
first-order poles and \cref{dyn:lem:essential} does not apply. A proof
discussing only simple-pole residues would miss
\cref{dyn:ex:multiplezero} entirely, even though the dichotomy remains
true there.
\end{remark}

The argument works over the larger base $\BB$ precisely because ordinary
deck transformations fix \emph{every rational coefficient function
separately} (\cref{dyn:warn:base}).

\begin{proposition}[Classification of algebraic leading coordinates]
\label{dyn:prop:leadingalgebraic}
For a nonlinear polynomial $V$ as above, $h_0$ is algebraic over $\C(w)$
if and only if all zeros of $V$ are simple and
\begin{equation}\label{dyn:eq:rationalresidues}
 \frac{\mu}{V'(a)}\in\Q\qquad\text{for every zero }a\text{ of }V ,
\end{equation}
in which case
\begin{equation}\label{dyn:eq:h0product}
 h_0(w)=w\prod_{a\ne0}(1-w/a)^{\mu/V'(a)} .
\end{equation}
Whenever these conditions hold, the first correction $h_1$ in
$H_\eps=h_0+\eps h_1+\cdots$ is \emph{already} transcendental over
$\C(w)$.
\end{proposition}
\begin{proof}
Multiple zeros are excluded by \cref{dyn:lem:essential}. With all zeros
simple, the leading part of \eqref{dyn:eq:fullfactorization} is
\eqref{dyn:eq:h0product}; if all exponents are rational, taking a common
denominator gives a polynomial relation, so $h_0$ is algebraic. Conversely
continuation of an algebraic function gives a finite orbit, so the local
multipliers of \eqref{dyn:eq:h0product} must have finite order, which for
an ordinary complex exponent is rationality. For the last assertion choose
$a$ with $\nu=V'(a)\ne\mu$ by \cref{dyn:lem:slopes}, put
$r_0=\mu/\nu\in\Q$ and a positive integer $m$ with $mr_0\in\Z$; after $m$
circuits about $a$ the function $h_0$ returns to itself, while by
\cref{dyn:prop:firstcorrection} the same continuation changes $h_1$ by
$2\pi im\frac\mu2(1-r_0)h_0\ne0$. Repeated blocks of $m$ circuits therefore
produce infinitely many distinct branches of $h_1$, and an algebraic
function has only finitely many.
\end{proof}

\begin{example}[Rational leading coordinate, logarithmic correction]
\label{dyn:ex:rational}
For $V=w+w^{2}$, $h_0=w/(1+w)$ and $h_1=\frac{w}{1+w}\log(1+w)$. At the
fixed point $-1$ the exact exponent is
$\Log(1+\eps)/\Log(1-\eps)=-1+\eps+O(\eps^{2})$, so the leading multiplier
is $1$ while the full multiplier begins $1+2\pi i\eps+O(\eps^{2})$: the
obstruction is present even without any branching in the leading function.
A rational leading coordinate does not even predict the \emph{branching
type} of the first correction.
\end{example}

\begin{example}[One trivial nonzero-root multiplier]\label{dyn:ex:trivialroot}
For $V=w(1-w)(1-2w)$ one has $\mu=1$, $V'(1)=1$, $V'(1/2)=-1/2$, and
\[
 h_0=\frac{w(1-w)}{(1-2w)^{2}},\qquad h_1=\tfrac32h_0\log(1-2w).
\]
The exact multiplier at $1$ is $1$, whereas at $1/2$ it is
$1+3\pi i\eps+O(\eps^{2})$ and has infinite order. Thus
\cref{dyn:thm:monodromy} does \emph{not} claim branching at every nonzero
fixed point; \cref{dyn:lem:slopes} guarantees at least one obstruction,
which suffices for transcendence over the whole base.
\end{example}

\begin{example}[Why multiple zeros need a separate argument]
\label{dyn:ex:multiplezero}
For $V=w(1+w)^{2}$, $\mu=1$ and
$h_0(w)=\frac{w}{1+w}\exp\bigl(\frac1{1+w}-1\bigr)$. The pole of $\mu/V$ at
$-1$ has order two, so the leading solution has an essential singularity
and is already transcendental --- \cref{dyn:rem:tworoutes-alg}.
\end{example}

\subsection{The resonant model: exact scale, isometry, continuation}
Fix an ordinary integer $q\ge1$ and let
$f(z)=(1+\eps)z+z^{q+1}$ with $\eps=t^{\delta}$, $s=t^{\delta/q}$
(divisibility of $\Gamma$ placing $s$ in $K_\Gamma$). Under $z=sw$,
\begin{equation}\label{dyn:eq:gmodel}
 g(w)=s^{-1}f(sw)=(1+\eps)w+\eps w^{q+1}=w+\eps V_q(w),
 \qquad V_q(w)=w(1+w^{q}).
\end{equation}
Let $\Psi_\eps$ be its normalized common-domain coordinate; $U=\{|w|<1\}$
is one admissible domain, and larger simply connected slit domains avoiding
the roots of $1+w^{q}$ are also possible. Since $g(\zeta w)=\zeta g(w)$ for
a $q$-th root of unity $\zeta$, uniqueness gives
$\Psi_\eps(\zeta w)=\zeta\Psi_\eps(w)$ and hence
\begin{equation}\label{dyn:eq:bseries}
 \Psi_\eps(w)=\sum_{m\ge0}b_m(\eps)w^{qm+1},\qquad b_0=1,\ b_m\in\C[[\eps]],
\end{equation}
initially as the Taylor series of each ordinary coefficient function; its
evaluation as a Hahn family at a specific $w$ is a separate question. The
leading solution is
\begin{equation}\label{dyn:eq:leadingmodel}
 \Psi_0(w)=w(1+w^{q})^{-1/q},\qquad
 b_m(0)=(-1)^{m}\frac{(1/q)_m}{m!}\ne0 ,
\end{equation}
with $(x)_m=x(x+1)\cdots(x+m-1)$.

\begin{proposition}[Exact finite recurrence]\label{dyn:prop:recurrence}
With $\lambda=1+\eps$, for $m\ge1$
\begin{equation}\label{dyn:eq:brecurrence}
 b_m=-\frac{\displaystyle\sum_{\substack{0\le j<m\\ m-j\le qj+1}}
 b_j\binom{qj+1}{m-j}\lambda^{qj+1-(m-j)}\eps^{m-j}}
 {\lambda^{qm+1}-\lambda},
\end{equation}
which uniquely constructs $b_m\in\C[[\eps]]$ and has leading recurrence
$qm\,b_m(0)=-\bigl(q(m-1)+1\bigr)b_{m-1}(0)$.
\end{proposition}
\begin{proof}
In $g(w)^{qj+1}=w^{qj+1}(\lambda+\eps w^{q})^{qj+1}$ the coefficient of
$w^{qm+1}$ is the binomial expression displayed; the term $j=m$ contributes
$\lambda^{qm+1}b_m$ and the right side of
$\Psi_\eps\circ g=\lambda\Psi_\eps$ contributes $\lambda b_m$. The
denominator is $qm\eps+O(\eps^{2})$ while every numerator term is divisible
by $\eps$, so induction gives $b_m\in\C[[\eps]]$; only $j=m-1$ survives
after division by $\eps$ and reduction modulo $\eps$, which also checks
\eqref{dyn:eq:leadingmodel} directly.
\end{proof}

The normalized formal coordinate of the original map is
\begin{equation}\label{dyn:eq:originalh}
 h(z)=s\Psi_\eps(z/s)=\sum_{m\ge0}\eps^{-m}b_m(\eps)z^{qm+1},
\end{equation}
a formal power series in $z$ over $K_\Gamma$ --- \emph{not} a Hahn family
already asserted summable at every $z$ --- whose nonzero coefficients have
the exact valuations $v(\eps^{-m}b_m(\eps))=-m\delta$.

\begin{theorem}[Sharp domain and valuative isometry]\label{dyn:thm:sharp}
The family of Taylor terms in \eqref{dyn:eq:originalh} is strongly summable
at $z\in K_\Gamma$ if and only if $z=0$ or
\begin{equation}\label{dyn:eq:sharpdomain}
 v(z)>\delta/q .
\end{equation}
On $D=\{v(z)>\delta/q\}$ the resulting function is a bijection
$h:D\to D$ satisfying $h\circ f=\lambda h$ and
\begin{equation}\label{dyn:eq:isometry}
 v\bigl(h(x)-h(y)\bigr)=v(x-y)\qquad(x,y\in D).
\end{equation}
No strictly larger centered valuative ball containing $D$ admits an
injective zero-fixing conjugacy to multiplication by $\lambda$.
\end{theorem}
\begin{proof}
Let $z\ne0$, $\beta=v(z)$: the $m$-th nonzero Taylor term has valuation
exactly
\begin{equation}\label{dyn:eq:termvaluation}
 v\bigl(\eps^{-m}b_mz^{qm+1}\bigr)=\beta+m(q\beta-\delta).
\end{equation}
\textbf{Two distinct failure mechanisms.} If $q\beta=\delta$, every term
has a nonzero coefficient at the \emph{same} exponent $\beta$, so
\emph{local finiteness} fails. If $q\beta<\delta$, the displayed leading
exponents form a strictly descending sequence in the union of supports, so
\emph{well-ordering} fails. This proves necessity with no convergence
terminology. If $q\beta>\delta$, put $w=z/s$ with $v(w)>0$: all
coefficients $b_m$ have nonnegative $\eps$-support, and a term of spatial
degree $qm+1$ contributes a word of at least that length from $\supp w$
together with a nonnegative number of $\delta$'s, so
\cref{dyn:lem:neumann} makes the double family strongly summable, and
multiplication by $s$ preserves this.

For invertibility, formal reversion of $\Psi_\eps=w+O(w^{q+1})$ gives
$\Theta_\eps\in\C[[\eps]][[w]]$ with linear coefficient $1$ and no constant
term, requiring no division except by a unit; the same positive-word
argument evaluates it at every $w\in\m_\Gamma$, and the formal identities
become strong identities there, both maps preserving $\m_\Gamma$ because
their difference from the identity has strictly greater valuation at a
nonzero infinitesimal input. For the isometry put $u=x/s$, $v_1=y/s$:
\[
 \Psi_\eps(u)-\Psi_\eps(v_1)=(u-v_1)
 \Bigl[1+\sum_{m\ge1}b_m(\eps)\sum_{j=0}^{qm}u^{qm-j}v_1^{j}\Bigr],
\]
whose bracket is a strong sum equal to $1$ plus a positive-valuation
element, each nonconstant summand having valuation at least
$q\min\{v(u),v(v_1)\}>0$ when both inputs are nonzero (a zero input is
handled with the vanishing terms omitted). Finally every nonzero solution
of $z^{q}=-\eps$ is a fixed point of $f$ of valuation $\delta/q$, lying in
$K_\Gamma$ because its constant factor is a complex $q$-th root of $-1$;
any strictly larger centered ball containing $D$ contains all elements at
valuation $\delta/q$, hence such a fixed point in addition to $0$, and an
injective conjugacy $k$ would satisfy $k(z)=\lambda k(z)$ at every fixed
point, mapping both to $0$ since $\lambda\ne1$.
\end{proof}

\begin{warning}[The maximality statement is about centered balls]
\label{dyn:warn:sharp}
\Cref{dyn:thm:sharp} concerns centered valuative balls defined by a
valuation threshold, under either boundary convention. It does \textbf{not}
exclude continuation to a non-ball domain avoiding the obstructing fixed
points, and does not assert that every coefficientwise continuation is
injective. Those distinctions matter at the critical scale, which is
exactly where \cref{dyn:subsec:boundary} operates.
\end{warning}

\subsection{Continuation at a scale where the Taylor family fails}
\label{dyn:subsec:boundary}

Let $w_0\in U$ be an ordinary nonzero point. At $z=sw_0$ we have
$v(z)=\delta/q$, so the Taylor family \eqref{dyn:eq:originalh} is
\emph{not} strongly summable. On the other hand
\begin{equation}\label{dyn:eq:resummedvalue}
 s\Psi_\eps(w_0)=s\sum_{n\ge0}\eps^{n}\Psi_n(w_0)
\end{equation}
\emph{is} a valid Hahn value, and more generally it is defined at
$z=s(w_0+\eta)$ for positive $\eta$ by
\eqref{dyn:eq:halo-evaluation}. There is no contradiction, and the
distinction is exactly where the monodromy lives: in
\eqref{dyn:eq:resummedvalue} each coefficient function $\Psi_n$ is first an
ordinary holomorphic function obtained by ordinary integration and
continuation, and summation in the Hahn exponent comes \emph{afterward};
in \eqref{dyn:eq:originalh} summation in spatial Taylor degree is required
to be strongly summable as a family of Hahn numbers, and it fails precisely
on this boundary. The boundary value is a \emph{coefficientwise analytic
continuation}, not an illegal regrouping of a strong sum; on the inner disk
both procedures are admissible and agree by finite-word Taylor identities.

For $V_q=w(1+w^{q})$ every nonzero zero satisfies $a^{q}=-1$ and
$V_q'(a)=-q$, so all these zeros share the exact exponent
\begin{equation}\label{dyn:eq:rhoq}
 \rho_q(\eps)=\frac{\Log(1+\eps)}{\Log(1-q\eps)}
 =-\frac1q+A_q\eps+B_q\eps^{2}+O(\eps^{3}),
\end{equation}
\begin{equation}\label{dyn:eq:ABC}
 A_q=\frac{q+1}{2q},\qquad B_q=\frac{(q+1)(q-4)}{12q},\qquad
 C_q=\frac{(q+1)(2q+1)}{12q},
\end{equation}
the ratio being legitimate because numerator and denominator both start in
degree one with the denominator's coefficient $-q\ne0$.

\begin{theorem}[All-orders factorization and finite-cover obstruction for
the model]\label{dyn:thm:modelfactor}
The rescaled coordinate has the all-orders form
\begin{equation}\label{dyn:eq:modelfactor}
 \Psi_\eps(w)=w(1+w^{q})^{\rho_q(\eps)}\ExpH(P_\eps(w)),
 \qquad P_\eps\in\eps^{2}\C[w][[\eps]],\ P_\eps(0)=0 .
\end{equation}
At every nonzero fixed point the monodromy multiplier
$M_q=\ExpH(2\pi i\rho_q)$ satisfies
\begin{equation}\label{dyn:eq:Mq}
 M_q^{q}=1+\pi i(q+1)\eps+O(\eps^{2}),
\end{equation}
and has infinite order. The leading coordinate becomes single-valued after
finite ramification, but no finite branched cover makes $\Psi_\eps$
single-valued at these points.
\end{theorem}
\begin{proof}
Apply \cref{dyn:cor:factorization}: all nonzero-root exponents coincide and
$\prod_{a^{q}=-1}(1-w/a)=1+w^{q}$, the branches of the logarithms agreeing
near $0$ where each is normalized to vanish, so their sum is
$\log(1+w^{q})$; this gives \eqref{dyn:eq:modelfactor} with the same
polynomial part. The coefficient at zero in $q\rho_q$ is $-1$ and the
coefficient at $\eps$ is $(q+1)/2$, so exponentiation gives
\eqref{dyn:eq:Mq}; no positive power is $1$ by \cref{dyn:lem:torsion} since
$A_q\ne0$. The leading exponent $-1/q$ is rational, so a $q$-fold local
ramification removes its branching, and \cref{dyn:thm:monodromy} shows that
the positive part of the full exponent survives every finite ramification.
\end{proof}

A circuit repeated $q$ times thus returns the leading branch to itself and
still changes the full answer by a nonzero infinitesimal amount: the
monodromy group records information lost under passage to the standard
part. This is the phenomenon whose linear-ODE ancestor is already recorded
by the companion report \emph{Differential Algebra and Differential
Equations over the Surreal and Surcomplex Numbers}
(\texttt{docs/surcomplex/differential-equations/}); see
\cref{dyn:subsec:diffeq} for the exact inheritance.

\begin{proposition}[Explicit second-order expansion]\label{dyn:prop:second}
With $L(w)=\log(1+w^{q})$ normalized by $L(0)=0$, on any admissible common
domain
\begin{equation}\label{dyn:eq:secondorder}
 \Psi_\eps(w)=w(1+w^{q})^{-1/q}\Bigl[1+A_q\eps L(w)
 +\eps^{2}\Bigl(B_qL(w)+\tfrac{A_q^{2}}{2}L(w)^{2}-C_qw^{q}\Bigr)
 +O(\eps^{3})\Bigr],
\end{equation}
equivalently $P_\eps(w)=-C_q\eps^{2}w^{q}+O(\eps^{3})$.
\end{proposition}
\begin{proof}
With $V=V_q$ and $\mu=1$, expansion using
\eqref{dyn:eq:modifiedvector} gives
$[\eps^{2}]\Omega_\eps=\bigl(1/3-V'/4-(V')^{2}/12-VV''/12\bigr)/V$;
substituting $V=w(1+w^{q})$ and collecting powers of $w^{q}$ yields
\[
 [\eps^{2}]\Omega_\eps=-\frac{(q+1)(2q+1)}{12}w^{q-1}
 +\frac{(q+1)(q-4)}{12}\frac{w^{q-1}}{1+w^{q}} ,
\]
whose first summand is the polynomial part with normalized primitive
$-C_qw^{q}$ and whose second has normalized primitive $B_qL(w)$; the
coefficient at $\eps$ integrates to $A_qL(w)$ and the coefficient at zero
gives $\Psi_0$.
\end{proof}

For $q=1$ this reads
\begin{equation}\label{dyn:eq:q1second}
 \Psi_\eps(w)=\frac{w}{1+w}\Bigl[1+\eps\log(1+w)
 +\frac{\eps^{2}}{2}\bigl(\log^{2}(1+w)-\log(1+w)-w\bigr)+O(\eps^{3})\Bigr].
\end{equation}
For $q=4$ the coefficient $B_q$ happens to vanish, but the first-order
logarithm and hence the infinite-order monodromy do not: an accidental
cancellation at one higher order cannot repair the obstruction.

\begin{remark}[Three representations of one inner-disk coordinate]
\label{dyn:rem:threereps}
The formal Taylor series \eqref{dyn:eq:originalh} supplies exact valuations
and the sharp disk; the common-domain coefficient series
$\Psi_\eps=\sum\eps^{n}\Psi_n$ supplies continuation at the rescaled
boundary; and the factorization \eqref{dyn:eq:modelfactor} isolates
monodromy in a single Hahn constant with a single-valued
polynomial-coefficient factor left over. They answer different questions:
to evaluate very close to $0$ the recurrence
\eqref{dyn:eq:brecurrence} is efficient and directly summable; to determine
branching near a boundary fixed point the exact ratio of logarithms avoids
computing many individual Taylor coefficients; and to see why algebraicity
fails even when the leading function is algebraic, the first correction or
the monodromy orbit suffices.
\end{remark}

\begin{warning}[Transcendental, but differentially algebraic]
\label{dyn:warn:diffalg}
Transcendence in \cref{dyn:thm:dichotomy} is relative to the specific base
$\BB$. Source 04 volunteers the counterweight: $H_\eps$ is
\emph{differentially} algebraic in a very explicit way, satisfying the
first-order linear equation $H_\eps'=\Omega_\eps H_\eps$ with
$\Omega_\eps\in\BB$. So ``transcendental'' must not be upgraded to
``differentially transcendental''. Nor is it a claim that any individual
value $H_\eps(w_0)$ lies outside $\No[i]$; the statement concerns a
\emph{function} in an ambient meromorphic Hahn field and is explicitly
consistent with $\No[i]$ being algebraically closed.

Everything above stays true of $H_\eps$, and the sentences above are not
retracted by what follows. What the companion report \emph{Tail-Spans and
Differential Transcendence in Surreal and Surcomplex Hahn Fields}
(\texttt{docs/surreal/tail-spans-and-differential-transcendence/}) shows is
that the refusal is about $H_\eps$ and not about the workspace: its
Theorem \texttt{tail:thm:mixed} exhibits explicit functions
$F_A(z,t)=\sum_{n\in A}t^{n}\sqrt{1-z/2^{n}}$, with the branch equal to
$1$ at $z=0$, one for each infinite $A\subseteq\N_{>0}$, all of whose
mixed jets under $\partial_z$ and $t\,d/dt$ are algebraically independent
over $\C(z)((t^{\Q}))$ --- the $\Gamma=\Q$ case of a base of the same
shape as $\BB$ --- so that these functions satisfy no algebraic
differential equation of any finite order over that field.
Its Theorem \texttt{tail:thm:continuumanalytic} does this for continuum
many such functions on the single fixed disk $|z|<1$, with no shrinking
of the disk, matching the fixed-domain discipline of
\cref{dyn:sec:periods}. Those functions are not solutions of
\cref{dyn:eq:schroederintro} and nothing about them transfers to
$H_\eps$; the pointer records that ``transcendental but differentially
algebraic'' is a property of $H_\eps$, not a ceiling on what a base of
this shape permits. Its Remark \texttt{tail:rem:dynupgrade} states
the same relation from the other side.
\end{warning}

\section{Parabolic dynamics II: time forms, periods, and moduli}
\label{dyn:sec:periods}

This section is also under $\pares$ and $\sdno$, and its maps are
near-identity, hence unipotent: there are \emph{no} eigenvalue ratios and
\emph{no} Diophantine hypotheses (\cref{dyn:warn:parabolic}). The single
defining choice is the coefficient algebra: everything happens in
$\AAA=\OO(D)((t^{\Gamma}))$ for \emph{one} fixed ordinary connected
$D\subset\C$, the fixed-polydisk row of \cref{dyn:subsec:rings}, and
\textbf{no shrinking of $D$ is permitted in any theorem}. The Hahn
monomials are constants, not functions of $z$; the only derivation is
$\partial_z$ applied coefficientwise. Write also
$\mathcal B=\MM(D)((t^{\Gamma}))$, a field, used only for division steps;
$\AAA$ itself is a commutative integral domain and generally not a field,
with $K=K_\Gamma$ embedded as the constants. All integrals below are
\emph{ordinary} complex line integrals taken coefficient by coefficient
over ordinary curves in $D$.

\begin{warning}[These periods are not surcomplex contour integrals]
\label{dyn:warn:periods}
A period below is an ordinary coefficientwise contour integral. It is
\textbf{not} a surcomplex contour integral, \textbf{not} a residue at an
arbitrary surcomplex point, and \textbf{not} a fine-topological notion.
Source 05 cites the companion report \emph{Contours and Stokes: Jordan
separation, Cauchy theory, and Stokes calculus on the surcomplex plane}
(\texttt{docs/surcomplex/contours-and-stokes/}) as background and
deliberately does \emph{not} use surcomplex contours; it carries a remark
forbidding exactly that identification. Furthermore, if one allows an
unrelated choice of primitive on each monad one has enlarged the function
class and can destroy the global obstruction: the theorems below are about
a single element of $\OO(D)((t^{\Gamma}))$, not a collection of unrelated
infinitesimal germs.
\end{warning}

\subsection{Primitives, time forms, and pullback exactness}

\begin{lemma}[Support-preserving primitives]\label{dyn:lem:primitive}
For $g\in\AAA$, a primitive $P\in\AAA$ with $P'=g$ exists if and only if
$\oint_\gamma g(z)\dd z=0$ (coefficientwise) for every ordinary closed
piecewise $C^{1}$ curve in $D$. If it exists, the primitive with
$P(z_{*})=0$ is unique and introduces no new support exponents.
\end{lemma}
\begin{proof}
Necessity holds coefficientwise by the ordinary fundamental theorem for
line integrals. Conversely define each coefficient primitive by path
integration from $z_{*}$; vanishing periods make this path-independent, and
on an ordinary disk the result is an ordinary holomorphic primitive, hence
holomorphic on all of $D$. Use zero for a coefficient that was zero, so
the support is a subset of $\supp g$. Two primitives differ by an element
of $K$, removed by the normalization.
\end{proof}

Assume now that $W=W_F=X_Fz$ has a nowhere-zero leading coefficient. By
\cref{dyn:lem:unit} the \emph{time form}
\[
 \Omega_F=\frac{\dd z}{W_F}
\]
lies in $\AAA\dd z$: ``holomorphic'' refers to all its ordinary coefficient
functions, and its \emph{valuation may be negative}. Its significance is
elementary: in a local primitive coordinate $P$ with $P'=1/W$ the vector
field $W\partial_z$ becomes a unit translation field; the issue is whether
such local primitives fit together globally, and whether one can compare
the primitives of two maps without leaving the common coefficient domain.
For $H=z+h\in\GG_+(D)$ the pullback of a one-form is
$H^{*}(g\dd z)=(g\circ H)H'\dd z$.

\begin{lemma}[Explicit exactness of an infinitesimal pullback]
\label{dyn:lem:homotopy}
For $g\in\AAA$ and $h\in\AAA_{>0}$ the series
\begin{equation}\label{dyn:eq:Q}
 Q_h(g)=\sum_{n\ge1}\frac{h^{n}}{n!}g^{(n-1)}
\end{equation}
is strongly summable and satisfies
\begin{equation}\label{dyn:eq:pullback-exact}
 (z+h)^{*}(g\dd z)-g\dd z=\dd Q_h(g).
\end{equation}
Consequently every coefficientwise period is unchanged by a near-identity
pullback.
\end{lemma}
\begin{proof}
For $S=\supp h$, the support in \eqref{dyn:eq:Q} lies in
$\supp g+(S^{*}\setminus\{0\})$ with finitely many contributions per
coefficient, so termwise differentiation is legitimate and gives
\[
 (Q_h(g))'=\sum_{n\ge1}\frac{h^{n}}{n!}g^{(n)}
 +h'\sum_{n\ge1}\frac{h^{n-1}}{(n-1)!}g^{(n-1)}
 =g(z+h)-g(z)+h'g(z+h)=g(z+h)(1+h')-g(z),
\]
which is \eqref{dyn:eq:pullback-exact}. Exact coefficient forms have zero
ordinary periods by \cref{dyn:lem:primitive}.
\end{proof}

\Cref{dyn:lem:homotopy} is a support-controlled version of the familiar
fact that an infinitesimal change of coordinate acts trivially on
one-dimensional holomorphic de Rham cohomology. It is included as
infrastructure; the homotopy principle itself is \emph{not} a novelty
claim, and there is no assertion that $s\mapsto z+sh$ is a fine-continuous
path in the surreal topology --- formula \eqref{dyn:eq:Q} is the
construction and its support certificate is what licenses it.

\begin{proposition}[Time forms under conjugacy]\label{dyn:prop:time-conjugacy}
If $W_{F_1},W_{F_2}$ are invertible in $\AAA$ then for $H\in\GG_+(D)$,
$H\circ F_1=F_2\circ H$ iff $H^{*}\Omega_{F_2}=\Omega_{F_1}$; in
particular conjugate maps have identical time-form periods on every
ordinary closed curve in $D$.
\end{proposition}
\begin{proof}
Equality of time forms reads $H'/W_{F_2}(H)=1/W_{F_1}$, equivalent to
\cref{dyn:prop:intertwine}; period invariance is
\cref{dyn:lem:homotopy}.
\end{proof}

\subsection{The discrete equation as a primitive problem}
Let $F\in\GG_+(D)$ be nonidentity, $W=W_F$, $\mathscr D=W\partial_z$, and
\begin{equation}\label{dyn:eq:E-B}
 E(X)=\frac{e^{X}-1}{X}=\sum_{n\ge0}\frac{X^{n}}{(n+1)!},\qquad
 B(X)=\frac{X}{e^{X}-1}=\sum_{n\ge0}\frac{B_n}{n!}X^{n},
\end{equation}
with $B_0=1$, $B_1=-1/2$. Substitution of $\mathscr D$ into either series
is strongly summable, and these are inverse automorphisms of the underlying
strongly linear space $\AAA$; \emph{no} multiplicativity of $B(\mathscr D)$
or $E(\mathscr D)$ is asserted. Write $\Delta_F\psi=\psi\circ F-\psi$.

\begin{theorem}[Discrete equation as a primitive problem]
\label{dyn:thm:difference}
For nonidentity $F\in\GG_+(D)$,
\begin{equation}\label{dyn:eq:images}
 \ker\Delta_F=K,\qquad \Delta_F(\AAA)=W\partial_z\AAA .
\end{equation}
For $g\in\AAA$ the equation $\psi\circ F-\psi=g$ has a solution
$\psi\in\AAA$ if and only if the meromorphic Hahn quotient $g/W$ belongs
to $\AAA$ \textbf{and}
\begin{equation}\label{dyn:eq:periodcondition}
 \oint_\gamma\frac{g(z)}{W(z)}\dd z=0
 \quad\text{for every ordinary closed curve }\gamma\subset D .
\end{equation}
If $P'=g/W$, all solutions are $\psi=B(\mathscr D)P+C$ with $C\in K$. When
the leading coefficient of $W$ is nowhere zero, the membership
$g/W\in\AAA$ is automatic.
\end{theorem}
\begin{proof}
By \cref{dyn:thm:exp-log} and the scalar identities in
\eqref{dyn:eq:E-B},
\begin{equation}\label{dyn:eq:factor-difference}
 \Delta_F=e^{\mathscr D}-I=\mathscr DE(\mathscr D)=E(\mathscr D)\mathscr D,
\end{equation}
and $E(\mathscr D)B(\mathscr D)=I=B(\mathscr D)E(\mathscr D)$ follows
coefficientwise: for each input $\psi\in\AAA$, apply
\cref{dyn:cor:ancestry} with $B=\supp\psi$ and $S=\supp W$ to justify
every rearrangement. Thus the image equals $\mathscr D\AAA$. The kernel of
$\mathscr D$ is $K$ since
$W\ne0$, $\AAA$ is a domain and $\ker\partial_z=K$, and both
$E(\mathscr D)$ and its inverse act as the identity on constants, so
\eqref{dyn:eq:factor-difference} also gives $\ker\Delta_F=K$. Since
division by $W$ is available in $\mathcal B$, membership in
$W\partial_z\AAA$ is equivalent to $g/W$ being the derivative of an element
of $\AAA$, which is \cref{dyn:lem:primitive}. Finally
$\Delta_FB(\mathscr D)P=\mathscr DP=WP'=g$, and the kernel statement gives
uniqueness up to $K$; the last assertion is \cref{dyn:lem:unit}.
\end{proof}

\begin{remark}[Why the uncorrected forcing supplies the periods]
\label{dyn:rem:bernoulli}
One might instead multiply by $B(\mathscr D)$ and integrate
$B(\mathscr D)g/W$. This gives the same obstruction: for $n\ge1$,
$\mathscr D^{n}g/W=(\mathscr D^{n-1}g)'$ is exact, so all Bernoulli
correction terms have zero periods. The factorization
\eqref{dyn:eq:factor-difference} makes the uncorrected quotient $g/W$ and
its obstruction especially transparent.
\end{remark}

\begin{corollary}[Exact global Abel criterion]\label{dyn:cor:abel}
Suppose the leading coefficient of $W_F$ is nowhere zero. Then
$\psi\circ F=\psi+1$ has a solution in $\AAA$ if and only if \emph{all} the
periods of $\Omega_F$ vanish, in which case any primitive $P'=1/W_F$
already satisfies the equation and all solutions differ from $P$ by a
constant in $K$.
\end{corollary}
\begin{proof}
Apply \cref{dyn:thm:difference} with $g=1$; also $\mathscr DP=1$ and
$\mathscr D^{2}P=0$, so $e^{\mathscr D}P=P+1$, while the Bernoulli formula
gives $B(\mathscr D)P=P-1/2$, differing only by a constant.
\end{proof}

\begin{warning}[An Abel function is not a globally injective coordinate]
\label{dyn:warn:abel}
The identity $P'=1/W_F$ guarantees a nonvanishing derivative on the halo
under the leading-unit hypothesis, but it does \textbf{not} make $P$
globally injective on a multiply covered domain: even in ordinary complex
analysis a locally univalent primitive need not be globally univalent.
\Cref{dyn:cor:abel} is an Abel-function theorem, not an unconditional
global conjugacy with translation on a subset of $K$.
\end{warning}

\begin{theorem}[Complete finite-dimensional discrete obstruction]
\label{dyn:thm:exact-sequence}
Let $D_r=\C\setminus\{a_1,\dots,a_r\}$ with distinct ordinary points and
positively oriented small ordinary circles $\gamma_j$ about $a_j$ enclosing
no other puncture, and put
\begin{equation}\label{dyn:eq:periodmap}
 (\mathcal P_Fg)_j=\frac1{2\pi i}\oint_{\gamma_j}\frac{g}{W_F}\dd z .
\end{equation}
If $F\in\GG_+(D_r)$ has a nowhere-zero leading displacement coefficient,
then
\begin{equation}\label{dyn:eq:exact-sequence}
 0\longrightarrow K\longrightarrow\AAA
 \xrightarrow{\ \Delta_F\ }\AAA
 \xrightarrow{\ \mathcal P_F\ }K^{r}\longrightarrow0
\end{equation}
is an exact sequence of $K$-vector spaces, with explicit right inverse
\begin{equation}\label{dyn:eq:rightinverse}
 (q_1,\dots,q_r)\longmapsto W_F(z)\sum_{j=1}^{r}\frac{q_j}{z-a_j},
\end{equation}
so that every forcing decomposes as
\begin{equation}\label{dyn:eq:forcing-decomposition}
 g=\Delta_F\psi+W_F\sum_{j=1}^{r}\frac{(\mathcal P_Fg)_j}{z-a_j},
\end{equation}
with $\psi$ determined up to a constant once the second term is fixed.
\end{theorem}
\begin{proof}
The initial kernel statement is \cref{dyn:thm:difference}. Every ordinary
closed curve in $D_r$ has, against a holomorphic one-form, the same
periods as the integer combination of small puncture circles given by its
winding numbers: remove small disks about the punctures inside a large disk
containing the curve and cut the remaining planar region along arcs joining
the boundary components, so that the cut region is simply connected and the
integrals along the two sides of each cut cancel. Applying this ordinary
argument coefficientwise shows that vanishing of the $r$ components of
\eqref{dyn:eq:periodmap} is equivalent to
\eqref{dyn:eq:periodcondition}, so middle exactness follows from
\cref{dyn:thm:difference}. The differential $\dd z/(z-a_j)$ has integral
$2\pi i$ around $\gamma_j$ and zero around the other chosen circles, so
\eqref{dyn:eq:rightinverse} is a right inverse; subtracting it from $g$ and
applying exactness gives \eqref{dyn:eq:forcing-decomposition}.
\end{proof}

No identification of the infinite-dimensional $\AAA$ with
$K\otimes_\C\OO(D_r)$ is used: the finite-dimensional quotient is computed
directly by coefficientwise periods and explicit representatives.

\begin{example}[A local obstruction where the leading coefficient vanishes]
\label{dyn:ex:localobstruction}
The quotient condition in \cref{dyn:thm:difference} must not be dropped
outside the nonsingular stratum. On a disk containing $0$ take
$W=\eps z^{2}$ and $F=z/(1-\eps z)$ with $\eps=t^{\delta}$, $\delta>0$;
this $F$ is the exponential of $W\partial_z$ in $\GG_+(D)$. The equation
$\psi\circ F-\psi=1$ has \emph{no} solution in $\AAA$, because
$1/W=\eps^{-1}z^{-2}$ is not holomorphic at $0$ --- equivalently $F$ fixes
$0$ and evaluating the equation there would give $0=1$. On a punctured
disk, by contrast, $-1/(\eps z)$ is an Abel function. This separates a
local divisibility obstruction from the global period obstruction.
\end{example}

\subsection{The complete global conjugacy classification}
Fix $\delta>0$, $\eps=t^{\delta}$, and $a\in\OO(D)$ \emph{nowhere zero},
and define the nonsingular stratum
\begin{equation}\label{dyn:eq:stratum}
 \GG_{\delta,a}(D)=\{F\in\GG_+(D):F=z+\eps a+r,\ v(r)>\delta\}
 =\{F:W_F=\eps(a+u),\ u\in\AAA_{>0}\},
\end{equation}
with normalized time form
\begin{equation}\label{dyn:eq:normalized-form}
 \beta_F=\eps\Omega_F=b_F(z)\dd z,\qquad
 b_F=a^{-1}+\text{positive Hahn terms} .
\end{equation}
The leading data $(\delta,a)$ are a conjugacy invariant under
$\GG_+(D)$: the identity $W_2(H)=H'W_1$ of \cref{dyn:prop:intertwine}
changes neither the leading coefficient nor the valuation, because
$H=z+$ positive and $H'=1+$ positive. It is therefore natural to classify
one stratum at a time.

\begin{lemma}[Implicit displacement lemma]\label{dyn:lem:displacement2}
Let $b=b_0+u\in\AAA_{\ge0}$ with $b_0\in\OO(D)$ nowhere zero and
$u\in\AAA_{>0}$. For each $p\in\AAA_{>0}$ there is a unique
$h\in\AAA_{>0}$ satisfying
\begin{equation}\label{dyn:eq:implicit-displacement}
 bh+\frac{b'}{2!}h^{2}+\frac{b''}{3!}h^{3}+\cdots=p .
\end{equation}
Every coefficient of $h$ is holomorphic on $D$; if
$S=\supp u\cup\supp p$ then $\supp h\subseteq S^{*}\setminus\{0\}$; and if
$p(z_{*})=0$ then $h(z_{*})=0$.
\end{lemma}
\begin{proof}
If $S=\emptyset$ take $h=0$. Otherwise let $\mathsf M=S^{*}$, well ordered
by \cref{dyn:lem:neumann}; set the coefficient of $h$ at $0$ to zero and
determine $h_\gamma$ for $\gamma\in\mathsf M\setminus\{0\}$ in increasing
order. The coefficient of $h_\gamma$ on the left of
\eqref{dyn:eq:implicit-displacement} is precisely $b_0$; any other linear
term comes from $u_\alpha h_\nu$ with $\alpha>0$ and
$\nu=\gamma-\alpha<\gamma$, and in a term with $n\ge2$ factors of $h$ every
factor exponent is positive and strictly smaller than $\gamma$. The
coefficient is a finite sum (expand the factors into words in $S$ and apply
\cref{dyn:lem:neumann}), so the recursion has the form
$h_\gamma=b_0^{-1}(p_\gamma-R_\gamma)$ with $R_\gamma$ a finite expression
in earlier coefficients and ordinary derivatives of $b$, defining a
holomorphic function on $D$. For unrestricted uniqueness, the difference of
the two left sides for two solutions $h_1,h_2\in\AAA_{>0}$ factors as
$(h_1-h_2)(b_0+v)$ with $v\in\AAA_{>0}$ --- factor each difference of
powers, the resulting sum being strongly summable over the finite union of
the positive supports of $u,h_1,h_2$ --- and the second factor is a unit by
\cref{dyn:lem:unit}. At $z_{*}$, evaluation gives the same scalar implicit
equation over $K$ with linear leading coefficient $b_0(z_{*})\ne0$, so if
$p(z_{*})=0$ then zero is a solution and the same factorization shows it is
the only positive one.
\end{proof}

\begin{warning}[Where the nowhere-zero hypothesis does its work]
\label{dyn:warn:nowherezero}
\Cref{dyn:lem:displacement2} is exactly the point at which the
classification proof stops applying when $a$ has zeros. Its
common-domain conclusion is stronger than solving a separate germ problem
at each point: \emph{every stage divides only by the same nowhere-zero
ordinary function $b_0$} and uses derivatives of functions already
holomorphic on $D$, so no radius is lost at any Hahn exponent. Source 05
states outright that the zeros case needs local zero divisors,
multiplicities and singular normal-form obstructions combined with global
periods, and that its classification is not claimed there.
\end{warning}

\begin{theorem}[Complete period invariant and marked conjugacy]
\label{dyn:thm:classification}
For $F_1,F_2\in\GG_{\delta,a}(D)$ with $a$ nowhere zero, the following are
equivalent:
\begin{enumerate}[label=(\roman*)]
\item there is $H\in\GG_+(D)$ with $H\circ F_1=F_2\circ H$;
\item $\oint_\gamma\Omega_{F_1}=\oint_\gamma\Omega_{F_2}$ for every
ordinary closed piecewise $C^{1}$ curve $\gamma\subset D$;
\item $\beta_{F_1}-\beta_{F_2}=\dd p$ for some $p\in\AAA_{>0}$.
\end{enumerate}
When they hold, for every specified $z_{*}\in D$ there is a \emph{unique}
conjugacy $H$ as in (i) with $H(z_{*})=z_{*}$, obtained by taking the
normalized primitive $p(z_{*})=0$ in (iii) and solving
\begin{equation}\label{dyn:eq:construct-conjugacy}
 \sum_{n\ge1}\frac{b_{F_2}^{(n-1)}}{n!}h^{n}=p,\qquad H=z+h .
\end{equation}
\textbf{No shrinking of $D$ is permitted anywhere in this statement.}
\end{theorem}
\begin{proof}
(i)$\Rightarrow$(ii) is \cref{dyn:prop:time-conjugacy}. Multiplication by
the nonzero constant $\eps$ shows (ii) equivalent to the vanishing of all
periods of $\beta_{F_1}-\beta_{F_2}$; both normalized time forms have the
same leading coefficient $a^{-1}$, so their difference has positive
support, and the normalized primitive of \cref{dyn:lem:primitive}, when it
exists, also has positive support --- this is the equivalence of (ii) and
(iii) including the normalization of $p$. Assume (iii): the leading
coefficient of $b_{F_2}$ is $a^{-1}$, nowhere zero, so
\cref{dyn:lem:displacement2} gives a unique positive $h$ in
\eqref{dyn:eq:construct-conjugacy}, and \cref{dyn:lem:homotopy} gives
$H^{*}\beta_{F_2}-\beta_{F_2}=\dd p=\beta_{F_1}-\beta_{F_2}$, i.e.\
$H^{*}\Omega_{F_2}=\Omega_{F_1}$, so $H$ is a conjugacy by
\cref{dyn:prop:time-conjugacy}; its inverse lies in $\GG_+(D)$ by
\cref{dyn:thm:exp-log} and $H(z_{*})=z_{*}$ by
\cref{dyn:lem:displacement2}. If $\widetilde H=z+\widetilde h$ is any
conjugacy fixing $z_{*}$, the same pullback identity yields
$\dd\bigl(Q_{\widetilde h}(b_{F_2})-p\bigr)=0$, so that expression is a
constant in $K$ with value zero at $z_{*}$, whence $\widetilde h$ satisfies
the same implicit equation and uniqueness in
\cref{dyn:lem:displacement2} gives $\widetilde H=H$.
\end{proof}

\begin{corollary}[Simply connected domains]\label{dyn:cor:simplyconnected}
If $D$ is simply connected, any two maps in $\GG_{\delta,a}(D)$ are
conjugate by a unique near-identity change of variable fixing a prescribed
ordinary basepoint, and each is conjugate to $\exp(\eps a\partial_z)z$.
\end{corollary}
\begin{proof}
Every ordinary holomorphic coefficient one-form on $D$ has a primitive, so
(iii) of \cref{dyn:thm:classification} holds.
\end{proof}

For $a=1$ the reference map is $z+\eps$, and this is a genuine global
near-identity conjugacy to that translation on $D^{\#}$; for general $a$
the reference map is the indicated vector-field exponential and no
univalence of a primitive of $\dd z/a$ is required.

\begin{warning}[Three hypotheses that do actual work]
\label{dyn:warn:classification}
$a$ must not vanish, because inversion and the implicit displacement
recursion divide by $a$ or $a^{-1}$
(\cref{dyn:warn:nowherezero}). The conjugating map must have identity
ordinary part, because otherwise it can act nontrivially on ordinary
homology and on the leading form --- source 05 records that extending the
classification to a larger group would require the ordinary automorphism
action on the leading form and its periods, and that equal period vectors
in a fixed basis would then no longer be the correct equivalence relation
by themselves. And all coefficient functions must be defined on the same
$D$: replacing global functions by independent germs removes the meaning of
a single period comparison. The theorem is not claimed after any of these
hypotheses is erased.
\end{warning}

\subsection{Explicit moduli}

\begin{theorem}[An explicit moduli space]\label{dyn:thm:moduli}
On $D_r=\C\setminus\{a_1,\dots,a_r\}$, for the fixed stratum
$(\delta,a)$ define
\begin{equation}\label{dyn:eq:moduli-map}
 q_j(F)=\frac1{2\pi i}\oint_{\gamma_j}
 \Bigl(\beta_F-\frac{\dd z}{a(z)}\Bigr)\in\m_K .
\end{equation}
Then $F\mapsto(q_1(F),\dots,q_r(F))$ induces a \emph{bijection}
\begin{equation}\label{dyn:eq:moduli-bijection}
 \GG_{\delta,a}(D_r)/\text{conjugacy by }\GG_+(D_r)
 \ \longleftrightarrow\ \m_K^{r},
\end{equation}
and for $\mathbf q\in\m_K^{r}$ a representative is
\begin{equation}\label{dyn:eq:moduli-representative}
 b_{\mathbf q}(z)=\frac1{a(z)}+\sum_{j=1}^{r}\frac{q_j}{z-a_j},
 \qquad W_{\mathbf q}=\frac{\eps}{b_{\mathbf q}},
 \qquad F_{\mathbf q}=\exp(W_{\mathbf q}\partial_z)z .
\end{equation}
No countability, finite generation or rank-one condition is imposed on the
supports of the $q_j$.
\end{theorem}
\begin{proof}
The finite union of the supports of the $q_j$ is positive and well ordered,
so $b_{\mathbf q}\in\AAA_{\ge0}$ with nowhere-zero leading coefficient
$a^{-1}$, invertible by \cref{dyn:lem:unit}; hence
$W_{\mathbf q}=\eps(a+\text{positive})$ and $F_{\mathbf q}$ lies in the
stated stratum with normalized time form exactly $b_{\mathbf q}\dd z$, and
the simple-pole representatives give $q_j(F_{\mathbf q})=q_j$. Conjugate
maps have the same vector by period invariance; conversely equal vectors
give equal periods around each puncture, hence, as in
\cref{dyn:thm:exact-sequence}, equality of every ordinary period, and
\cref{dyn:thm:classification} supplies a conjugacy.
\end{proof}

The parametrization is algebraic over the Hahn constants in its displayed
one-form representatives. Its basis is the \emph{specified} puncture loops;
arbitrary changes of ordinary coordinates are not quotiented out.

\subsection{Nonconjugate dynamics beyond every one-scale jet}
\label{dyn:subsec:flat}

Fix $\delta>0$, $\eps=t^{\delta}$, and define the \emph{flat ideal}
\begin{equation}\label{dyn:eq:flat-ideal}
 J_\delta=\{0\}\cup\{c\in K:v(c)>n\delta\ \text{for every integer }n\ge1\},
\end{equation}
an ideal of the valuation ring $K_{\ge0}$, nonzero \emph{precisely} when
some $\eta\in\Gamma$ exceeds all the $n\delta$ (the forward implication by
taking the valuation of a nonzero member, the reverse from
$t^{\eta}\in J_\delta$). In particular $J_\delta=0$ when the multiples of
$\delta$ are cofinal, as in any Archimedean value group. For coefficient
functions the corresponding ideal is
\begin{equation}\label{dyn:eq:flat-functions}
 I_\delta(D)=\bigcap_{n\ge1}\eps^{n}\AAA_{\ge0}
 =\{0\}\cup\{f\in\AAA:v(f)>n\delta\ \forall n\ge1\},
\end{equation}
the equivalence of weak and strict inequalities following by replacing $n$
with $n+1$.

\begin{warning}[The ideals must be formed in $\AAA_{\ge0}$]
\label{dyn:warn:flatideal}
It is essential that \eqref{dyn:eq:flat-ideal},
\eqref{dyn:eq:flat-functions} are formed in $K_{\ge0}$ and
$\AAA_{\ge0}$, \textbf{not} in $K$ or $\AAA$, where $\eps$ is a
\emph{unit} and its powers generate everything. Consequently the map to
one-scale jets
\begin{equation}\label{dyn:eq:jet-map}
 \AAA_{\ge0}\longrightarrow
 \varprojlim_n\AAA_{\ge0}/\eps^{n}\AAA_{\ge0}
\end{equation}
has kernel exactly $I_\delta(D)$, and in the noncofinal case it loses
nonzero information. The next theorem shows it can lose a complete family
of global conjugacy obstructions, not merely some irrelevant constants.
\end{warning}

\begin{theorem}[Flatly indistinguishable but globally nonconjugate]
\label{dyn:thm:flat}
Let $r\ge1$, $D_r=\C\setminus\{a_1,\dots,a_r\}$, and suppose
$J_\delta\ne0$. For $\mathbf c\in J_\delta^{r}$ put
\begin{equation}\label{dyn:eq:flat-family}
 R_{\mathbf c}(z)=\sum_{j=1}^{r}\frac{c_j}{z-a_j},\qquad
 W_{\mathbf c}=\frac{\eps}{1+\eps R_{\mathbf c}},\qquad
 F_{\mathbf c}=\exp(W_{\mathbf c}\partial_z)z .
\end{equation}
All these maps lie in $\GG_{\delta,1}(D_r)$ and satisfy:
\begin{enumerate}[label=(\alph*)]
\item $F_{\mathbf0}=z+\eps$ and
$F_{\mathbf c}-F_{\mathbf0}\in I_\delta(D_r)$ for every $\mathbf c$; if
$\mathbf c\ne0$ and $\eta=\min_{c_j\ne0}v(c_j)$ then \emph{exactly}
\begin{equation}\label{dyn:eq:flat-valuation}
 v\bigl(F_{\mathbf c}-F_{\mathbf0}\bigr)=2\delta+\eta ;
\end{equation}
\item their time forms are \emph{exactly}
\begin{equation}\label{dyn:eq:exact-flat-form}
 \Omega_{F_{\mathbf c}}=\eps^{-1}\dd z+\sum_{j=1}^{r}\frac{c_j\dd z}{z-a_j};
\end{equation}
\item $F_{\mathbf c}$ and $F_{\mathbf d}$ are conjugate by an element of
$\GG_+(D_r)$ if and only if $\mathbf c=\mathbf d$;
\item a global Abel function in $\AAA$ exists for $F_{\mathbf c}$ if and
only if $\mathbf c=\mathbf0$.
\end{enumerate}
\end{theorem}
\begin{proof}
The finite union of the supports of the $c_j$ is well ordered and positive,
so $R_{\mathbf c}\in\AAA_{>0}$, and the denominator in
\eqref{dyn:eq:flat-family} has leading coefficient $1$ and is a unit, so
$W_{\mathbf c}=\eps+$ terms of valuation greater than $\delta$;
\cref{dyn:thm:exp-log} places $F_{\mathbf c}$ in the stated stratum, and
for $\mathbf c=0$ the vector field is constant with exponential exactly
$z+\eps$. If $\mathbf c\ne0$ then $v(R_{\mathbf c})=\eta$, since a
nontrivial linear combination of the $(z-a_j)^{-1}$ cannot vanish
identically --- ordinary residues at the punctures recover its
coefficients. The geometric expansion gives
\begin{equation}\label{dyn:eq:Wflat-difference}
 W_{\mathbf c}-\eps=-\eps^{2}R_{\mathbf c}+\eps^{3}R_{\mathbf c}^{2}-\cdots,
 \qquad v(W_{\mathbf c}-\eps)=2\delta+\eta .
\end{equation}
Put $V=W_{\mathbf c}-\eps$ and $\lambda_0=2\delta+\eta$. In the exponential
expansion of $(\eps\partial_z+V\partial_z)^{n}z$ for $n\ge2$, the word with
no factor $V\partial_z$ vanishes because $(\eps\partial_z)^{2}z=0$, and
every other word has valuation at least $\lambda_0+(n-1)\delta>\lambda_0$
since derivatives do not lower Hahn valuation; all such words are summable
by \cref{dyn:lem:neumann}. So the leading difference comes from $V$ itself,
proving \eqref{dyn:eq:flat-valuation}, and $\eta>n\delta$ for all $n$ puts
the difference in \eqref{dyn:eq:flat-functions}. The logarithm of
$F_{\mathbf c}$ is $W_{\mathbf c}$ by the inverse identity in
\cref{dyn:thm:exp-log}, and taking its reciprocal gives
\eqref{dyn:eq:exact-flat-form}; hence
$(2\pi i)^{-1}\oint_{\gamma_j}\Omega_{F_{\mathbf c}}=c_j$, and for the
normalized form the corresponding value is $\eps c_j$. Equality of periods
is therefore equivalent to $\mathbf c=\mathbf d$, so
\cref{dyn:thm:classification} proves (c) and \cref{dyn:cor:abel} proves
(d).
\end{proof}

\begin{remark}[A failure of one-scale information, not of Hahn equality]
\label{dyn:rem:flat}
All the maps of \cref{dyn:thm:flat} have the same image under
\eqref{dyn:eq:jet-map}. They do \emph{not} have the same full Hahn
coefficients: the missing coefficients occur at exponents larger than every
$n\delta$. The result therefore contradicts neither uniqueness of Hahn
expansions nor separation by arbitrary value-group cutoffs. It shows
exactly why replacing arbitrary-rank support control by a single
``small-parameter completion'' can erase a global obstruction. This is the
interpretive reading source 05 attaches to its own theorem, and is a
framing argument supported by the theorem rather than itself a theorem.
\end{remark}

\begin{example}[A concrete surcomplex instance]\label{dyn:ex:flatsurcomplex}
Take $\Gamma=\Q\oplus\Q$ lexicographically, $\delta=(0,1)$,
$\eta=(1,0)$, so that $\eta>n\delta$ for every integer $n\ge1$ because
comparison is decided by the first coordinate; realize this inside the
surreals as $\Q\omega+\Q$ via $(q,r)\mapsto q\omega+r$. Under
\eqref{dyn:eq:embedding}, $\eps=\omega^{-1}$ and
$c=t^{\eta}=\omega^{-\omega}$. On $D=\C^{*}$ the field
\begin{equation}\label{dyn:eq:surreal-flat-example}
 W_c(z)=\frac{\omega^{-1}}{1+\omega^{-(\omega+1)}/z}
\end{equation}
defines an explicit surcomplex map $F_c=\exp(W_c\partial_z)z$ with the same
truncation as $z+\omega^{-1}$ modulo every positive integer power of
$\omega^{-1}$, yet
$(2\pi i)^{-1}\oint_{|z|=1}\Omega_{F_c}=\omega^{-\omega}\ne0$: it is not
globally near-identity conjugate to that translation and has no
single-valued Hahn-coherent Abel function on the punctured plane. This is a
global effect at a strictly smaller \emph{scale}, and is explicitly
\textbf{not} a claim about ordinary exponentially small asymptotics of
complex numbers --- $\omega^{-\omega}$ is an exact Hahn monomial in a
rank-two workspace. A nonzero ordinary multiple of it already gives a
different member of the family, and the full parameter space $J_\delta$ is
much larger.
\end{example}

\subsection{Worked parabolic examples and the iterative residue}
\label{dyn:subsec:eulerexamples}

For a map written in powers of one positive monomial,
$F=z+\eps v_1+\eps^{2}b+\eps^{3}c+O(\eps^{4})$ with
$v_1,b,c\in\OO(D)$, direct expansion of \eqref{dyn:eq:explog} gives
\begin{equation}\label{dyn:eq:log-first-terms}
 W_F=\eps v_1+\eps^{2}\Bigl(b-\tfrac12v_1v_1'\Bigr)
 +\eps^{3}\Bigl(c-\tfrac12(v_1b'+bv_1')+\tfrac1{12}v_1^{2}v_1''
 +\tfrac13v_1(v_1')^{2}\Bigr)+O(\eps^{4}),
\end{equation}
obtained from $\Delta z=h$, $\Delta^{2}z=hh'+\frac12h^{2}h''+\cdots$,
$\Delta^{3}z=h(h')^{2}+h^{2}h''+\cdots$ with $h=F-z$ and the logarithm
coefficients $1,-\frac12,\frac13$. Here $O(\eps^{4})$ means membership in
$\eps^{4}\OO(D)[[\eps]]$ and is used only in examples already specified as
one-parameter series; it is never a substitute for arbitrary-rank support
conditions. If $v_1$ is nowhere zero, the time form begins
\begin{equation}\label{dyn:eq:time-first-terms}
 \Omega_F=\Bigl[\frac1{\eps v_1}+\frac{v_1'}{2v_1}-\frac{b}{v_1^{2}}
 +O(\eps)\Bigr]\dd z ,
\end{equation}
so already at the next order the ordinary logarithmic derivative of the
leading vector field can create a nonzero period.

\begin{proposition}[Exact residue of the Euler time form]
\label{dyn:prop:euler-residue}
On $D=\C^{*}$ let $F_p(z)=z+\eps z^{p+1}$, $p\ge1$. Then
\begin{align}
 W_{F_p}(z)&=\sum_{n\ge1}c_{p,n}\eps^{n}z^{np+1},
 & c_{p,1}&=1,\quad c_{p,2}=-\frac{p+1}{2},\label{dyn:eq:euler-grading}\\
 \frac1{W_{F_p}(z)}&=\eps^{-1}z^{-p-1}+\frac{p+1}{2}z^{-1}
 +\sum_{n\ge1}d_{p,n}\eps^{n}z^{np-1},\label{dyn:eq:euler-reciprocal}
\end{align}
so that
\begin{equation}\label{dyn:eq:euler-residue}
 \Res_{z=0}\Omega_{F_p}=\frac{p+1}{2}
\end{equation}
\emph{exactly}, with no omitted positive-valuation correction. Hence there
is no global Abel function for $F_p$ in $\OO(\C^{*})((t^{\Gamma}))$.
\end{proposition}
\begin{proof}
Taylor substitution by $F_p$ preserves the grading in which
$\eps^{n}z^{np+1}$ has grade $n$, since
$(z+\eps z^{p+1})^{kp+1}=z^{kp+1}(1+\eps z^{p})^{kp+1}$; every term of the
finite difference and its iterates therefore has the displayed form, and
taking the logarithm preserves it, the first two coefficients following
from \eqref{dyn:eq:log-first-terms}. Factor
$W_{F_p}=\eps z^{p+1}C_p(\eps z^{p})$ with
$C_p(x)=1-(p+1)x/2+\cdots$ and invert $C_p$ to obtain
\eqref{dyn:eq:euler-reciprocal}; for $n\ge1$ the exponent $np-1$ is
nonnegative, so none of the last terms has a $z^{-1}$ coefficient, and
\eqref{dyn:eq:euler-residue} is exact. The period about $0$ is nonzero, so
\cref{dyn:cor:abel} forbids a global Abel function.
\end{proof}

The leading field $\eps z^{p+1}$ itself has zero period on $\C^{*}$ and its
exponential has Abel function $-1/(p\eps z^{p})$; so passing from that
exact flow to its Euler map changes the global obstruction even though
their first displacements agree. For $p=1$,
\begin{equation}\label{dyn:eq:euler-known}
 W_{F_1}=\eps z^{2}-\eps^{2}z^{3}+\tfrac32\eps^{3}z^{4}
 -\tfrac83\eps^{4}z^{5}+\tfrac{31}6\eps^{5}z^{6}
 -\tfrac{157}{15}\eps^{6}z^{7}+O(\eps^{7}).
\end{equation}

\begin{warning}[These coefficients are known, and nothing is claimed about
their convergence]\label{dyn:warn:hairer}
The coefficients \eqref{dyn:eq:euler-known} agree, after multiplication by
the step size, with the standard modified equation in Hairer's worked
example \cite{Hairer}; they are \textbf{known examples, not new
coefficients}, and are not proposed as discoveries. Nothing is asserted
about convergence of the ordinary step-size series in a real step size;
Hairer's warning that it generally diverges is cited and accepted. The
constant \eqref{dyn:eq:euler-residue} is the familiar iterative-residue
phenomenon, and the assertion here is its exact interpretation in the
stated common-domain Hahn category, not the discovery of that classical
invariant. The \emph{same} degree-six expansion appears independently in
source 03 as the iterative logarithm of $z+tz^{2}$
(\cref{dyn:ex:quadlog}); it is one calculation and is printed once, in
\eqref{dyn:eq:euler-known}.
\end{warning}

\begin{corollary}[A global coherent normal form for the Euler family]
\label{dyn:cor:euler-normal}
$F_p$ is conjugate on $(\C^{*})^{\#}$, by a near-identity Hahn-coherent
change of variable, to the time-one map of
\begin{equation}\label{dyn:eq:euler-normal}
 V_p(z)=\frac{\eps z^{p+1}}{1+\frac{p+1}{2}\eps z^{p}} ,
\end{equation}
the conjugacy fixing $z=1$ being unique. It is \textbf{not} a conjugacy of
$F_p$ to the time-one map of $\eps z^{p+1}$.
\end{corollary}
\begin{proof}
The time form of \eqref{dyn:eq:euler-normal} is
$\eps^{-1}z^{-p-1}\dd z+\frac{p+1}2\dd z/z$, with the same period as
\eqref{dyn:eq:euler-reciprocal}; their difference consists of nonnegative
powers of $z$ coefficientwise, hence is exact on $\C^{*}$, so
\cref{dyn:thm:classification} applies. The last assertion holds because the
leading vector-field flow has zero period while $F_p$ does not.
\end{proof}

On any simply connected slit domain in $\C^{*}$ the normal form has the
local Abel function $-1/(p\eps z^{p})+\frac{p+1}2\log z$, whose logarithmic
monodromy is precisely the obstruction to a single-valued global Abel
function on $\C^{*}$.

\begin{example}[A difference equation with an explicit Bernoulli correction]
\label{dyn:ex:bernoulli}
Take $W=\eps z$, so $F=e^{\eps}z$, on $D=\C^{*}$. The forcing $g=z^{m}$
with nonzero integer $m$ has zero period after division by $W$, and
$\psi(z)=z^{m}/(e^{m\eps}-1)$ is a solution; the denominator has valuation
$\delta$ so its reciprocal is a legitimate Laurent--Hahn constant, and
expanding gives
$\psi(z)=\bigl(\frac1{m\eps}-\frac12+\frac{m\eps}{12}-\cdots\bigr)z^{m}$,
exactly the Bernoulli solver of \cref{dyn:thm:difference} applied to
$P=z^{m}/(m\eps)$. For the constant forcing $g=1$, by contrast,
$g\dd z/W=\eps^{-1}\dd z/z$ has nonzero period and no global solution
exists, displaying the one-dimensional cokernel for the punctured plane.
\end{example}

\begin{example}[Local straightening does not erase global periods]
\label{dyn:ex:straighten}
$F(z)=(1+\eps)z$ and $G(z)=e^{\eps}z$ have the same leading displacement
$\eps z$ on $\C^{*}$, with logarithmic vector fields
$\log(1+\eps)z$ and $\eps z$ and periods
$2\pi i/\log(1+\eps)$ and $2\pi i/\eps$, which differ --- so there is no
global near-identity conjugacy there. On a simply connected domain with a
branch of $\log z$ there is one:
\[
 H(z)=z\exp\Bigl(\Bigl[\frac{\eps}{\log(1+\eps)}-1\Bigr]\log z\Bigr),
\]
the bracketed coefficient having positive valuation, so the exponential is
a strongly summable Hahn expansion with coefficients holomorphic on the
chosen branch domain, and direct substitution gives $H\circ F=G\circ H$.
The failure to make that formula single-valued on $\C^{*}$ is exactly what
the period theorem detects.
\end{example}

\subsection{Ordinary slow time on one fixed holomorphic flow chart}
\label{dyn:subsec:slowtime}

\Cref{dyn:thm:exp-log} concerns an infinitesimal vector field. One may also
ask what happens at the larger time scale on which the leading ordinary
vector field moves a finite distance; a direct exponential is generally not
a strongly summable Hahn expression there. The correct construction solves
the leading ordinary flow first and then lifts all positive Hahn
coefficients along it. Write
\begin{equation}\label{dyn:eq:slow-vector}
 W=\eps X,\qquad X(z)=a(z)+\sum_{\beta\in S}t^{\beta}X_\beta(z),
\end{equation}
with $S\subset\Gamma_{>0}$ well ordered and all coefficient functions
holomorphic on $D$; in this subsection $a$ \emph{may} have zeros. Let
$\Sigma\subset\C$ be an ordinary disk centered at $0$, $D_0\subset D$ an
ordinary open domain, and assume a holomorphic ordinary flow chart
\begin{equation}\label{dyn:eq:ordinary-flow}
 \phi:\Sigma\times D_0\to D,\qquad
 \partial_s\phi(s,z)=a(\phi(s,z)),\qquad \phi(0,z)=z
\end{equation}
is \emph{specified}, with $J(s,z)=\partial_z\phi(s,z)$ nowhere vanishing on
that chart.

\begin{theorem}[Common-chart Hahn lifting of the slow flow]
\label{dyn:thm:slowtime}
There is \emph{exactly one} series
\begin{equation}\label{dyn:eq:slow-solution}
 \Psi(s,z)=\phi(s,z)+\sum_{\gamma\in S^{*}\setminus\{0\}}
 t^{\gamma}\psi_\gamma(s,z)
\end{equation}
with coefficients holomorphic on the \emph{same} product domain
$\Sigma\times D_0$ satisfying
\begin{equation}\label{dyn:eq:slow-equation}
 \partial_s\Psi=X(\Psi),\qquad \Psi(0,z)=z,
\end{equation}
$X(\Psi)$ meaning Taylor expansion at the ordinary point $\phi(s,z)$.
Uniqueness holds among \emph{all} positive-Hahn perturbations of $\phi$ on
this product domain, not only those with the displayed support bound. Each
coefficient is obtained by one ordinary integral,
\begin{equation}\label{dyn:eq:variation-formula}
 \psi_\gamma(s,z)=J(s,z)\int_0^{s}J(u,z)^{-1}R_\gamma(u,z)\dd u ,
\end{equation}
with $R_\gamma$ a finite expression in previously determined coefficients
and in ordinary derivatives of the $X_\beta$ and $a$ along $\phi$.
Finally Taylor evaluation at $s=\eps$ is strongly summable and
\begin{equation}\label{dyn:eq:slow-epsilon}
 \Psi(\eps,z)=\exp(W\partial_z)z\qquad(z\in D_0^{\#}),
\end{equation}
and more generally, for $c\in K$ with $c=0$ or $v(c)+\delta>0$, evaluation
at $s=c\eps$ is defined and gives $\Psi(c\eps,z)=\exp(cW\partial_z)z$.
\end{theorem}
\begin{proof}
Let $\mathsf M=S^{*}$ and seek $q=\Psi-\phi$ with support in
$\mathsf M\setminus\{0\}$. Expanding
$X(\phi+q)=a(\phi)+a'(\phi)q+\sum_{n\ge2}a^{(n)}(\phi)q^{n}/n!
+\sum_{\beta\in S}t^{\beta}\sum_{n\ge0}X_\beta^{(n)}(\phi)q^{n}/n!$ --- all
compositions with $\phi$ being holomorphic on the specified product because
its image lies in $D$ --- the coefficient equation at
$\gamma\in\mathsf M\setminus\{0\}$ is
\begin{equation}\label{dyn:eq:variational}
 \partial_s\psi_\gamma=a'(\phi)\psi_\gamma+R_\gamma,
 \qquad \psi_\gamma(0,z)=0,
\end{equation}
every factor in $R_\gamma$ involving a positive exponent smaller than
$\gamma$ except for specified coefficients of $X$, with finitely many
contributions by \cref{dyn:lem:neumann} exactly as in
\cref{dyn:lem:displacement2}. Differentiating
\eqref{dyn:eq:ordinary-flow} in $z$ gives
$\partial_sJ=a'(\phi)J$, $J(0,z)=1$, so variation of constants solves
\eqref{dyn:eq:variational} by \eqref{dyn:eq:variation-formula}; since $J$
is nowhere zero and $\Sigma$ is simply connected the integral is
independent of its path in $\Sigma$ and is jointly holomorphic in $(s,z)$,
locally by integrating the holomorphic power series in $s$ with
coefficients holomorphic in $z$ and patching the uniquely normalized
primitives. There is therefore \emph{no} successive domain shrinkage and
\emph{no} requirement of a uniform bound on these coefficients as $\gamma$
varies. The support lies in $\mathsf M$ and the expansion of $X(\Psi)$ is
strongly summable by \cref{dyn:lem:neumann}, so the construction solves
\eqref{dyn:eq:slow-equation}. For uniqueness, at the least exponent where
two positive-Hahn perturbations of $\phi$ satisfying that equation differ,
their difference solves the homogeneous version of
\eqref{dyn:eq:variational} with zero initial condition and is zero; this
argument does not require either support to lie in $\mathsf M$. For the
last assertion expand every coefficient at $s=0$ and substitute
$s=\eps$: the possible exponents lie in the monoid generated by
$S\cup\{\delta\}$ and a fixed exponent has finitely many word
representations, so the evaluation is strongly summable. Repeated
differentiation of \eqref{dyn:eq:slow-equation} gives
$\partial_s^{n}\Psi(0,z)=(X\partial_z)^{n}z$, whence
$\Psi(\eps,z)=\sum_n\eps^{n}(X\partial_z)^{n}z/n!=\exp(W\partial_z)z$. For
$s=c\eps\in\m_K$ replace $\{\delta\}$ by the well-ordered positive support
of $c\eps$.
\end{proof}

\begin{warning}[Exactly what the slow-time theorem licenses]
\label{dyn:warn:slowtime}
At an ordinary time $s_0\in\Sigma$, evaluating
\eqref{dyn:eq:slow-solution} produces an element with leading ordinary part
$\phi(s_0,z)$ and positive Hahn corrections, the ordinary coefficient
integrals having \emph{already been performed}: that is a separate analytic
operation from an unregrouped Hahn exponential. For example, with $X=z^{2}$
and $W=\eps z^{2}$ the slow flow is the ordinary rational function
$\phi(s,z)=z/(1-sz)$ on any product chart avoiding $sz=1$, whereas at a
nonzero ordinary $s$ the raw formal terms of
$\exp(sz^{2}\partial_z)z=\sum_{n\ge0}s^{n}z^{n+1}$ all have Hahn exponent
$0$ with infinitely many nonzero, and are \emph{not} a strongly summable
Hahn family. The rational function comes from ordinary complex analysis on
its chart, not from declaring that family a Hahn sum.
\Cref{dyn:thm:slowtime} therefore licenses one particular passage to the
time scale $\eps^{-1}$ --- solve the leading ordinary motion, then
integrate each positive correction along it --- and licenses \textbf{no}
arbitrary exchange of ordinary with Hahn summation, \textbf{no} global
continuation past an ordinary singularity of the supplied chart, and
\textbf{no} interpretation of an arbitrary surreal-valued time as an actual
continuous trajectory. It is also explicitly \emph{not} analytic
convergence of the modified series in a real step size.
\end{warning}

\section{Hamiltonian normal forms, integrability, and certified domains}
\label{dyn:sec:hamiltonian}

This section is under $\lnot\fres$ and, according to the case,
$\sdz$ or $\sdfr$. The frequencies are \emph{ordinary complex constants},
never surreal or Hahn quantities; $\Gamma$ needs neither divisibility nor
algebraic closure; and the two coefficient algebras
\begin{equation}\label{dyn:eq:coeffalgebra}
 \PP_n=\C[z_1,\dots,z_n],\qquad \EE_n=\OO(\C^{n}),\qquad
 \HH_\Gamma(\PP_n),\quad \HH_\Gamma(\EE_n)
\end{equation}
are the load-bearing distinction: \textbf{no} uniform norm bound as
$\gamma$ varies in the entire case, \textbf{no} uniform degree bound in the
polynomial case, and results about actual finite-halo evaluation are
\textbf{not} transferred to $\HH_\Gamma(\C\{z\})$, whose coefficients are
unrelated convergent germs. In $2d$ phase variables
$q_1,\dots,q_d,p_1,\dots,p_d$ use the fixed sign convention
\begin{equation}\label{dyn:eq:pb}
 \{F,G\}=\sum_{j=1}^{d}(F_{q_j}G_{p_j}-F_{p_j}G_{q_j}),
 \qquad \ad_\chi F=\{\chi,F\},
\end{equation}
and write $I_j=q_jp_j$,
\begin{equation}\label{dyn:eq:H0intro}
 H_0=\sum_{j=1}^{d}\omega_jI_j,\qquad
 \omega\cdot k\ne0\quad(0\ne k\in\Z^{d}).
\end{equation}

\subsection{Supported substitutions and Lie maps}

\begin{proposition}[Positive near-identity maps]\label{dyn:prop:inverse}
Let $U\in\HH_\Gamma(\EE_n)_{>0}^{n}$. Then $\Phi(z)=z+U(z)$ has a unique
inverse $z+V(z)$ with $V\in\HH_\Gamma(\EE_n)_{>0}^{n}$; both are bijections
of $\OO_\Gamma^{n}$ preserving standard parts; and if $U$ has polynomial
coefficients, so does $V$.
\end{proposition}
\begin{proof}
Let $S$ be the union of the supports of the components of $U$ and
$\mathsf M=M(S)$. The inverse equation is $V(z)=-U(z+V(z))$; expanding the
right side by Taylor's formula, its term independent of $V$ is $-U(z)$ and
every other term has at least one positive input from $U$ and at least one
from $V$, so at exponent $\gamma>0$ all occurrences of coefficients of $V$
have exponents strictly smaller than $\gamma$, and recursion on the
well-ordered $\mathsf M$ determines $V_\gamma$ uniquely --- including at
limit positions of $\mathsf M$, where \emph{no topological limit is
requested}. At a fixed exponent only finitely many Taylor terms and input
words occur by \cref{dyn:lem:neumann}, their coefficients using only
derivatives, products and finite sums of the coefficients of $U$, hence
entire, or polynomial in the polynomial case. This gives a right inverse;
the same construction for a left inverse, with associativity of supported
substitution or the least-exponent uniqueness argument, shows the two
agree, and \cref{dyn:prop:eval} gives the actual bijections.
\end{proof}

\begin{proposition}[Supported Hamiltonian exponential]\label{dyn:prop:exp}
If $\chi\in\HH_\Gamma(\EE_{2d})_{>0}$ then for every
$F\in\HH_\Gamma(\EE_{2d})$ the family $(\ad_\chi^{n}F/n!)_{n\ge0}$ is
strongly summable; its sum $T_\chi F=e^{\ad_\chi}F$ is an algebra and
Poisson automorphism with inverse $T_{-\chi}$; and with
$\Phi=(T_\chi q_1,\dots,T_\chi q_d,T_\chi p_1,\dots,T_\chi p_d)$ one has
$T_\chi F=F\circ\Phi$, the realized $\Phi$ being a symplectic bijection of
the finite halo. Polynomial coefficients are preserved.
\end{proposition}
\begin{proof}
With $A=\supp F$ and $S=\supp\chi$, each term of $\ad_\chi^{n}F$ has
support in $A+nS$, differentiation does not change Hahn weights, and a
fixed exponent receives contributions for finitely many words in $S$
attached to an exponent in $A$: \cref{dyn:lem:neumann} gives strong
summability and every regrouping used below. Leibniz gives
$\ad_\chi^{n}(FG)=\sum_j\binom nj(\ad_\chi^{j}F)(\ad_\chi^{n-j}G)$, so
after division by $n!$ and supported summation $T_\chi$ is multiplicative;
the Jacobi identity gives the identical binomial rule for the bracket, so
$T_\chi\{F,G\}=\{T_\chi F,T_\chi G\}$; and the binomial identity for
$\exp(\mathscr E)\exp(-\mathscr E)$ gives the inverse. For an ordinary
entire function the formal chain rule for the derivation $\ad_\chi$
identifies $T_\chi F$ with Taylor substitution by $\Phi$ --- check this at
each power of an auxiliary formal scalar multiplying $\chi$, where it is
the usual formal flow identity, then evaluate that scalar at $1$ using the
proved word summability --- and applying it coefficientwise gives the
assertion for general $F$. The coordinate Poisson relations are preserved,
proving symplecticity, and $\Phi=\id+$ positive terms, so
\cref{dyn:prop:inverse,dyn:prop:eval} give the actual maps.
\end{proof}

\begin{definition}[The finite-expression certificate]\label{dyn:def:trees}
Label each input perturbation coefficient occurrence by its positive Hahn
exponent. A term in a normalizing coefficient is built by a \emph{finite
expression tree} whose leaves are labeled inputs and whose internal
operations are Poisson brackets, linear projections, scalar division, or
fixed linear solution operators; its weight is the sum of its leaf weights.
For a fixed total weight there are finitely many leaf words by
\cref{dyn:lem:neumann}; every genuinely nonlinear internal step has at
least two positive inputs, linear operators only decorating those steps or
the leaves, so a fixed leaf word admits only finitely many relevant tree
shapes. This is a finite certificate for every coefficient even when the
set of smaller weights is infinite, and it is also the reason a later
reordering of weights under rescaling is legitimate
(\cref{dyn:thm:rescale}).
\end{definition}

\subsection{The homological operator and a sharp divisor threshold}
Define $LF=\{F,H_0\}$, so that
\begin{equation}\label{dyn:eq:hamhomological}
 L(q^{\alpha}p^{\beta})=\omega\cdot(\alpha-\beta)\,q^{\alpha}p^{\beta}
\end{equation}
and $\lnot\fres$ makes the kernel exactly the monomials with
$\alpha=\beta$. For a Taylor expansion at the origin set
\begin{equation}\label{dyn:eq:projection}
 \Pi\Bigl(\sum_{\alpha,\beta}c_{\alpha\beta}q^{\alpha}p^{\beta}\Bigr)
 =\sum_\alpha c_{\alpha\alpha}I^{\alpha},\qquad \Pi+Q_\Pi=1 .
\end{equation}

\begin{lemma}[Entire action projection]\label{dyn:lem:projection}
$\Pi$ maps $\EE_{2d}$ to functions $g(I)$ with $g\in\EE_d$, uniquely; it
preserves polynomial coefficients, does not decrease phase vanishing order,
and commutes with $L$. On $Q_\Pi\PP_{2d}$ the inverse $L^{-1}$ is
\emph{always} defined and polynomial, with no arithmetic estimate.
\end{lemma}
\begin{proof}
For entire $f$, Cauchy's estimate on a phase polydisk of radius $R$ gives
$|c_{\alpha\beta}|\le C_RR^{-|\alpha|-|\beta|}$, so
$|c_{\alpha\alpha}|\le C_R(R^{2})^{-|\alpha|}$; as $R$ is arbitrary the
action series is entire, and distinct action monomials are distinct phase
monomials, giving uniqueness. The remaining assertions follow termwise, and
for a polynomial the inverse divides finitely many nonresonant coefficients
by nonzero complex constants.
\end{proof}

One may also realize $\Pi$ by averaging over the ordinary torus action
$q_j\mapsto e^{i\theta_j}q_j$, $p_j\mapsto e^{-i\theta_j}p_j$; \emph{no}
surcomplex contour integration is needed for that observation.

\begin{theorem}[Sharp homological criterion]\label{dyn:thm:homological}
For nonresonant $\omega$ the following are equivalent:
(i) $\Theta(\omega)<\infty$; (ii) the coefficientwise $L^{-1}$ maps every
nonresonant \emph{entire} function to an entire function; (iii) the
coefficientwise $L^{-1}$ maps every nonresonant convergent germ at the
origin to a convergent germ, with a possible loss of radius. If these fail,
there is a nonresonant entire $F$ of phase order at least three for which
$L^{-1}F$ has no positive radius of convergence.
\end{theorem}
\begin{proof}
If $M_D\le CA^{D}$ and the coefficients of $F$ obey
$|f_{\alpha\beta}|\le C_RR^{-D}$, those of $L^{-1}F$ obey
$|(L^{-1}F)_{\alpha\beta}|\le CC_R(A/R)^{D}$; the number of monomials of
total degree $D$ in $2d$ variables is $\binom{D+2d-1}{2d-1}$, which grows
polynomially, so this bound proves absolute convergence on every polydisk
of radius less than $R/A$ --- one admissible $R$ for a germ, arbitrarily
large $R$ for an entire function. Conversely if $\Theta=\infty$, select
distinct degrees $D_j\to\infty$ with $D_j\ge\max(j,3)$ and nonresonant
monomials $m_j=q^{\alpha_j}p^{\beta_j}$ of those degrees with
$\delta_j=|\omega\cdot(\alpha_j-\beta_j)|$ and
$\delta_j^{-1/D_j}\ge j^{4}$, and put
\begin{equation}\label{dyn:eq:sparseF}
 F=\sum_j\delta_j^{1/2}m_j .
\end{equation}
On a polydisk of radius $R$ the $j$-th term has modulus at most
$(R/j^{2})^{D_j}$ and the tail is bounded by a geometric series since
$D_j\ge j$, so $F$ is entire; in the putative inverse the absolute
coefficient of $m_j$ is $\delta_j^{-1/2}$, whose $D_j$-th root is at least
$j^{2}$, violating every Cauchy bound on a positive-radius polydisk. Since
$F$ is itself a germ, both (ii) and (iii) fail, and its order and
nonresonance were built into the selection.
\end{proof}

\begin{warning}[What $\Theta$ measures, and what it must not be advertised as]
\label{dyn:warn:theta}
A polynomial lower bound on $|\omega\cdot k|$ implies $\Theta=1$ but is far
stronger than necessary here. Finite exponential growth is the exact
condition for this \emph{single coefficientwise homological operator} to
preserve all entire functions. The nonlinear Hahn theorem below applies
that operator repeatedly and requires \emph{no} uniform norm estimate
across Hahn layers; this is precisely why the criterion must
\textbf{not} be advertised as an ordinary complex-analytic convergence
criterion of Bruno or KAM type. Sources 06 repeats this disclaimer in three
separate places, and the merge preserves it.
\end{warning}

\subsection{Arbitrary-support Hamiltonian normalization}

\begin{theorem}[Exact positive-Hahn normal form]\label{dyn:thm:normalform}
Let $\omega$ satisfy \eqref{dyn:eq:H0intro} and let
\begin{equation}\label{dyn:eq:Hinput}
 H=H_0+R,\qquad R=\sum_{s\in S}t^{s}R_s(q,p),\qquad
 S\subseteq\Gamma_{>0}\text{ well ordered},\quad \ord_0R_s\ge3 .
\end{equation}
Assume either
\textbf{(a)} every $R_s$ is a polynomial, with \emph{no} restriction on
$\omega$ beyond nonresonance (tag $\sdz$), or
\textbf{(b)} every $R_s$ is entire and $\Theta(\omega)<\infty$ (tag
$\sdfr$). Then there is a \emph{unique} positive generator $\chi$ in the
corresponding coefficient algebra with
\begin{equation}\label{dyn:eq:gauge}
 \Pi\chi=0,\qquad \ord_0\chi_\gamma\ge3,\qquad
 e^{\ad_\chi}H=H_0+B(I),
\end{equation}
where $B$ has positive Hahn support and coefficients of action vanishing
order at least two, and $\supp\chi,\supp B\subseteq M(S)\setminus\{0\}$.
The coordinate map $\Phi=e^{\ad_\chi}(q,p)$ and its inverse are symplectic
bijections of $\OO_\Gamma^{2d}$, fixing the origin, tangent to the identity
there, preserving standard parts everywhere, and satisfying the actual
identity
\begin{equation}\label{dyn:eq:actualnormal}
 H^{\#}\circ\Phi^{\#}=\bigl(H_0+B(I)\bigr)^{\#}.
\end{equation}
All statements commute with enlargement of the Hahn workspace.
\end{theorem}
\begin{proof}
Expanding the Lie exponential,
\begin{equation}\label{dyn:eq:normalexpansion}
 e^{\ad_\chi}(H_0+R)=H_0+R+L\chi+\mathcal B(\chi,R),\qquad
 \mathcal B(\chi,R)=\sum_{n\ge1}\frac{\ad_\chi^{n}R}{n!}
 +\sum_{n\ge2}\frac{\ad_\chi^{n}H_0}{n!},
\end{equation}
every term of $\mathcal B$ involving at least two positive factors,
counting occurrences of $R$ and $\chi$. Since $\Pi L=0$ and $\Pi\chi=0$,
the nonresonant normal-form equation is exactly
\begin{equation}\label{dyn:eq:chifixed}
 \chi=-L^{-1}Q_\Pi\bigl(R+\mathcal B(\chi,R)\bigr).
\end{equation}
Put $\mathsf M=M(S)$: at weight $\gamma\in\mathsf M\setminus\{0\}$ the
coefficient of $\mathcal B$ depends only on $\chi_\mu$ with
$0<\mu<\gamma$, since an occurrence of $\chi_\gamma$ would be accompanied
by another positive factor and so have total weight greater than $\gamma$.
Well-founded recursion on $\mathsf M$ therefore determines
\begin{equation}\label{dyn:eq:coefficientrecursion}
 \chi_\gamma=-L^{-1}Q_\Pi\bigl(R_\gamma
 +[t^{\gamma}]\mathcal B(\chi_{<\gamma},R)\bigr),
\end{equation}
a formula that remains meaningful when $\gamma$ has infinitely many
predecessors, because \cref{dyn:lem:neumann} allows only finitely many of
them to participate in a word of total weight $\gamma$. More explicitly,
each coefficient is a finite sum of the labeled trees of
\cref{dyn:def:trees}: their leaves are input coefficients $R_s$, their
total weight is $\gamma$, and their internal nonlinear nodes have at least
two positive inputs, so the finitely many words of total weight $\gamma$
bound the number and choice of leaves and there are finitely many tree
shapes and bracket orderings. Hence every coefficient uses finitely many
applications of derivatives, multiplication, $Q_\Pi$ and $L^{-1}$: the
claimed support bound holds and no unproved infinite operation occurs
within a Hahn coefficient. In case (a) these operations preserve
polynomials by \cref{dyn:lem:projection}; in case (b) they preserve entire
functions by \cref{dyn:thm:homological}. Every application of
$L^{-1}Q_\Pi$ gives $\Pi\chi_\gamma=0$, and the estimate
$\ord_0\{F,G\}\ge\ord_0F+\ord_0G-2$ shows inductively that each
$\chi_\gamma$ has order at least three while terms in $\mathcal B$ have
order at least four; hence $B=\Pi(R+\mathcal B(\chi,R))$ has phase order at
least four, since a resonant monomial has even phase degree and $R$ has
order at least three, and by \cref{dyn:lem:projection} it is an entire or
polynomial function of the actions at every Hahn weight, of action order at
least two. Equation \eqref{dyn:eq:chifixed} removes the nonresonant part of
\eqref{dyn:eq:normalexpansion}, so \eqref{dyn:eq:gauge} holds as a Hahn
identity, and \cref{dyn:prop:exp} defines the exponential. For uniqueness,
at the least exponent $\gamma$ of the difference of two positive
gauge-normalized generators, all terms coming from $\mathcal B$ have weight
strictly greater than $\gamma$, leaving
$L(\chi_\gamma-\widetilde\chi_\gamma)=0$ with a nonresonant difference, on
which $L$ is injective by \eqref{dyn:eq:hamhomological}; this also rules
out additional support outside $\mathsf M$. Finally
\cref{dyn:prop:exp,dyn:prop:eval} give the symplectic maps, their inverse,
and \eqref{dyn:eq:actualnormal}, and no construction used an exponent
outside $M(S)$.
\end{proof}

\begin{remark}[What has become exact]\label{dyn:rem:becameexact}
In a classical analytic problem the generator can be formal in the phase
variables with zero radius of convergence. In case (a) each Hahn
coefficient of the generator is instead a \emph{finite polynomial}: its
coefficients may be enormous ordinary complex numbers, but they are finite
constants at their Hahn layer, and an arbitrarily large ordinary number
cannot cancel the positive valuation of $t^{\gamma}$
(\cref{dyn:rem:invisible}). Evaluation on the finite halo is therefore
controlled by supports rather than by a numerical norm of the normalizing
coefficients. In case (b) infinitely many phase monomials occur within a
single Hahn layer, so their ordinary convergence must be settled
\emph{before} Hahn summability can be invoked; the homological threshold
does exactly that, and requires no common numerical bound on the resulting
entire coefficients because the common domain is already all of
$\C^{2d}$.
\end{remark}

\begin{corollary}[Finite surcomplex phase space]\label{dyn:cor:hamsurcomplex}
For polynomial layers, and for entire layers satisfying the threshold,
\cref{dyn:thm:normalform} realizes the normalization at every finite point
of $\No[i]^{2d}$, and the maps obtained in different set-sized workspaces
agree on their common domain.
\end{corollary}
\begin{proof}
Given a phase tuple, enlarge the workspace to contain all of its exponents
and those of the input and apply \cref{dyn:thm:normalform}; its coefficient
formulas and the Taylor evaluation are unchanged by the enlargement.
\end{proof}

\begin{warning}[Uniqueness has a gauge]
\label{dyn:warn:gauge}
The unique object in \cref{dyn:thm:normalform} is the \emph{single
exponential} generator with $\Pi\chi=0$, and the normal form it produces in
that gauge. Arbitrary symplectic changes of variables normalizing a
Hamiltonian need \textbf{not} be unique: in particular transformations
generated by functions of the actions can preserve a normal form. Source 06
declines to suppress this freedom by an unstated uniqueness claim, and
\cref{dyn:app:algorithm} distinguishes the single-exponential gauge
recursion from the different algorithm that composes successive
exponentials without taking a logarithm.
\end{warning}

\subsection{Exact integrability and the whole Poisson centralizer}
Write $T=e^{\ad_\chi}$, $N=TH$, and $J_j=T^{-1}I_j=e^{-\ad_\chi}I_j$.

\begin{theorem}[Exact integrability]\label{dyn:thm:integrability}
The $J_1,\dots,J_d$ are entire-coefficient Hahn functions --- with
polynomial coefficients in case (a) --- and
\begin{equation}\label{dyn:eq:integrals}
 \{H,J_j\}=0,\qquad \{J_j,J_k\}=0,\qquad H=N(J_1,\dots,J_d).
\end{equation}
They are functionally independent wherever, in the normalized coordinates,
none of the pairs $(q_j,p_j)$ is $(0,0)$. These hold as realized identities
on the finite halo.
\end{theorem}
\begin{proof}
The $I_j$ commute with one another and with every function of the actions,
including $N$; applying the Poisson automorphism $T^{-1}$ gives the first
two identities and the algebra substitution $T^{-1}$ applied to $N(I)$ the
third. The $d$ differentials $\mathrm dI_j=p_j\mathrm dq_j+q_j\mathrm dp_j$
have disjoint coordinate pairs and are linearly independent exactly when
each is nonzero; pullback by the invertible differential of $\Phi^{-1}$
preserves that rank, and \cref{dyn:prop:eval} respects all these
identities.
\end{proof}

\begin{theorem}[Poisson centralizer]\label{dyn:thm:pcentralizer}
In $\HH_\Gamma(\EE_{2d})$,
\begin{equation}\label{dyn:eq:pcentralizer}
 \Cent(H):=\{F:\{F,H\}=0\}
 =\{G(J_1,\dots,J_d):G\in\HH_\Gamma(\EE_d)\},
\end{equation}
the representing $G$ being unique. Its Hahn support need \textbf{not} be
positive. Thus \eqref{dyn:eq:pcentralizer} describes all first integrals
that admit an entire-coefficient Hahn representation on the finite halo.
\end{theorem}
\begin{proof}
It suffices to classify the centralizer of $N(I)=H_0+B(I)$, since $T$ is a
Poisson automorphism. Let $\widetilde F=TF$ with
$\{\widetilde F,N\}=0$ and apply $\Pi$ coefficientwise: by
\cref{dyn:lem:projection}, $\Pi\widetilde F=G(I)$ for a unique
$G\in\HH_\Gamma(\EE_d)$, and every action function commutes with $N$, so
$U:=\widetilde F-\Pi\widetilde F$ is coefficientwise nonresonant and also
commutes with $N$. If $U\ne0$, let $a_0$ be the least exponent of its
support: the coefficient at $t^{a_0}$ in $\{U,B\}$ is zero, all exponents
of $B$ being positive and all of $U$ at least $a_0$, so the coefficient at
$t^{a_0}$ in $\{U,N\}=0$ is $LU_{a_0}=0$ with $U_{a_0}$ nonresonant and
nonzero, contradicting \eqref{dyn:eq:hamhomological}. Hence $U=0$ and
$\widetilde F=G(I)$, so $F=G(J)$. Conversely every $G(J)$ commutes with
$H$; its definition is legitimate even for arbitrary well-ordered support
of $G$ --- write $J=I+\Delta$ with positive $\Delta$, use the ordinary
Taylor expansion of each entire coefficient of $G$ at $I$, and apply
\cref{dyn:lem:neumann}. Uniqueness follows by applying $T$ or by the
injectivity of $g(I)\mapsto g$; faithfulness of \cref{dyn:prop:eval} means
no additional identity appears on passing to finite-halo functions.
\end{proof}

\begin{warning}[The coefficient class is essential]\label{dyn:warn:pcent}
\Cref{dyn:thm:pcentralizer} does \textbf{not} classify first integrals
defined only near a particular point, multivalued functions, or arbitrary
functions that happen to be locally constant in a fine topology, and does
not say that all first integrals have polynomial coefficients. The
entire-coefficient class is large enough to contain arbitrary entire
functions of the actions and rigid enough for the leading nonresonant
coefficient argument to work; that is exactly the balance the theorem
exploits.
\end{warning}

\begin{proposition}[Exact ordinary-time evolution]\label{dyn:prop:flow}
Let $\Omega_j(I)=\partial N/\partial I_j=\omega_j+\Delta_j(I)$ with
$v(\Delta_j)>0$ (positive Hahn support). For ordinary complex $\tau$ put
\begin{equation}\label{dyn:eq:qflow}
 q_j(\tau)=e^{\tau\omega_j}\ExpH(\tau\Delta_j(I))q_j(0),\qquad
 p_j(\tau)=e^{-\tau\omega_j}\ExpH(-\tau\Delta_j(I))p_j(0).
\end{equation}
These give an exact flow of $N$ on the finite halo for every $\tau\in\C$;
conjugation by $\Phi$ gives the flow of $H$; the coefficients depend
holomorphically on $\tau$; and the flow law holds as an identity of
Hahn-coherent functions.
\end{proposition}
\begin{proof}
The actions are constant along \eqref{dyn:eq:qflow} because the two factors
for each $j$ are inverse, so $\Delta_j(I)$ is constant along the curve, and
coefficientwise differentiation in ordinary time gives
$\dot q_j=\Omega_j(I)q_j=\partial N/\partial p_j$ and
$\dot p_j=-\Omega_j(I)p_j=-\partial N/\partial q_j$. The Hahn exponentials
are strongly summable by positivity, each coefficient is entire in $\tau$
and the phase variables, and their constant Hahn terms are nonzero ordinary
exponentials, so finite phase points remain finite; the exponential
identities, justified coefficientwise by finite Hahn regroupings, give the
flow law, and symplectic conjugacy transports the equations to $H$.
\end{proof}

\begin{warning}[This is a coefficientwise holomorphic-time identity]
\label{dyn:warn:hamflow}
\Cref{dyn:prop:flow} is not an assertion about continuous ordinary-time
paths in the fine topology of the proper-class field. No infinite surreal
time, no global surreal exponential of an arbitrary surcomplex argument,
and no intrinsic coefficient-field derivation is used
(\cref{dyn:warn:derivative}). ``Analytic'' in this section means exactly
the entire-coefficient Taylor-compatible realization of
\cref{dyn:prop:eval}.
\end{warning}

\subsection{Failure for entire layers: an exact converse}

\begin{theorem}[Sharp universal normalization criterion]\label{dyn:thm:universal}
For fixed nonresonant $\omega$ the following are equivalent:
(i) $\Theta(\omega)<\infty$; (ii) \emph{every} positive Hahn perturbation
with entire coefficients of phase order at least three admits the
entire-coefficient gauge-normalized generator and finite-halo realization
of \cref{dyn:thm:normalform}, for \emph{every} set-sized ordered value
group; (iii) every one-parameter Hamiltonian $H_0+\eps_0F$ with $F$ entire
of order at least three admits a symplectic formal change
$\Psi=\id+O(\eps_0)$ with \emph{convergent-germ} coefficients at the origin
transforming it into a function of the actions. If (i) fails, failure of
(iii) is already forced at the coefficient of $\eps_0$, and there is no
positive gauge-normalized entire Hahn generator for the same input in any
workspace containing its scale.
\end{theorem}
\begin{proof}
(i)$\Rightarrow$(ii) is \cref{dyn:thm:normalform}; (ii)$\Rightarrow$(iii)
follows by taking $\Gamma=\Z$ and $\eps_0=t$. Suppose $\Theta=\infty$ and
take the entire nonresonant $F$ of \cref{dyn:thm:homological}. If a
symplectic change as in (iii) exists, write
$\Psi(z)=z+\eps_0V(z)+O(\eps_0^{2})$ with $V$ a holomorphic vector-field
germ; the first-order symplectic condition says $V$ is a symplectic vector
field, so its contraction with the constant symplectic form is a closed
holomorphic one-form, and a closed holomorphic one-form germ is exact (a
primitive is given on a small star-shaped polydisk by the radial homotopy
integral). Hence, with the sign matching \eqref{dyn:eq:pb}, there is a
holomorphic germ $\chi_1$ with $V(G)=\{\chi_1,G\}$ for every holomorphic
germ $G$. The coefficient of $\eps_0$ in
$(H_0+\eps_0F)\circ\Psi$ is $F+L\chi_1$, and if the transformed Hamiltonian
is a function of the actions its nonresonant part must vanish:
$L(Q_\Pi\chi_1)=-F$. The projection $Q_\Pi$ preserves convergent germs by
the same Cauchy estimate as in \cref{dyn:lem:projection}, and
coefficientwise injectivity of $L$ on nonresonant series forces
$Q_\Pi\chi_1=-L^{-1}F$, which is not a convergent germ. For the final
assertion, a putative positive gauge-normalized entire Hahn generator for
$H_0+t^{\eta}F$ cannot have a least nonzero coefficient below $\eta$ (its
normal-form equation there would be $L\chi_\gamma=0$, impossible in the
gauge), and at weight $\eta$ the nonlinear terms have larger weight, so
again $\chi_\eta=-L^{-1}F$, not entire. This argument does not depend on
the rank or divisibility of the enlarged value group.
\end{proof}

\begin{proposition}[An explicit super-Liouville obstruction]
\label{dyn:prop:explicitbad}
Define integers by
\begin{equation}\label{dyn:eq:Liouv}
 a_1=1,\quad a_{n+1}=10^{4a_n},\quad Q_n=10^{a_n},\quad
 \alpha=\sum_{n\ge1}10^{-a_n},
\end{equation}
put $P_n/Q_n=\sum_{j=1}^{n}10^{-a_j}$ (so $P_n\in\Z$) and
$D_n=P_n+Q_n<2Q_n$. Then
\begin{equation}\label{dyn:eq:Liouvbound}
 0<\alpha Q_n-P_n<2Q_n10^{-a_{n+1}}=2Q_n10^{-Q_n^{4}}<e^{-D_n^{2}},
\end{equation}
$\alpha$ is irrational, and $\omega=(1,\alpha)$ is nonresonant. In two
degrees of freedom,
\begin{equation}\label{dyn:eq:explicitF}
 F(q_1,q_2,p_1,p_2)=\sum_{n\ge1}e^{-D_n^{2}/2}q_1^{P_n}p_2^{Q_n},
 \qquad H=I_1+\alpha I_2+\eps_0F,
\end{equation}
where $F$ is entire of phase order at least three ($D_1=11$), the
Hamiltonian has \emph{no} symplectic normalization $\id+O(\eps_0)$ with
convergent-germ coefficients to an action-only normal form, and its first
nonresonant homological solution has zero radius of convergence.
\end{proposition}
\begin{proof}
The first bound in \eqref{dyn:eq:Liouvbound} dominates the very rapidly
decreasing tail by twice its first term; for the last, $Q_n\ge10$ and
$\log(2Q_n)-Q_n^{4}\log10<-4Q_n^{2}<-D_n^{2}$. If $\alpha=A/B$ were
rational then $\alpha Q_n-P_n$ would be at least $1/B$ for every $n$,
contradicting \eqref{dyn:eq:Liouvbound}; hence $\omega$ is nonresonant. On
a phase polydisk of radius $R$ the $n$-th monomial of $F$ is bounded by
$e^{-D_n^{2}/2}R^{D_n}$ and the distinct integers $D_n$ tend to infinity,
so these bounds are summable for every finite $R$ and $F$ is entire. By
\eqref{dyn:eq:hamhomological} the homological divisor of the $n$-th
monomial is $P_n-\alpha Q_n$, so the absolute coefficient after division
exceeds $e^{D_n^{2}/2}$ by \eqref{dyn:eq:Liouvbound} and its $D_n$-th root
tends to infinity, excluding every positive analytic radius; the
first-order argument of \cref{dyn:thm:universal} then applies.
\end{proof}

\begin{remark}[The obstruction is one homological inverse, not a long
iteration]\label{dyn:rem:oneinverse}
Both $F$ in \eqref{dyn:eq:explicitF} and its formal homological solution
depend only on $q_1$ and $p_2$, and any two functions of those variables
have zero Poisson bracket. Therefore in the purely formal phase category
$\chi=-\eps_0L^{-1}F$ removes the perturbation with \emph{no} higher Lie
corrections at all. The obstruction is thus not divergence of a long
nonlinear iteration; it is already the failure of a single homological
inverse to preserve the specified analytic coefficient class. This
observation is specific to source 06.
\end{remark}

\begin{remark}[Why polynomial layers escape the same example]
\label{dyn:rem:whypoly}
Assign the same monomials distinct increasing positive Hahn weights, for
instance
$\widehat H=H_0+\sum_{n\ge1}t^{n}e^{-D_n^{2}/2}q_1^{P_n}p_2^{Q_n}$: then
every Hahn layer is a polynomial and \cref{dyn:thm:normalform}(a) applies,
despite the \emph{same} frequency vector, because dividing an individual
coefficient by its tiny ordinary divisor still gives an ordinary complex
number rather than a failure of a phase Taylor series within that layer.
This is no contradiction: $\widehat H$ is a \emph{different} Hamiltonian,
with a different distribution of monomials among the infinitesimal scales,
and source 06 records explicitly that ``which scale carries which
monomial'' is an input choice with mathematical consequences and not a
physical fact. The divergent formal generator of
\cref{dyn:prop:explicitbad} also remains evaluable at infinitesimal phase
points --- its ordinary complex Taylor coefficients are summed against
positive phase supports by \cref{dyn:lem:neumann} --- so the counterexample
does not exclude normalization on an infinitesimal monad; its obstruction
is to the ordinary analytic coefficient \emph{germs} required by the global
finite-halo theorem. This is the Hamiltonian instance of
\cref{dyn:prop:monad}.
\end{remark}

\subsection{Support-and-degree control beyond the finite halo}
\label{dyn:subsec:domains}

Assume polynomial coefficients in \eqref{dyn:eq:Hinput}, split each nonzero
coefficient into its finitely many monomials, and regard a monomial
occurrence as an input \emph{label}
$\ell=(\gamma,\alpha,\beta)$ with coefficient $c_\ell\ne0$ and
$d_\ell=|\alpha|+|\beta|\ge3$. The label set is a set; infinitely many
labels are allowed. For $\rho\in\Gamma$ put
\begin{equation}\label{dyn:eq:weights}
 w_\rho(\ell)=\gamma+(d_\ell-2)\rho,\qquad
 D_\rho=(t^{\rho}\OO_\Gamma)^{2d} .
\end{equation}

\begin{definition}[Admissible rescaling]\label{dyn:def:admissible}
The input monomial family is \emph{admissible at $\rho$} if
(1) $w_\rho(\ell)>0$ for every label; (2) the set
$W_\rho=\{w_\rho(\ell)\}$ is well ordered; and (3) each weight in $W_\rho$
comes from only finitely many labels. These are exactly strong summability
of the individually rescaled monomials into a positive Hahn series with
polynomial coefficients.
\end{definition}

The point $\rho=0$ is always admissible; a \emph{negative} $\rho$ gives a
domain containing phase points of \emph{infinite} magnitude. For infinitely
many input monomials, condition (1) alone implies neither (2) nor (3)
(\cref{dyn:ex:badorder,dyn:ex:badfiber}). The Hamiltonian rescaling
operator is
\begin{equation}\label{dyn:eq:S}
 \mathscr S_\rho F(z)=t^{-2\rho}F(t^{\rho}z),
\end{equation}
which is a Poisson \emph{Lie} algebra homomorphism wherever its monomial
expansion is summable,
\begin{equation}\label{dyn:eq:scalepb}
 \{\mathscr S_\rho F,\mathscr S_\rho G\}=\mathscr S_\rho\{F,G\},
\end{equation}
but is \textbf{not} a homomorphism for ordinary products: the factor
$t^{-2\rho}$ is exactly what compensates the two phase derivatives in a
bracket. It fixes $H_0$, preserves resonant monomials, and commutes with
$L$ and $\Pi$.

\begin{theorem}[Compatible support-certified domains]\label{dyn:thm:rescale}
Suppose the polynomial input of \cref{dyn:thm:normalform} is admissible at
$\rho$. Then its canonical normalizing map and inverse have strongly
summable monomial realizations on $D_\rho$ and are mutually inverse
symplectic bijections of that domain, and the Hamiltonian identity and the
commuting action identities hold there. If the normalizing construction is
performed independently after rescaling at two admissible values, the
realized maps agree on the overlap after undoing the scalings: these
domains carry \emph{one} compatible normalization, not a choice of
unrelated local conjugacies.
\end{theorem}
\begin{proof}
Admissibility makes
$H_\rho:=\mathscr S_\rho H
=H_0+\sum_\ell c_\ell t^{w_\rho(\ell)}q^{\alpha_\ell}p^{\beta_\ell}$
a positive Hahn perturbation with polynomial coefficients, so
\cref{dyn:thm:normalform}(a) gives $\chi_\rho,\Phi_\rho$ on the finite
halo, and undoing the phase scaling gives
\begin{equation}\label{dyn:eq:Phirescaled}
 \Psi_\rho(z)=t^{\rho}\Phi_\rho(t^{-\rho}z)\quad\text{of }D_\rho,
\end{equation}
symplectic in the original bracket because an equal scalar dilation in all
phase coordinates is conformally symplectic and the dilation and its
inverse cancel. That this is the \emph{same} normalization uses the labeled
trees of \cref{dyn:def:trees} with each input coefficient expanded into its
monomial labels: every nonzero monomial contribution to a generator or
Hamiltonian coefficient with leaves $\ell_1,\dots,\ell_r$ has phase degree
\begin{equation}\label{dyn:eq:degreeham}
 2+\sum_{j=1}^{r}(d_{\ell_j}-2),
\end{equation}
and every contribution to a nonlinear coordinate coefficient has phase
degree
\begin{equation}\label{dyn:eq:degreecoord}
 1+\sum_{j=1}^{r}(d_{\ell_j}-2),
\end{equation}
both inductively, since a Poisson bracket subtracts two from the sum of
phase degrees while $L$, $L^{-1}$ and $\Pi$ preserve degree, the coordinate
construction starting at degree one. Under Hamiltonian scaling the exponent
of the tree in \eqref{dyn:eq:degreeham} becomes precisely
$\sum_j\gamma_{\ell_j}+\rho\sum_j(d_{\ell_j}-2)=\sum_jw_\rho(\ell_j)$, and
for coordinates the same sum is obtained using
$t^{-\rho}\Phi(t^{\rho}z)$. Since $W_\rho$ is positive and well ordered,
\cref{dyn:lem:neumann} gives finitely many words at a fixed rescaled
weight, the finite-fiber hypothesis gives finitely many choices of original
labels per word, and there are then finitely many trees: the rescaled
generator, Hamiltonian and coordinate expressions are strongly summable
\emph{before} any regrouping. Hence \eqref{dyn:eq:scalepb} applies to every
finite tree and may then be summed; the rescaled generator is
$\mathscr S_\rho\chi$, satisfying the same nonresonant gauge and
normal-form equation as $\chi_\rho$, and uniqueness identifies them, so on
coordinates $\Phi_\rho(z)=t^{-\rho}\Phi(t^{\rho}z)$. The same argument
applies to the inverse exponential and to the actions, and Taylor
evaluation at finite rescaled points is \cref{dyn:prop:eval}. For two
admissible rescalings both maps arise from the same labeled trees, and at a
point of the overlap those tree evaluations are strongly summable by each
certificate with every regrouping agreeing coefficientwise --- which proves
compatibility \emph{without} assuming that rescaling preserves the order of
the original Hahn exponents; in general it does not.
\end{proof}

\begin{warning}[Scope of the extension]\label{dyn:warn:rescale}
\Cref{dyn:thm:rescale} extends the \emph{polynomial}-coefficient
normalizer, Hamiltonian and constructed actions only. It does
\textbf{not} license evaluating an arbitrary entire function at an infinite
phase or action argument; a centralizer statement on $D_\rho$ must instead
be made in the corresponding \emph{rescaled} entire-coefficient algebra,
and transporting all original entire coefficients to infinite arguments
without a summability certificate would be an unjustified enlargement of
the function class. The theorem gives \emph{guaranteed}, not maximal,
domains.
\end{warning}

\begin{corollary}[A degree-dependent radius]\label{dyn:cor:homogeneous}
Let $H=H_0+t^{\eta}P_m$ with $\eta>0$ and $P_m$ homogeneous of phase degree
$m\ge3$, and write $\chi=\sum_{n\ge1}t^{n\eta}\chi_n$,
$\Phi=z+\sum_{n\ge1}t^{n\eta}\Phi_n$. Every nonzero $\chi_n$ is homogeneous
of degree $2+n(m-2)$ and every nonzero component of $\Phi_n$ has degree
$1+n(m-2)$, and the normalizer and its inverse are defined on every
$D_\rho$ with
\begin{equation}\label{dyn:eq:homogeneousdomain}
 \eta+(m-2)\rho>0 .
\end{equation}
\end{corollary}
\begin{proof}
Every tree at exponent $n\eta$ has exactly $n$ leaves, an ordered abelian
group being torsion-free and $\eta$ the only input weight; apply
\eqref{dyn:eq:degreeham}--\eqref{dyn:eq:degreecoord}. The finitely many
input monomials all have the same shifted weight $\eta+(m-2)\rho$, whose
positivity supplies all three admissibility conditions.
\end{proof}

\subsection{A sharp cubic boundary, and two independent summability failures}

\begin{example}[The cubic boundary, with a coordinate-independent
obstruction]\label{dyn:ex:cubic}
In one degree of freedom take
\begin{equation}\label{dyn:eq:cubicH}
 H(q,p)=qp+\eps_0q^{2}p,\qquad \eps_0=t^{\eta},\ \eta>0,
\end{equation}
where there is no arithmetic issue at all because $\omega_1=1$. With
$\chi=-\eps_0q^{2}p$ one has $\ad_\chi q=\eps_0q^{2}$,
$\ad_\chi p=-2\eps_0qp$, and the coordinate exponentials give
\begin{equation}\label{dyn:eq:cubicmap}
 \Phi(q,p)=\Bigl(\frac{q}{1-\eps_0q},\ p(1-\eps_0q)^{2}\Bigr),
 \qquad H\circ\Phi=qp ,
\end{equation}
verifiable also by direct rational substitution; its Jacobian determinant
is one, so it is symplectic, and $\Pi\chi=0$ makes it the canonical
normalizer of \cref{dyn:thm:normalform}. For $\eta+\rho>0$ the expansion
$q/(1-\eps_0q)=\sum_{n\ge0}\eps_0^{n}q^{n+1}$ is strongly summable on
$D_\rho$; at the boundary point $q=\eps_0^{-1}$, $p=0$, every term of that
family has exponent $-\eta$ and the rational function has a \emph{pole}.
Hence replacing the strict inequality in
\eqref{dyn:eq:homogeneousdomain} by a non-strict one cannot be valid in
general.

There is also a \emph{coordinate-independent} obstruction on the whole
boundary polydisk $D_{-\eta}$: from $H_p=q(1+\eps_0q)$ and
$H_q=p(1+2\eps_0q)$, \eqref{dyn:eq:cubicH} has exactly two critical points
there, $(0,0)$ and $(-\eps_0^{-1},0)$, whereas $qp$ has exactly one. An
everywhere invertible symplectic change on that domain with
$H\circ\Psi=qp$ would give a bijection between these critical-point sets by
the chain rule, which is impossible. The boundary failure is therefore not
merely a defect in one geometric-series proof.
\end{example}

\begin{example}[Positive shifted weights, but no well-order]
\label{dyn:ex:badorder}
Over $\Gamma=\Q$ consider the valid polynomial-coefficient perturbation
$R(q,p)=\sum_{n\ge2}t^{n+1/n}q^{n+2}$, whose original exponents strictly
increase. At $\rho=-1$ the shifted weights are $1/n$, all strictly
positive, but forming a \emph{descending} sequence; at the rescaled point
$q=1$ the rescaled perturbation would be $\sum_{n\ge2}t^{1/n}$, which is
not a Hahn series. Positivity alone is therefore insufficient even to
\emph{define} the rescaled input.
\end{example}

\begin{example}[Positive well-ordered weights, but infinite fibers]
\label{dyn:ex:badfiber}
Take $R(q,p)=\sum_{n\ge2}t^{n+1}q^{n+2}$. At $\rho=-1$ all shifted weights
equal $1$: the weight set is positive and well ordered, but infinitely many
labels contribute to that one weight. The rescaled expression is
$t\sum_{n\ge2}q^{n+2}$, whose evaluation at the ordinary point $q=1$ is
defined neither by finite coefficient addition nor by ordinary convergence;
it cannot supply a polynomial-coefficient Hahn perturbation on the whole
rescaled finite halo.
\end{example}

\Cref{dyn:ex:badorder,dyn:ex:badfiber} show why none of the three
conditions of \cref{dyn:def:admissible} can be dropped from a general
theorem of this form. They do \textbf{not} assert that finite fibers are
necessary for every possible analytic resummation in some larger,
differently specified class.

\subsection{A genuinely multiscale worked example}
\label{dyn:subsec:multiscale}

Take $\omega=(1,\sqrt2)$, $\kappa=2-\sqrt2$, and
$A=q_1^{2}p_2$, $B_{\mathrm h}=p_1^{2}q_2$, $C=I_1I_2$. For arbitrary
positive scales $\eta,\theta\in\Gamma$ put
\begin{equation}\label{dyn:eq:workedH}
 H=I_1+\sqrt2I_2+u(A+B_{\mathrm h})+vC,\qquad u=t^{\eta},\ v=t^{\theta},
\end{equation}
nonresonance following from the irrationality of $\sqrt2$. The two scales
need \emph{not} be Archimedean-comparable; the lexicographic example of
\cref{dyn:warn:contraction} allows $n\eta<\theta$ for every $n$. Retaining
$u,v$ as independent \emph{labels}, and writing $O_{\mathrm{lab}}(3)$ for
terms involving at least three input labels --- \textbf{not} terms above a
valuation cutoff --- direct homological calculation gives
\begin{align}
 \chi&=\frac{u}{\kappa}(B_{\mathrm h}-A)
 +\frac{uv}{\kappa^{2}}(A-B_{\mathrm h})(2I_2-I_1)+O_{\mathrm{lab}}(3),
 \label{dyn:eq:workedchi}\\
 N&=I_1+\sqrt2I_2+vI_1I_2
 +\frac{u^{2}}{\kappa}(I_1^{2}-4I_1I_2)+O_{\mathrm{lab}}(3),
 \label{dyn:eq:workedN}
\end{align}
the coefficients of $v$, $u^{2}$ and $v^{2}$ in $\chi$ vanishing and those
of $uv$ and $v^{2}$ in the displayed normal form vanishing. The key bracket
is $\{A,B_{\mathrm h}\}=4I_1I_2-I_1^{2}$; since $LA=\kappa A$ and
$LB_{\mathrm h}=-\kappa B_{\mathrm h}$ the first generator is
$u(B_{\mathrm h}-A)/\kappa$, and the second normal term is
$\frac12\Pi\{u(B_{\mathrm h}-A)/\kappa,\,u(A+B_{\mathrm h})\}
=\frac{u^{2}}{\kappa}(I_1^{2}-4I_1I_2)$. Also
$\{A,C\}=A(2I_2-I_1)$ and $\{B_{\mathrm h},C\}=-B_{\mathrm h}(2I_2-I_1)$,
which gives the mixed generator by $-L^{-1}Q_\Pi$; these identities fix
both the sign convention and the placement of the inverse transformation.

To first label degree the coordinate map and actions are
\[
 \Phi(q_1)=q_1-\frac{2u}{\kappa}p_1q_2+O_{\mathrm{lab}}(2),\quad
 \Phi(p_1)=p_1-\frac{2u}{\kappa}q_1p_2+O_{\mathrm{lab}}(2),
\]
\[
 \Phi(q_2)=q_2+\frac{u}{\kappa}q_1^{2}+O_{\mathrm{lab}}(2),\quad
 \Phi(p_2)=p_2+\frac{u}{\kappa}p_1^{2}+O_{\mathrm{lab}}(2),
\]
\[
 J_1=I_1+\frac{2u}{\kappa}(A+B_{\mathrm h})+O_{\mathrm{lab}}(2),\qquad
 J_2=I_2-\frac{u}{\kappa}(A+B_{\mathrm h})+O_{\mathrm{lab}}(2),
\]
the full supported series satisfying $H=N(J)$ and commuting exactly, the
displayed remainders serving only to show a manageable initial segment.
There are two input degrees --- the $u$ terms cubic, the $v$ term quartic
--- so every generator or Hamiltonian coefficient with label $u^{a}v^{b}$
is homogeneous of phase degree $2+a+2b$ while the corresponding coordinate
coefficient has degree $1+a+2b$, and a sufficient domain certificate is
simply
\begin{equation}\label{dyn:eq:exampledomain}
 \eta+\rho>0,\qquad \theta+2\rho>0 .
\end{equation}
The input label set being finite, positivity already supplies well-ordering
and finite fibers here. In a higher-rank group, writing these as valuation
inequalities rather than numerical radii avoids introducing a nonexistent
real-valued absolute value. If two label weights coincide, for instance
$2\eta=\theta$, their terms are added at that Hahn exponent and the
construction remains valid; label-degree truncation and Hahn-weight
truncation are still different operations, and the verification code
deliberately uses the former.

\subsection{A discrete one-variable parallel}
\label{dyn:subsec:discrete}

The same support mechanism is visible in a simpler setting. It is included
as a bridge to the established non-Archimedean literature
\cite{Lindahl}, \emph{not} as a claim that polynomial perturbative
linearization is new in rank one.

\begin{proposition}[Polynomial-layer linearization]\label{dyn:prop:discrete}
Let $\lambda\in\C^{\times}$ not be a root of unity and let
$f(z)=\lambda z+R(z)$ with $R\in\HH_\Gamma(\PP_1)_{>0}$ and
$\ord_0R_s\ge2$. There is a unique $h(z)=z+U(z)$ with positive polynomial
Hahn coefficients, $\ord_0U_\gamma\ge2$, satisfying
\begin{equation}\label{dyn:eq:discreteconj}
 f\circ h=h\circ(\lambda\,\id),
\end{equation}
it is a bijection of the finite halo, and
\eqref{dyn:eq:discreteconj} holds there as an actual identity. \emph{No}
bound on $|\lambda^{n}-\lambda|^{-1}$ is required.
\end{proposition}
\begin{proof}
On $z^{n}$ with $n\ge2$ the operator $L_\lambda U=U(\lambda z)-\lambda U(z)$
is multiplication by the nonzero constant $\lambda^{n}-\lambda$, hence
invertible on polynomials of order at least two, and
\eqref{dyn:eq:discreteconj} becomes $U=L_\lambda^{-1}R(z+U)$. At a fixed
Hahn exponent every nonlinear occurrence of $U$ on the right is accompanied
by a positive coefficient of $R$, hence uses strictly smaller exponents of
$U$, so the supported recursion and finite-expression argument of
\cref{dyn:prop:inverse} apply; each coefficient remains a polynomial, the
least-difference argument proves uniqueness, and
\cref{dyn:prop:inverse,dyn:prop:eval} give the actual conjugacy.
\end{proof}

For $f(z)=\lambda z+\eps_0z^{2}$, writing
$h(z)=\sum_{n\ge1}b_n\eps_0^{n-1}z^{n}$ with $b_1=1$ gives the exact
recurrence
\begin{equation}\label{dyn:eq:discreterec}
 (\lambda^{n}-\lambda)b_n=\sum_{j=1}^{n-1}b_jb_{n-j}\qquad(n\ge2),
\end{equation}
with one phase monomial per Hahn layer; at finite $z$ it is the
\emph{support}, not the complex growth of the $b_n$, that makes this a
valid Hahn identity. This is the same quadratic Schr\"oder recursion that
sources 02, 03, 04 and 05 each also compute, in four different
normalizations; \cref{dyn:subsec:quadratic} prints it once with the
translation between them.

\section{Quasi-periodic divisors: a finite flag for a surreal frequency}
\label{dyn:sec:flag}

Everything up to this point has been \emph{germ and fixed-point} dynamics:
a map $F$ with a distinguished fixed point, its multiplier, its return
iterates, its periods, or a Hamiltonian at the origin of a finite phase
halo. Source 07 is the other classical half of small-divisor theory. Its
object is a \emph{vector field on a real torus} whose frequency vector has
surreal components, its divisors are the infinitely many scalars
$k\cdot\Upsilon$ indexed by $k\in\Z^{d}$, and its equation is the
cohomological equation. None of the six germ sources contains a resonance
lattice, a cohomological equation, a Fourier divisor or a conjugacy to a
constant vector field; conversely nothing in this part concerns a
multiplier, a return iterate or a period. What the two halves share is the
machinery of \cref{dyn:sec:framework} --- the set-sized workspace, strong
summability, the positive-support word lemma, halo evaluation --- which is
therefore not restated here, and the architecture of
\cref{dyn:rem:oneinvariant}, which is instantiated a fourth time in
\cref{dyn:sec:qplinear}.

The headline of this part has no counterpart in the other six. A frequency
vector with $d$ surreal components can carry a transfinite hierarchy of
scales, and there are infinitely many Fourier divisors; yet only
\textbf{at most $d$} distinct divisor valuations are arithmetically
visible, and a canonical descending flag of integer lattices identifies
them. That finite reduction is what turns an apparently transfinite problem
into a finite one, and it is the reason this material occupies three
sections rather than a remark.

\subsection{The torus, the strip, and the notation of this part}
\label{dyn:subsec:qpnotation}

Before anything else, read
\cref{dyn:warn:twofrequencies,dyn:warn:qpspectrum,dyn:warn:qptwosizes}.
They record that this report now carries \emph{two} frequency vectors of
different kinds, that $\kups$ below is an element of $\Gamma$ and not a
rate, that $\Upsilon$ is never the exponent $\Omega$ of
\cref{dyn:ex:ranktwo} nor the time form $\Omega_F$ of
\cref{dyn:sec:periods}, and that the finitely many visible divisor
valuations are not a spectrum in any sense this collection uses.

Fix $d\ge1$ and write $\TT^{d}=(\R/2\pi\Z)^{d}$. Fix $\rho>0$ and let
\begin{equation}\label{dyn:eq:strip}
 \TT^{d}_\rho=\{z\in\C^{d}/(2\pi\Z)^{d}:|\operatorname{Im}z_j|<\rho\
 \text{for all }j\},\qquad \strip
\end{equation}
be the \emph{open} complex strip of half-width $\rho$ around the real torus
and the algebra of $2\pi\Z^{d}$-periodic functions holomorphic on it.
Functions in $\strip$ need not be bounded at the boundary of the strip.
With $|k|=|k_1|+\cdots+|k_d|$, write
\begin{equation}\label{dyn:eq:seminorm}
 f(\theta)=\sum_{k\in\Z^{d}}\widehat f(k)e^{ik\cdot\theta},
 \qquad
 \|f\|_{r}=\sum_{k\in\Z^{d}}|\widehat f(k)|e^{r|k|}\quad(0<r<\rho).
\end{equation}
A periodic holomorphic function lies in $\strip$ exactly when all these
seminorms are finite: boundedness on a narrower closed strip and Cauchy's
Fourier estimate give exponential decay there, and inserting an
intermediate strip gives the weighted $\ell^{1}$ bound; conversely those
bounds give locally uniform convergence of the Fourier series and its
termwise derivatives on every narrower strip. Hence $\strip$ is closed
under finite sums, products and all orders of differentiation, and Fourier
projection onto \emph{any} subset of $\Z^{d}$ preserves it, since it only
deletes terms from $\|f\|_{r}$.

The torus derivatives $\partial/\partial\theta_j$ act coefficientwise and
annihilate $K_\Gamma$: they are external derivations, and
\cref{dyn:warn:derivative} applies verbatim --- they are not the
Berarducci--Mantova derivation, and $\partial_\theta\sin\theta=\cos\theta$
here says nothing about whether a surreal $y$ solves $\partial^{2}y+y=0$
for that internal derivation. Source 07 states this disclaimer
independently and in its own words; it is the same disclaimer, printed
once.

\subsection{The fixed-strip coefficient category}
\label{dyn:subsec:qpstrip}

\begin{definition}[The common-strip Hahn algebra]\label{dyn:def:qpstrip}
$\HH_\Gamma(\strip)=\strip((t^{\Gamma}))$ is the Hahn algebra over the
coefficient algebra \eqref{dyn:eq:strip}: elements are
$f=\sum_{\beta\in B}t^{\beta}f_\beta$ with $B$ well ordered and every
$f_\beta\in\strip$. Fourier coefficients and the torus mean are taken
coefficientwise,
\[
 \widehat f(k)=\sum_\beta t^{\beta}\widehat{f_\beta}(k)\in K_\Gamma,
 \qquad \av(f)=\widehat f(0),\qquad \Qz f=f-\av(f),
\]
and $\HH_\Gamma(\strip)_{>0}$ denotes positive support, zero included. For
a vector the support is the union of the component supports and the
valuation their minimum.
\end{definition}

This is the \emph{fixed-domain} row of \cref{dyn:subsec:rings} in its torus
instantiation: one and the same $\rho$ serves every Hahn exponent $\beta$,
and separate analytic germs for separate coefficients do not suffice. The
distinction is stated once, in \cref{dyn:conv:rings}, and is not re-erected
here; source 07 imports it from the collection's catalogue explicitly, as
a hypothesis of each of its main theorems rather than as a convention. As
in the entire row of \cref{dyn:subsec:rings}, \textbf{no} bound on
$\|f_\beta\|_{r}$ uniform in $\beta$ is imposed --- a deliberate choice,
and the reason nothing here specializes $t$ to a nonzero real number
(\cref{dyn:subsec:qplimits}).

Elements of $\HH_\Gamma(\strip)$ are genuine functions on infinitesimal
thickenings: for ordinary $z_0\in\TT^{d}_\rho$ and $\xi\in\m_\Gamma^{d}$,
\eqref{dyn:eq:halo-evaluation} applies verbatim and
\cref{dyn:prop:eval} makes $f\mapsto f^{\#}(z_0+\xi)$ a faithful
$K_\Gamma$-algebra homomorphism into $\SC$. No continuity from an ordinary
torus into the fine topology of $\No[i]$ is claimed.

\begin{lemma}[Subexponential Fourier multipliers]\label{dyn:lem:qpmultiplier}
Let $m:\Z^{d}\to\C$ satisfy $|m(k)|\le C_\eps e^{\eps|k|}$ for every
$\eps>0$. Then $\widehat{Tf}(k)=m(k)\widehat f(k)$ maps $\strip$ into
itself, and so does any finite product of such multipliers with polynomial
functions of $k$.
\end{lemma}
\begin{proof}
For $r<r'<\rho$ apply the estimate with $\eps=r'-r$ to get
$\|Tf\|_{r}\le C_{r'-r}\|f\|_{r'}$. A fixed polynomial in $k$ is bounded by
a constant times $e^{\eta|k|}$ for every $\eta>0$; split an arbitrary
exponent budget among finitely many factors.
\end{proof}

\begin{lemma}[Taylor substitution on the torus, with its Lipschitz
estimate]\label{dyn:lem:qpsubstitution}
Let $h\in\HH_\Gamma(\strip)^{d}$ have positive valuation. Then
\begin{equation}\label{dyn:eq:qpsubstitution}
 S_hf=f\circ(\id+h)
 :=\sum_\beta t^{\beta}\sum_{\alpha\in\N^{d}}
 \frac{\partial^{\alpha}f_\beta}{\alpha!}h^{\alpha}
\end{equation}
is well defined in $\HH_\Gamma(\strip)$, with
$\supp(S_hf)\subseteq\supp(f)+C^{*}$ for $C\supseteq\supp h$ well ordered
and positive. Substitution is a ring homomorphism, obeys the chain rule,
and is associative for infinitesimal substitutions; these are the torus
instances of \cref{dyn:prop:subgroup} and are not reproved. In addition,
for positive-valuation $h,g$,
\begin{equation}\label{dyn:eq:qplipschitz}
 v(S_hf)\ge v(f),\qquad
 v(S_hf-S_gf)\ge v(f)+v(h-g).
\end{equation}
\end{lemma}
\begin{proof}
Well-definedness, the support bound and the algebraic identities are
\cref{dyn:lem:neumann,dyn:cor:ancestry,dyn:prop:subgroup} applied with
$U=\TT^{d}_\rho$: at each exponent the word lemma permits only finitely
many Taylor monomials and finitely many shifted word pairs in
$\word_{\supp f,C}(\gamma)$, and every resulting coefficient is a finite
sum of products of derivatives of functions in $\strip$, hence lies in $\strip$. For the
second estimate in \eqref{dyn:eq:qplipschitz} subtract corresponding Taylor
monomials and telescope the products: every surviving term carries one
factor from $h-g$ and all its other factors have positive valuation, and
the multi-index-zero terms cancel.
\end{proof}

\Cref{dyn:eq:qplipschitz} is the one part of the torus substitution
calculus that is not already in \cref{dyn:prop:subgroup}; it is what
forces uniqueness in \cref{dyn:thm:qpnormalform} and controls the remainder
in \cref{dyn:thm:qpnecessity}.

\begin{warning}[Two summations, never interchanged]\label{dyn:warn:qptwosums}
The Fourier series of an ordinary coefficient is an ordinary complex
analytic sum; the series in $t^{\Gamma}$ is a Hahn sum. Every argument
below groups all Fourier modes \emph{inside each analytic coefficient}
first. It is nowhere asserted that the unrestricted two-index family of
Fourier and Hahn terms is Hahn-summable, and indeed a family with
infinitely many nonzero terms at one Hahn exponent violates
\cref{dyn:def:summable} however well its complex coefficients converge.
This is the torus form of the discipline of \cref{dyn:def:summable}.
\end{warning}

\subsection{The flag, and one support for every reciprocal divisor}
\label{dyn:subsec:qpflag}

Let $\Upsilon\in F_\Gamma^{d}\setminus\{0\}$ and collect the component
Conway/Hahn normal forms into
\begin{equation}\label{dyn:eq:qpfrequency}
 \Upsilon=\sum_{\beta\in S}t^{\beta}\varpi_\beta,
 \qquad \varpi_\beta\in\R^{d}\setminus\{0\},\quad S\ \text{well ordered},
\end{equation}
so the \emph{exact resonance lattice} is
\begin{equation}\label{dyn:eq:qpexactlattice}
 \qlat=\{k\in\Z^{d}:k\cdot\Upsilon=0\}
 =\bigcap_{\beta\in S}\ker(k\mapsto k\cdot\varpi_\beta).
\end{equation}
This is \emph{exact vanishing of a Hahn scalar}, never smallness; that is
the content of the tag $\qres$.

Set $\mathsf L_0=\Z^{d}$. While $\mathsf L_{j-1}\ne\qlat$, let $\gamma_j$ be
the \emph{least} exponent of $S$ whose coefficient form does not vanish
identically on $\mathsf L_{j-1}$, and put
\begin{equation}\label{dyn:eq:qpflagdef}
 \fla_j(k)=k\cdot\varpi_{\gamma_j},\qquad
 \mathsf L_j=\{k\in\mathsf L_{j-1}:\fla_j(k)=0\},\qquad
 \Sigma_j=\mathsf L_{j-1}\setminus\mathsf L_j .
\end{equation}
The construction stops as soon as the current lattice equals $\qlat$. The
$\Sigma_j$ are the \emph{nonresonant strata}.

\begin{theorem}[Finite flag and one common reciprocal support]
\label{dyn:thm:flag}
The construction \eqref{dyn:eq:qpflagdef} terminates after
\begin{equation}\label{dyn:eq:qprbound}
 r\le d
\end{equation}
steps, and:
\begin{enumerate}[label=(\roman*),leftmargin=*]
\item each $\mathsf L_j$ is a saturated subgroup of $\Z^{d}$ whose rank
strictly decreases at each step, $\gamma_1<\cdots<\gamma_r$,
$\mathsf L_r=\qlat$, and the strata partition $\Z^{d}\setminus\qlat$;
\item for $k\in\Sigma_j$,
\begin{equation}\label{dyn:eq:qpleading}
 k\cdot\Upsilon=t^{\gamma_j}\fla_j(k)(1+\varrho_k),\qquad
 v(\varrho_k)>0,\qquad v(k\cdot\Upsilon)=\gamma_j,
\end{equation}
with $\varrho_k=0$ permitted; in particular the divisor valuation takes at
most $d$ distinct values on $\Z^{d}\setminus\qlat$;
\item writing
\begin{equation}\label{dyn:eq:qpE}
 \mathsf E_j=\{\beta-\gamma_j:\beta\in S,\ \beta>\gamma_j\},\qquad
 \mathsf E=\bigcup_{j=1}^{r}\mathsf E_j,\qquad
 \Tups=\bigcup_{j=1}^{r}\bigl(-\gamma_j+\mathsf E^{*}\bigr),
\end{equation}
the set $\mathsf E$ is well ordered and positive, $\Tups$ is well ordered,
and
\begin{equation}\label{dyn:eq:qpTsupport}
 \supp\bigl((k\cdot\Upsilon)^{-1}\bigr)\subseteq\Tups
 \qquad\text{for every }k\notin\qlat .
\end{equation}
With $\kups=\gamma_r$ one has $\Tups\subseteq[-\kups,+\infty)$.
\end{enumerate}
\end{theorem}
\begin{proof}
Each $\mathsf L_j$ is a common kernel of real additive homomorphisms on
$\Z^{d}$ and is therefore saturated: if a nonzero integer multiple of $k$
lies in the kernel, so does $k$. At a nonterminal step $\fla_j$ is a
nonzero homomorphism on the finitely generated free abelian group
$\mathsf L_{j-1}$ with torsion-free nonzero image, so $\rank\mathsf L_j<\rank
\mathsf L_{j-1}$; at most $d$ such decreases occur, which is
\eqref{dyn:eq:qprbound}.

By the minimal choice of $\gamma_j$, every coefficient form below
$\gamma_j$ vanishes on $\mathsf L_{j-1}$, and the form at $\gamma_j$
vanishes on $\mathsf L_j$; hence the next selected exponent is strictly
larger. Termination occurs exactly when every coefficient form vanishes on
the current lattice, i.e.\ at $\qlat$. The same observation identifies the
leading term for $k\in\Sigma_j$: explicitly
\begin{equation}\label{dyn:eq:qptail}
 \varrho_k=\sum_{\substack{\beta\in S\\ \beta>\gamma_j}}
 \frac{k\cdot\varpi_\beta}{\fla_j(k)}\,t^{\beta-\gamma_j},
\end{equation}
whose support lies in $\mathsf E_j$ and is positive. This is (i) and (ii).

A tail of a well-ordered set translated by a fixed element is well ordered,
and a finite union of well-ordered subsets of a linear order is well
ordered; this gives $\mathsf E$. \Cref{dyn:lem:neumann} then gives
$\mathsf E^{*}$ and $\Tups$. Finally
\begin{equation}\label{dyn:eq:qpreciprocal}
 \frac1{k\cdot\Upsilon}
 =\frac{t^{-\gamma_j}}{\fla_j(k)}\sum_{n\ge0}(-\varrho_k)^{n}
 \qquad(k\in\Sigma_j)
\end{equation}
is a Hahn identity by \cref{dyn:lem:neumann} --- not a limit of partial
sums (\cref{dyn:def:summable,dyn:warn:contraction}) --- with the asserted
support. Every element of $\mathsf E^{*}$ is nonnegative and
$\gamma_j\le\kups$, giving the lower bound.
\end{proof}

\begin{remark}[$r$ counts visible scales, not monomials]
\label{dyn:rem:qpvisible}
The number $r$ counts \emph{arithmetically visible} scales, not the
monomials of $\Upsilon$. Exponents of $S$ whose forms vanish on the current
lattice are skipped, possibly across a long transfinite interval. The
theorem does \textbf{not} say that only $d$ monomials of $\Upsilon$ matter
to the eventual solution: every later coefficient can still move higher
Hahn coefficients of the inverse, through \eqref{dyn:eq:qptail}. What is
finite is the list of divisor valuations and the list of real forms on
which all small-divisor arithmetic is tested.
\end{remark}

\begin{proposition}[Coefficient structure of the inverse]
\label{dyn:prop:qpjets}
Fix $\eta\in\Gamma$ and a stratum $\Sigma_j$. On that stratum the function
$k\mapsto[t^{\eta}](k\cdot\Upsilon)^{-1}$ is a finite sum of expressions
\begin{equation}\label{dyn:eq:qpmultform}
 c\,\frac{\prod_{\ell=1}^{n}(k\cdot\varpi_{\beta_\ell})}
 {\fla_j(k)^{n+1}},
 \qquad \sum_{\ell=1}^{n}(\beta_\ell-\gamma_j)=\eta+\gamma_j,
 \quad \beta_\ell>\gamma_j,
\end{equation}
with $c$ real and empty products allowed at $n=0$. The number of terms and
the denominator powers depend on $\eta$ but are finite for each fixed
$\eta$.
\end{proposition}
\begin{proof}
Expand \eqref{dyn:eq:qpreciprocal} and extract the coefficient. All letters
lie in $\mathsf E_j$, and \cref{dyn:lem:neumann} permits finitely many
words of all lengths at a fixed weight; $\beta\mapsto\beta-\gamma_j$ is
injective, so the coefficient choices are also finitely many.
\end{proof}

\begin{remark}[A common support is not summability over the modes]
\label{dyn:rem:qpnotsummable}
The family $\bigl((k\cdot\Upsilon)^{-1}\bigr)_{k\notin\qlat}$ is in general
\textbf{not} a Hahn-summable family: infinitely many leading terms may
occupy the exponent $-\gamma_j$. \Cref{dyn:thm:flag} supplies a
\emph{common support}, not a Hahn sum over $k$, exactly as
\cref{dyn:def:summable} requires. \Cref{dyn:prop:qpjets} is what permits
ordinary Fourier summation at each Hahn coefficient; that is the additional
analytic step, and it is where arithmetic enters.
\end{remark}

\begin{proposition}[An arithmetic-free valuation bound]
\label{dyn:prop:qpvalbound}
Let $\Cups=\Upsilon_1\partial_{\theta_1}+\cdots+\Upsilon_d\partial_{\theta_d}$.
If $h\in\HH_\Gamma(\strip)$ has no Fourier modes in $\qlat$ and
$\Cups h=g$, then $v(h)\ge v(g)-\kups$. No small-divisor hypothesis is used.
\end{proposition}
\begin{proof}
For $h\ne0$ put $\beta=v(h)$. The nonzero $h_\beta$ has a nonzero Fourier
coefficient at some $k\notin\qlat$, so $v(\widehat h(k))=\beta$ and
$v(\widehat g(k))=\beta+v(k\cdot\Upsilon)\le\beta+\kups$ by
\eqref{dyn:eq:qpleading}. Every Fourier coefficient of $g$ has valuation at
least $v(g)$; rearrange. The zero case is immediate.
\end{proof}

\section{The fixed-strip solvability threshold, and an arbitrary-order
lifting obstruction}
\label{dyn:sec:qplinear}

\subsection{The stratified condition, and the category table instantiated a
second time}
\label{dyn:subsec:qpss}

\begin{definition}[Stratified subexponential condition $\sdss$]
\label{dyn:def:qpss}
$\Upsilon$ satisfies $(\sdss)$ if for each flag stratum $\Sigma_j$ and
every real $\eps>0$ there is a real $C_{j,\eps}>0$ with
\begin{equation}\label{dyn:eq:qpss}
 \frac1{|\fla_j(k)|}\le C_{j,\eps}e^{\eps|k|}\qquad(k\in\Sigma_j),
\end{equation}
equivalently
$\limsup_{|k|\to\infty,\,k\in\Sigma_j}\log^{+}(1/|\fla_j(k)|)/|k|=0$.
\end{definition}

\begin{remark}[The one invariant, a fourth divisor family, and two
instantiations of one table]\label{dyn:rem:qponeinvariant}
\Cref{dyn:rem:oneinvariant} defines one quantity --- the exponential rate
of the reciprocals of a graded family of nonzero ordinary complex divisors
--- and evaluates it on three families. $(\sdss)$ is that same quantity on
a \emph{fourth} family, evaluated stratum by stratum:
\begin{center}
\begin{tabular}{@{}lll@{}}
\toprule
Divisor family & Grading & Rate\\
\midrule
$\lambda^{n}-\lambda$ (drifted multiplier) & $n$ & $\sigma$ ($\sddr$)\\
$\lambda^{\alpha}-\lambda_j$ (exact multiplier) & $|\alpha|$ & $\tau$ ($\sdex$)\\
$\omega\cdot(\alpha-\beta)$ (ordinary frequency) & $|\alpha|+|\beta|$
 & $\log\Theta$ ($\sdfr$)\\
$\fla_j(k)=k\cdot\varpi_{\gamma_j}$ (leading real forms of $\Upsilon$,
 on $\Sigma_j$) & $|k|$ & $0$ required on every $j$ ($\sdss$)\\
\bottomrule
\end{tabular}
\end{center}
So $(\sdss)$ is the vanishing-rate condition, imposed on each of the
finitely many strata; it is the exact analogue of $\tau=0$ and of
$\sigma=0$, and it is \emph{not} the analogue of $\tau<\infty$. That is
consistent with the category table of \cref{dyn:rem:oneinvariant}, which is
stated there once and is instantiated twice in this report:
\begin{center}
\begin{tabular}{@{}P{0.30\textwidth}P{0.29\textwidth}P{0.31\textwidth}@{}}
\toprule
Row of \cref{dyn:rem:oneinvariant} & Germ instantiation, divisors
$\lambda^{\alpha}-\lambda_j$ or $\lambda^{q}-1$ & Quasi-periodic
instantiation, divisors $k\cdot\Upsilon$\\
\midrule
one fixed unchanged domain, rate $0$
& fixed polydisk $\OO(\mathbb D_R^{d})((t^{\Gamma}))$, $\tau=0$
 (row one of \cref{dyn:thm:main})
& fixed strip $\HH_\Gamma(\strip)$, $(\sdss)$
 (\cref{dyn:thm:qplinear})\\[3pt]
polynomial / formal coefficients, no arithmetic
& \cref{dyn:thm:main} rows four and five; $\sdz$
& \textbf{not established here}: source 07 proves no polynomial or formal
row for the torus, and none is claimed\\[3pt]
radius-free germs / entire coefficients, rate finite
& \cref{dyn:thm:main} rows two and three; \cref{dyn:thm:homological}
& \textbf{not established here}; the fixed strip is the only category
source 07 classifies\\
\bottomrule
\end{tabular}
\end{center}
Only the first row is instantiated on both sides. Reading across the other
two rows, in either direction, would assert theorems no source proves.
\end{remark}


\subsection{Universal solvability on a fixed strip}

Let $\Pi_{\qlat}$ denote coefficientwise Fourier projection onto $\qlat$;
it preserves $\HH_\Gamma(\strip)$ because arbitrary Fourier projection
preserves $\strip$.

\begin{theorem}[Universal fixed-strip solvability]\label{dyn:thm:qplinear}
For a nonzero $\Upsilon\in F_\Gamma^{d}$ the following are equivalent.
\begin{enumerate}[label=(\roman*),leftmargin=*]
\item $\Upsilon$ satisfies $(\sdss)$.
\item For every $\rho>0$ and every $f\in\HH_\Gamma(\strip)$ with
$\Pi_{\qlat}f=0$ there is a unique $u\in\HH_\Gamma(\strip)$ with
$\Cups u=f$ and $\Pi_{\qlat}u=0$.
\end{enumerate}
When these hold, the solution operator is
\begin{equation}\label{dyn:eq:qpgreen}
 \widehat{\Gups f}(k)=
 \begin{cases}
 \widehat f(k)/(i\,k\cdot\Upsilon), & k\notin\qlat,\\
 0,& k\in\qlat,
 \end{cases}
\end{equation}
and it satisfies
\begin{equation}\label{dyn:eq:qpgreensupport}
 \supp(\Gups f)\subseteq\supp(f)+\Tups,\qquad
 v(\Gups f)\ge v(f)-\kups .
\end{equation}
Without the normalization the kernel consists exactly of the elements of
$\HH_\Gamma(\strip)$ with Fourier support inside $\qlat$. If $(\sdss)$
fails, there is already an ordinary real analytic mean-compatible forcing,
supported at a single Hahn exponent and on one stratum, for which a forced
coefficient of any putative solution is not a distribution on the real
torus.
\end{theorem}
\begin{proof}[Sufficiency]
For $\eta\in\Tups$ define the scalar multiplier
$m_\eta(k)=[t^{\eta}](i\,k\cdot\Upsilon)^{-1}$ for $k\notin\qlat$ and
$m_\eta(k)=0$ on $\qlat$. On each stratum \cref{dyn:prop:qpjets} writes
$m_\eta$ as a finite sum of expressions with polynomial numerator in $k$
and a finite power of $\fla_j(k)$ in the denominator; $(\sdss)$ and
\cref{dyn:lem:qpmultiplier} make each a subexponential multiplier. There
are finitely many strata --- this is where \eqref{dyn:eq:qprbound} is used
--- and projection onto a stratum does not affect the estimate, so $m_\eta$
defines $G_\eta:\strip\to\strip$.

Write $f=\sum_{\beta\in B}t^{\beta}f_\beta$ and put
\begin{equation}\label{dyn:eq:qpgreencoefficient}
 u_\gamma=\sum_{\substack{\beta\in B,\ \eta\in\Tups\\ \beta+\eta=\gamma}}
 G_\eta f_\beta .
\end{equation}
By \cref{dyn:lem:neumann} this sum is finite at each $\gamma$ and
$B+\Tups$ is well ordered, so $u=\sum_\gamma t^{\gamma}u_\gamma$ lies in
$\HH_\Gamma(\strip)$: the common domain and the common support are
\emph{proved}, not inferred from the modewise formula. Fourier extraction
gives \eqref{dyn:eq:qpgreen}, and multiplying modewise by
$i\,k\cdot\Upsilon$ gives the equation. Uniqueness and the kernel
description follow from the same modewise calculation, $K_\Gamma$ being a
field. The support estimate is the construction and
$\Tups\ge-\kups$ gives the valuation estimate.
\end{proof}
\begin{proof}[Necessity]
Suppose $(\sdss)$ fails on $\Sigma_j$. Choose $\eps>0$ and modes
$k_n\in\Sigma_j$, no two equal up to sign, with $|k_n|\to\infty$ and
$|\fla_j(k_n)|^{-1}>e^{2\eps|k_n|}$; this is possible after decreasing
$\eps$ if necessary. The real function
\begin{equation}\label{dyn:eq:qpbadforcing}
 f(\theta)=2\sum_{n\ge1}e^{-\eps|k_n|}\cos(k_n\cdot\theta)
\end{equation}
lies in $\OO(\TT^{d}_\eps)$, has no resonant modes and sits at Hahn
exponent zero. Any solution must have
$\widehat u(k_n)=e^{-\eps|k_n|}/(i\,k_n\cdot\Upsilon)$, so by
\eqref{dyn:eq:qpleading} the coefficient of $t^{-\gamma_j}$ has Fourier
coefficients of modulus greater than $e^{\eps|k_n|}$. Fourier coefficients
of a distribution on a compact torus grow at most polynomially --- apply
its finite-order continuity estimate to the Fourier characters --- so that
forced coefficient is not a distribution, and in particular not holomorphic
on any strip of positive width. Adding resonant modes cannot alter these
nonresonant coefficients.
\end{proof}

\begin{corollary}[The divisor loss $\kups$ is attained]\label{dyn:cor:qpsharp}
For any $k\in\Sigma_r$ and any $s\in\Gamma$ the solution for
$f=t^{s}e^{ik\cdot\theta}$ has valuation $s-\kups$. Hence the valuation
loss in \eqref{dyn:eq:qpgreensupport} cannot be improved uniformly.
\end{corollary}
\begin{proof}
Its unique nonzero Fourier mode is divided by a scalar of valuation
$\kups$.
\end{proof}

\begin{remark}[A stratified Diophantine sufficient condition]
\label{dyn:rem:qpdiophantine}
A convenient sufficient hypothesis is a stratified Diophantine bound
$|\fla_j(k)|\ge c_j(1+|k|)^{-\tau_j}$ on $\Sigma_j$. A term of length $n$
in \eqref{dyn:eq:qpmultform} then has a multiplier growing at most
polynomially with exponent $n+(n+1)\tau_j$, so every individual coefficient
operator has explicit finite loss estimates between narrower-strip
seminorms. It is \textbf{not} asserted that those losses stay bounded as
$\eta$ varies. The fixed-strip theorem is available precisely because the
category demands a common open domain and not a single bounded analytic
norm for every coefficient --- the same mechanism as in
\cref{dyn:warn:nouniform}.
\end{remark}

\begin{corollary}[Higher tails do not change the arithmetic test]
\label{dyn:cor:qptail}
If $v(\widetilde\Upsilon-\Upsilon)>\kups$ then $\widetilde\Upsilon$ has the
same flag, the same leading real forms, and --- whenever $\qlat=\{0\}$ ---
the same exact resonance lattice; in that nonresonant case
$\widetilde\Upsilon$ satisfies $(\sdss)$ exactly when $\Upsilon$ does.
\end{corollary}
\begin{proof}
No coefficient at or below $\kups$ changes. Each nonzero integer vector
already has a nonvanishing leading divisor coefficient at one of those
levels, so no resonance can appear and its leading form is unchanged.
\end{proof}

The nonresonance qualification is not decoration: when $\qlat\ne\{0\}$ a
new higher tail can split some exact resonances and introduce additional
flag levels, and the criterion must then be applied to the extended flag.

\subsection{Arbitrarily long analytic jets with no analytic lift}
\label{dyn:subsec:qpjet}

The next construction is a failure that no valuation-only argument can
detect, and it lives entirely at integer exponents: the obstruction is not
caused by exotic uncountable supports.

Define integers and a real number by
\begin{equation}\label{dyn:eq:qpliouvilleconstant}
 b_1=2,\qquad b_{j+1}=10^{2b_j},\qquad
 \mathfrak l=\sum_{j\ge1}10^{-b_j},
\end{equation}
and put $Q_j=10^{b_j}$, $P_j=Q_j\sum_{\ell\le j}10^{-b_\ell}$,
$\mathfrak d_j=Q_j\mathfrak l-P_j$, $k_j=(-P_j,Q_j)$, $K_j=|k_j|$. Here
$P_j$ is an integer with $0<P_j<Q_j$ and
\begin{equation}\label{dyn:eq:qpliouville}
 0<\mathfrak d_j\le2Q_j10^{-Q_j^{2}},\qquad Q_j\le K_j<2Q_j ,
\end{equation}
by bounding the tail of the rapidly increasing sequence of decimal places;
the same estimate proves $\mathfrak l$ irrational, since a fixed
denominator $q$ would force $|\mathfrak d_j|\ge1/q$. Consider the
infinitesimally detuned frequency and operator
\begin{equation}\label{dyn:eq:qpdetuned}
 \Upsilon=(1,\mathfrak l+t),\qquad
 \Cups=\partial_x+(\mathfrak l+t)\partial_y=L_0+t\partial_y,
 \quad L_0=\partial_x+\mathfrak l\,\partial_y .
\end{equation}
The real leading vector $(1,\mathfrak l)$ is already nonresonant, so there
is one flag level, at $\gamma_1=0$, and $(\sdss)$ fails very strongly.

\begin{theorem}[Arbitrary finite analytic lifting, but no full lift]
\label{dyn:thm:qpjet}
For every integer $N\ge1$ and every $\rho>0$ there is an ordinary real
entire periodic $f_N$ such that
\begin{enumerate}[label=(\roman*),leftmargin=*]
\item $\Cups u=f_N$ has a mean-zero solution modulo $t^{N+1}$ with
coefficients $u_0,\dots,u_N\in\strip$;
\item $u_0,\dots,u_{N-1}$ are entire periodic functions;
\item the forced coefficient $u_{N+1}$ is not a distribution on the real
torus, so there is no solution in $\HH_\Gamma(\OO(\TT^{2}_{\rho'}))$ for
any $\rho'>0$, \emph{and none after an order-preserving extension of the
value group}.
\end{enumerate}
One may take
\begin{equation}\label{dyn:eq:qpjetforcing}
 f_N(x,y)=2\sum_{j\ge1}\mathfrak d_j^{\,N+1}e^{-\rho K_j}
 \cos\bigl(k_j\cdot(x,y)\bigr).
\end{equation}
\end{theorem}
\begin{proof}
By \eqref{dyn:eq:qpliouville} the Fourier coefficients in
\eqref{dyn:eq:qpjetforcing} decay faster than $e^{-RK_j}$ for every fixed
$R>0$, so $f_N$ is entire, real and of zero mean. At the mode $k_j$ the
divisor is $\mathfrak d_j+Q_jt$; expanding its inverse as a Hahn geometric
series (legitimate by \cref{dyn:lem:neumann}) gives
\begin{equation}\label{dyn:eq:qpjetcoefficients}
 \widehat{u_n}(k_j)=\frac{(-1)^{n}}{i}\,
 Q_j^{\,n}\mathfrak d_j^{\,N-n}e^{-\rho K_j}\qquad(n\ge0),
\end{equation}
with complex conjugates at $-k_j$ and no other nonzero modes; set every
mean to zero. For $n<N$ the remaining positive power of $\mathfrak d_j$
gives superexponential decay, so $u_n$ is entire. For $n=N$ the
coefficients are the polynomial factor $Q_j^{N}$ times $e^{-\rho K_j}$,
whose weighted Fourier sums are finite on every strip of width $r<\rho$, so
$u_N\in\strip$. The scalar identities give $L_0u_0=f_N$ and
$L_0u_n=-\partial_yu_{n-1}$ for $1\le n\le N$, which is the finite-jet
statement in $\strip[[t]]/(t^{N+1})$. At the next coefficient
\[
 |\widehat{u_{N+1}}(k_j)|
 =Q_j^{\,N+1}\mathfrak d_j^{-1}e^{-\rho K_j}
 \ge\tfrac12Q_j^{\,N}10^{Q_j^{2}}e^{-2\rho Q_j},
\]
which exceeds every polynomial in $K_j$ along the sequence and so is not a
distributional Fourier coefficient sequence.

For the value-group clause: at each nonzero Fourier mode the scalar
equation has the \emph{unique} solution
$\widehat{f_N}(k)/(i\,k\cdot\Upsilon)$ in the Hahn field, and the inverse
of that divisor is the same geometric series in any larger ordered-group
Hahn field. Therefore the coefficient of $t^{N+1}$ remains exactly
\eqref{dyn:eq:qpjetcoefficients}, and admitting negative or noninteger Hahn
exponents in a proposed solution changes nothing; extra constants affect
only the zero Fourier mode.
\end{proof}

\begin{warning}[The robustness under value-group extension is argued, not
assumed]\label{dyn:warn:qpextension}
This report's standing hazard is the temptation to assert that some norm or
estimate is invariant under enlarging $\Gamma$. \Cref{dyn:thm:qpjet}(iii)
does \textbf{not} do that. It cites no invariance of any norm across a
value-group extension; it argues from \emph{uniqueness of the modewise Hahn
inverse} in the larger field, which is an algebraic statement about a
field, together with the fact that the same geometric expansion computes
it. Source 07 states the argument in exactly that form, and the
corresponding invariance in the positive direction is
\cref{dyn:prop:workspace}, which is likewise proved rather than assumed.
\end{warning}

This is a finite-versus-infinite lifting obstruction in a precise analytic
ring. It does not say that the scalar algebraic inverses fail to exist ---
they all exist --- but that infinitely many modewise inverses do not
assemble into even one analytic coefficient at a specified order. The same
operator \eqref{dyn:eq:qpdetuned} works for every $N$; only the entire
forcing changes. In particular \textbf{no test consisting of a fixed finite
number of analytic coefficient equations can replace the arithmetic
hypothesis} in a general lifting theorem. The construction adapts the
classical Liouville mechanism \cite{Kamienski}, which is not claimed as
new; what is offered is the arbitrarily prescribed order at which it first
bites.

\section{Torus normal form: conjugacy to a constant vector field}
\label{dyn:sec:qpnonlinear}

We pass from the scalar equation to vector fields. Multiplying the whole
field by a scalar monomial normalizes
\begin{equation}\label{dyn:eq:qpnormalized}
 \gamma_1=0,\qquad 0\le\kups=\gamma_r ,
\end{equation}
which subtracts $\gamma_1$ from both the frequency and the perturbation
valuations, so the eventual condition $v(\mathcal V)>\kups$ is unchanged
when translated back. Throughout this section $\Gamma$ is nontrivial,
$\qlat=\{0\}$ (tag $\lnot\qres$: no exact resonances) and
$(\sdss)$ holds. The operator $\Gups$ then inverts $\Cups$ on all mean-zero
elements and is applied componentwise. Let
\begin{equation}\label{dyn:eq:qpvectorfield}
 \mathcal X(\theta)=\Upsilon+\mathcal V(\theta),\qquad
 \mathcal V\in\HH_\Gamma(\strip)^{d},\qquad
 s_0=v(\mathcal V)>\kups .
\end{equation}

\begin{theorem}[Hahn-analytic constant normal form]
\label{dyn:thm:qpnormalform}
Under the preceding assumptions there is a \emph{unique} pair
\[
 h\in\HH_\Gamma(\strip)^{d},\qquad \mathfrak c\in K_\Gamma^{d},
 \qquad v(h)>0,\qquad \av(h)=0,
\]
such that, with $\Phi_h=\id+h$,
\begin{equation}\label{dyn:eq:qpconjugacy}
 D\Phi_h(\theta)\,(\Upsilon+\mathfrak c)=\mathcal X(\Phi_h(\theta)).
\end{equation}
Moreover:
\begin{enumerate}[label=(\roman*),leftmargin=*]
\item $v(h)\ge s_0-\kups$ and $v(\mathfrak c)\ge s_0$;
\item with $B=\supp\mathcal V$ and
\begin{equation}\label{dyn:eq:qpmonoid}
 P_0=\bigcup_{j=1}^{r}(B-\gamma_j),\qquad
 P=\mathsf E\cup P_0,\qquad \mathsf M=P^{*},
\end{equation}
these are positive generators in a well-ordered set, and
$\supp h\subseteq\mathsf M\setminus\{0\}$,
$\supp\mathfrak c\subseteq B+\mathsf M$;
\item $\Phi_h$ has an infinitesimal Taylor inverse on
$\HH_\Gamma(\strip)$, and all nonzero integer divisors of
$\Upsilon+\mathfrak c$ have the same leading real coefficients and the same
valuations as those of $\Upsilon$;
\item if $\mathcal V$ has real values coefficientwise on the real torus,
then $h$ and $\mathfrak c$ are real in the same sense, with
$\mathfrak c\in F_\Gamma^{d}$.
\end{enumerate}
The case $\mathcal V=0$ is included, with $h=\mathfrak c=0$.
\end{theorem}

\subsection{Reduction to a fixed-point equation}

Taking the torus mean of \eqref{dyn:eq:qpconjugacy} and using
$\av(Dh)=0$ gives
\begin{equation}\label{dyn:eq:qpdrift}
 \mathfrak c=\mathfrak c(h):=\av\bigl(\mathcal V\circ(\id+h)\bigr),
\end{equation}
and the remaining equation becomes
\begin{equation}\label{dyn:eq:qpfixedpoint}
 h=\mathcal T(h):=
 \Gups\Bigl(\Qz\bigl(\mathcal V\circ(\id+h)\bigr)
 -Dh\,\mathfrak c(h)\Bigr).
\end{equation}
The argument of $\Gups$ has zero mean, since $\mathfrak c(h)$ is constant
and every column of $Dh$ has zero mean, so the inverse is used only where
it is defined.

All terms of $B-\gamma_j$ are positive because $s_0>\kups\ge\gamma_j$. The
sets in \eqref{dyn:eq:qpmonoid} are well ordered and \cref{dyn:lem:neumann}
applies. If $h$ has positive support in $\mathsf M$ then
\cref{dyn:lem:qpsubstitution} gives
$\supp(\mathcal V\circ(\id+h))\subseteq B+\mathsf M$ and the same for
$\mathfrak c(h)$; $Dh\,\mathfrak c(h)$ adds $\mathsf M$ once more; and
applying an inverse coefficient at $-\gamma_j+\mathsf E^{*}$ shifts these
supports into $B-\gamma_j+\mathsf M\subseteq\mathsf M$. Every resulting
term contains at least one letter from $P_0$, so no constant exponent
appears. This proves preservation of a support certificate. Existence still
requires proof that infinitely many iteration steps assemble into a Hahn
series.

\subsection{Existence by marked support words}
\label{dyn:subsec:qpmarked}

\begin{proof}[Existence in \cref{dyn:thm:qpnormalform}]
Put $h^{(0)}=0$ and $h^{(n+1)}=\mathcal T(h^{(n)})$. Each iterate lies in
$\HH_\Gamma(\strip)^{d}$, has zero mean and positive support in
$\mathsf M$, its coefficients obtained by finitely many analytic operations
at each exponent.

Call a generator in $P_0$ a \emph{marked} letter and a generator in
$\mathsf E$ \emph{unmarked}; an exponent lying in both keeps both labels.
There are only finitely many origins for a given letter --- $j$ ranges over
finitely many values and the translation determines the original exponent
of $B$ or $S$ --- so the finite-word conclusion of \cref{dyn:lem:neumann}
is unchanged by the labels.

A term of $\mathcal T(h)$ consists of a coefficient of $\mathcal V$, an
inverse coefficient, and finitely many factors of $h$ or its derivatives.
The combined exponent of the $\mathcal V$ coefficient and the leading shift
$-\gamma_j$ of that inverse supplies one marked letter $\beta-\gamma_j$;
the higher inverse terms supply unmarked letters in $\mathsf E$; all other
letters come from the factors of $h$. By induction, the difference
\begin{equation}\label{dyn:eq:qpiteratedifference}
 h^{(n+1)}-h^{(n)}
\end{equation}
can have a nonzero coefficient only at an exponent representable by a word
with at least $n+1$ \emph{marked} letters. For $n=0$ there is at least one
$\mathcal V$ factor. For the induction step, subtract the two Taylor
expressions and telescope the finite products at each coefficient: every
difference term contains a factor from $h^{(n)}-h^{(n-1)}$ or its
derivative together with one additional $\mathcal V$ coefficient. The term
$Dh\,\mathfrak c(h)$ has the same property, by writing its difference as
$D(h-g)\mathfrak c(h)+Dg(\mathfrak c(h)-\mathfrak c(g))$ and expanding the
second difference by the same telescoping. Applying $\Gups$ creates one new
marked shift and unmarked inverse-tail letters, exactly as above.

At a fixed exponent $\gamma$, \cref{dyn:lem:neumann} permits only finitely
many labelled word representations, so their numbers of marked letters are
bounded; hence \eqref{dyn:eq:qpiteratedifference} has zero coefficient at
$\gamma$ for all large $n$. The coefficients $[t^{\gamma}]h^{(n)}$
therefore stabilize \emph{exactly}, not approximately. The union of the
difference supports lies in the well-ordered $\mathsf M$ and each exponent
occurs in finitely many differences, so the differences form a
Hahn-summable family in the sense of \cref{dyn:def:summable}; let $h$ be
their sum. Each stabilized coefficient lies in $\strip$; no new analytic
limit has been taken. At a fixed exponent $\mathcal T(h)$ uses only
finitely many certified word decompositions and finitely many stabilized
input coefficients, so passing to the stabilized coefficients gives
$h=\mathcal T(h)$. Define $\mathfrak c$ by \eqref{dyn:eq:qpdrift};
\eqref{dyn:eq:qpfixedpoint} and \eqref{dyn:eq:qpdrift} give
\eqref{dyn:eq:qpconjugacy}.

Since substitution does not lower valuation, $v(\mathfrak c)\ge s_0$; the
right side inside $\Gups$ has valuation at least $s_0$ because $Dh$ has
positive valuation; and \cref{dyn:prop:qpvalbound} or
\eqref{dyn:eq:qpgreensupport} gives $v(h)\ge s_0-\kups$. The support
statements are the construction.
\end{proof}

\begin{warning}[The word argument is the existence step, and the obvious
substitute is invalid]\label{dyn:warn:qpcontraction}
The weaker estimate
$v(h^{(n+1)}-h^{(n)})\ge(n+1)(s_0-\kups)$ does \textbf{not} suffice in a
general ordered group: that sequence of lower bounds need not be cofinal.
With $\Gamma=\Z+\Z\Omega$ and $v(h)=1$ the bounds never pass $\Omega$, so a
sequential valuation argument would prove nothing at or beyond that
exponent, and \cref{dyn:ex:qpthree} exhibits a solution with coefficients
exactly there. No Banach fixed-point theorem and no sequential
completeness of $\No[i]$ is used. This is the same refusal as
\cref{dyn:warn:contraction}, reached independently by source 07 and
implemented by a \emph{labelled} refinement of the word certificate ---
counting marked letters --- which is the one mechanism in this report that
\cref{dyn:cor:ancestry} does not already supply.
\end{warning}

\subsection{Uniqueness and the inverse coordinate map}

\begin{proof}[Uniqueness in \cref{dyn:thm:qpnormalform}]
Let $h,g$ be positive-valuation mean-zero solutions with $h\ne g$ and put
$\delta=v(h-g)>0$; their drifts are forced by \eqref{dyn:eq:qpdrift}. By
\eqref{dyn:eq:qplipschitz},
$v(\mathfrak c(h)-\mathfrak c(g))\ge s_0+\delta$, and subtracting the
expressions inside \eqref{dyn:eq:qpfixedpoint} gives valuation at least
$s_0+\delta$, since $\mathfrak c(h)$ has valuation at least $s_0$ and
differentiation does not lower valuation. Hence
$v(\mathcal T(h)-\mathcal T(g))\ge s_0+\delta-\kups>\delta$, contradicting
$h-g=\mathcal T(h)-\mathcal T(g)$. Uniqueness holds among \emph{all}
positive-valuation mean-zero solutions, not only those carrying the
constructed support certificate.
\end{proof}

\begin{lemma}[Infinitesimal coordinate inverses on the torus]
\label{dyn:lem:qpcoordinverse}
Every $\Phi_h=\id+h$ with $v(h)>0$ has a unique inverse $\id+g$ of positive
valuation in the Taylor sense, inducing an automorphism of
$\HH_\Gamma(\strip)$.
\end{lemma}
\begin{proof}
Solve $g=-h\circ(\id+g)$ by iteration from zero: the same marked-word
stabilization, now marking each occurrence of a coefficient of $h$, gives
coefficient stabilization inside the positive monoid generated by
$\supp h$. Uniqueness follows from
$v(h\circ(\id+g_1)-h\circ(\id+g_2))\ge v(h)+v(g_1-g_2)>v(g_1-g_2)$, which
is \eqref{dyn:eq:qplipschitz}. For two-sidedness write
$\Psi\circ\Phi_h=\id+q$; associativity gives
$\Phi_h\circ(\id+q)=\Phi_h$, hence $q+h\circ(\id+q)-h=0$, and if $q\ne0$
the second difference has valuation strictly greater than $v(q)$, which is
impossible. This is the torus counterpart of \cref{dyn:lem:inverse} and
\cref{dyn:prop:subgroup}, proved here by the labelled word certificate
rather than by the triangular recursion, because the unknown is a full
Hahn-analytic coefficient family rather than a Taylor coefficient list.
\end{proof}

\begin{proof}[Completion of \cref{dyn:thm:qpnormalform}]
\Cref{dyn:lem:qpcoordinverse} gives the inverse assertion. Since
$v(\mathfrak c)\ge s_0>\kups$, \cref{dyn:cor:qptail} applies and no leading
coefficient of a nonzero integer divisor changes; for complex data the same
leading real coefficients persist although some higher coefficients can be
complex. Complex conjugation together with $k\mapsto-k$ preserves all
scalar Fourier inversions for a real frequency, and Taylor substitution,
mean and differentiation preserve coefficientwise reality, so real
$\mathcal V$ gives real $h,\mathfrak c$; alternatively reality follows from
uniqueness.
\end{proof}

\subsection{Necessity, and the sharp scale boundary}
\label{dyn:subsec:qpnecessity}

\begin{theorem}[Non-linearizable perturbations beyond any prescribed
valuation]\label{dyn:thm:qpnecessity}
Assume $\gamma_1=0$, $\qlat=\{0\}$ and $\Gamma$ nontrivial. If $(\sdss)$
fails, then for every prescribed $H\in\Gamma$ there are $s>H$, a real
$\rho>0$ and a real perturbation $\mathcal V\in\HH_\Gamma(\strip)^{d}$ with
$v(\mathcal V)=s$ such that $\Upsilon+\mathcal V$ admits \emph{no}
positive-valuation mean-zero Hahn-analytic conjugacy to a constant vector
field, even on a narrower positive-width strip. One may require
$s>2\kups$ and take $\mathcal V$ with a single Hahn exponent.
\end{theorem}
\begin{proof}
Choose $s>\max\{H,2\kups\}$, possible in a nontrivial ordered abelian
group. Take the real $f$ of \eqref{dyn:eq:qpbadforcing} on a failing
stratum $\Sigma_j$ and put $\mathcal V=t^{s}fe_1$. Suppose
$\Phi_h=\id+h$ were such a conjugacy on some strip; restrict to a common
narrower strip if necessary. Taking means still forces
$\mathfrak c=\av(\mathcal V\circ\Phi_h)$ with $v(\mathfrak c)\ge s$, and
since $v(h)>0$ the equation
$\Cups h=\Qz(\mathcal V\circ\Phi_h)-Dh\,\mathfrak c$ has right side of
valuation at least $s$, so \cref{dyn:prop:qpvalbound} gives
$v(h)\ge s-\kups$. Subtracting the uncomposed forcing, write
$\Cups h=t^{s}fe_1+R$ with
$R=\Qz(\mathcal V\circ\Phi_h-\mathcal V)-Dh\,\mathfrak c$; Taylor
telescoping \eqref{dyn:eq:qplipschitz} gives $v(R)\ge s+v(h)\ge2s-\kups$.
At every nonzero Fourier mode the scalar equation therefore gives
\[
 \widehat h(k)=\frac{t^{s}\widehat f(k)e_1}{i\,k\cdot\Upsilon}
 +\frac{\widehat R(k)}{i\,k\cdot\Upsilon},\qquad
 v\Bigl(\frac{\widehat R(k)}{i\,k\cdot\Upsilon}\Bigr)\ge2s-2\kups>s .
\]
For $k=k_n$ in the failing stratum the coefficient of $t^{s-\gamma_j}$ in
the first component of $h$ is therefore forced to be
$e^{-\eps|k_n|}/(i\,\fla_j(k_n))$ --- here $s-\gamma_j\le s$, so the
remainder cannot contribute --- and those coefficients grow faster than
$e^{\eps|k_n|}$ and cannot belong to a distribution, exactly as in the
necessity proof of \cref{dyn:thm:qplinear}.
\end{proof}

\begin{corollary}[Exact universal infinitesimal normal-form criterion]
\label{dyn:cor:qpequivalence}
For normalized $\Upsilon$ with $\qlat=\{0\}$ over nontrivial $\Gamma$,
$(\sdss)$ is equivalent to the assertion that for every $\rho>0$ every
$\mathcal V\in\HH_\Gamma(\strip)^{d}$ of valuation greater than $\kups$
admits the normalized conjugacy of \cref{dyn:thm:qpnormalform}. When
$(\sdss)$ fails, restricting perturbations to arbitrarily higher valuations
does \textbf{not} rescue a universal theorem.
\end{corollary}

\begin{proposition}[Failure exactly at the threshold, with no zero of the
slow speed]\label{dyn:prop:qpboundary}
Let $\kups>0$ lie in $\Gamma$ and take
\[
 \Upsilon=(1,t^{\kups}),\qquad
 \mathcal V(\theta_1,\theta_2)=\bigl(0,\tfrac12t^{\kups}\cos\theta_2\bigr).
\]
Then $\Upsilon$ satisfies $(\sdss)$ and its last flag level is $\kups$, yet
this perturbation of valuation \emph{exactly} $\kups$ admits no
positive-valuation mean-zero conjugacy to a constant vector field.
\end{proposition}
\begin{proof}
The two strata have leading forms $k_1$ and $k_2$, so $(\sdss)$ is
immediate. Suppose a conjugacy existed. The mean equation forces
$\mathfrak c_1=0$; since substitution by a positive-valuation $h$ does not
change the coefficient at exponent $\kups$ of $\mathcal V$, that
coefficient is $(0,\tfrac12\cos\theta_2)$, of mean zero, so
$v(\mathfrak c_2)>\kups$ (with $\mathfrak c_2=0$ allowed). At exponent
$\kups$ in the second component of \eqref{dyn:eq:qpconjugacy},
$t^{\kups}\partial_2h_2$ contributes nothing because $[t^{0}]h_2=0$, and
neither do the drift or $Dh\,\mathfrak c$. What remains is
$\partial_1\bigl([t^{\kups}]h_2\bigr)=\tfrac12\cos\theta_2$; averaging in
$\theta_1$ makes the left side zero and leaves the right side unchanged.
\end{proof}

The slow component of this field is
$t^{\kups}\bigl(1+\tfrac12\cos\theta_2\bigr)$, whose leading real factor is
everywhere positive: the example does \textbf{not} rely on introducing a
stationary point. It says that a perturbation of the \emph{same} size as a
slow frequency may require a coordinate correction of order one rather than
an infinitesimal one. The strict threshold $v(\mathcal V)>\kups$ is
therefore sharp \emph{as a uniform sufficient bound}; particular
perturbations at the boundary can of course still be linearizable.

\subsection{A canonical invariant density}
\label{dyn:subsec:qpdensity}

The claim here concerns coefficientwise integration of analytic functions
on the torus. It is \textbf{not} a countably additive $\No$-valued
probability measure on arbitrary subsets, and not a real-time ergodic
theorem; \cref{dyn:warn:periods} and \cref{dyn:warn:hamflow} record the
same refusal elsewhere in this report.

\begin{lemma}[Formal change of variables]\label{dyn:lem:qpchangevar}
For $\Phi_h=\id+h$ with $v(h)>0$ and every $f\in\HH_\Gamma(\strip)$,
\begin{equation}\label{dyn:eq:qpchangevar}
 \av\bigl((f\circ\Phi_h)\det D\Phi_h\bigr)=\av(f).
\end{equation}
\end{lemma}
\begin{proof}
Introduce an ordinary formal auxiliary variable $z$ and
$\Phi_z=\id+zh$. At each Hahn exponent Taylor substitution and the
determinant produce a \emph{polynomial} in $z$, there being finitely many
support words at that exponent (\cref{dyn:lem:neumann}). With $A=D\Phi_z$,
the finite determinant identities and the chain rule give
\begin{equation}\label{dyn:eq:qppiola}
 \frac{\mathrm d}{\mathrm dz}\bigl((f\circ\Phi_z)\det A\bigr)
 =\diver_\theta\bigl((f\circ\Phi_z)\operatorname{adj}(A)h\bigr):
\end{equation}
use $A\operatorname{adj}(A)=(\det A)I$ for the derivative of $f$, the
derivative-of-determinant formula for the $Dh$ terms, and the elementary
Piola identity --- the remaining cofactor derivatives cancel in pairs
because mixed partials of $\Phi_z$ commute. Each coefficient on the right
of \eqref{dyn:eq:qppiola} is a divergence of a periodic analytic vector
field, hence of mean zero, so the mean on the left is independent of $z$;
evaluate the coefficient polynomials at $z=0$ and $z=1$. There is no
infinite integration in $z$ and no ordinary convergence of a Hahn sum is
assumed.
\end{proof}

\begin{theorem}[Invariant density and its uniqueness]\label{dyn:thm:qpdensity}
Let $\mathcal X$, $\Phi_h$ and $\mathfrak c$ be as in
\cref{dyn:thm:qpnormalform} and put $\Psi=\Phi_h^{-1}$. Then
\begin{equation}\label{dyn:eq:qpdensity}
 \mu_{\mathcal X}=\det D\Psi\in1+\HH_\Gamma(\strip)_{>0}
\end{equation}
is the unique density in $\HH_\Gamma(\strip)$ with
$\diver(\mu_{\mathcal X}\mathcal X)=0$ and $\av(\mu_{\mathcal X})=1$, and
the $K_\Gamma$-linear functional
\begin{equation}\label{dyn:eq:qpfunctional}
 I_{\mathcal X}(f)=\av(f\circ\Phi_h)=\av(f\mu_{\mathcal X})
\end{equation}
satisfies $I_{\mathcal X}(1)=1$ and
$I_{\mathcal X}(\mathcal X\cdot\nabla f)=0$. For real data
$\mu_{\mathcal X}$ is real, with positive surcomplex values on the real
torus, being one plus an infinitesimal.
\end{theorem}
\begin{proof}
The inverse is infinitesimal, so its Jacobian determinant is one plus an
element of positive valuation, and the chain rule gives
$(\mu_{\mathcal X}\circ\Phi_h)\det D\Phi_h=1$. Applying
\cref{dyn:lem:qpchangevar} to $f\mu_{\mathcal X}$ gives the equality in
\eqref{dyn:eq:qpfunctional}, and $f=1$ gives the normalization.
Conjugacy and the chain rule yield
$(\mathcal X\cdot\nabla f)\circ\Phi_h=(\Upsilon+\mathfrak c)\cdot
\nabla(f\circ\Phi_h)$, whose mean is zero; coefficientwise integration by
parts against every Fourier character gives
$\diver(\mu_{\mathcal X}\mathcal X)=0$. For uniqueness let
$\widetilde\mu$ be a normalized solution and set
$g=(\widetilde\mu\circ\Phi_h)\det D\Phi_h$; the same Jacobian identities
transport its divergence equation to
$(\Upsilon+\mathfrak c)\cdot\nabla g=0$, and since $\Upsilon+\mathfrak c$
has no nonzero integer resonances (\cref{dyn:thm:qpnormalform}(iii)),
Fourier uniqueness makes $g$ constant in $\theta$. Then
\eqref{dyn:eq:qpchangevar} gives $\av(g)=\av(\widetilde\mu)=1$, so $g=1$
and $\widetilde\mu=\mu_{\mathcal X}$. Reality and pointwise positivity
follow from \cref{dyn:thm:qpnormalform}(iv) and the positive-valuation form
of \eqref{dyn:eq:qpdensity}.
\end{proof}

The identity $I_{\mathcal X}(\mathcal X\cdot\nabla f)=0$ is the precise
invariance assertion. No real-time flow taking values continuously in the
fine topology of $\No[i]$ is constructed, and no ergodic theorem for such a
flow is claimed.

\subsection{Worked examples with genuinely different scales}
\label{dyn:subsec:qpexamples}

\begin{example}[Three visible scales and a reciprocal support of order type
$\omega^{2}$]\label{dyn:ex:qpthree}
Take $\Gamma=\Z+\Z\Omega$ with $\Omega$ greater than every ordinary
positive integer --- the same use of the symbol $\Omega$ as in
\cref{dyn:ex:ranktwo,dyn:ex:rank2}, an element of $\Gamma$ and not a
frequency (\cref{dyn:warn:twofrequencies}) --- and
\begin{equation}\label{dyn:eq:qpthreefreq}
 \Upsilon=(1,t,t^{\Omega}).
\end{equation}
The flag is $\Z^{3}\supset\{k_1=0\}\supset\{k_1=k_2=0\}\supset\{0\}$ with
levels $0,1,\Omega$ and leading forms $k_1,k_2,k_3$ on their strata. Each
nonzero leading form is a nonzero integer, so its reciprocal has modulus at
most one and $(\sdss)$ holds with no delicate ordinary arithmetic at all.
Nevertheless the mode $k=(1,1,1)$ already has
\begin{equation}\label{dyn:eq:qpranktwoinverse}
 \frac1{1+t+t^{\Omega}}
 =\sum_{n,m\ge0}(-1)^{n+m}\binom{n+m}{n}t^{\,n\Omega+m},
\end{equation}
whose support has order type $\omega^{2}$ --- the ordinal $\omega$, not a
frequency --- since every exponent of the $n=0$ layer precedes every
exponent of the $n=1$ layer, and so on. Formula
\eqref{dyn:eq:qpranktwoinverse} follows by grouping words according to
their numbers of letters $1$ and $\Omega$, and is confirmed by the finite
recurrence $C_{n,m}+C_{n-1,m}+C_{n,m-1}=\delta_{(n,m),(0,0)}$ for
$C_{n,m}=(-1)^{n+m}\binom{n+m}{n}$ (negative index meaning zero), which the
verification program checks.

On the second stratum
$\bigl(k_2t+k_3t^{\Omega}\bigr)^{-1}
 =\frac{t^{-1}}{k_2}\sum_{n\ge0}(-k_3/k_2)^{n}t^{\,n(\Omega-1)}$, and on
the last stratum the inverse is simply $t^{-\Omega}/k_3$: a large support
coexists with a three-level divisor classification. The forcing
$f=t^{\Omega+1}(\sin\theta_1+\sin\theta_2+\sin\theta_3)$ has exact
mean-zero solution
\begin{equation}\label{dyn:eq:qpthreesolution}
 u=-t^{\Omega+1}\cos\theta_1-t^{\Omega}\cos\theta_2-t\cos\theta_3 ,
\end{equation}
of valuation $1=v(f)-\Omega$, so the last scale really does control the
worst inverse loss (\cref{dyn:cor:qpsharp}). A vector perturbation of
valuation $\Omega+1$ falls under \cref{dyn:thm:qpnormalform} with a
coordinate correction of valuation at least one. This is also the example
that exposes the invalidity of a sequential contraction proof
(\cref{dyn:warn:qpcontraction}): corrections of order $1,2,3,\dots$ never
pass $\Omega$, while the coefficient-stabilization proof constructs the
exact transformation including its coefficients at and beyond that scale.
\end{example}

\begin{example}[Real arithmetic on the leading stratum, and an arbitrary
tail]\label{dyn:ex:qpsqrt2}
Take $\Upsilon=(1,\sqrt2,t)$. Its first stratum is all $(p,q,r)$ with
$(p,q)\ne(0,0)$, with leading form $p+\sqrt2q$; its second is $(0,0,r)$
with $r\ne0$, at level one. Since
$|p+\sqrt2q|^{-1}=|p-\sqrt2q|/|p^{2}-2q^{2}|\le|p|+\sqrt2|q|$ for
$(p,q)\ne(0,0)$, the first stratum is Diophantine and the second has the
integer leading form $r$; the universal conjugacy threshold is
$v(\mathcal V)>1$. One may now append an arbitrary set-supported real
vector tail of valuation greater than one: it can have transfinite support
and change every higher coefficient of the solution, yet by
\cref{dyn:cor:qptail} it changes neither the arithmetic test nor the
threshold. This separates ordinary approximation properties from surreal
support complexity.
\end{example}

\begin{example}[An explicit slow-circle normal form]\label{dyn:ex:qpcircle}
Let $s>\kups>0$ and consider
$\mathcal X(\theta_1,\theta_2)=(1,\,t^{\kups}+t^{s}\sin\theta_2)$, and put
$\varrho=t^{\,s-\kups}$. The mean-zero conjugacy has the form
$\Phi_h=(\theta_1,\theta_2+h(\theta_2))$ and carries $\mathcal X$ to
$(1,\nu)$, where
\begin{equation}\label{dyn:eq:qpcircleconjugacy}
 (1+h'(\theta))\nu=t^{\kups}+t^{s}\sin(\theta+h(\theta)),
\end{equation}
and the new frequency is exactly
\begin{equation}\label{dyn:eq:qpcirclefrequency}
 \nu=t^{\kups}\sqrt{1-\varrho^{2}}
 =t^{\kups}\Bigl(1-\tfrac12\varrho^{2}-\tfrac18\varrho^{4}
 -\tfrac1{16}\varrho^{6}-\cdots\Bigr),
\end{equation}
the square root being the Hahn binomial branch with constant term one. To
see this, use the inverse coordinate in the second variable: the inverse
equation gives $\Psi'(\theta)=\nu\,t^{-\kups}(1+\varrho\sin\theta)^{-1}$,
whose mean is one, so
$t^{\kups}/\nu=\av\bigl((1+\varrho\sin\theta)^{-1}\bigr)
 =\sum_{m\ge0}\binom{2m}{m}\varrho^{2m}4^{-m}=(1-\varrho^{2})^{-1/2}$,
using the geometric expansion, the vanishing of odd sine moments and
$\av(\sin^{2m}\theta)=\binom{2m}{m}4^{-m}$. All of these are Hahn
identities in the positive monomial $\varrho$. The first terms of the
unique mean-zero $h$ are
\begin{equation}\label{dyn:eq:qpcircleh}
 h(\theta)=-\varrho\cos\theta-\frac{\varrho^{2}}{4}\sin2\theta
 +\varrho^{3}\Bigl(\frac1{12}\cos3\theta-\frac14\cos\theta\Bigr)
 +O_H(\varrho^{4}),
\end{equation}
where $O_H(\varrho^{4})$ means that the remaining integer powers of
$\varrho$ have exponent at least four and is \textbf{not} an assertion
about convergence in a real parameter. In particular the lower bound
$v(h)\ge s-\kups$ of \cref{dyn:thm:qpnormalform}(i) is attained, and the
first frequency drift occurs later, at valuation $2s-\kups$. The invariant
density of \cref{dyn:thm:qpdensity} is equally explicit,
\begin{equation}\label{dyn:eq:qpcircledensity}
 \mu_{\mathcal X}(\theta_2)
 =\frac{\sqrt{1-\varrho^{2}}}{1+\varrho\sin\theta_2},
\end{equation}
with mean one, making $\mu_{\mathcal X}\cdot
(t^{\kups}+t^{s}\sin\theta_2)$ independent of $\theta_2$. The verification
program reconstructs $h$ and $\nu$ independently by exact Fourier
arithmetic and checks \eqref{dyn:eq:qpcircleconjugacy} through degree ten
in $\varrho$.
\end{example}

\subsection{What this part does not reach}
\label{dyn:subsec:qplimits}

Three delimitations belong here rather than only in
\cref{dyn:sec:nonclaims}, because they bound the reading of every statement
above.

First, all of it is a \emph{coefficientwise Hahn} theory. A formal series
$\sum_{n\ge0}t^{n}h_n(\theta)$ lies in $\HH_\Gamma(\strip)$ whenever every
$h_n$ is holomorphic on the same strip, \emph{however fast} $\|h_n\|_{r}$
grows --- faster than every exponential in $n$ is permitted --- so
substituting a nonzero real number for $t$ need not be meaningful. Nothing
here improves classical real analytic KAM or Brjuno arithmetic: those
theorems must control Fourier decay and the accumulation of losses through
infinitely many nonlinear steps simultaneously \cite{Bounemoura,Chevallier},
whereas $(\sdss)$ only guarantees that each \emph{individual} coefficient
operation stays holomorphic, the Hahn support argument handling how the
coefficients are assembled. This is the torus instance of
\textbf{N55}, \textbf{N56} and \cref{dyn:warn:refuse}.

Second, the torus derivatives are external derivations annihilating
$K_\Gamma$ (\cref{dyn:warn:derivative}); nothing above speaks to internal
surreal differential equations, and source 07 says so in order to keep
itself separate from the companion report on differential equations
(\cref{dyn:subsec:diffeq}).

Third, the nonlinear theorem requires $\qlat=\{0\}$. The linear theorem
allows $\qlat\ne\{0\}$, but the resonant nonlinear case is open, and
source 07 warns explicitly that replacing the mean projection by
$\Pi_{\qlat}$ is \emph{not} justified, since the retained normal form is
then no longer a constant vector field and creates additional terms in the
conjugacy equation (\cref{dyn:q:qpresonant}). The flag itself is finite but
not automatically computable: detecting exact integer relations among
arbitrary named reals may be unavailable, so existence of the flag is
separated from decidability of its entries (\cref{dyn:q:qpeffective}).

\section{Finite-order multipliers and resonant splitting}
\label{dyn:sec:resonant-normal}

Nonresonance fails for multipliers of finite order. There is nevertheless
an exact normal form retaining precisely the resonant dynamics, with
\emph{no} small-divisor estimates anywhere (tag $\sdz$). Note that the
hypothesis $\Lambda^{q}=I$ here is the multiplier \emph{itself} of finite
order, which is neither $\resq$ --- where $\lambda$ is explicitly not a
root of unity --- nor mere $\mres$ (\cref{dyn:warn:resonance}).

\begin{theorem}[Same-domain finite-order decomposition]
\label{dyn:thm:finite-order}
Let $\Lambda$ be diagonal unitary with $\Lambda^{q}=I$ for an integer
$q\ge1$, and let $F=\Lambda z+f$ with
$f\in\HH_\Gamma(\OO(\mathbb D_R^{d}))_+^{d}$ vanishing coefficientwise to
order at least two. There exist a normalized positive coordinate change
$H$ and a positive vector field $Y$ vanishing to order at least two, both
on the \emph{same} $\mathbb D_R^{d}$, with
\begin{equation}\label{dyn:eq:resonant-normal}
 H\circ F\circ H^{-1}=\Lambda\circ\Phi_Y,\qquad Y(\Lambda z)=\Lambda Y(z),
\end{equation}
the rotation and flow in \eqref{dyn:eq:resonant-normal} commuting. Every
monomial $z^{\alpha}e_j$ of $Y$ satisfies $\lambda^{\alpha}=\lambda_j$. All
coefficients have finite support-word certificates from $\supp f$, and the
construction commutes with workspace enlargement. Moreover
\begin{equation}\label{dyn:eq:finiteorder-equivalence}
 F^{\circ q}=\id\iff Y=0\iff
 F\text{ is positively conjugate to }\Lambda .
\end{equation}
The same statement holds for radius-free germs, entire coefficients,
polynomials and formal coefficients.
\end{theorem}
\begin{proof}
The reduction of $P=F^{\circ q}$ is the identity, it fixes zero and has
exact Jacobian $I$ there, so $X=\log C_P$ is a positive vector field whose
coefficients vanish to order at least two (\cref{dyn:thm:exp-log}). Since
$F$ commutes with $P$, its pullback commutes with $\log C_P$, so $F$
preserves $X$ and commutes with all its admissible ordinary flow times.
Set $A_{\mathrm f}=F\circ\Phi_{-X/q}$: the factors commute, so
$A_{\mathrm f}^{\circ q}=P\circ\Phi_{-X}=\id$, and its ordinary reduction
is $\Lambda$. Define the finite average
\begin{equation}\label{dyn:eq:average}
 H(z)=\frac1q\sum_{j=0}^{q-1}\Lambda^{-j}A_{\mathrm f}^{\circ j}(z).
\end{equation}
The leading term of every summand after multiplication by $\Lambda^{-j}$ is
$z$, so $H=\id+$ positive with value zero and derivative $I$ at zero, with
a positive inverse on the same coefficient domain; reindexing the finite sum
and using $A_{\mathrm f}^{q}=\id$, $\Lambda^{q}=I$ gives
$H\circ A_{\mathrm f}=\Lambda H$. With
$(H_{*}X)(w)=\mathrm dH(H^{-1}(w))\,a(H^{-1}(w))$ for $X=a\cdot\partial$,
Taylor substitution shows $H_{*}X$ is positive of the same analytic
category and that $H\circ\Phi_X\circ H^{-1}=\Phi_{H_{*}X}$, either by the
coefficientwise chain rule or by conjugating pullback operators. Put
$Y=\frac1qH_{*}X$: since $A_{\mathrm f}$ preserves $X$, the conjugated map
$\Lambda$ preserves $Y$, which is exactly $Y(\Lambda z)=\Lambda Y(z)$ and
gives the monomial resonance condition; and $F=A_{\mathrm f}\circ\Phi_{X/q}$
gives \eqref{dyn:eq:resonant-normal} with its commuting factors. The
operations used are finite compositions with the ordinary rotation,
positive logarithms and exponentials, substitution inverses and the finite
sum \eqref{dyn:eq:average}, each preserving the original coefficient domain
and the generated support monoid. Finally $P=\id$ iff $X=0$ iff $Y=0$; if
$Y=0$ the displayed normal form is a conjugacy to $\Lambda$, and
conversely such a conjugacy forces $F^{q}=\id$.
\end{proof}

In one variable with $\lambda$ of exact order $q$, the resonant vector field
has the form $Y(z)=\sum_{n\ge2,\ n\equiv1\!\pmod q}a_nz^{n}\partial_z$,
interpreted coefficientwise when the $a_n$ are themselves Hahn series; the
coefficient \emph{family}, rather than an unrestricted rearrangement in
$n$, remains the defining object. For $q=1$ the theorem reduces to
\cref{dyn:thm:exp-log}.

\begin{warning}[Credit and non-uniqueness]\label{dyn:warn:finiteorder}
The finite-factor/flow splitting follows the classical
semisimple--unipotent pattern \cite{Ribon}, and neither the finite
averaging idea nor that abstract pattern is claimed new; what is
established here is that the entire construction is valid for arbitrary
positive Hahn support and retains the same ordinary coefficient domain.
Uniqueness of resonant normal forms is \textbf{not} claimed: further
coordinate changes commuting with $\Lambda$ can transform $Y$, and formula
\eqref{dyn:eq:average} specifies the particular normalized coordinate used.
For resonant multipliers of \emph{infinite} order no universal analytic
Poincar\'e--Dulac statement is proved; see \cref{dyn:q:resonant}.
\end{warning}

\section{Critical scales: four different things that fail}
\label{dyn:sec:scales}

Four sources each exhibit a scale at which something stops working, and the
four failures are \emph{not} the same failure. Collecting them is what
prevents a reader from concluding that the report has one scale obstruction
with four proofs.

\begin{proposition}[Quadratic scaling and the boundary of the halo theorem]
\label{dyn:prop:quadratic-scale}
For every non-root-of-unity $\lambda\in\C^{*}$ and every $\eta>0$, the map
$F(z)=\lambda z+t^{\eta}z^{2}$ has a normalized conjugacy with
\emph{polynomial} Hahn coefficients
\begin{equation}\label{dyn:eq:scaled-H}
 H(z)=z+\sum_{n\ge1}b_nt^{n\eta}z^{n+1},
\end{equation}
which evaluates on every finite halo irrespective of the growth of the
$b_n$. For every nonzero ordinary $w$, substituting $z=t^{-\eta}w$ into the
individual terms of \eqref{dyn:eq:scaled-H} is \emph{not} a strongly
summable family. If $\lambda$ is a non-Brjuno irrational rotation, the
ordinary series $w+\sum_{n\ge1}b_nw^{n+1}$ has zero radius of convergence
as well.
\end{proposition}
\begin{proof}
Let $L_\lambda(w)=w+\sum_{n\ge1}b_nw^{n+1}$ be the unique formal solution
of $L_\lambda(\lambda w+w^{2})=\lambda L_\lambda(w)$, whose coefficients
nonresonance gives recursively; substituting $w=t^{\eta}z$ and dividing by
$t^{\eta}$ gives exactly \eqref{dyn:eq:scaled-H} with
$H\circ F=\lambda H$. Each Hahn coefficient is polynomial, the support is a
subset of $\{n\eta:n\ge1\}$, and evaluation at a finite point is justified
by \cref{dyn:lem:neumann} with no condition on the ordinary magnitudes.
Infinitely many $b_n$ are nonzero: otherwise $L_\lambda$ would be a
nonconstant polynomial and the degrees of $L_\lambda\circ Q_\lambda$ and
$\lambda L_\lambda$ would be $2\deg L_\lambda$ and $\deg L_\lambda$. At the
critical scale every nonzero summand becomes
$b_nt^{n\eta}(t^{-\eta}w)^{n+1}=t^{-\eta}b_nw^{n+1}$, so infinitely many
members share the Hahn exponent $-\eta$, violating
\cref{dyn:def:summable}: this is a failure of \textbf{local finiteness},
not of a real-valued norm estimate. Finally, if $L_\lambda$ had a positive
ordinary radius it would be a local holomorphic coordinate with derivative
one, analytically linearizing $Q_\lambda$, contradicting Brjuno--Yoccoz for
a non-Brjuno rotation.
\end{proof}

Even for a Brjuno rotation, ordinary convergence of the coefficient series
would not by itself make the repeated-weight family a strong Hahn sum: it
would be an additional ordinary analytic summation convention. The
non-Brjuno example rules out even that local analytic reinterpretation of
the formal linearizer. No assertion is made about sectorial or resurgent
summation under additional hypotheses.

\begin{proposition}[A precise rescaled-germ obstruction]\label{dyn:prop:critical}
For the same family, $H_t(z)=t^{-1}h_\lambda(tz)$ with
$h_\lambda(w)=w+\sum_{n\ge2}b_nw^{n}$ the ordinary normalized formal
linearizer of $p_\lambda(w)=\lambda w+w^{2}$, the change of variable $H_t$
and its inverse are well defined by Hahn sums on the \emph{full} set
\begin{equation}\label{dyn:eq:subcritical}
 \{z\in K_\Gamma:v(tz)>0\}
\end{equation}
and give a conjugacy there. Rescaling to the critical scale $z=t^{-1}w$
gives $tH_t(t^{-1}w)=h_\lambda(w)$, and this formal germ is the Taylor germ
of a common-domain Hahn-holomorphic section at $w=0$ \emph{if and only if}
$\rad(h_\lambda)>0$.
\end{proposition}
\begin{proof}
$H_t(z)=t^{-1}h_\lambda(tz)$ is coefficientwise algebra, and
$tF_t(z)=p_\lambda(tz)$ then gives $H_t\circ F_t=\lambda H_t$ by formal
substitution. At any argument with $tz\in\m_\Gamma$ the ordinary formal
series $h_\lambda(tz)$ is Hahn summable --- scalar coefficients of
valuation zero, powers of a positive-support argument, and
\cref{dyn:lem:neumann} --- as is its ordinary formal inverse, and both
preserve the infinitesimal monad because their linear term is the identity
and higher terms have strictly higher valuation; multiplying by $t^{-1}$
preserves \eqref{dyn:eq:subcritical}. For the last assertion, suppose a
section $A(w)=\sum_{\gamma\ge0}t^{\gamma}a_\gamma(w)$ with all $a_\gamma$
holomorphic on one $\mathbb D_{r}$ has Taylor series $h_\lambda(w)$ in
$K_\Gamma[[w]]$: the coefficient at Hahn exponent zero has ordinary Taylor
series $h_\lambda$, which must therefore converge at least on
$\mathbb D_{r}$, all other coefficient germs having zero Taylor series.
Conversely, if $h_\lambda$ converges on a disk, the section with only
exponent-zero coefficient $h_\lambda$ is the desired one.
\end{proof}

For $t=\omega^{-1}$ the domain \eqref{dyn:eq:subcritical} includes every
finite surcomplex number and many infinite ones, and excludes exactly the
critical scale where $tz$ has nonzero ordinary standard part. This is
consistent with the familiar non-Archimedean linearization phenomenon over
a trivially valued coefficient field, and is not presented as a replacement
for the ultrametric results of Lindahl \cite{Lindahl}. Even when
$h_\lambda$ converges ordinarily, substituting a nonzero ordinary $w$
directly into the displayed Hahn family is \emph{not} a Hahn summation:
infinitely many nonzero terms would occupy exponent zero. One must first
use the separately justified ordinary holomorphic germ and then its
canonical Taylor lift, and this is where an illegitimate exchange of
ordinary with Hahn summation would occur.

\begin{warning}[The four critical-scale phenomena are distinct]
\label{dyn:warn:fourscales}
\begin{enumerate}[label=(\arabic*)]
\item \Cref{dyn:prop:quadratic-scale} (source 03): at the critical scale
\emph{local finiteness} fails --- infinitely many terms occupy one exponent
--- and this is what stops the finite-halo theorems from being all-scale
theorems.
\item \Cref{dyn:thm:sharp} (source 04): the sharp disk
$v(z)>\delta/q$ has \textbf{two} distinct failure mechanisms, local
finiteness \emph{at} equality and \emph{well-ordering} strictly below
(\eqref{dyn:eq:termvaluation}), and the boundary nevertheless carries a
legitimate coefficientwise continuation
(\cref{dyn:subsec:boundary}) --- so here the scale is where one summation
stops and a \emph{different} legal operation begins.
\item \Cref{dyn:prop:critical} (source 02): the failure is at the level of
\emph{representability}, an exact iff between positive ordinary radius of
the rescaled germ and existence of a common-domain section, plus the
observation that a coefficientwise-entire section need not be globally
analytic on all of $\No[i]$ under every spatial rescaling.
\item \Cref{dyn:ex:cubic} (source 06): the boundary polydisk carries
\emph{two} independent obstructions --- a pole of an explicit rational
normalizer, and a coordinate-independent \emph{critical-point count} --- so
the failure is not a defect of one geometric-series proof, and
\cref{dyn:ex:badorder,dyn:ex:badfiber} separate well-ordering from finite
fibers as independent requirements.
\end{enumerate}
Sources 02 and 05 each add the further remark that the existence of a Hahn
conjugacy does \emph{not} imply ordinary holomorphic dependence on $t$ near
$t=0$: for a divergent $h_\lambda$ the formal dependence encoded in
\eqref{dyn:eq:scaled-H} need not converge as an ordinary two-variable
series in $(t,z)$; and source 05's \cref{dyn:warn:slowtime} gives the
positive counterpart, one licensed passage to the scale $\eps^{-1}$.
\end{warning}

\section{Worked examples and exact coefficient calculations}
\label{dyn:sec:examples}

\subsection{The quadratic Schr\"oder recursion, once, in four normalizations}
\label{dyn:subsec:quadratic}

Four sources compute the coefficients of the normalized formal linearizer
of an ordinary quadratic. They are one calculation; the translation between
their normalizations is the following.

For the ordinary equation
$L_\lambda(\lambda w+w^{2})=\lambda L_\lambda(w)$ with
$L_\lambda(w)=w+\sum_{n\ge1}b_nw^{n+1}$, the first coefficients are
\begin{align}\label{dyn:eq:quadratic-bs}
 b_1&=-\frac1{\lambda(\lambda-1)},\qquad
 b_2=\frac2{\lambda(\lambda-1)^{2}(\lambda+1)},\\
 b_3&=-\frac{5\lambda^{2}+1}
 {\lambda^{2}(\lambda-1)^{3}(\lambda+1)(\lambda^{2}+\lambda+1)},\notag\\
 b_4&=\frac{2(7\lambda^{3}+2\lambda+3)}
 {\lambda(\lambda-1)^{4}(\lambda+1)^{2}(\lambda^{2}+1)
 (\lambda^{2}+\lambda+1)},\notag
\end{align}
obtained by equating successive powers of $w$, the denominators at stage
$n$ arising from $\lambda^{n+1}-\lambda$. In the shifted indexing
$h_\lambda(w)=w+\sum_{n\ge2}\beta_nw^{n}$, the same recursion reads
\begin{equation}\label{dyn:eq:quadraticrecurrence}
 \beta_1=1,\qquad
 \beta_n=-\frac1{\lambda^{n}-\lambda}\sum_{m=\lceil n/2\rceil}^{n-1}
 \beta_m\binom{m}{n-m}\lambda^{2m-n}\quad(n\ge2),
\end{equation}
the binomial being the coefficient of $w^{n}$ in $(\lambda w+w^{2})^{m}$,
with $\beta_2=-1/(\lambda^{2}-\lambda)$ and
$\beta_3=2\lambda/\bigl((\lambda^{2}-\lambda)(\lambda^{3}-\lambda)\bigr)$;
in the normalization $h(z)=\sum_{n\ge1}b_n\eps_0^{n-1}z^{n}$ of
\cref{dyn:subsec:discrete} it reads
\eqref{dyn:eq:discreterec}. Multiplying $b_nz^{n+1}$ by $t^{n\eta}$ turns
these ordinary formal data into the strongly summable surcomplex linearizer
of \eqref{dyn:eq:scaled-H}, \emph{even when the ordinary series has zero
radius} --- which is the whole point of \cref{dyn:prop:quadratic-scale}.
None of these is a convergence assertion: they are coefficient identities.

\subsection{The iterative logarithm of a quadratic perturbation}
\label{dyn:ex:quadlog}
Take $\Gamma=\Q$, write $t=t^{1}$, and consider $F(z)=z+tz^{2}$ with
entire polynomial coefficients. Directly from \eqref{dyn:eq:explog},
\begin{equation}\label{dyn:eq:quadratic-log}
 X_F=\Bigl(tz^{2}-t^{2}z^{3}+\tfrac32t^{3}z^{4}-\tfrac83t^{4}z^{5}
 +\tfrac{31}6t^{5}z^{6}-\tfrac{157}{15}t^{6}z^{7}+O(t^{7})\Bigr)\partial_z ,
\end{equation}
where $O(t^{7})$ refers only to this single-parameter formal expansion.
Exponentiating half of this vector field produces the unique positive
half-iterate of $F$, with polynomial coefficients at every Hahn weight, and
the ordinary size and growth of the rational coefficients in
\eqref{dyn:eq:quadratic-log} have no bearing on strong summability in $t$.
For a general map, with a bookkeeping parameter $\varrho$ and
$F_\varrho=z+\varrho u$, expanding the logarithm through degree three gives
\begin{equation}\label{dyn:eq:general-third-log}
 X_{F_\varrho}z=\varrho u-\frac{\varrho^{2}}{2}(Du)u
 +\varrho^{3}\Bigl(\frac1{12}D^{2}u(u,u)+\frac13(Du)^{2}u\Bigr)
 +O(\varrho^{4}),
\end{equation}
since $\Delta^{2}z=\varrho^{2}(Du)u+\frac12\varrho^{3}D^{2}u(u,u)+O(\varrho^{4})$
and $\Delta^{3}z=\varrho^{3}(D^{2}u(u,u)+(Du)^{2}u)+O(\varrho^{4})$; with
$u=z^{2}$ this gives the first three terms of
\eqref{dyn:eq:quadratic-log}. \Cref{dyn:eq:quadratic-log} is
\eqref{dyn:eq:euler-known} with $\eps\mapsto t$ and $p=1$; it is the same
calculation, recorded once in each of the two notations its sources use and
identified here (\cref{dyn:warn:hairer}).

\subsection{A complete centralizer and a persistent double fixed point}
For $X=tz^{2}\partial_z$ the flow has the especially simple expression
\begin{equation}\label{dyn:eq:mobius-flow}
 \Phi_c(z)=\frac{z}{1-ctz}=z+\sum_{n\ge1}(ct)^{n}z^{n+1},
 \qquad c=0\text{ or }v(c)>-1,
\end{equation}
the last series being a strongly summable expansion on finite halos. The
time-one map here is $z/(1-tz)$, \emph{not} the quadratic map of
\cref{dyn:ex:quadlog}. By \cref{dyn:thm:centralizer} the full positive
centralizer of this time-one map on the entire-coefficient halo consists
exactly of \eqref{dyn:eq:mobius-flow}, and for every nonzero admissible
time $\Phi_c(z)-z=ctz^{2}/(1-ctz)$, whose factor $ct/(1-ctz)$ is a unit in
the Hahn coefficient algebra. Hence the fixed-point ideal is \emph{exactly}
$(z^{2})$ for every such time: the point at zero retains its
scheme-theoretic double structure, not merely its location --- which is
what \cref{dyn:thm:fixed-ideal} buys over a zero-set statement
(\cref{dyn:rem:tworoutes-fixed}).

\subsection{Two independent perturbation scales, and a resonant splitting}
Let $t$ and $u$ represent $t^{1}$ and $t^{\Omega}$ in
$\Gamma=\Q+\Q\Omega$, and consider the map with multiplier $-1$,
$F(z)=-z+tz^{2}+uz^{3}$, illustrating \cref{dyn:thm:finite-order} with
$q=2$. The averaging coordinate and the resonant vector field begin
\begin{align}\label{dyn:eq:resonant-example}
 H(z)&=z-\frac{t}{2}z^{2}+\frac{t^{2}}{2}z^{3}-\frac{5tu}{4}z^{4}
 +O((t,u)^{3}),\\
 Y(z)&=\Bigl(-uz^{3}-t^{2}z^{3}-\tfrac32u^{2}z^{5}
 -\tfrac{25}4t^{2}uz^{5}-\tfrac72u^{3}z^{7}+O((t,u)^{4})\Bigr)\partial_z,
 \notag
\end{align}
and $HF=(-\Phi_Y)H$. Every displayed monomial of $Y$ has odd coordinate
degree, exactly as $Y(-z)=-Y(z)$ requires.

\begin{warning}[$O((t,u)^{N})$ is a label cutoff, not a valuation cutoff]
\label{dyn:warn:labelcutoff}
In \eqref{dyn:eq:resonant-example}, $O((t,u)^{N})$ means discarded
monomials of \emph{total parameter degree} at least $N$ in an auxiliary
two-parameter jet algebra. It is \textbf{not} a valuation cutoff: for
example $t^{100}$ has much smaller Hahn valuation than $u$, although its
total parameter degree is larger. The same caution applies to
$O_{\mathrm{lab}}(\cdot)$ in \cref{dyn:subsec:multiscale} and to
\emph{every} verification program in \cref{dyn:sec:verification}, all of
which truncate by parameter or label degree. Sources 03 and 06 each state
this independently, and source 06 adds the sharpest form: if
$\theta>n\eta$ for every $n$, infinitely many powers of $u$ lie below the
first $v$ layer, and no finite label truncation pretends to compute that
entire valuation initial segment.
\end{warning}

The verification program of source 03 also uses the two-dimensional
identity-reducing map
$G(x,y)=(x+tx^{2}+uy,\ y+txy+ux)$, checking at total parameter depth three
the logarithm, the inverse, the exponential, the Leibniz identity and the
fixed-point factorization \eqref{dyn:eq:fixedfactor}. That example includes
infinitesimal \emph{linear} terms, which are permitted in the
identity-reducing flow theorem \cref{dyn:thm:exp-log} even though they are
excluded from the exact-multiplier linearization theorem
\cref{dyn:thm:main} --- an instance of the tag distinction of
\cref{dyn:subsec:smalldivisor} showing up in a computation.

\subsection{An arbitrary-support example beyond a finite grid}
\label{dyn:ex:Omega}
Use $S$ and $\Gamma$ from \cref{dyn:ex:ranktwo} and let
$F=\Lambda z+\sum_{\sigma\in S}t^{\sigma}f_\sigma(z)$ with each
$f_\sigma$ vanishing to order at least two. Under a valid analytic
hypothesis from \cref{dyn:thm:main}, the higher-rank coefficients obey
\begin{align}\label{dyn:eq:Omega-example}
 h_\Omega&=-\Li^{-1}f_\Omega,\\
 h_{\Omega+1}&=\Li^{-1}\Bigl(
 D(\Li^{-1}f_\Omega)(\Lambda z)f_1(z)
 +D(\Li^{-1}f_1)(\Lambda z)f_\Omega(z)\Bigr).\notag
\end{align}
Only the word $(\Omega)$ contributes in the first line, and only
$(\Omega,1)$ and $(1,\Omega)$ can contribute in the second. Infinitely many
input coefficients at rational weights below $\Omega$ enter neither
formula. This is a direct use of finite word dependence
(\cref{dyn:prop:linearization}), not a rank-one approximation to the value
group.

\section{Relation to the companion reports in this collection}
\label{dyn:sec:companions}

Each entry names the report by title and repository path. Companion reports
are cited this way and never by their internal labels.

\subsection{Extends: differential equations over the surreals}
\label{dyn:subsec:diffeq}
\emph{Differential Algebra and Differential Equations over the Surreal and
Surcomplex Numbers} (\texttt{docs/surcomplex/differential-equations/})
already contains, in its section on monodromy as an ordinary-domain
construction, the phenomenon that
\cref{dyn:sec:parabolic-monodromy} builds on. Its example takes the scalar
equation $\partial_zY=((\lambda+\eps)/z)Y$, obtains the monodromy
$M=e^{2\pi i\lambda}\operatorname{Exp}(2\pi i\eps)$, and concludes that for
$\lambda\in\Z$ the \emph{ordinary} monodromy is $1$ while the full
monodromy is not, with leading change exactly $2\pi i\eps$: reduction to
the ordinary coefficient can erase a genuine global obstruction.

\textbf{Inheritance.} The two reports stay separate --- fundamental
matrices of linear coherent ODEs are not conjugacies of iterated maps, and
the algebraicity dichotomy \cref{dyn:thm:dichotomy} over
$\C(w)((t^{\Gamma}))$ has no counterpart there. But that example is
\emph{the} antecedent of the monodromy exponent $\rho_a$ of
\eqref{dyn:eq:monodromy}, and is cited as such at the point where $\rho_a$
is introduced. What is new here is not the invisibility phenomenon but the
package around it: the torsion lemma
\cref{dyn:lem:torsion}, the finite-cover obstruction
\cref{dyn:thm:monodromy}, the classification of algebraic leading
coordinates \cref{dyn:prop:leadingalgebraic} with its
already-transcendental first correction, and the dichotomy itself. That
report's cross-category warning is also adopted verbatim, once, as
\cref{dyn:warn:derivative}, in place of re-deriving the
coordinate-versus-surreal-derivation distinction seven times.

\textbf{Source 07 separates itself from that report by name.} Its torus
derivatives $\partial/\partial\theta_j$ are external derivations
annihilating $K_\Gamma$, and it says in as many words that this convention
``keeps the present theory separate from the repository's internal
differential-equation report''. The cohomological equation
\eqref{dyn:eq:qpgreen} is therefore not a surreal differential equation for
the Berarducci--Mantova derivation, and no theorem of either report
transfers to the other. That is the same boundary
\cref{dyn:warn:derivative} draws for the germ half, stated once and
applying to both. Source 07 also records that a targeted repository code
search for ``cohomological'' returned no match, while stating that this is
not an exhaustive audit of every manuscript or every synonym.

\subsection{Adjacent: surcomplex analysis}
\emph{Surcomplex analysis: holomorphic functions on $K=\No[i]$}
(\texttt{docs/surcomplex/analysis/}) is the foundation that fixes the
no-topological-convergence discipline and states the Neumann support
material. All seven sources restate that discipline and all seven state
the support lemma, with sources 02, 03 and 05 each reproducing a
minimal-bad-sequence Higman proof.

\textbf{Inheritance.} Inherited, not restated. Exactly one proof of the
support lemma is kept --- \cref{dyn:lem:neumann} with
\cref{dyn:app:higman}, the cleanest finite-word form --- and the
summability discipline is cited to that report in
\cref{dyn:def:summable}. The other five statements are deleted. This is
also what keeps the present report honest about which parts of its
machinery are classical.

\subsection{Adjacent: infinitesimal analytic geometry}
\emph{Infinitesimal Analytic Geometry over the Surcomplex Numbers}
(\texttt{docs/\allowbreak surcomplex/\allowbreak analytic-\allowbreak
geometry/}) owns the four-ring table and its separating witness. Sources
02, 03 and 04 each re-erect that
separation from scratch, with source 02's $\sum_{m\ge1}t^{m}/(1-mz)$ and
source 04's $\sum_{n\ge1}t^{n\delta}/(1-nw)$ being that report's witness
with $\eta=1$ and $\eta=\delta$.

\textbf{Inheritance.} Neither the ring table nor the witness is
reprinted. Every theorem here states its ring in that report's names
(\cref{dyn:conv:rings}), and exactly two rows are added ---
entire coefficients and polynomial coefficients --- with the explicit note
that they are new \emph{rows}, not new rings
(\cref{dyn:subsec:rings}). Several sources already say they are continuing
that distinction; this merge makes the continuation literal.

\subsection{Adjacent: contours and Stokes}
\emph{Contours and Stokes: Jordan separation, Cauchy theory, and Stokes
calculus on the surcomplex plane}
(\texttt{docs/surcomplex/contours-and-stokes/}) is background for
\cref{dyn:sec:periods} and is deliberately \emph{not} used there: every
integral in that section is an ordinary complex line integral taken
coefficient by coefficient over ordinary curves in a fixed $D\subset\C$.
\Cref{dyn:warn:periods} states the resulting prohibition, which source 05
carries as a remark of its own.

\subsection{Keep separate, with a cross-note: rank-one Berkovich geometry}
\label{dyn:subsec:berkovich}
\emph{Rank-One Non-Archimedean Analytic Geometry inside the Surcomplex
Numbers}
(\texttt{docs/\allowbreak surcomplex/\allowbreak rank-one-berkovich/}) is
the one report in the collection where convergence is genuine convergence:
it works inside
a \emph{fixed rank-one} $\C((t^{\R}))$ with $|a|=\exp(-v(a))$ and spherical
completeness, where the absolute value is a real-valued one and the field
is complete and algebraically closed.

\textbf{Cross-note, and it is load-bearing.} This report forbids that
reading throughout. It works at arbitrary valuation rank, with a
deliberately rank-two example (\cref{dyn:ex:rank2}), and cites Lindahl's
rank-one non-Archimedean linearization as an antecedent it does not
reprove. A Lindahl-type ultrametric disk radius and source 01's maximal
centered valuation ball are \emph{different categories}
(\cref{dyn:warn:ball}), and the exact-ball and shell-periodicity
conclusions of \cref{dyn:thm:exact,dyn:thm:shell} must \textbf{not} be
transported to $\C_p$: source 01 says so itself and cites $p$-adic
quadratic-like examples whose bounding sphere carries no periodic point
(\cref{dyn:warn:charzero}). Without this note a reader will read the two
reports' ``disks'' as the same object. No result of this report and no
result of that one is transferred in either direction.

\section{Computational record and reproducibility}
\label{dyn:sec:verification}

The seven verification programs are preserved unmodified in \texttt{code/}
and their recorded outputs in \texttt{data/}. They supply \emph{finite
exact symbolic} checks. None of them proves an infinite statement in
this report. Separately, the repository's Lean modules
\texttt{Surreal/HahnSeries/Neumann.lean} and
\texttt{Surreal/HahnSeries/Ancestry.lean} prove the finite-word support
lemma, the corrected input-relative combinatorial ancestry bounds, and the
counterexample to an input-independent bound. This does not yet formalize
the corollary's full operator-summability assertion or the analytic
theorems built on it. The exact coverage is recorded in
\texttt{docs/FORMALIZATION.md}.

\subsection{What each suite checks, and what it does not}

\begin{longtable}{@{}P{0.12\textwidth}P{0.10\textwidth}P{0.42\textwidth}P{0.29\textwidth}@{}}
\toprule
Source & Checks & Content & Explicit disclaimer\\
\midrule
\endhead
01 & 54 &
Nine coefficients of the general cubic second iterate against a closed
form; the tuned factorization \eqref{dyn:eq:tuned-factor} and the
degree-$3$/degree-$4$ cancellation at $e=0$; the three phase regimes with
\emph{integer} exponent representatives $\delta=8$, $\sigma_0=2,4,6$
(radii $3,2,2$), their thresholds, shell polynomials and squarefreeness;
multiplier ratios by polynomial remainder; the three coherent leading
coefficients through their \emph{logarithmic derivatives} (deliberately
avoiding a numeric branch choice of fractional powers); Schr\"oder
recurrence to degree $9$ for $\lambda=3/2,-4/3,5/4,1+i$; the tangent
binomial identity \eqref{dyn:eq:binomial-mult} for $m=2..7$;
$\Disc_Y(1+\kappa Y-2Y^{2})=\kappa^{2}+8$; integrality and residues of
$h_2..h_9$ for the normalized tuned cubic at $e=t^{4}$, $r=1$, including
$[X^{5}]H_0=1/2$ and $[X^{9}]H_0=5/8$. All exact/symbolic, no floating
point.
& Does not establish the arbitrary-rank support lemma, algebraic closure,
maximality for every polynomial, the transfinite coherent lifting, or
novelty; not an implementation of general Hahn fields; not proof-assistant
verification.\\[4pt]
02 & 1{,}109 &
Reciprocal-jet formula \eqref{dyn:eq:reciprocaljet} versus the inverse
recurrence ($133$) and truncated product of a denominator with its
reciprocal ($133$), both for $2\le n\le20$, $0\le k\le6$; ordinary
quadratic Schr\"oder identity through degree $12$ ($13$); Hahn-scaled
quadratic identity ($156$); nonlinear recursion with multiplier drift and
cubic forcing, $F(z)=(\lambda+t)z+t(z^{2}+z^{3})$, checked after each
parameter step through $t^{4}$ and spatial degree ten ($154$);
two-variable multiplicative-detuning reciprocal jets with exact unit
phases $(3+4i)/5$ and $(5+12i)/13$, total degrees $2..9$, both components,
jet orders $0..4$ ($520$). Exact in $\Q(i)$ via
\texttt{fractions.Fraction}; test multiplier $\lambda=3/5+4i/5$, not a root
of unity because a root of unity in $\Q(i)$ is an algebraic integer.
& Tests algebraic identities, signs, normalizations and truncation only.
Does not prove Hahn summability, well-founded recursion, domain
preservation, the small-divisor limsups, the exact radii, or novelty. The
multivariable test verifies nonzero denominators only in its finite range
and does not certify nonresonance for all multi-indices.\\[4pt]
03 & 21 &
$\exp(\log F)=F$ and two-sided inversion; half and third iterates; time
addition; Leibniz; identification of the logarithm with its coordinate
vector field; the two-variable fixed-point matrix factorization
\eqref{dyn:eq:fixedfactor}; normalized nonresonant conjugacy, its inverse
and the degree bound \eqref{dyn:eq:degree-bound}; the finite-order
averaging and resonant flow identities. One-parameter flow checks through
degree six; two-parameter multivariable, nonresonant and resonant checks
through total parameter degree three. Sparse exact jet algebra over
$\Q(i)$ with a symbolic $\lambda$ for the four coefficients
\eqref{dyn:eq:quadratic-bs}.
& Finite polynomial jet identities truncated by \emph{total parameter
degree}, explicitly not a Hahn valuation cutoff in a higher-rank group
(\cref{dyn:warn:labelcutoff}). Does not test arbitrary well-ordered
supports, analytic radius claims, continued-fraction asymptotics, the
universal quantifiers, or priority. No Lean verification.\\[4pt]
04 & 2{,}764 &
Koenigs recurrence \eqref{dyn:eq:brecurrence} and the leading Pochhammer
coefficients for $q=1..8$, $m\le24$, through $\eps^{8}$; the \emph{directly
substituted} Schr\"oder residual through $\eps^{9}$ (an independent test,
not merely agreement with a displayed formula); the first and second
coefficient functions through $w^{24q+1}$; the logarithm ratio $\rho_q$ of
\eqref{dyn:eq:rhoq} at orders $0,1,2$ for $q=1..8$; the derivation identity
$\log C_g(p)=(\log C_g(w))p'$ for $q=1..6$, $p=w^{n}$, $0\le n\le6$,
through $\eps^{6}$; generator coefficients and fixed-point derivatives,
reduced using $w^{q}=-1$ to avoid approximate complex roots. Standard
library only, \texttt{fractions.Fraction}, no floating point.
& No finite test proves Neumann summability, existence on an arbitrary
ordinary domain, compatibility for arbitrary value groups, or the
algebraicity dichotomy. Not a symbolic surcomplex CAS and not a numerical
evaluator at arbitrary surreal inputs.\\[4pt]
05 & 245 &
Euler families $p=1,2,3$: $\exp$ of $\log$, the Julia identity
\cref{dyn:cor:julia} in the form $W_F(F)=F'W_F$, inverse via negative
time, the homogeneous grading coefficientwise, residues equal $(p+1)/2$ at
degree one and zero elsewhere; a mixed-coefficient map (round trip,
log-equals-vector-derivation, Leibniz, Bernoulli difference solver); the
flow group law (times $2$ and $3$ compose to $5$); the unique-root
formula; an explicit conjugacy with pushforward of the vector field,
pullback of the normalized time form and the basepoint marking
$h_k(1)=0$; untruncated symbolic identities for the flat family
($1/W_c=1/\eps+c/z$, residue $c$, leading correction $-1/z$); and $50$
lexicographic order comparisons. Exact in $\Q(z)[\eps]/(\eps^{8})$.
& Transcription and algebra checking only. Proves nothing about
arbitrary-rank support finiteness, global primitive existence, the
classification, or the slow-time theorem. \textbf{Fifty of the $245$
assertions are finite lexicographic illustrations for $n=1,\dots,50$; the
all-$n$ inequality is proved in the text, not inferred from them.} The
verifier also loses one top coefficient on division by $\eps$, so
reciprocal and pullback comparisons are checked only through degree
$N-1$.\\[4pt]
06 & 119 &
Exact arithmetic in $\Q(\sqrt2)$, with a hand-rolled bivariate label
series ring truncated at total label degree four: the single-exponential
gauge recursion for \eqref{dyn:eq:workedH}; resonance of every normal
coefficient and the nonresonant gauge of every generator coefficient; the
four displayed coefficients; all $16$ symplectic coordinate brackets; four
inverse Lie maps; direct substitution $H(\Phi)=N$ and $J(\Phi)=I$; the two
first integrals, their involution, and one nonlinear function of the
actions being an integral; $51$ degree-defect identities; the cubic
boundary map \eqref{dyn:eq:cubicmap}; the support counterexample's
finite-sample order pattern; and \eqref{dyn:eq:discreterec} for
$\lambda=2$ through $n=10$.
& Truncates by \emph{total label degree} in two independent labels, not by
Hahn valuation, and only to degree four. Explicitly not a proof of
arbitrary-support summability, nor of
\cref{dyn:thm:homological,dyn:thm:universal}, nor of
\cref{dyn:thm:pcentralizer}, nor of the Liouville construction
\cref{dyn:prop:explicitbad}.\\[4pt]
07 & --- &
No printed total; two named checks, every assertion exact. The slow-circle
conjugacy \cref{dyn:ex:qpcircle} is reconstructed recursively from
$(1+h')\nu=1+\varrho\sin(\theta+h)$ with $\nu(0)=1$: at each order the mean
is determined first, giving the new coefficient of $\nu$, then the
remaining zero-mean Fourier polynomial is integrated. It asserts at every
order that the mean is real and that $h$ is Hermitian, that the full
conjugacy residual vanishes through degree ten, that $\nu^{2}=1-\varrho^{2}$
coefficientwise through the same degree, and that the three displayed terms
of \eqref{dyn:eq:qpcircleh} are correct; the recorded frequency
coefficients are $1,0,-\tfrac12,0,-\tfrac18,0,-\tfrac1{16},0,
-\tfrac5{128},0,-\tfrac7{256}$. It then checks $117$ instances
($n\in0..8$, $m\in0..12$) of the coefficient recurrence for
\eqref{dyn:eq:qpranktwoinverse}, and classifies all $1{,}330$ nonzero
integer modes in $[-5,5]^{3}$ into the three strata of
\cref{dyn:ex:qpthree}, obtaining $1{,}210/110/10$. Exact $\Q(i)$
arithmetic by pairs of \texttt{fractions.Fraction}; every check is an
\texttt{assert}; no floating point and no tolerance anywhere.
& Finite algebraic checks only. They verify neither the irrationality of
\eqref{dyn:eq:qpliouvilleconstant}, nor any infinite-mode subexponential
estimate, nor the word lemma, nor the infinite-support conjugacy theorem;
those rest on the proofs in the text. No Lean or Wolfram verification of
any infinite theorem is claimed.\\
\bottomrule
\end{longtable}

\subsection{Independently verified facts about the suites}
\label{dyn:subsec:orchestrator}

The following were checked directly against the delivered files and by
running the code, and override any contrary statement in a source README.

\begin{enumerate}[label=(\arabic*)]
\item \textbf{All seven suites pass.} Run on Python 3.14.4 with SymPy 1.14.0
where SymPy is used, every one exits $0$. Reported totals where a script
prints one: $54$ (01), $1{,}109$ (02), $21$ (03), $2{,}764$ (04), $245$
(05), $119$ (06); source 07 prints no total. Source 04's count reproduced
the recorded $2{,}764$ exactly; source 05 reproduced its recorded
transcript line for line; source 03 reproduced its four recorded example
expansions byte for byte; and the regenerated JSON of sources 06 and 07
matched the shipped ones up to line endings only.
\item \textbf{Four of the seven scripts rewrite their own evidence.} Source
03's \texttt{verify.py} writes \texttt{data/verification.json} by default;
source 04's writes \texttt{data/verification.json} \emph{and}
\texttt{data/coefficients.csv}; source 06's writes
\texttt{data/verification.json}, and does so only after all checks pass;
source 07's takes \texttt{-{}-output} with \emph{default}
\texttt{verification.json} resolved against the current working directory
and writes it with an unconditional \texttt{write\_text} --- no existence
check, no backup, no dry run --- and the invocation documented in its own
README is exactly the one that overwrites the delivered record.
Source 01's script prints to stdout, but its \texttt{Makefile}'s
\texttt{verify} target redirects that stdout into
\texttt{verification.txt} in the same directory, overwriting the delivered
record. \textbf{Every suite must therefore be run on a copy, never in the
delivered tree}: a byte-identity check performed after an in-place run
compares two equally modified copies and passes while proving nothing.
\item \textbf{Every suite prints its own scope disclaimer.} That is a real
and unusual discipline in this material, and those disclaimers are exactly
the limitations carried in \cref{dyn:sec:nonclaims}.
\item Source 04's README lists \texttt{SHA256SUMS.txt} among its package
contents, and no such file exists in the delivered archive. This is a
packaging gap, recorded rather than repaired.
\item Source 03's recorded console transcript names its report path as
\texttt{/mnt/\allowbreak data/\allowbreak surcomplex\_\allowbreak
dynamics/\allowbreak data/\allowbreak verification.\allowbreak json},
i.e.\ it
was recorded in a different container from the delivered one. The check
contents match the article's prose regardless.
\item Source 07's README likewise lists \texttt{SHA256SUMS.txt} among its
package contents, and no such file exists in the delivered archive. Like
item~(4) this is recorded rather than repaired.
\item \textbf{The finiteness bound of \cref{dyn:thm:flag} was checked
independently of source 07.} Modelling
$\Upsilon_i=\sum_jc_{ij}t^{\gamma_j}$ with $\gamma_1<\gamma_2<\cdots$ gives
$k\cdot\Upsilon=\sum_jL_j(k)t^{\gamma_j}$ with
$L_j(k)=\sum_ik_ic_{ij}$, so $v(k\cdot\Upsilon)$ is the least $\gamma_j$
with $L_j(k)\ne0$ and
$\{k:\text{level}\ge\gamma_{m+1}\}=\ker L_1\cap\cdots\cap\ker L_m$, whose
saturation has rank $d-\rank(L_1,\dots,L_m)$. A new level appears exactly
when that rank increases, so the kernel rank strictly drops, and it can
drop at most $d$ times: that is both the bound and the flag. Brute-force
testing over rational coefficient matrices with deliberate resonances
forced in $40$ per cent of the cases gave \textbf{zero violations of the
bound in $400$ random trials}, and a worked instance at $d=3$ had kernel
ranks $2\to1\to0$ with achieved levels $\{0,1,3\}$ --- three levels,
attaining the bound, with two indices skipped exactly where the rank did
not move (\cref{dyn:rem:qpvisible}). The theorem and its stated mechanism
both hold. This is an independent check of a finite mechanism, not a proof
of any infinite statement.
\end{enumerate}

\subsection{Building this report}
The report is standalone LaTeX with an internal bibliography and no
external \texttt{.bib} file, graphics, shell escape or network access:
\begin{verbatim}
latexmk -pdf -interaction=nonstopmode article.tex
\end{verbatim}
Reproducing the seven suites, each \emph{on a copy}, requires Python 3.10 or
later and, for four of them, SymPy; the pinned versions recorded by the
sources are \texttt{sympy==1.14.0}.

\section{Open questions}
\label{dyn:sec:questions}

The questions below are exactly those the sources leave open, with the
categories in which they are open stated explicitly.

\begin{question}[Common-domain germs at a finite positive rate, for a fixed
multiplier]\label{dyn:q:commongerm-restate}
This is \cref{dyn:q:commongerm}, and \cref{dyn:warn:resolution} is its
status: settled negatively in the \emph{drifted} category $\sddr$ by
\cref{dyn:thm:collapse,dyn:thm:trichotomy}, and \textbf{open} in the
exact-multiplier category $\sdex$. An exact characterization would need to
control interactions across unbounded word depths rather than a single
homological inverse.
\end{question}

\begin{question}[The fixed cyclic value group]\label{dyn:q:cyclic}
The negative direction of \cref{dyn:thm:collapse} is established over
$\Gamma_{*}=\Q+\sqrt2\,\Q$, hence for the universal class of workspaces,
and it embeds in any surcomplex workspace containing that ordered subgroup.
It is \emph{not} asserted for a fixed cyclic value group such as $\Z$,
where different parameter blocks can collide at one Hahn exponent.
Resolving the effect of those collisions is a separate problem, not a
consequence of the two-independent-exponent proof, and it concerns the
existence of \emph{some} smaller common neighborhood for every input, not
the weaker statement about retaining the original disk.
\end{question}

\begin{question}[Resonant multipliers of infinite order]\label{dyn:q:resonant}
For a diagonal unitary multiplier with resonances but \emph{without} finite
order, find a support-controlled analytic normal form with precise
necessary and sufficient coefficient-category hypotheses. Nonresonant
monomials can be treated by the homological inverse, but resonant terms
create a coupled normal-form problem and remove the global inverse used in
\eqref{dyn:eq:linearizer-series}, and a finite cyclic averaging operator is
no longer available. The Hamiltonian analogue is open too: for resonant
$\omega$ the first reduction retains a nontrivial resonant Hamiltonian
algebra, the gauge equation changes because $L^{-1}$ has a larger kernel,
and whether the dynamics of that retained algebra can itself be integrated
under effective support conditions is not answered by the nonresonant
centralizer proof.
\end{question}

\begin{question}[Higher-dimensional centralizers]\label{dyn:q:centralizers}
In several variables, classification of the positive centralizer reduces
exactly to classification of the commuting positive vector fields, and the
one-variable meromorphic-ratio argument of \cref{dyn:thm:centralizer}
cannot classify those higher-dimensional Lie centralizers; additional first
integrals or commuting geometric structures can occur. Independently,
source 05 records that in several variables there is no single reciprocal
time one-form encoding a vector field and commuting fields need not be
constant multiples, so the period normal form does not transport either.
\end{question}

\begin{question}[Moving fixed points, and several variables, for the
parabolic theory]\label{dyn:q:moving}
Characterize the monodromy of $w+\eps V(w)+\eps^{2}R(w)$ when its fixed
points move. An analogue of the local derivative identity
\cref{dyn:lem:fixedderivative} at moving formal fixed points is a natural
intermediate target; translating it into a common-domain continuation
theorem would also require controlling those moving singularities. The
Euler family \emph{avoids} that difficulty rather than concealing it. In
several variables the logarithm-of-substitution construction still suggests
a vector field, but the scalar quotient $\alpha_{\mathrm c}/W$ has
\emph{no} direct replacement, and resonant homological equations with
noncommuting monodromy would replace the one-dimensional integration
argument.
\end{question}

\begin{question}[The period classification where the leading coefficient
vanishes]\label{dyn:q:zeros}
When $a$ has zeros, \cref{dyn:thm:difference} still controls the discrete
equation, but the conjugacy classification needs more: local zero divisors,
multiplicities and singular normal-form obstructions must be combined with
global periods. \Cref{dyn:warn:nowherezero} identifies the precise point
where the proof stops. Separately, extending
\cref{dyn:thm:classification} to conjugators with nonidentity ordinary part
requires the ordinary automorphism action on the leading form and on its
periods, after which equal period vectors in a fixed basis would no longer
be the right equivalence relation by themselves.
\end{question}

\begin{question}[Maximal domains at infinite scale]\label{dyn:q:maximal}
\Cref{dyn:thm:rescale} gives \emph{guaranteed}, not maximal, domains.
\Cref{dyn:ex:cubic} detects a sharp boundary through both a pole and an
extra critical point; a general procedure combining support certificates
with algebraic or analytic obstructions could yield substantially better
maximal-domain information for finitely presented perturbations. Relatedly,
for entire layers failing the universal threshold, particular inputs may
still normalize because their coefficients vanish at the dangerous
monomials; an input-specific necessary and sufficient condition would have
to remain stable under the nonlinear brackets created by normalization, and
a first-order divisibility test alone need not be enough.
\end{question}

\begin{question}[Ordinary analytic realization]\label{dyn:q:realization}
Which of the Hahn normal forms constructed here admit compatible
\emph{ordinary} analytic realizations? Ordinary analytic specialization of
a Hahn or formal parameter requires growth estimates or an additional
summation theory, and this report supplies neither. Sectorial analytic
moduli of ordinary parabolic germs are not removed, and the analytic
embedding obstruction for ordinary parabolic germs is not resolved.
Likewise, adding weighted coefficient norms, Gevrey bounds in the
parameter, or ordinary convergence in $(u,v)$ produces strictly stronger
problems whose arithmetic thresholds could differ from the thresholds
established here; \eqref{dyn:eq:compactestimate} is offered only as a
starting point.
\end{question}

\begin{question}[The resonant quasi-periodic normal form]
\label{dyn:q:qpresonant}
\Cref{dyn:thm:qplinear} allows $\qres$, but the constant normal form
\cref{dyn:thm:qpnormalform} requires $\qlat=\{0\}$. In the resonant case a
natural target is a normal form retaining the Fourier modes in $\qlat$.
Source 07 warns explicitly that simply replacing the mean projection $\Qz$
by $\Pi_{\qlat}$ in the proof is \textbf{not} justified: the retained
normal form is then no longer a constant vector field and creates
additional terms in the conjugacy equation. Products and substitutions
interact with those modes, and higher perturbations may split the remaining
resonances, so the extended flag of \cref{dyn:cor:qptail} would have to be
carried through the nonlinear argument. This is the quasi-periodic
counterpart of \cref{dyn:q:resonant} and is open for the same structural
reason --- a larger kernel for the inverse operator --- but it is a
different problem on a different divisor family.
\end{question}

\begin{question}[Quantitative coefficient-growth classes on the torus]
\label{dyn:q:qpgrowth}
One may add bounds linking $\|f_\beta\|_{r}$ to a weight on the support.
The argument of \cref{dyn:sec:qpnonlinear} does \emph{not} prove
preservation of such a weighted class, because inverse denominator powers
and derivative losses grow along support words
(\cref{dyn:rem:qpdiophantine}). Identifying hypotheses that preserve these
weights, and that would permit a meaningful specialization of selected
monomials to small real parameters, is open; that is exactly where genuine
non-linear convergence arithmetic would have to re-enter, and it is the
torus form of \cref{dyn:q:realization}.
\end{question}

\begin{question}[Effective arithmetic certificates for the flag]
\label{dyn:q:qpeffective}
\Cref{dyn:thm:flag} makes the flag finite but not automatically computable
from arbitrary real coefficients: detecting an exact integer relation among
arbitrary named real numbers may itself be unavailable. For rational and
suitable algebraic coefficient data, explicit lattice and approximation
certificates can make the construction effective. The theorems as stated
separate \emph{existence} of a finite flag from \emph{decidability} of its
entries and its arithmetic bounds, and no complexity claim is made. This is
the same distinction as \textbf{N64} and \cref{dyn:app:certificate} draw
for the finite-return certificate.
\end{question}

\section{Limitations and non-claims}
\label{dyn:sec:nonclaims}

This section is the consolidated record required by
\cref{dyn:conv:merge}(iv). It carries \textbf{seventy-eight} distinct
limitations stated by the seven sources --- eleven from source 01, nine from
source 02, ten from source 03, fourteen from source 04, twelve from source
05, twelve from source 06 and ten from source 07 --- with the originating
source named in each case. None has been merged away, softened, or absorbed
into a stronger statement. Where two sources state overlapping limitations,
both wordings are preserved because their scopes differ.

\subsection{Status, refereeing, and formalization}
The following items record the status of the preserved source packages.
For the later, partial repository formalization of the support results,
see \cref{dyn:sec:verification}.
\begin{enumerate}[label=\textbf{N\arabic*.},leftmargin=3.2em]
\item (01) An AI-assisted, unrefereed research draft; not refereed and not
formally verified. Its PDF metadata records the author as ``Research
manuscript prepared with ChatGPT''. Priority is explicitly not certified
and no named published conjecture is claimed solved.
\item (02) An AI-assisted, unrefereed research draft; not refereed, not
proof-assistant verified. \textbf{No Lean theorem corresponding to the new
analytic results was compiled} as part of the package.
\item (03) Self-described AI-assisted, unrefereed, not formally verified.
Its literature search is described as targeted and auditable but
\emph{not} a priority determination; no named published conjecture is
claimed solved; independent expert review of both the proofs and the
literature comparison is called appropriate.
\item (04) An unrefereed, AI-generated research draft whose proofs are
``offered for scrutiny'', not proofs checked by a proof assistant. No Lean,
Wolfram or other proof-assistant formalization is included or claimed.
\item (05) Not independently refereed and not formally verified; none of
the central claims is an independently checked Lean theorem.
\item (06) Not peer-reviewed and not proof-assistant verified; no Lean or
other formalization is supplied or claimed.
\end{enumerate}

\subsection{Priority and novelty}
\begin{enumerate}[label=\textbf{N\arabic*.},leftmargin=3.2em,start=7]
\item (01) The repository and literature comparisons are both explicitly
\emph{targeted}: three catalogue README files plus a search for
``linearization'', not a line-by-line audit of every research package, and
the novelty assessment is a ``proposed-originality assessment, not an
exhaustive priority claim''.
\item (02) Priority is not independently verified. Only three theorems ---
here \cref{dyn:thm:collapse,dyn:thm:trichotomy,dyn:thm:multicollapse} ---
are named as candidate-new, and the source explicitly concedes that ``a
specialist could identify an existing formal-family theorem that subsumes
part of the present formulation; that possibility should be checked before
any publication claim''. The support lemma, the homological inverse, formal
nonresonant conjugacy, the quadratic scaling trick and the centralizer
argument have classical antecedents and are not offered as independent
discoveries.
\item (02) The repository inspection is explicitly \emph{not} a
line-by-line audit of every manuscript or every Lean file; only the
catalogue and the opening foundation definitions were read.
\item (03) Classical credit is stated plainly and is not claimed as new:
formal logarithms, flow embedding and semisimple--unipotent splitting
\cite{Ribon}; the small-divisor problem and formal-versus-analytic
linearization \cite{CarlettiMarmi,Yoccoz,BuffCheritat}; and the support
lemmas \cite{Neumann,Higman,Poonen}.
\item (04) Priority is \emph{not} claimed: the literature search is
described as targeted and bounded and ``did not certify priority''; no
named published open problem is claimed solved; novelty is ``to the
author's present knowledge''.
\item (04) Acknowledged precedents, not contributions: Neumann's support
lemma, formal embedding and formal logarithms, and rank-one ultrametric
radius theory with the quadratic disk scale \cite{Lindahl}. In that
source's own words, ``we do not present a disk estimate alone as a new
theorem of classical dynamics''.
\item (04) The repository inspection was ``a targeted inspection, not an
exhaustive line-by-line audit'' at one frozen revision.
\item (05) Novelty is assessed by a \emph{targeted}, not exhaustive,
literature search; the statements are called candidate new results; an
unlocated result may exist in the literature on formal flows, deformation
of one-forms, or generalized series; and no named published conjecture is
claimed settled.
\item (06) Priority is not established: the contributions are labelled
``proposed contributions with proofs'', not certified bibliographic firsts,
on the basis of a non-exhaustive literature search; no named published open
conjecture is claimed solved. Classical formal Birkhoff normalization,
non-Archimedean linearization and general formal-sum Lie exponentiation are
credited rather than claimed
\cite{PerezMarco,Krikorian,Lindahl,BKKPS}.
\item (04) For the Neumann 1949 reference, the publisher's direct PDF
endpoint was not accessible in the session; the metadata and the standard
lemma were cross-checked instead.
\item (04) Packaging gap: the source README lists \texttt{SHA256SUMS.txt}
among the archive contents, and no such file exists in the delivered
archive (see \cref{dyn:subsec:orchestrator}(4)).
\end{enumerate}

\subsection{Scope of the hypotheses}
\begin{enumerate}[label=\textbf{N\arabic*.},leftmargin=3.2em,start=18]
\item (01) \Cref{dyn:thm:exact} and its package cover nonlinear
polynomials in \emph{one} variable, with unit multiplier not a root of
unity, over complex Hahn fields of arbitrary ordered-group rank, in equal
characteristic zero. Explicitly \emph{not} covered: general infinite
Hahn-analytic power series --- where a common support certificate, a finite
or attained dominant face, and analytic composition would each need
separate proof; actual root-of-unity multipliers; attracting or repelling
multipliers; and higher-dimensional simultaneous resonances.
\item (01) Equal characteristic zero is a \emph{genuine} hypothesis, not
cosmetic: the divisor cancellation \eqref{dyn:eq:denomcancel} leaves a
factor $n-1$ in the denominator, a valuation unit here but not in mixed
characteristic. The formula and the shell-periodicity conclusion must
\textbf{not} be transported to $\C_p$; \cite[Cor.~C]{LindahlPadic} gives
$p$-adic quadratic-like examples whose linearization bounding sphere
contains no periodic point (\cref{dyn:warn:charzero}).
\item (02) Resonant multipliers are excluded \emph{entirely} from
\cref{dyn:sec:drift}; resonant normal forms need a different target
equation, not the linear normal form used there.
\item (02) The several-variable results are stated only for diagonal
unitary $L$ and \emph{equal-radius} polydisks; no optimal general
Reinhardt-domain claim is made or needed.
\item (03) The finite-order theorem \cref{dyn:thm:finite-order} requires
$\Lambda^{q}=I$, not merely the presence of resonances, and resonant normal
forms are \emph{not} claimed unique: further $\Lambda$-commuting changes
can transform $Y$.
\item (04) The most consequential hypotheses of
\cref{dyn:sec:parabolic-monodromy} are not cosmetic: the simple zero at the
normalization point permits the regular quotient
\eqref{dyn:eq:regularR}; a common ordinary domain permits coefficientwise
primitives without an uncontrolled sequence of shrinking radii; positive
support makes the operator logarithm strongly meaningful; the polynomial
Euler form fixes all zeros and gives the exact factor
$W_\eps=VB_\eps$, whereas a general perturbation can move zeros or create
different coefficient pole patterns; and divisibility of $\Gamma$ is used
to keep the resonant rescaling inside the chosen workspace.
\item (05) \Cref{dyn:thm:classification,dyn:thm:moduli} require the
leading coefficient $a$ to be \emph{nowhere zero} on $D$. The unit
hypothesis in \cref{dyn:lem:displacement2} is exactly where the proof stops
applying, and the zeros case needs local zero divisors, multiplicities and
singular normal-form obstructions combined with global periods.
\item (05) Conjugators are restricted to identity ordinary part. A
conjugacy with nonidentity ordinary part can change the leading vector
field and the homology basis, so equal period vectors in a fixed basis
would no longer be the right equivalence relation; that extension is not
done.
\item (05) No converse characterization of the units of $\AAA$ is claimed,
only the sufficient criterion \cref{dyn:lem:unit}.
\item (05) The $\AAA$-membership condition $g/W\in\AAA$ in
\cref{dyn:thm:difference} is indispensable outside the nonsingular stratum;
period vanishing alone is insufficient there
(\cref{dyn:ex:localobstruction}).
\item (06) Nonresonant frequencies only. The gauge equation changes for
resonant $\omega$, because the homological inverse then has a larger
kernel; the resonant case is listed as open
(\cref{dyn:q:resonant}).
\item (06) No result is claimed for Hahn series with unrelated
convergent-germ coefficients, nor for infinite-dimensional Hamiltonian
systems.
\item (06) Uniqueness is \emph{gauge-relative}: only the single
exponential generator with $\Pi\chi=0$ and the normal form it produces are
unique; general normalizing symplectic changes are not
(\cref{dyn:warn:gauge}).
\end{enumerate}

\subsection{Exactly what ``exact'' and ``analytic'' mean}
\begin{enumerate}[label=\textbf{N\arabic*.},leftmargin=3.2em,start=31]
\item (01) ``Exact'' in \cref{dyn:thm:exact} means the maximal
\emph{centered valuation ball} only --- not the full pointwise summability
locus of the series (exceptional evaluations on the bounding shell are
allowed), not a maximal non-ball invariant domain, and not every analytic
continuation of the conjugacy.
\item (01) The bounding shell $S_r$ is an algebraic valuation-shell
expression: \emph{not} a topological boundary (valuation balls are clopen)
and \emph{not} an ordered-modulus circle. The common-domain disk $D$ is
\emph{not} identified with $B_r$: all points of the normalized monad have
ordinary residue zero while $D$ contains many nonzero ordinary points
(\cref{dyn:warn:ball}).
\item (01) \Cref{dyn:thm:coherent} gives \emph{no} uniform ordinary norm
bound on the coefficient family $(H_\gamma)$; these may be arbitrarily
large with arbitrarily large derivatives. Only the domain of holomorphy and
the support certificate are uniform (\cref{dyn:warn:nouniform}).
\item (01) The coherent conclusion is \emph{resonant-only}, and
\cref{dyn:prop:noncoherent} shows failure of a \emph{universal}
common-domain statement in the nonresonant case --- explicitly \textbf{not}
an if-and-only-if classification, and not every nonresonant case is
noncoherent.
\item (01) The stability certificate \cref{dyn:thm:stability} is
sufficient, \emph{not} necessary: a perturbation may change the first
return face and still preserve its maximal slope.
\item (02) The positive theorem gives \emph{no} uniform bound on
$\|h_\gamma\|_r$ as $\gamma$ varies, no norm, no Gevrey or weighted
coefficient estimates, and no ordinary holomorphic dependence on $t$ near
$t=0$; those are strictly stronger problems whose arithmetic thresholds
could differ, and \eqref{dyn:eq:compactestimate} is offered only as a
starting point.
\item (03) ``Finite dependence'' bounds the number of input
\emph{coefficient functions} and analytic operators at a specified weight.
It is \emph{not} an unconditional finite-time algorithm: one application of
$\Li^{-1}$ to an arbitrary analytic function can involve infinitely many of
its ordinary Taylor coefficients, and only polynomial input makes each
requested coefficient a finite exact algebraic calculation
(\cref{dyn:warn:finitedep}).
\item (04) Transcendence is relative to the specific base
$\C(w)((t^{\Gamma}))$. The source volunteers the counterweight that
$H_\eps$ is \emph{differentially algebraic} in a very explicit way,
satisfying $H_\eps'=\Omega_\eps H_\eps$ with $\Omega_\eps\in\BB$, so
``transcendental'' must not be upgraded to ``differentially
transcendental'' (\cref{dyn:warn:diffalg}).
\item (04) ``Transcendental'' concerns a \emph{function} in an ambient
meromorphic Hahn field. It is not a claim that any individual value
$H_\eps(w_0)$ lies outside $\No[i]$, and is explicitly consistent with
$\No[i]$ being algebraically closed.
\item (04) \Cref{dyn:thm:common} does \emph{not} assert that the global
evaluated map $H^{\#}$ is injective on all of $U^{\#}_\Gamma$; injectivity
is proved only on the infinitesimal disk for the resonant model
(\cref{dyn:warn:embedding}).
\item (04) The sharp-domain maximality of \cref{dyn:thm:sharp} is about
\emph{centered valuative balls} under either boundary convention. It does
not exclude continuation to a non-ball domain avoiding the obstructing
fixed points, and does not assert that every coefficientwise continuation
is injective (\cref{dyn:warn:sharp}).
\item (05) An Abel function is not automatically a globally injective
coordinate: a locally univalent primitive need not be globally univalent
(\cref{dyn:warn:abel}).
\item (06) ``Analytic'' is a \emph{defined} term meaning the
entire-coefficient Taylor-compatible realization of
\cref{dyn:prop:eval}. The source does not assert fine-topological
continuity of nonconstant ordinary-time paths in $\No[i]$, and does not
identify phase differentiation with an intrinsic derivation of the surreal
coefficient field (\cref{dyn:warn:derivative,dyn:warn:hamflow}).
\item (06) \Cref{dyn:thm:rescale} extends the polynomial-coefficient
normalizer, Hamiltonian and constructed actions \emph{only}. It does
\textbf{not} license evaluating an arbitrary entire function at an infinite
phase or action argument; a centralizer statement on $D_\rho$ must be made
in the rescaled entire-coefficient algebra
(\cref{dyn:warn:rescale}).
\item (06) Rescaling gives \emph{guaranteed}, not maximal, domains. The
cubic example shows a sharp boundary, but the general maximal-domain
problem is left open (\cref{dyn:q:maximal}).
\item (06) \Cref{dyn:ex:badorder,dyn:ex:badfiber} show that none of the
three admissibility conditions can be dropped from a theorem of this form.
They do \emph{not} assert that finite fibers are necessary for every
resummation in some larger, differently specified class.
\end{enumerate}

\subsection{No global, all-scale, or topological claims}
\begin{enumerate}[label=\textbf{N\arabic*.},leftmargin=3.2em,start=47]
\item (02) No global conjugacy on all of $\No[i]$, no surreal-time flow, no
class-sized analytic construction, and no coherence between arbitrary
infinite-scale charts is claimed; \cref{dyn:prop:critical} supplies a
concrete critical-scale test that any such extension must pass.
\item (03) No all-scale theorem. The common-domain results concern finite
halos and the germ/formal results the infinitesimal monad;
\cref{dyn:prop:quadratic-scale} explicitly blocks automatic extension to
infinite surcomplex arguments. No global analytic atlas on $\No[i]$, no
ordinary contour theory on its fine topology, and no fine-topologically
continuous complex-time flow is constructed.
\item (03) There is no assertion that every set-theoretic or algebraic
conjugacy of evaluated functions must arise from a coefficient family;
failure in one specified analytic category prohibits no set-theoretic
conjugacy, and the formal-coefficient conjugacy always exists under
nonresonance and can be evaluated on the monad
(\cref{dyn:prop:monad}). No interchange of an infinite ordinary coefficient
sum with a non-locally-finite Hahn family is permitted.
\item (03) No claim of compatibility with the Berarducci--Mantova surreal
derivation or with any global composition operator on $\No$; no claim that
these coefficient categories fall outside every framework of
Cluckers--Lipshitz analytic structures \cite{CL}; and no sectorial or
resurgent summation is asserted.
\item (04) No global Julia or Fatou theory on the proper class $\No[i]$, no
classification of arbitrary surcomplex dynamical systems, and no ordinary
topological convergence in the full surcomplex plane is used anywhere.
\item (04) The formal logarithm is \emph{not} identified with a convergent
autonomous flow at a numerical parameter; \cite{GV} is cited precisely to
keep formal embedding and analytic embedding apart
(\cref{dyn:warn:embedding}).
\item (05) The slow-time theorem gives nothing past the ordinary
singularities of the supplied flow chart, licenses no general exchange of
ordinary with Hahn summation, and gives no interpretation of an arbitrary
surreal time as an actual continuous trajectory
(\cref{dyn:warn:slowtime}).
\item (05) No ordinary analytic specialization: growth estimates and a
summation theory are neither supplied nor claimed; sectorial analytic
moduli of ordinary parabolic germs are not removed; the analytic embedding
obstruction for ordinary parabolic germs is not resolved; and which Hahn
normal forms admit compatible ordinary analytic realizations is a distinct
problem (\cref{dyn:q:realization}).
\item (06) No ordinary analytic convergence in a complex perturbation
parameter is inferred: substituting $t^{\eta}=\eps_0\in\C^{\times}$ need
not exist on the coefficient algebra and can turn a valid Hahn series into
a divergent ordinary series.
\item (06) No fixed real-valued norm and no sequential completeness of
$K_\Gamma$ is assumed or available; all control is by supports, and the
proper-class field $\No[i]$ is reached only through compatible set-sized
workspaces, never as a single object.
\end{enumerate}

\subsection{Unresolved problems recorded as such}
\begin{enumerate}[label=\textbf{N\arabic*.},leftmargin=3.2em,start=57]
\item (02) The negative common-domain result \emph{requires} multiplier
drift $F'(0)=\lambda+u$ and two rationally independent positive exponents.
The source does \textbf{not} settle the fixed-multiplier problem when
$0<\sigma<\infty$ --- ``eliminating drift removes the direct reciprocal-jet
mechanism'' --- and does not claim the necessity direction for a fixed
cyclic value group (\cref{dyn:warn:drift,dyn:q:cyclic}).
\item (03) The intermediate regime $0<\tau<\infty$ for
\emph{common-domain} germs is explicitly unresolved
(\cref{dyn:q:commongerm}), and the radius bound
\eqref{dyn:eq:radius-depth} is a lower bound only, stated \emph{not} to
prove that radius collapse occurs (\cref{dyn:warn:lowerbound}).
\item (03) No universal analytic Poincar\'e--Dulac statement for resonant
multipliers of infinite order is proved
(\cref{dyn:q:resonant}).
\item (03) Higher-dimensional positive centralizers are reduced to
commuting positive vector fields but \emph{not} classified; the
one-variable meromorphic-ratio argument does not extend
(\cref{dyn:q:centralizers}).
\item (04) The moving-fixed-point case $w+\eps V+\eps^{2}R$ is open and
named as future work; the Euler family ``avoids that additional difficulty
rather than conceals it''. Several variables is open, the scalar quotient
$\alpha_{\mathrm c}/W$ having no direct replacement
(\cref{dyn:q:moving}).
\item (05) Several variables: no single reciprocal time one-form encodes a
vector field and commuting fields need not be constant multiples, so
\cref{dyn:thm:centralizer} and the finite-period normal form are explicitly
\emph{not} to be transported (\cref{dyn:q:centralizers}).
\end{enumerate}

\subsection{What the computations do and do not establish}
\begin{enumerate}[label=\textbf{N\arabic*.},leftmargin=3.2em,start=63]
\item (01) The $54$ finite checks do not establish the arbitrary-rank
support lemma, algebraic closure, maximality for every polynomial, the
transfinite coherent lifting, or novelty; they are not an implementation of
general Hahn fields and not proof-assistant verification.
\item (01) The finite procedure of \cref{dyn:app:certificate} is relative
to oracles: equality and valuation of arbitrary Hahn expressions need not
be effective, root-of-unity recognition for an arbitrary complex
coefficient needs a presentation allowing that decision, $d^{q}$ may be
very large, and treating a truncated zero as an exact zero would invalidate
the maximum. It is an exact characterization and a certificate, not a
polynomial-time complexity theorem.
\item (02) The finite checks test algebraic identities, signs,
normalizations and truncation only. They do \emph{not} prove Hahn
summability, well-founded recursion, ordinary analytic domain preservation,
an infinite small-divisor limsup, the exact convergence radii, or novelty;
and the multivariable test verifies nonzero denominators only in its finite
range, not nonresonance for all multi-indices of its sample tuple.
\item (03) The $21$ checks are finite polynomial jet identities truncated
by \emph{total parameter degree}, explicitly not a Hahn valuation cutoff in
a higher-rank group. They do not test arbitrary well-ordered supports,
analytic radius claims, continued-fraction asymptotics, the universal
quantifiers, or priority, and no proof-assistant verification is claimed
(\cref{dyn:warn:labelcutoff}).
\item (04) The $2{,}764$ exact assertions are finite consistency checks
only. No finite test proves Neumann summability, existence on an arbitrary
ordinary domain, compatibility for arbitrary value groups, or the
algebraicity dichotomy; the program is neither a complete symbolic
surcomplex CAS nor a numerical evaluator at arbitrary surreal inputs.
\item (05) \texttt{verify.py} is transcription and algebra checking only:
it proves nothing about arbitrary-rank support finiteness, global primitive
existence, the classification, or the slow-time theorem.
\textbf{Fifty of its $245$ assertions are finite lexicographic
illustrations for $n=1,\dots,50$; the all-$n$ inequality is proved in the
text, not inferred from them.}
\item (05) The verifier loses one top coefficient on division by $\eps$
(its division routine sets the new top coefficient to zero), so
reciprocal and pullback comparisons are checked only through degree $N-1$.
\item (06) The $119$ checks truncate by \emph{total label degree} in the
two independent labels, not by Hahn valuation, and only to degree four.
They are explicitly not a proof of arbitrary-support summability, nor of
\cref{dyn:thm:homological,dyn:thm:universal}, nor of
\cref{dyn:thm:pcentralizer}, nor of the Liouville construction
\cref{dyn:prop:explicitbad}.
\end{enumerate}

\subsection{The quasi-periodic source}
\begin{enumerate}[label=\textbf{N\arabic*.},leftmargin=3.2em,start=71]
\item (07) Self-declared as proposed original results with detailed proofs,
unrefereed, and not formalized in any proof assistant. Its README and title
page both state that the exact theorem package was not located in a
\emph{focused} literature review and that this does not establish priority;
independent mathematical review is called for before correctness or
originality is treated as settled. Its classical inputs --- ordinary
Fourier division with its subexponential criterion, the Liouville
construction method, positive-support Hahn inversion, and the finite
Jacobian identities --- are named as classical and are not claimed
\cite{Bounemoura,Kamienski,Neumann,Poonen}.
\item (07) A \emph{coefficientwise Hahn} theorem only. It explicitly does
\textbf{not} assert convergence after substituting a nonzero real parameter
for a Hahn monomial, and does not improve classical real analytic KAM or
Brjuno arithmetic \cite{Bounemoura,Chevallier}. No bound uniform in the
Hahn exponent on the analytic seminorms $\|f_\beta\|_{r}$ is supplied, and
growth faster than every exponential in the exponent index is permitted
(\cref{dyn:subsec:qplimits}).
\item (07) The torus derivatives are external derivations annihilating the
Hahn scalar field, deliberately \textbf{not} the Berarducci--Mantova
derivation on $\No$ \cite{BM}, so nothing in
\cref{dyn:sec:flag,dyn:sec:qplinear,dyn:sec:qpnonlinear} speaks to internal
surreal differential equations (\cref{dyn:warn:derivative}).
\item (07) \textbf{A support bound is not summability.} The family of
reciprocal divisors indexed by the Fourier modes is stated to be in general
\emph{not} Hahn-summable; only a common support is proved
(\cref{dyn:thm:flag}(iii), \cref{dyn:rem:qpnotsummable}), with ordinary
Fourier summation performed inside each Hahn coefficient and never
interchanged with it (\cref{dyn:warn:qptwosums}).
\item (07) No class-sized objects. Everything is a set-sized workspace,
with the full surreal field entering only through the embedding
\eqref{dyn:eq:embedding}; no continuity into the fine topology of $\No[i]$
and no class-indexed analytic limit is claimed. The invariant density
\cref{dyn:thm:qpdensity} is coefficientwise integration of analytic
functions --- \textbf{not} a countably additive $\No$-valued measure on
arbitrary subsets and \textbf{not} a real-time ergodic theorem.
\item (07) The nonlinear constant normal form requires a trivial exact
resonance lattice. The resonant case is left open, and the source warns
that replacing the mean projection by the resonant projection is not
justified, because the retained normal form is then no longer a constant
vector field (\cref{dyn:q:qpresonant}).
\item (07) The flag is finite but \emph{not} automatically computable:
detecting exact integer relations among arbitrary named reals may be
unavailable, so existence of the flag is explicitly separated from
decidability of its entries (\cref{dyn:q:qpeffective}).
\item (07) Three extensions are listed as needing further work ---
quantitative coefficient-growth classes preserved under the construction,
the resonant normal form, and effective arithmetic certificates --- and the
source states that these arise from the article and are \emph{not} claimed
to be catalogued open problems
(\cref{dyn:q:qpgrowth,dyn:q:qpresonant,dyn:q:qpeffective}).
\item (07) The finite checks verify neither the irrationality of the
Liouville constant \eqref{dyn:eq:qpliouvilleconstant}, nor any infinite-mode
subexponential estimate, nor the word lemma, nor the infinite-support
conjugacy theorem.
\item (07) An instance of this report's standing hazard, but handled: the
source cites \emph{no} invariance of any norm across a value-group
extension. The robustness of \cref{dyn:thm:qpjet} under enlarging $\Gamma$
is argued from uniqueness of the modewise Hahn inverse rather than assumed
(\cref{dyn:warn:qpextension}).
\end{enumerate}

\noindent The list above contains $80$ numbered items for the $78$ distinct
source-level limitations. The two extra entries arise because source 03's
limitation on resonant normal forms is recorded in two places ---
\textbf{N22}, the finite-order hypothesis $\Lambda^{q}=I$, and
\textbf{N59}, the openness of the infinite-order case --- and because
source 04's hypothesis audit is given its own entry
(\textbf{N23}) in addition to its fourteen listed non-claims. The
distribution of \emph{distinct} source-level limitations is
$11/9/10/14/12/12/10$ across sources $01$ through $07$; the corresponding
numbered-item counts are $11/9/11/15/12/12/10$.

\begin{warning}[The one sentence this report refuses]
\label{dyn:warn:refuse}
No statement of this report says that Hahn or surcomplex coefficients
\emph{abolish} the classical small-divisor problem. What is true is stated
in \cref{dyn:rem:invisible,dyn:rem:oneinvariant}: passing to positive Hahn
support \emph{separates} two questions that ordinary analytic convergence
forces one to solve at once --- algebraic dependence across perturbation
orders, and ordinary analyticity within each coefficient --- and the sharp
thresholds describe exactly that separation. Support alone suffices for
polynomial and formal coefficients at arbitrary rank
($\sdz$); arithmetic enters exactly when an ordinary domain or entirety is
demanded ($\sddr$, $\sdex$, $\sdfr$, $\sdss$); and under $\pares$ there is
no divisor family at all ($\sdno$). The quasi-periodic part makes the same
point in its sharpest form: \cref{dyn:prop:qpvalbound} is a genuine
theorem with \emph{no} arithmetic hypothesis whatever, and
\cref{dyn:thm:qplinear} shows that the moment an analytic strip is demanded
the arithmetic becomes not merely sufficient but necessary. The two
positions are complementary, not opposed, and stating them without their
coefficient categories is what would make them look contradictory.
\end{warning}

\appendix

\section{A short proof of the word-order input}
\label{dyn:app:higman}

We record the special case of Higman's lemma used in
\cref{dyn:lem:neumann}, so that the support proof can be audited without a
separate combinatorial argument hidden in the notation. It is classical
\cite{Higman}. Sources 02 and 05 each give the same argument; this is the
one copy.

Let $S$ be well ordered. A word $(a_1,\dots,a_m)$ \emph{embeds} in a word
$(b_1,\dots,b_n)$ if there are indices $i_1<\cdots<i_m$ with
$a_j\le b_{i_j}$ for every $j$. A sequence of words is \emph{bad} if no
earlier word embeds in a later one. Suppose an infinite bad sequence
exists. Choose a bad sequence $w_0,w_1,\dots$ minimally by word length at
each stage: after choosing a prefix, choose a shortest next word among
those admitting a bad infinite extension. The choices take place in a set,
so ordinary choice suffices, and no $w_n$ is empty.

Write $w_n=v_na_n$, separating the last letter. Any infinite sequence in a
well-ordered set has an infinite nondecreasing subsequence --- select a
position attaining the minimum of the current infinite tail, then repeat on
the tail strictly after that position; the successive minima cannot
decrease. So choose $n_0<n_1<\cdots$ with
$a_{n_0}\le a_{n_1}\le\cdots$ and consider
\[
 w_0,\dots,w_{n_0-1},v_{n_0},v_{n_1},v_{n_2},\dots .
\]
This is bad: an earlier $w_i$ embedding in a later $v_{n_j}$ would also
embed in $w_{n_j}$, contradicting the original badness; if $v_{n_i}$ embeds
in $v_{n_j}$ for $i<j$, appending the dominated last letters embeds
$w_{n_i}$ in $w_{n_j}$, again a contradiction; and comparisons within the
unchanged prefix were already forbidden. The new bad sequence agrees with
the chosen one until position $n_0$ and has the strictly shorter word
$v_{n_0}$ there, contradicting the minimal choice. Hence no bad sequence
exists, which is the required word well-quasi-order.

When $S\subset\Gamma_{>0}$, a proper word embedding \emph{strictly} raises
the sum: either an unused positive letter is present in the larger word, or
a matched letter is strictly larger. That observation converts the word
theorem into both the well-ordering of $S^{*}$ and the finiteness of the
words of one weight, as claimed in \cref{dyn:lem:neumann}.

\section{The finite certificate for the resonant threshold}
\label{dyn:app:certificate}

For polynomial inputs with effective coefficient arithmetic, the geometric
certificate of \cref{dyn:cor:finite-certificate} can be computed without
constructing $H$ at all.
\begin{enumerate}[label=\arabic*.]
\item Determine $\alpha=\bar\lambda$. If it is not a root of unity, compute
the finite maximum $\eta(f)$ and return $r=\eta(f)$.
\item If its order is $q$, compute the finite composition $g=f^{\circ q}$,
$c=\lambda^{q}-1$ and $\delta=v(c)$.
\item Compute every nonzero nonlinear return coefficient $b_n$
\emph{including actual cancellations}, and take the maximum
\eqref{dyn:eq:exact-r}.
\item Scale by $\rho=t^{r}$ and reduce \eqref{dyn:eq:Q-shell} to obtain
$P$. Its degree gives the shell multiplicity; factoring $P$ determines the
residue directions; and $XP'(X)$ gives the leading multiplier ratios at
simple roots.
\end{enumerate}
This is a finite algebraic procedure \emph{relative to the required
oracles}, and not a claim of universal computability for arbitrary
surcomplex inputs: see \textbf{N64} of \cref{dyn:sec:nonclaims}.

\section{A coefficientwise recipe for a certified linearizer coefficient}
\label{dyn:app:recipe}

Suppose one wishes to compute $h_\gamma$ for polynomial input in
\cref{dyn:thm:main}, with an effective presentation of the finite word set
$\word_S(\gamma)$. The following specifies what must be certified.

Collect the finite letter set $D_S(\gamma)$ and the finite set of partial
sums of all words of weight $\gamma$: only input coefficient polynomials at
these letters can matter. The maximum word length gives the bound
$L=\ell_S(\gamma)$. In \eqref{dyn:eq:linearizer-series} retain only
$k=0,\dots,L-1$, and in each Taylor expansion of $\NN_f$ retain only
products whose support words can complete an allowed word of total weight
$\gamma$; polynomial degrees can be bounded in advance by
\eqref{dyn:eq:degree-bound}. Apply $\Li^{-1}$ to each retained monomial by
the exact scalar division \eqref{dyn:eq:Linverse}. Nonresonance of the
finitely many divisors actually used must be established, or supplied as an
exact symbolic condition. Sum the resulting finite expressions and check
Schr\"oder's equation at the requested weight; the word certificate proves
that discarded expressions cannot affect it. This is a finite exact
algorithm \emph{relative to the support certificate and the polynomial
coefficient data}.

For analytic input the same support certificate still applies, but
$\Li^{-1}$ is an operator on an entire coefficient function, not merely a
finite list of Taylor coefficients; a radius certificate from
\cref{dyn:thm:radius-depth} is then part of the analytic output. A finite
jet calculation alone must \emph{not} be described as computing the entire
analytic coefficient, nor as proving its convergence
(\cref{dyn:warn:finitedep}).

The parabolic counterpart, for \cref{dyn:thm:classification}, is
explicit in the same sense. Given $F_1,F_2$ in one nonsingular stratum,
form their logarithmic vector fields by \eqref{dyn:eq:explog}, compute the
normalized reciprocal forms $\beta_i=b_i\dd z$ by the geometric unit
expansion of \cref{dyn:lem:unit}, and test the periods of
$(b_1-b_2)\dd z$; on a finitely punctured plane that is a finite list of
Hahn periods, though computing an arbitrary coefficient integral is not
assumed decidable. When the periods vanish, integrate coefficientwise from
$z_{*}$ to get $p$ with $p'=b_1-b_2$, $p(z_{*})=0$; with $b_2=b_0+u$ and
$S$ the positive union of the supports of $p$ and $u$, the desired
$h=H-z$ has support in $S^{*}\setminus\{0\}$ and, for each $\gamma$ in that
well-ordered set,
\begin{equation}\label{dyn:eq:explicit-recursion}
 h_\gamma=b_0^{-1}\left(p_\gamma-[t^{\gamma}]
 \Bigl[uh+\sum_{n\ge2}\frac{b_2^{(n-1)}}{n!}h^{n}\Bigr]\right),
\end{equation}
the bracket using only earlier coefficients of $h$. The number of
operations at each exponent is finite, but the total support may be
transfinite and need not have a finite presentation. For an ordinary
single-parameter expansion $b_2=b_0+\eps b_1+\eps^{2}b_2^{[2]}+\cdots$ and
$p=\eps p_1+\eps^{2}p_2+\cdots$ the first coefficients are
\[
 h_1=\frac{p_1}{b_0},\qquad
 h_2=\frac1{b_0}\Bigl(p_2-b_1h_1-\tfrac12b_0'h_1^{2}\Bigr),
\]
displaying both sources of nonlinear correction: the tail of the target
time form, and the change of its ordinary leading coefficient along the
displacement. The support certificate is part of the result --- it is what
makes each coefficient defined even when the multiples of the smallest
input exponent are not cofinal --- and it is \emph{not} a promise that
arbitrary Hahn-series equality, primitive existence or period evaluation
can be decided by a finite program.

\section{The single-exponential gauge recursion as an algorithm}
\label{dyn:app:algorithm}

For a finite collection of input labels $x_1,\dots,x_k$, an implementation
of \cref{dyn:thm:normalform} can work over the ordinary formal label ring
through a chosen total label degree $N$:
\begin{enumerate}[label=\arabic*.]
\item set $\chi=0$ and retain $H=H_0+R$ with independent labels attached to
its input layers;
\item for $n=1,\dots,N$ compute $e^{\ad_\chi}H$ through total label degree
$n$, using the current $\chi$ of label degree less than $n$;
\item for each coefficient $E_\nu$ with $|\nu|=n$, append
$-L^{-1}Q_\Pi E_\nu$ to the coefficient of $x^{\nu}$ in $\chi$;
\item recompute $e^{\ad_\chi}H$ through the requested order, and compute
coordinates by the same exponential and actions by the inverse exponential.
\end{enumerate}
Adding a degree-$n$ generator changes the degree-$n$ transformed Hamiltonian
only by $L\chi_n$, so step 3 is the \emph{single-exponential gauge}
recursion --- \textbf{not} the different algorithm that composes a list of
successive exponential maps without taking their logarithm
(\cref{dyn:warn:gauge}). For an infinite input label set this is a
conceptual construction rather than a finite executable loop; homogeneous
leaf degree still organizes the universal expression trees of
\cref{dyn:def:trees}, and at a specified Hahn weight
\cref{dyn:lem:neumann} permits only finitely many leaf words, hence
finitely many homogeneous leaf degrees and finite expressions. That is the
exact connection between the finite-label algorithm and the
arbitrary-support proof. Total label degree is not a valuation cutoff
(\cref{dyn:warn:labelcutoff}).

\section{Antecedent ledger}
\label{dyn:app:antecedents}

For audit, the external inputs and their exact roles. None is claimed as a
contribution of this report.
\begin{longtable}{@{}P{0.30\textwidth}P{0.62\textwidth}@{}}
\toprule
Input & Role here\\
\midrule
\endhead
Neumann's positive-support lemma \cite{Neumann}, via Higman's word
well-quasi-order \cite{Higman}, with Poonen's account of
Hahn--Mal'cev--Neumann units \cite{Poonen}
& \Cref{dyn:lem:neumann,dyn:cor:ancestry,dyn:app:higman}: the certificate
behind \emph{every} infinite operator expression in this report. Classical;
used as a black box for Higman's lemma beyond the reproved special case.\\[3pt]
Algebraic closure of a maximally complete Hahn field
\cite[Cor.~4]{Poonen}
& \Cref{dyn:lem:root-count,dyn:thm:workspace} only. Not needed for
\cref{dyn:sec:hamiltonian}.\\[3pt]
Conway normal forms and surreal normal-form arithmetic \cite{Gonshor}
& \eqref{dyn:eq:embedding} and the workspace discipline. No set of all
surreals is formed.\\[3pt]
Formal logarithms, flow embedding, Jordan--Chevalley /
semisimple--unipotent splitting \cite{Ribon}
& The classical pattern behind
\cref{dyn:thm:exp-log,dyn:thm:finite-order}. What is established here is
the arbitrary-positive-Hahn-support, same-coefficient-domain
implementation, not the pattern.\\[3pt]
Formal interpolation of near-identity maps versus analytic embedding at a
numerical parameter \cite{GV}
& Cited exactly to mark the line that
\cref{dyn:warn:embedding} draws.\\[3pt]
Exponentiating strongly summable derivations of generalized-series algebras
\cite{BKKPS}; automorphisms of the surreals \cite{KKS}
& The general algebraic mechanism behind
\cref{dyn:thm:exp-log,dyn:prop:exp}. The existence of a formal logarithm is
credited, not claimed.\\[3pt]
Julia's equation and differential transcendence of iterative logarithms
\cite{AB}
& Background for \cref{dyn:cor:julia} and
\cref{dyn:sec:periods}. The parameter-Hahn construction here is a
different category; the analytic embedding obstructions for ordinary
parabolic germs are not removed.\\[3pt]
Ordinary homological equation, small divisors, formal-versus-analytic
linearization \cite{CarlettiMarmi,Marmi}; the Brjuno--Yoccoz quadratic
theorem \cite{Yoccoz,BuffCheritat}
& \Cref{dyn:sec:divisors}; and Brjuno--Yoccoz used as an
\emph{external} theorem, not reproved, in
\cref{dyn:prop:noncoherent,dyn:rem:nocontra,dyn:prop:quadratic-scale}.\\[3pt]
Regular dependence of ordinary linearization on the multiplier
\cite{CarminatiMarmi}; explicit tree expansions
\cite{FauvetMenousSauzin}; linear versus nonlinear arithmetic conditions in
regularity classes \cite{Bounemoura}
& Natural parameter-regularity context for the reciprocal jets
\eqref{dyn:eq:reciprocaljet}. No Whitney estimates, no joint analytic
parameter convergence, and no novelty claim for tree-organized formal
coefficients.\\[3pt]
Equal-characteristic-zero ultrametric linearization, sharp disk radii, and
$p$-adic bounding spheres without periodic points
\cite{Lindahl,LindahlPadic}
& The rank-one antecedent of \cref{dyn:sec:return} and the
source of the cross-note \cref{dyn:subsec:berkovich} and the
characteristic warning \cref{dyn:warn:charzero}.\\[3pt]
Modified differential equations in backward error analysis
\cite{Hairer}
& The origin of the coefficients \eqref{dyn:eq:euler-known}, which are
known examples (\cref{dyn:warn:hairer}); its warning that the ordinary
step-size series generally diverges is cited and accepted.\\[3pt]
Convergence and generic divergence of Birkhoff normal forms
\cite{PerezMarco,Krikorian}
& Classical context for \cref{dyn:sec:hamiltonian}; ordinary analytic
convergence is never inferred from a Hahn identity.\\[3pt]
Transseries as germs of surreal functions \cite{BM}; fields with analytic
structure \cite{CL}; Hahn-meromorphic function theory
\cite{MuellerStrohmaier}
& Adjacent frameworks, kept distinct: see
\cref{dyn:warn:derivative} and \textbf{N50} of
\cref{dyn:sec:nonclaims}. No theorem is transferred between categories
without a separate proof.\\[3pt]
Linear versus nonlinear arithmetic conditions for torus vector fields
\cite{Bounemoura}; the one-frequency cohomological equation and its
Liouville obstructions \cite{Kamienski}; the Brjuno condition through best
approximations \cite{Chevallier}
& The classical small-divisor background of
\cref{dyn:sec:qplinear,dyn:sec:qpnonlinear}: ordinary Fourier division and
its subexponential criterion, and the Liouville mechanism adapted in
\cref{dyn:thm:qpjet}. Neither the ingredient nor any strengthening of
ordinary analytic KAM arithmetic is claimed
(\textbf{N72}, \cref{dyn:subsec:qplimits}).\\[3pt]
The motivating repository at its pinned snapshot \cite{Repo}
& Sources 01--06 cite it at commit
\texttt{39f2be6667ade51bca2b45daa47e289d69c09764} and source 07 at
\texttt{aa846271b4dcae2c055b216126a87210292ec19b}, for the
coefficient-category distinction only; none uses an unrefereed repository
theorem as a premise. The inspections were targeted, not exhaustive
(\textbf{N9}, \textbf{N13}, \textbf{N71}).\\
\bottomrule
\end{longtable}

\section{Notation at a glance}
\label{dyn:app:notation}

\begin{longtable}{@{}P{0.22\textwidth}P{0.71\textwidth}@{}}
\toprule
Symbol or condition & Meaning and role\\
\midrule
\endhead
$\Gamma$, $K_\Gamma$, $t^{\gamma}$ & Nonzero set-sized ordered abelian
group; $K_\Gamma=\C((t^{\Gamma}))$; $t^{\gamma}=\omega^{-\gamma}$ for the
surcomplex reading \eqref{dyn:eq:embedding}. Divisibility and algebraic
closure are assumed only where flagged.\\[3pt]
$v$, $\OO_\Gamma$, $\m_\Gamma$ & $v=\min\supp$ with \emph{increasing}
supports (\cref{dyn:warn:orientation}); $v\ge0$ and $v>0$. Trivial on
nonzero complex constants (\cref{dyn:rem:invisible}).\\[3pt]
$+$, $>0$ subscript & All nonzero support exponents strictly positive. Not
positivity of ordinary complex coefficients.\\[3pt]
$B_r$, $S_r$, $D_\rho$ & Valuation ball, shell and polydisk; larger
threshold means \emph{smaller} domain (\cref{dyn:subsec:radius}).
$S_r$ is not a topological boundary.\\[3pt]
$\mathbb D_R$, $\mathbb D_R^{d}$, $D$, $U$ & Ordinary disk, polydisk or
open set; larger $R$ means \emph{larger} domain. A domain for coefficient
functions only (\cref{dyn:warn:ball}).\\[3pt]
$U^{\#}_\Gamma$, $D^{\#}$, $\m_\Gamma^{d}$ & Finite halo and infinitesimal
monad; evaluation \eqref{dyn:eq:halo-evaluation} is a faithful algebra
homomorphism (\cref{dyn:prop:eval}).\\[3pt]
$\rad$, $\radiso$ & Ordinary Taylor radius and isotropic Taylor radius of
one coefficient function; larger is better.\\[3pt]
$S^{*}=M(S)$, $\word_S(\gamma)$, $\ell_S(\gamma)$, $D_S(\gamma)$ &
Additive monoid, ordered words of weight $\gamma$, maximum word length,
occurring letters (\cref{dyn:lem:neumann}). $\ell_S(\gamma)$ is finite even
when infinitely many support elements precede $\gamma$.\\[3pt]
$\word_{B,S}(\gamma)$, $\ell_{B,S}(\gamma)$,
$E_{B,S}(\gamma)$, $D_{B,S}(\gamma)$ &
Input-exponent/word pairs reaching $\gamma$ from $B$, their maximum word
length (zero if there are no pairs), their starting exponents, and their
positive letters. These are the relative data for arbitrary Hahn inputs
in \cref{dyn:cor:ancestry}.\\[3pt]
$\BB_{\mathsf M}$ & Uniform-monoid support algebra
(\cref{dyn:def:BM}); the uniformity across Taylor degrees is
load-bearing.\\[3pt]
$\Li$, $\Li_L$, $L$ & Homological operators
\eqref{dyn:eq:homological}, \eqref{dyn:eq:hamhomological}; $\Li^{-1}$
preserves Hahn exponents and total degree.\\[3pt]
$\sigma(\lambda)$, $\tau(\Lambda)$, $\Theta(\omega)$ & The three divisor
rates: drifted, exact-multiplier, frequency
(\cref{dyn:subsec:smalldivisor}). $\sigma,\tau$ are log-rates in
$[0,\infty]$; $\Theta\in[1,\infty]$ is the exponential of a rate.\\[3pt]
$\resq$, $\mres$/$\nres$, $\fres$, $\pares$, $\qres$ & The five resonance
notions (\cref{dyn:subsec:resonance}). A threshold proved under one is
never quoted under another.\\[3pt]
$\sdz$, $\sddr$, $\sdex$, $\sdfr$, $\sdno$, $\sdss$ & The six small-divisor
notions (\cref{dyn:subsec:smalldivisor}).\\[3pt]
$\Upsilon$, $\varpi_\beta$, $\qlat$ & The \emph{Hahn} frequency vector in
$F_\Gamma^{d}$, its real coefficient vectors, and its exact resonance
lattice $\{k:k\cdot\Upsilon=0\}$. Never $\omega\in\C^{d}$ of
\cref{dyn:sec:hamiltonian}, and never the exponent $\Omega$ of
\cref{dyn:ex:ranktwo} (\cref{dyn:warn:twofrequencies}).\\[3pt]
$\mathsf L_j$, $\Sigma_j$, $\gamma_j$, $\fla_j$, $r$ & The resonance
lattices of the flag, its strata, its visible levels, its leading
\emph{real} forms and its length $r\le d$
(\cref{dyn:thm:flag}). $\mathsf L_j$ is not the multiplier $\Lambda$ and
not the homological $\Li$; $\Sigma_j$ is not the ordinary time disk
$\Sigma$ of \cref{dyn:thm:slowtime}.\\[3pt]
$\mathsf E$, $\Tups$, $\kups$ & The shifted tail generators, the common
reciprocal support $\bigcup_j(-\gamma_j+\mathsf E^{*})$, and the last
visible level $\gamma_r$: an \emph{element of $\Gamma$}, never a rate
(\cref{dyn:warn:twofrequencies}).\\[3pt]
$\TT^{d}_\rho$, $\strip$, $\HH_\Gamma(\strip)$ & The fixed open complex
strip, its periodic holomorphic functions, and the fixed-strip Hahn
algebra: the torus instance of the fixed-domain row of
\cref{dyn:subsec:rings}, with one $\rho$ for every Hahn exponent and no
uniform seminorm bound.\\[3pt]
$\Cups$, $\Gups$, $\Pi_{\qlat}$, $\av$, $\Qz$ & The cohomological operator
$\sum_j\Upsilon_j\partial_{\theta_j}$, its normalized inverse
\eqref{dyn:eq:qpgreen}, resonant Fourier projection, coefficientwise torus
mean, and mean-zero projection. $\Qz$ is not the Hamiltonian $Q_\Pi$.\\[3pt]
$\mathcal X$, $\mathcal V$, $\Phi_h$, $\mathfrak c$, $\mu_{\mathcal X}$ &
The perturbed torus field $\Upsilon+\mathcal V$, its perturbation, the
torus coordinate change $\id+h$ (\emph{not} the flow $\Phi_X$ or $\Phi_c$),
the constant frequency correction, and the invariant density
\eqref{dyn:eq:qpdensity}.\\[3pt]
$r$, $q$, $c$, $\delta$, $P$, $M$ & Finite-return threshold
\eqref{dyn:eq:exact-r}, residue order, $\lambda^{q}-1$, $v(c)$, shell
polynomial \eqref{dyn:eq:P}, its degree.\\[3pt]
$\eps=t^{\delta}$, $s=t^{\delta/q}$ & A single Hahn monomial of positive
valuation; \emph{not} a small complex number. All $O(\eps^{k})$ mean
divisibility in a formal series ring.\\[3pt]
$X_F$, $W_F$, $\Phi_X$, $\Phi_c$ & Logarithmic vector field, its one-variable
coefficient, time-one and time-$c$ maps
(\cref{dyn:thm:exp-log,dyn:cor:flow}).\\[3pt]
$\Omega_F=\dd z/W_F$, $\beta_F=\eps\Omega_F$ & Time form and its
normalization; ``holomorphic'' refers to the coefficient functions, and the
valuation may be negative.\\[3pt]
$J_\delta$, $I_\delta(D)$ & Flat ideal and one-scale-jet kernel, formed in
$K_{\ge0}$ and $\AAA_{\ge0}$, \emph{never} in $K$ or $\AAA$
(\cref{dyn:warn:flatideal}).\\[3pt]
$\BB=\C(w)((t^{\Gamma}))$, $\LL$ & The algebraicity base with arbitrary
rational-function coefficients (\emph{not} $K_\Gamma(w)$;
\cref{dyn:warn:base}) and the meromorphic ambient field
\eqref{dyn:eq:ambientfield}.\\[3pt]
$\rho_a$, $M_a$, $\rho_q$ & Exact monodromy exponent and multiplier
\eqref{dyn:eq:monodromy}, and the model's common exponent
\eqref{dyn:eq:rhoq}.\\[3pt]
$\Pi$, $Q_\Pi$, $\chi$, $J_j$, $\Theta(\omega)$ & Action projection and its
complement, the gauge-normalized generator, the commuting actions, and the
frequency divisor invariant.\\[3pt]
$w_\rho(\ell)$, $\mathscr S_\rho$ & Shifted monomial weight
\eqref{dyn:eq:weights} and the Hamiltonian rescaling operator
\eqref{dyn:eq:S}, a Poisson Lie but \emph{not} a ring homomorphism.\\[3pt]
$I_{\Fix(F)}$ & Ideal generated by $F_j-z_j$, compared by the explicit
invertible matrix \eqref{dyn:eq:fixedmatrix}, not by its zero set.\\
\bottomrule
\end{longtable}

\section{Per-theorem hypothesis audit}
\label{dyn:app:audit}

\begin{longtable}{@{}P{0.26\textwidth}P{0.67\textwidth}@{}}
\toprule
Statement & Indispensable hypotheses, and the boundary of the conclusion\\
\midrule
\endhead
\cref{dyn:lem:neumann} &
Well-ordered \emph{positive} support. No Archimedean hypothesis, no
rank-one, no countability, no finitely generated grid. Gives finite words,
not a valuation-contraction estimate
(\cref{dyn:warn:contraction}).\\[3pt]
\cref{dyn:prop:eval} &
One common ordinary domain for all coefficients; positive support of the
infinitesimal increment. Faithful, but polynomial/entire coefficients do
not automatically evaluate at infinite surcomplex points.\\[3pt]
\cref{dyn:thm:exp-log} &
Ordinary reduction the identity; positive perturbative support; arbitrary
open $U$. No ordinary parameter convergence, no uniform coefficient
bound.\\[3pt]
\cref{dyn:thm:fixed-ideal} &
Positive vector field; nonzero admissible time. Ideals compared by
\eqref{dyn:eq:fixedfactor}. Says nothing about maps whose reduction is not
the identity.\\[3pt]
\cref{dyn:thm:centralizer} &
$U\subset\C$ \emph{connected}, one variable, $F\ne\id$; admissible times
$v(c)+\alpha>0$. Not a classification in $d>1$
(\cref{dyn:rem:centralizer-limits}); the centralizer can \emph{grow} under
workspace extension (\cref{dyn:warn:workspace}).\\[3pt]
\cref{dyn:thm:main} &
$\nres$ with the multiplier \emph{exactly} $\Lambda$; positive support;
no degree-$0$ or degree-$1$ coordinate terms at any exponent; a fixed
polydisk of finite positive radius in row one. Rows are not
interchangeable; the sixth category is \cref{dyn:q:commongerm}.\\[3pt]
\cref{dyn:thm:radius-depth} &
$\tau<\infty$. \textbf{Lower bound only}
(\cref{dyn:warn:lowerbound}).\\[3pt]
\cref{dyn:thm:degrees} &
Polynomial input. Holds for any invertible diagonal nonresonant $\Lambda$,
without $|\lambda_j|=1$.\\[3pt]
\cref{dyn:thm:lifting} &
$\nres$; drift permitted; well-ordered positive $A$; $f_a(0)=0$. No
uniform bound, no Gevrey estimate, no dependence on $t$ at $t=0$.\\[3pt]
\cref{dyn:thm:collapse,dyn:thm:trichotomy,dyn:thm:multicollapse} &
\textbf{Multiplier drift required}, plus two rationally independent
positive exponents; several-variable version needs multiplicative detuning
and equal-radius polydisks. Not asserted for a fixed cyclic value group
(\cref{dyn:warn:drift,dyn:q:cyclic}).\\[3pt]
\cref{dyn:thm:exact,dyn:thm:shell,dyn:thm:phase} &
$\resq$ with $\lambda$ \emph{not} a root of unity; one variable;
polynomial $f$; \emph{equal characteristic zero}
(\cref{dyn:warn:charzero}); divisible $\Gamma$ for
\cref{dyn:lem:root-count}. ``Exact'' means the maximal centered valuation
ball only.\\[3pt]
\cref{dyn:thm:coherent} &
$\resq$. Gives one fixed ordinary disk on which $P$ has no zero; no uniform
norm bound on $(H_\gamma)$; resonant-only, and not an iff
(\cref{dyn:warn:noncoherent}).\\[3pt]
\cref{dyn:thm:stability} &
Same multiplier $\lambda$; inequalities on the \emph{return} coefficients;
squarefree $P$ for the matching. Sufficient, not necessary
(\cref{dyn:warn:stability}).\\[3pt]
\cref{dyn:thm:common,dyn:thm:dichotomy} &
$\pares$; simply connected $U$ containing $0$; $V$ nonvanishing on
$U\setminus\{0\}$ with a \emph{simple} zero at $0$; polynomial Euler form
for the dichotomy; divisible $\Gamma$. Transcendence over
$\C(w)((t^{\Gamma}))$ only, and not differential transcendence
(\cref{dyn:warn:diffalg}).\\[3pt]
\cref{dyn:thm:monodromy} &
A \emph{simple} zero $a$; ordinary loops on the ordinary universal cover.
Finite order exactly when $V'(a)=\mu$; the obstruction is visible already
at order $\eps^{1}$.\\[3pt]
\cref{dyn:thm:sharp} &
$\pares$; the model $f=(1+\eps)z+z^{q+1}$; divisible $\Gamma$. Maximality
for \emph{centered} valuative balls; two distinct failure mechanisms at and
below the threshold (\cref{dyn:warn:sharp}).\\[3pt]
\cref{dyn:thm:difference} &
$F\ne\id$. Outside the nonsingular stratum the quotient $g/W_F$ must still
lie in $\AAA$; period vanishing alone is insufficient
(\cref{dyn:ex:localobstruction}).\\[3pt]
\cref{dyn:thm:classification,dyn:thm:moduli} &
Same leading data $(\delta,a)$ with $a$ \emph{nowhere zero}; near-identity
conjugator on the \emph{original} common domain; ordinary coefficientwise
periods; the specified puncture-loop basis
(\cref{dyn:warn:classification,dyn:warn:nowherezero}).\\[3pt]
\cref{dyn:thm:flat} &
$\delta$ \emph{noncofinal}, i.e.\ $J_\delta\ne0$; jets formed in
$\AAA_{\ge0}$ (\cref{dyn:warn:flatideal}). Not a statement about ordinary
exponentially small asymptotics.\\[3pt]
\cref{dyn:thm:slowtime} &
A \emph{specified} ordinary holomorphic flow chart with image in $D$ and
nowhere-vanishing variational factor $J$. Nothing past the chart's ordinary
singularities (\cref{dyn:warn:slowtime}).\\[3pt]
\cref{dyn:thm:homological,dyn:thm:universal} &
Nonresonant \emph{ordinary complex} $\omega$. Measures one coefficientwise
operator; not a Bruno- or KAM-type convergence criterion
(\cref{dyn:warn:theta}).\\[3pt]
\cref{dyn:thm:normalform} &
$\lnot\fres$; positive well-ordered $S$; $\ord_0R_s\ge3$; case (a)
polynomial layers with no arithmetic condition, case (b) entire layers with
$\Theta<\infty$. Uniqueness is gauge-relative
(\cref{dyn:warn:gauge}).\\[3pt]
\cref{dyn:thm:pcentralizer} &
The \emph{entire-coefficient} Hahn algebra. Does not classify local,
multivalued or fine-topologically locally constant integrals
(\cref{dyn:warn:pcent}).\\[3pt]
\cref{dyn:thm:rescale} &
Polynomial input admissible at $\rho$: positivity, well-ordering
\emph{and} finite fibers, all three independent
(\cref{dyn:ex:badorder,dyn:ex:badfiber}). Guaranteed, not maximal, domains
(\cref{dyn:warn:rescale}).\\[3pt]
\cref{dyn:thm:finite-order} &
$\Lambda^{q}=I$ --- the multiplier itself of finite order, which is neither
$\resq$ nor mere $\mres$. Resonant normal forms not unique
(\cref{dyn:warn:finiteorder}).\\[3pt]
\cref{dyn:thm:flag} &
Finite dimension $d$; \emph{one} well-ordered support for the whole
frequency vector. Without finite $d$ the kernel rank need not reach its
minimum in finitely many steps; independent unbounded scalar choices do not
supply a shared support. Gives a \emph{common support}, not summability of
the reciprocal family (\cref{dyn:rem:qpnotsummable}); and $r$ counts
visible scales, not monomials of $\Upsilon$
(\cref{dyn:rem:qpvisible}).\\[3pt]
\cref{dyn:prop:qpvalbound} &
No Fourier modes in $\qlat$. \textbf{No divisor estimate of any kind}; this
is the one quantitative statement of the quasi-periodic part that is
arithmetic-free.\\[3pt]
\cref{dyn:thm:qplinear} &
$(\sdss)$ on \emph{every} nonresonant stratum; \emph{all} input
coefficients holomorphic on \emph{one} open strip. Without the first, a
forced coefficient can fail to be a distribution; without the second,
independent germs need not give a common analytic domain
(\cref{dyn:conv:rings}). The loss $\kups$ is attained
(\cref{dyn:cor:qpsharp}); no loss bound uniform in the Hahn exponent is
asserted (\cref{dyn:rem:qpdiophantine}).\\[3pt]
\cref{dyn:thm:qpjet} &
Integer exponents only; the explicit Liouville constant
\eqref{dyn:eq:qpliouvilleconstant}. The obstruction is to assembling
infinitely many modewise inverses into \emph{one} analytic coefficient ---
the scalar inverses all exist. Robustness under enlarging $\Gamma$ is
argued from uniqueness of the modewise Hahn inverse, not from any norm
invariance (\cref{dyn:warn:qpextension}).\\[3pt]
\cref{dyn:thm:qpnormalform} &
$\Gamma$ nontrivial; $\lnot\qres$; $(\sdss)$; normalization $\gamma_1=0$;
$v(\mathcal V)>\kups$ \emph{strictly} --- \cref{dyn:prop:qpboundary} fails
at equality. Existence needs finitely many labelled support words per
exponent; increasing valuations alone do \textbf{not} produce a Hahn sum
(\cref{dyn:warn:qpcontraction}). Uniqueness needs the mean-zero
normalization and positive valuation of $h$; constant translations remove
the normalization.\\[3pt]
\cref{dyn:thm:qpnecessity,dyn:cor:qpequivalence} &
$(\sdss)$ \emph{fails}; $\lnot\qres$; $\Gamma$ nontrivial. Refutes a
\emph{universal} theorem at arbitrarily high valuation; says nothing about
any particular perturbation.\\[3pt]
\cref{dyn:thm:qpdensity} &
Coefficientwise torus mean and finite Jacobian calculus, on the conjugacy
of \cref{dyn:thm:qpnormalform}. No measure on arbitrary subsets and no
real-time ergodic theorem follows (\textbf{N75}).\\
\bottomrule
\end{longtable}

\clearpage
\phantomsection
\addcontentsline{toc}{section}{References}
\begin{thebibliography}{99}
\raggedright

\bibitem{Repo}
Vladimir Reshetnikov,
\emph{Surreal}: surreal and surcomplex research reports and formalization,
GitHub repository. Documentation catalogue \texttt{docs/README.md},
\texttt{docs/manifest.tex}, and the packages
\texttt{docs/surcomplex/analysis/},
\texttt{docs/surcomplex/analytic-geometry/},
\texttt{docs/surcomplex/contours-and-stokes/}. Snapshot commit
\texttt{39f2be6667ade51bca2b45daa47e289d69c09764}, inspected
21 September 2026. The reports are AI-assisted, unrefereed research drafts,
not independent peer-reviewed validation.
\url{https://github.com/VladimirReshetnikov/Surreal}.

\bibitem{Neumann}
B.~H. Neumann,
On ordered division rings,
\emph{Transactions of the American Mathematical Society} \textbf{66}
(1949), 202--252.
\href{https://doi.org/10.1090/S0002-9947-1949-0032593-5}{doi:10.1090/S0002-9947-1949-0032593-5}.

\bibitem{Higman}
G.~Higman,
Ordering by divisibility in abstract algebras,
\emph{Proceedings of the London Mathematical Society} (3) \textbf{2}
(1952), 326--336.
\href{https://doi.org/10.1112/plms/s3-2.1.326}{doi:10.1112/plms/s3-2.1.326}.

\bibitem{Poonen}
B.~Poonen,
Maximally complete fields,
\emph{L'Enseignement Math\'ematique} (2) \textbf{39} (1993), 87--106
(in particular Corollary~4); and
\emph{Units in Hahn--Mal'cev--Neumann rings}, manuscript dated
14 January 2026, 4~pp.
\url{https://math.mit.edu/~poonen/papers/malcev.pdf}.

\bibitem{Gonshor}
H.~Gonshor,
\emph{An Introduction to the Theory of Surreal Numbers},
London Mathematical Society Lecture Note Series 110,
Cambridge University Press, 1986.

\bibitem{Ribon}
J.~Rib\'on,
Embedding smooth and formal diffeomorphisms through the Jordan--Chevalley
decomposition,
\emph{Journal of Differential Equations} \textbf{253} (2012), no.~12,
3211--3231.
\href{https://doi.org/10.1016/j.jde.2012.09.009}{doi:10.1016/j.jde.2012.09.009};
preprint \href{https://arxiv.org/abs/1107.3601}{arXiv:1107.3601}.

\bibitem{GV}
V.~Gelfreich and A.~Vieiro,
Interpolating vector fields for near identity maps and averaging,
\href{https://arxiv.org/abs/1711.01983}{arXiv:1711.01983} (2017);
published version
\href{https://doi.org/10.1088/1361-6544/aacb8e}{doi:10.1088/1361-6544/aacb8e}.

\bibitem{BKKPS}
V.~Bagayoko, L.~S. Krapp, S.~Kuhlmann, D.~Panazzolo, and M.~Serra,
\emph{Automorphisms and derivations on algebras endowed with formal
infinite sums}, \href{https://arxiv.org/abs/2403.05827}{arXiv:2403.05827},
version~2, 2025 (first posted 2024).

\bibitem{KKS}
E.~Kaplan, L.~S. Krapp, and M.~Serra,
\emph{Decomposing the automorphism group of the surreal numbers},
\href{https://arxiv.org/abs/2509.22374v3}{arXiv:2509.22374v3}, 2026
revision.

\bibitem{AB}
M.~Aschenbrenner and W.~Bergweiler,
Julia's equation and differential transcendence,
\emph{Illinois Journal of Mathematics} \textbf{59} (2015), 277--294.
Preprint \url{https://arxiv.org/abs/1307.6381}.

\bibitem{CarlettiMarmi}
T.~Carletti and S.~Marmi,
Linearization of analytic and non-analytic germs of diffeomorphisms of
$(\C,0)$,
\emph{Bulletin de la Soci\'et\'e Math\'ematique de France} \textbf{128}
(2000), no.~1, 69--85.
\href{https://doi.org/10.24033/bsmf.2363}{doi:10.24033/bsmf.2363}.

\bibitem{Marmi}
S.~Marmi,
\emph{An Introduction to Small Divisors Problems},
Istituti Editoriali e Poligrafici Internazionali, 2000.
\href{https://arxiv.org/abs/math/0009232}{arXiv:math/0009232}.

\bibitem{Yoccoz}
J.-C. Yoccoz,
Th\'eor\`eme de Siegel, nombres de Bruno et polyn\^omes quadratiques,
in \emph{Petits diviseurs en dimension 1},
\emph{Ast\'erisque} \textbf{231} (1995), 3--88.
\url{https://www.numdam.org/item/AST_1995__231__1_0/}.

\bibitem{BuffCheritat}
X.~Buff and A.~Ch\'eritat,
\emph{The Brjuno function continuously estimates the size of quadratic
Siegel disks}, 2004 preprint,
\href{https://arxiv.org/abs/math/0401044}{arXiv:math/0401044}.

\bibitem{CarminatiMarmi}
C.~Carminati and S.~Marmi,
Linearization of germs: regular dependence on the multiplier,
\emph{Bulletin de la Soci\'et\'e Math\'ematique de France} \textbf{136}
(2008), no.~4, 533--564.
\href{https://doi.org/10.24033/bsmf.2565}{doi:10.24033/bsmf.2565};
\href{https://arxiv.org/abs/0801.2844}{arXiv:0801.2844}.

\bibitem{FauvetMenousSauzin}
F.~Fauvet, F.~Menous, and D.~Sauzin,
Explicit linearization of one-dimensional germs through tree-expansions,
\emph{Bulletin de la Soci\'et\'e Math\'ematique de France} \textbf{146}
(2018), no.~2, 241--285.
\href{https://doi.org/10.24033/bsmf.2757}{doi:10.24033/bsmf.2757}.

\bibitem{Bounemoura}
A.~Bounemoura,
Optimal linearization of vector fields on the torus in non-analytic
Gevrey classes,
\emph{Annales de l'Institut Henri Poincar\'e C, Analyse Non Lin\'eaire}
\textbf{39} (2022), no.~3, 501--528.
\href{https://doi.org/10.4171/AIHPC/12}{doi:10.4171/AIHPC/12};
preprint \href{https://arxiv.org/abs/2006.05673}{arXiv:2006.05673} (2020).

\bibitem{Kamienski}
P.~Kamie\'nski,
\emph{The one-frequency cohomological equation, Brjuno-like functions and
Khintchine--L\'evy numbers}, preprint, 2019.
\href{https://arxiv.org/abs/1903.00311}{arXiv:1903.00311}.

\bibitem{Chevallier}
N.~Chevallier, J.~Lopes Dias, J.~P. Gaiv\~ao, A.~Marnat, and
N.~Moshchevitin,
\emph{Brjuno condition through best approximations and the linearization
problem}, July 2026.
\href{https://arxiv.org/abs/2607.25610}{arXiv:2607.25610v1}.

\bibitem{Lindahl}
K.-O. Lindahl,
Linearization in ultrametric dynamics in fields of characteristic
zero---equal characteristic case,
\emph{p-Adic Numbers, Ultrametric Analysis and Applications} \textbf{1}
(2009), no.~4, 307--316.
\href{https://doi.org/10.1134/S2070046609040049}{doi:10.1134/S2070046609040049};
\href{https://arxiv.org/abs/1111.1993}{arXiv:1111.1993}.

\bibitem{LindahlPadic}
K.-O. Lindahl,
\emph{The size of quadratic $p$-adic linearization disks},
\href{https://arxiv.org/abs/1311.4477}{arXiv:1311.4477} (2013),
version~1; in particular Corollary~C.

\bibitem{Hairer}
E.~Hairer,
\emph{Geometric Numerical Integration, Lecture 3: Backward error
analysis}, TU M\"unchen lectures, January--February 2010.
\url{https://www.unige.ch/~hairer/poly_geoint/week3.pdf}.

\bibitem{PerezMarco}
R.~P\'erez-Marco,
\emph{Convergence or generic divergence of Birkhoff normal form},
\href{https://arxiv.org/abs/math/0009028}{arXiv:math/0009028} (2000).

\bibitem{Krikorian}
R.~Krikorian,
\emph{On the divergence of Birkhoff normal forms},
\href{https://arxiv.org/abs/1906.01096}{arXiv:1906.01096} (2019).

\bibitem{BM}
A.~Berarducci and V.~Mantova,
Transseries as germs of surreal functions,
\emph{Transactions of the American Mathematical Society} \textbf{371}
(2019), 3549--3592.
\href{https://doi.org/10.1090/tran/7428}{doi:10.1090/tran/7428};
\href{https://arxiv.org/abs/1703.01995}{arXiv:1703.01995}.

\bibitem{CL}
R.~Cluckers and L.~Lipshitz,
Fields with analytic structure,
\emph{Journal of the European Mathematical Society} \textbf{13} (2011),
no.~4, 1147--1223.
\href{https://doi.org/10.4171/JEMS/278}{doi:10.4171/JEMS/278}.

\bibitem{MuellerStrohmaier}
J.~M\"uller and A.~Strohmaier,
The theory of Hahn-meromorphic functions, a holomorphic Fredholm theorem,
and its applications,
\emph{Analysis \& PDE} \textbf{7} (2014), no.~3, 745--770.
\href{https://doi.org/10.2140/apde.2014.7.745}{doi:10.2140/apde.2014.7.745};
\href{https://arxiv.org/abs/1205.0236}{arXiv:1205.0236}.

\end{thebibliography}
\end{document}
